\chapter{Communication Complexity --- Func Disjointness Lower Bound}
%%% generated by the blueprint pipeline from TCSlib.CommunicationComplexity.NewmanTheorem.FuncDisjointnessLowerBound
\section{Overview}

This module formalises a linear lower bound on the public-coin randomised communication
complexity of the $n$-bit Disjointness function, following the information-theoretic
argument of Rao and Yehudayoff (Chapter~6 of \emph{Communication Complexity}).
The key objects are the hard input distribution on \texttt{CommunicationComplexity.Functions.Disjointness.RandomizedLowerBound.HardSample},
the coarse conditioning variable $Z = (M, T, X_{<T}, Y_{>T})$, and two corrected mutual
information terms whose sum is bounded above by $O(\ell/n)$ and shown to be bounded below
by a constant whenever the protocol error on the hard distribution is at most $1/32$.

\section{Declarations}

\begin{definition}[Xbit]
\lean{CommunicationComplexity.Functions.Disjointness.RandomizedLowerBound.DisjointCoordinate.xBit}
\label{CommunicationComplexity.Functions.Disjointness.RandomizedLowerBound.DisjointCoordinate.xBit}
\leanok
Alice's bit in a disjoint coordinate-pair.
\end{definition}

\begin{definition}[Ybit]
\lean{CommunicationComplexity.Functions.Disjointness.RandomizedLowerBound.DisjointCoordinate.yBit}
\label{CommunicationComplexity.Functions.Disjointness.RandomizedLowerBound.DisjointCoordinate.yBit}
\leanok
Bob's bit in a disjoint coordinate-pair.
\end{definition}

\begin{theorem}[Number of disjoint-coordinate states]
\lean{CommunicationComplexity.Functions.Disjointness.RandomizedLowerBound.DisjointCoordinate.card}
\label{CommunicationComplexity.Functions.Disjointness.RandomizedLowerBound.DisjointCoordinate.card}
\leanok
\difficulty{1}
Consider the type of possible roles a single coordinate can play in the disjointness
communication problem: it belongs to neither of the two sets, to the left set only, or
to the right set only. This is a finite type, and it has exactly three elements.
\end{theorem}

\begin{theorem}[Alice's and Bob's bits are never both set]
\lean{CommunicationComplexity.Functions.Disjointness.RandomizedLowerBound.DisjointCoordinate.not_xBit_and_yBit}
\label{CommunicationComplexity.Functions.Disjointness.RandomizedLowerBound.DisjointCoordinate.not_xBit_and_yBit}
\leanok
\uses{CommunicationComplexity.Functions.Disjointness.RandomizedLowerBound.DisjointCoordinate.xBit, CommunicationComplexity.Functions.Disjointness.RandomizedLowerBound.DisjointCoordinate.yBit}
\difficulty{2}
Consider a coordinate in the disjointness distribution, whose state is one of three
cases — belonging to neither player, to Alice only, or to Bob only — and let
$\mathrm{xBit}(c)$ and $\mathrm{yBit}(c)$ denote Alice's and Bob's bits in that
coordinate. Then for every such coordinate $c$ it is never the case that both
$\mathrm{xBit}(c)$ and $\mathrm{yBit}(c)$ equal true.
\end{theorem}

\begin{definition}[Swap]
\lean{CommunicationComplexity.Functions.Disjointness.RandomizedLowerBound.DisjointCoordinate.swap}
\label{CommunicationComplexity.Functions.Disjointness.RandomizedLowerBound.DisjointCoordinate.swap}
\leanok
Swap Alice's and Bob's sides in a disjoint coordinate.
\end{definition}

\begin{theorem}[Double swap is the identity]
\lean{CommunicationComplexity.Functions.Disjointness.RandomizedLowerBound.DisjointCoordinate.swap_swap}
\label{CommunicationComplexity.Functions.Disjointness.RandomizedLowerBound.DisjointCoordinate.swap_swap}
\leanok
\uses{CommunicationComplexity.Functions.Disjointness.RandomizedLowerBound.DisjointCoordinate.swap}
\difficulty{2}
A disjoint coordinate records, for a single index, which of the two players' inputs
contains that index: neither, Alice only, or Bob only. Writing $\mathrm{swap}$ for the
operation that exchanges the two players' sides — fixing "neither" and interchanging
"Alice only" and "Bob only" — one has $\mathrm{swap}(\mathrm{swap}(c)) = c$ for every
disjoint coordinate $c$; that is, $\mathrm{swap}$ is an involution.
\end{theorem}

\begin{theorem}[Swapping sides exchanges Alice's and Bob's bits]
\lean{CommunicationComplexity.Functions.Disjointness.RandomizedLowerBound.DisjointCoordinate.xBit_swap}
\label{CommunicationComplexity.Functions.Disjointness.RandomizedLowerBound.DisjointCoordinate.xBit_swap}
\leanok
\uses{CommunicationComplexity.Functions.Disjointness.RandomizedLowerBound.DisjointCoordinate.swap, CommunicationComplexity.Functions.Disjointness.RandomizedLowerBound.DisjointCoordinate.xBit, CommunicationComplexity.Functions.Disjointness.RandomizedLowerBound.DisjointCoordinate.yBit}
\difficulty{2}
For every disjoint coordinate $c$, Alice's bit of the swapped coordinate
$\mathrm{swap}(c)$ equals Bob's bit of the original coordinate $c$.
\end{theorem}

\begin{definition}[Hardsample]
\lean{CommunicationComplexity.Functions.Disjointness.RandomizedLowerBound.HardSample}
\label{CommunicationComplexity.Functions.Disjointness.RandomizedLowerBound.HardSample}
\leanok
The sample space for the hard distribution.  The \texttt{other} coordinate at $T$ is ignored; keeping it in the sample space makes the ambient type a simple product-like finite type.
\end{definition}

\begin{definition}[Equivprod]
\lean{CommunicationComplexity.Functions.Disjointness.RandomizedLowerBound.HardSample.equivProd}
\label{CommunicationComplexity.Functions.Disjointness.RandomizedLowerBound.HardSample.equivProd}
\leanok
\uses{CommunicationComplexity.Functions.Disjointness.RandomizedLowerBound.HardSample}
Technical auxiliary result for \texttt{equivProd}.
\end{definition}

\begin{theorem}[Cardinality of the hard-distribution sample space]
\lean{CommunicationComplexity.Functions.Disjointness.RandomizedLowerBound.HardSample.card}
\label{CommunicationComplexity.Functions.Disjointness.RandomizedLowerBound.HardSample.card}
\leanok
\uses{CommunicationComplexity.Functions.Disjointness.RandomizedLowerBound.DisjointCoordinate.card, CommunicationComplexity.Functions.Disjointness.RandomizedLowerBound.HardSample, CommunicationComplexity.Functions.Disjointness.RandomizedLowerBound.HardSample.equivProd}
\difficulty{3}
A hard sample over $n \ge 1$ coordinates consists of a distinguished coordinate $T \in
\{0,1,\dots,n-1\}$, two bits $x_T, y_T \in \{0,1\}$, and an assignment giving each of
the $n$ coordinates one of the three tags *neither*, *left only*, or *right only*. The
number of such samples is exactly $n \cdot 4 \cdot 3^{n}$.
\end{theorem}

\begin{definition}[Uniformfin]
\lean{CommunicationComplexity.Functions.Disjointness.RandomizedLowerBound.uniformFin}
\label{CommunicationComplexity.Functions.Disjointness.RandomizedLowerBound.uniformFin}
\leanok
Uniform law on the special coordinate.
\end{definition}

\begin{theorem}[Singleton mass under the uniform law on a finite set]
\lean{CommunicationComplexity.Functions.Disjointness.RandomizedLowerBound.uniformFin_singleton}
\label{CommunicationComplexity.Functions.Disjointness.RandomizedLowerBound.uniformFin_singleton}
\leanok
\uses{CommunicationComplexity.Functions.Disjointness.RandomizedLowerBound.uniformFin}
\difficulty{2}
Let $n$ be a positive integer, and let the uniform probability law on the $n$-element
set $\{0,1,\dots,n-1\}$ (the finite type $\mathrm{Fin}\,n$) assign equal mass to each
point. Then for every element $i$, the mass of the singleton $\{i\}$ equals
$\dfrac{1}{n}$.
\end{theorem}

\begin{theorem}[Probability of a subset under the uniform law on a finite set]
\lean{CommunicationComplexity.Functions.Disjointness.RandomizedLowerBound.uniformFin_real}
\label{CommunicationComplexity.Functions.Disjointness.RandomizedLowerBound.uniformFin_real}
\leanok
\uses{CommunicationComplexity.Functions.Disjointness.RandomizedLowerBound.uniformFin, CommunicationComplexity.uniformOn_univ_measureReal_eq_card_subtype}
\difficulty{2}
Let $n \ge 1$ and equip the $n$-element set $\{0, 1, \dots, n-1\}$ with the uniform
probability law, under which each point carries mass $1/n$. Then for every subset $S$ of
this set, its probability is $\abs{S}/n$, the number of elements of $S$ divided by $n$.
\end{theorem}

\begin{theorem}[Total mass of the uniform law on a finite set]
\lean{CommunicationComplexity.Functions.Disjointness.RandomizedLowerBound.uniformFin_univ}
\label{CommunicationComplexity.Functions.Disjointness.RandomizedLowerBound.uniformFin_univ}
\leanok
\uses{CommunicationComplexity.Functions.Disjointness.RandomizedLowerBound.uniformFin, CommunicationComplexity.Functions.Disjointness.RandomizedLowerBound.uniformFin_real}
\difficulty{2}
Let $n$ be a positive integer, and let $\mu$ be the uniform probability law on the
$n$-element set $\{0, 1, \dots, n-1\}$, which assigns to each point equal weight. Then
the total mass of $\mu$ is one; that is, $\mu(\{0,1,\dots,n-1\}) = 1$.
\end{theorem}

\begin{definition}[Uniformbool]
\lean{CommunicationComplexity.Functions.Disjointness.RandomizedLowerBound.uniformBool}
\label{CommunicationComplexity.Functions.Disjointness.RandomizedLowerBound.uniformBool}
\leanok
Uniform law on one bit.
\end{definition}

\begin{theorem}[Singleton mass under the uniform law on a bit]
\lean{CommunicationComplexity.Functions.Disjointness.RandomizedLowerBound.uniformBool_singleton}
\label{CommunicationComplexity.Functions.Disjointness.RandomizedLowerBound.uniformBool_singleton}
\leanok
\uses{CommunicationComplexity.Functions.Disjointness.RandomizedLowerBound.uniformBool}
\difficulty{2}
Under the uniform probability law on the two Boolean values, each individual outcome $b$
carries mass exactly $\tfrac{1}{2}$; that is, the singleton event $\{b\}$ has
probability $\tfrac{1}{2}$ for either choice of $b$.
\end{theorem}

\begin{theorem}[Uniform measure of a subset of the Booleans]
\lean{CommunicationComplexity.Functions.Disjointness.RandomizedLowerBound.uniformBool_real}
\label{CommunicationComplexity.Functions.Disjointness.RandomizedLowerBound.uniformBool_real}
\leanok
\uses{CommunicationComplexity.Functions.Disjointness.RandomizedLowerBound.uniformBool, CommunicationComplexity.uniformOn_univ_measureReal_eq_card_subtype}
\difficulty{2}
Let $\mu$ denote the uniform probability law on the two-element type $\mathrm{Bool}$.
Then for every subset $S \subseteq \mathrm{Bool}$, the real-valued measure of $S$ equals
$|S|/2$, where $|S|$ is the number of Boolean values lying in $S$.
\end{theorem}

\begin{theorem}[Total mass of the uniform law on one bit]
\lean{CommunicationComplexity.Functions.Disjointness.RandomizedLowerBound.uniformBool_univ}
\label{CommunicationComplexity.Functions.Disjointness.RandomizedLowerBound.uniformBool_univ}
\leanok
\uses{CommunicationComplexity.Functions.Disjointness.RandomizedLowerBound.uniformBool, CommunicationComplexity.Functions.Disjointness.RandomizedLowerBound.uniformBool_real}
\difficulty{2}
The uniform law on a single bit, viewed as a measure on the two-element set of Boolean
values, assigns real-valued mass $1$ to the whole space; that is, the measure of the
universal set equals $1$.
\end{theorem}

\begin{definition}[Uniformdisjointcoordinatevector]
\lean{CommunicationComplexity.Functions.Disjointness.RandomizedLowerBound.uniformDisjointCoordinateVector}
\label{CommunicationComplexity.Functions.Disjointness.RandomizedLowerBound.uniformDisjointCoordinateVector}
\leanok
The uniform law on generated disjoint coordinate vectors.
\end{definition}

\begin{definition}[Uniformdisjointcoordinate]
\lean{CommunicationComplexity.Functions.Disjointness.RandomizedLowerBound.uniformDisjointCoordinate}
\label{CommunicationComplexity.Functions.Disjointness.RandomizedLowerBound.uniformDisjointCoordinate}
\leanok
The uniform law on one disjoint coordinate.
\end{definition}

\begin{theorem}[Uniform mass of a single disjoint-coordinate vector]
\lean{CommunicationComplexity.Functions.Disjointness.RandomizedLowerBound.uniformDisjointCoordinateVector_singleton}
\label{CommunicationComplexity.Functions.Disjointness.RandomizedLowerBound.uniformDisjointCoordinateVector_singleton}
\leanok
\uses{CommunicationComplexity.Functions.Disjointness.RandomizedLowerBound.DisjointCoordinate.card, CommunicationComplexity.Functions.Disjointness.RandomizedLowerBound.uniformDisjointCoordinateVector, CommunicationComplexity.uniformOn_univ_measureReal_eq_card_subtype}
\difficulty{4}
Let $n \ge 1$, and consider vectors of length $n$ whose entries are drawn from the
three-element set $\{\text{neither}, \text{leftOnly}, \text{rightOnly}\}$ of
disjoint-coordinate labels. Under the uniform law on these $3^n$ vectors, every
individual vector $c$ carries probability mass \[ \frac{1}{3^n}. \]
\end{theorem}

\begin{theorem}[Uniform law on disjointness-coordinate vectors]
\lean{CommunicationComplexity.Functions.Disjointness.RandomizedLowerBound.uniformDisjointCoordinateVector_real}
\label{CommunicationComplexity.Functions.Disjointness.RandomizedLowerBound.uniformDisjointCoordinateVector_real}
\leanok
\uses{CommunicationComplexity.Functions.Disjointness.RandomizedLowerBound.DisjointCoordinate.card, CommunicationComplexity.Functions.Disjointness.RandomizedLowerBound.uniformDisjointCoordinateVector, CommunicationComplexity.uniformOn_univ_measureReal_eq_card_subtype}
\difficulty{2}
Let $n \ge 1$, and consider the finite set of vectors $x \colon \{1,\dots,n\} \to
\{\text{neither}, \text{left only}, \text{right only}\}$ assigning to each coordinate
one of three labels, of which there are $3^n$ in total. Under the uniform probability
law on this set, every subset $S$ of such vectors has measure
\[
\frac{\abs{S}}{3^{n}},
\]
where $\abs{S}$ denotes the number of vectors lying in $S$.
\end{theorem}

\begin{theorem}[Total mass of the uniform law on disjoint-coordinate vectors]
\lean{CommunicationComplexity.Functions.Disjointness.RandomizedLowerBound.uniformDisjointCoordinateVector_univ}
\label{CommunicationComplexity.Functions.Disjointness.RandomizedLowerBound.uniformDisjointCoordinateVector_univ}
\leanok
\uses{CommunicationComplexity.Functions.Disjointness.RandomizedLowerBound.DisjointCoordinate.card, CommunicationComplexity.Functions.Disjointness.RandomizedLowerBound.uniformDisjointCoordinateVector, CommunicationComplexity.Functions.Disjointness.RandomizedLowerBound.uniformDisjointCoordinateVector_real}
\difficulty{2}
Fix a positive integer $n$, and consider the finite set of vectors $(x_1,\dots,x_n)$ of
length $n$ in which each coordinate $x_i$ takes one of the three values *neither*,
*left-only*, or *right-only*. Under the uniform probability law on this set, the total
mass — the measure assigned to the set of all such vectors — equals $1$.
\end{theorem}

\begin{theorem}[Uniform law assigns each disjoint coordinate probability one third]
\lean{CommunicationComplexity.Functions.Disjointness.RandomizedLowerBound.uniformDisjointCoordinate_singleton}
\label{CommunicationComplexity.Functions.Disjointness.RandomizedLowerBound.uniformDisjointCoordinate_singleton}
\leanok
\uses{CommunicationComplexity.Functions.Disjointness.RandomizedLowerBound.DisjointCoordinate.card, CommunicationComplexity.Functions.Disjointness.RandomizedLowerBound.uniformDisjointCoordinate, CommunicationComplexity.uniformOn_univ_measureReal_eq_card_subtype}
\difficulty{4}
Let a *disjoint coordinate* be one of the three symbols $\{\text{neither},\
\text{leftOnly},\ \text{rightOnly}\}$, and let $\mu$ be the uniform law on this
three-element set. Then for every disjoint coordinate $c$, the measure of the singleton
$\{c\}$ satisfies $\mu(\{c\}) = \tfrac{1}{3}$.
\end{theorem}

\begin{definition}[Xbit]
\lean{CommunicationComplexity.Functions.Disjointness.RandomizedLowerBound.xBit}
\label{CommunicationComplexity.Functions.Disjointness.RandomizedLowerBound.xBit}
\leanok
\uses{CommunicationComplexity.Functions.Disjointness.RandomizedLowerBound.DisjointCoordinate.xBit, CommunicationComplexity.Functions.Disjointness.RandomizedLowerBound.HardSample}
Alice's bit at coordinate $i$ under a hard-distribution sample.
\end{definition}

\begin{definition}[Ybit]
\lean{CommunicationComplexity.Functions.Disjointness.RandomizedLowerBound.yBit}
\label{CommunicationComplexity.Functions.Disjointness.RandomizedLowerBound.yBit}
\leanok
\uses{CommunicationComplexity.Functions.Disjointness.RandomizedLowerBound.DisjointCoordinate.yBit, CommunicationComplexity.Functions.Disjointness.RandomizedLowerBound.HardSample}
Bob's bit at coordinate $i$ under a hard-distribution sample.
\end{definition}

\begin{definition}[X]
\lean{CommunicationComplexity.Functions.Disjointness.RandomizedLowerBound.X}
\label{CommunicationComplexity.Functions.Disjointness.RandomizedLowerBound.X}
\leanok
\uses{CommunicationComplexity.Functions.Disjointness.RandomizedLowerBound.HardSample, CommunicationComplexity.Functions.Disjointness.RandomizedLowerBound.xBit}
Alice's input set under a hard-distribution sample.
\end{definition}

\begin{definition}[Y]
\lean{CommunicationComplexity.Functions.Disjointness.RandomizedLowerBound.Y}
\label{CommunicationComplexity.Functions.Disjointness.RandomizedLowerBound.Y}
\leanok
\uses{CommunicationComplexity.Functions.Disjointness.RandomizedLowerBound.HardSample, CommunicationComplexity.Functions.Disjointness.RandomizedLowerBound.yBit}
Bob's input set under a hard-distribution sample.
\end{definition}

\begin{definition}[Xvector]
\lean{CommunicationComplexity.Functions.Disjointness.RandomizedLowerBound.xVector}
\label{CommunicationComplexity.Functions.Disjointness.RandomizedLowerBound.xVector}
\leanok
\uses{CommunicationComplexity.Functions.Disjointness.RandomizedLowerBound.HardSample, CommunicationComplexity.Functions.Disjointness.RandomizedLowerBound.xBit}
Alice's full input vector under a hard-distribution sample.
\end{definition}

\begin{definition}[Yvector]
\lean{CommunicationComplexity.Functions.Disjointness.RandomizedLowerBound.yVector}
\label{CommunicationComplexity.Functions.Disjointness.RandomizedLowerBound.yVector}
\leanok
\uses{CommunicationComplexity.Functions.Disjointness.RandomizedLowerBound.HardSample, CommunicationComplexity.Functions.Disjointness.RandomizedLowerBound.yBit}
Bob's full input vector under a hard-distribution sample.
\end{definition}

\begin{definition}[Reverseboolvector]
\lean{CommunicationComplexity.Functions.Disjointness.RandomizedLowerBound.reverseBoolVector}
\label{CommunicationComplexity.Functions.Disjointness.RandomizedLowerBound.reverseBoolVector}
\leanok
Reverse the coordinate order of a boolean vector.
\end{definition}

\begin{theorem}[Reversing a boolean vector is an involution]
\lean{CommunicationComplexity.Functions.Disjointness.RandomizedLowerBound.reverseBoolVector_reverseBoolVector}
\label{CommunicationComplexity.Functions.Disjointness.RandomizedLowerBound.reverseBoolVector_reverseBoolVector}
\leanok
\uses{CommunicationComplexity.Functions.Disjointness.RandomizedLowerBound.reverseBoolVector}
\difficulty{2}
For any boolean vector $v = (v_0, v_1, \ldots, v_{n-1})$ of length $n \ge 1$, reversing
its coordinate order twice returns the original vector: if $v'$ denotes the vector
obtained by reversing the coordinate order of $v$, then reversing the coordinate order
of $v'$ recovers $v$ exactly.
\end{theorem}

\begin{theorem}[Injectivity of Boolean vector reversal]
\lean{CommunicationComplexity.Functions.Disjointness.RandomizedLowerBound.reverseBoolVector_injective}
\label{CommunicationComplexity.Functions.Disjointness.RandomizedLowerBound.reverseBoolVector_injective}
\leanok
\uses{CommunicationComplexity.Functions.Disjointness.RandomizedLowerBound.reverseBoolVector, CommunicationComplexity.Functions.Disjointness.RandomizedLowerBound.reverseBoolVector_reverseBoolVector}
\difficulty{3}
For a positive integer $n$, let $\rho \colon \{0,1\}^n \to \{0,1\}^n$ be the map that
reverses the coordinate order of a Boolean vector, sending $v = (v_0, v_1, \dots,
v_{n-1})$ to $(v_{n-1}, \dots, v_1, v_0)$. Then $\rho$ is injective: whenever two
Boolean vectors have the same reversal, they are equal.
\end{theorem}

\begin{definition}[Reverseset]
\lean{CommunicationComplexity.Functions.Disjointness.RandomizedLowerBound.reverseSet}
\label{CommunicationComplexity.Functions.Disjointness.RandomizedLowerBound.reverseSet}
\leanok
Reverse the coordinate order of a subset of coordinates.
\end{definition}

\begin{theorem}[Coordinate reversal is an involution on subsets]
\lean{CommunicationComplexity.Functions.Disjointness.RandomizedLowerBound.reverseSet_reverseSet}
\label{CommunicationComplexity.Functions.Disjointness.RandomizedLowerBound.reverseSet_reverseSet}
\leanok
\uses{CommunicationComplexity.Functions.Disjointness.RandomizedLowerBound.reverseSet}
\difficulty{2}
Fix a positive integer $n$, and for a subset $S \subseteq \{0, 1, \dots, n-1\}$ let
$\overline{S} = \{\, i : (n-1) - i \in S \,\}$ be the subset obtained by reversing the
coordinate order. Then applying this reversal twice returns the original set:
$\overline{\overline{S}} = S$ for every $S \subseteq \{0, 1, \dots, n-1\}$.
\end{theorem}

\begin{definition}[Input]
\lean{CommunicationComplexity.Functions.Disjointness.RandomizedLowerBound.input}
\label{CommunicationComplexity.Functions.Disjointness.RandomizedLowerBound.input}
\leanok
\uses{CommunicationComplexity.Functions.Disjointness.RandomizedLowerBound.HardSample, CommunicationComplexity.Functions.Disjointness.RandomizedLowerBound.X, CommunicationComplexity.Functions.Disjointness.RandomizedLowerBound.Y}
The input pair generated by a hard-distribution sample.
\end{definition}

\begin{definition}[Specialintersect]
\lean{CommunicationComplexity.Functions.Disjointness.RandomizedLowerBound.specialIntersect}
\label{CommunicationComplexity.Functions.Disjointness.RandomizedLowerBound.specialIntersect}
\leanok
\uses{CommunicationComplexity.Functions.Disjointness.RandomizedLowerBound.HardSample}
The event that the special coordinate is an actual intersection.
\end{definition}

\begin{definition}[Specialbitsevent]
\lean{CommunicationComplexity.Functions.Disjointness.RandomizedLowerBound.specialBitsEvent}
\label{CommunicationComplexity.Functions.Disjointness.RandomizedLowerBound.specialBitsEvent}
\leanok
\uses{CommunicationComplexity.Functions.Disjointness.RandomizedLowerBound.HardSample}
The event that the two special-coordinate bits have prescribed values.
\end{definition}

\begin{definition}[Specialcoordinateevent]
\lean{CommunicationComplexity.Functions.Disjointness.RandomizedLowerBound.specialCoordinateEvent}
\label{CommunicationComplexity.Functions.Disjointness.RandomizedLowerBound.specialCoordinateEvent}
\leanok
\uses{CommunicationComplexity.Functions.Disjointness.RandomizedLowerBound.HardSample}
The event that the special coordinate has a prescribed value.
\end{definition}

\begin{definition}[Specialzerozero]
\lean{CommunicationComplexity.Functions.Disjointness.RandomizedLowerBound.specialZeroZero}
\label{CommunicationComplexity.Functions.Disjointness.RandomizedLowerBound.specialZeroZero}
\leanok
\uses{CommunicationComplexity.Functions.Disjointness.RandomizedLowerBound.HardSample, CommunicationComplexity.Functions.Disjointness.RandomizedLowerBound.specialBitsEvent}
The event that both special-coordinate bits are zero.
\end{definition}

\begin{definition}[Disjointevent]
\lean{CommunicationComplexity.Functions.Disjointness.RandomizedLowerBound.disjointEvent}
\label{CommunicationComplexity.Functions.Disjointness.RandomizedLowerBound.disjointEvent}
\leanok
\uses{CommunicationComplexity.Functions.Disjointness.RandomizedLowerBound.HardSample, CommunicationComplexity.Functions.Disjointness.RandomizedLowerBound.X, CommunicationComplexity.Functions.Disjointness.RandomizedLowerBound.Y}
The event that the generated input pair is disjoint.
\end{definition}

\begin{definition}[Specialcoordinate]
\lean{CommunicationComplexity.Functions.Disjointness.RandomizedLowerBound.specialCoordinate}
\statementsource{roughgarden-cc-bootcamp}{pdf-d54d402b16e2-p016-b002}
\label{CommunicationComplexity.Functions.Disjointness.RandomizedLowerBound.specialCoordinate}
\leanok
\uses{CommunicationComplexity.Functions.Disjointness.RandomizedLowerBound.HardSample}
The special coordinate selected by the hard distribution.
\end{definition}

\begin{definition}[Specialx]
\lean{CommunicationComplexity.Functions.Disjointness.RandomizedLowerBound.specialX}
\label{CommunicationComplexity.Functions.Disjointness.RandomizedLowerBound.specialX}
\leanok
\uses{CommunicationComplexity.Functions.Disjointness.RandomizedLowerBound.HardSample}
Alice's bit at the special coordinate.
\end{definition}

\begin{definition}[Specialy]
\lean{CommunicationComplexity.Functions.Disjointness.RandomizedLowerBound.specialY}
\label{CommunicationComplexity.Functions.Disjointness.RandomizedLowerBound.specialY}
\leanok
\uses{CommunicationComplexity.Functions.Disjointness.RandomizedLowerBound.HardSample}
Bob's bit at the special coordinate.
\end{definition}

\begin{definition}[Xbeforespecial]
\lean{CommunicationComplexity.Functions.Disjointness.RandomizedLowerBound.xBeforeSpecial}
\label{CommunicationComplexity.Functions.Disjointness.RandomizedLowerBound.xBeforeSpecial}
\leanok
\uses{CommunicationComplexity.Functions.Disjointness.RandomizedLowerBound.HardSample, CommunicationComplexity.Functions.Disjointness.RandomizedLowerBound.xBit}
Alice's input bits before the special coordinate, padded by \texttt{false} elsewhere.
\end{definition}

\begin{definition}[Xlespecial]
\lean{CommunicationComplexity.Functions.Disjointness.RandomizedLowerBound.xLeSpecial}
\label{CommunicationComplexity.Functions.Disjointness.RandomizedLowerBound.xLeSpecial}
\leanok
\uses{CommunicationComplexity.Functions.Disjointness.RandomizedLowerBound.HardSample, CommunicationComplexity.Functions.Disjointness.RandomizedLowerBound.xBit}
Alice's input bits up to and including the special coordinate, padded by \texttt{false} elsewhere.
\end{definition}

\begin{definition}[Yafterspecial]
\lean{CommunicationComplexity.Functions.Disjointness.RandomizedLowerBound.yAfterSpecial}
\label{CommunicationComplexity.Functions.Disjointness.RandomizedLowerBound.yAfterSpecial}
\leanok
\uses{CommunicationComplexity.Functions.Disjointness.RandomizedLowerBound.HardSample, CommunicationComplexity.Functions.Disjointness.RandomizedLowerBound.yBit}
Bob's input bits after the special coordinate, padded by \texttt{false} elsewhere.
\end{definition}

\begin{definition}[Ygespecial]
\lean{CommunicationComplexity.Functions.Disjointness.RandomizedLowerBound.yGeSpecial}
\label{CommunicationComplexity.Functions.Disjointness.RandomizedLowerBound.yGeSpecial}
\leanok
\uses{CommunicationComplexity.Functions.Disjointness.RandomizedLowerBound.HardSample, CommunicationComplexity.Functions.Disjointness.RandomizedLowerBound.yBit}
Bob's input bits from the special coordinate onward, padded by \texttt{false} elsewhere.
\end{definition}

\begin{definition}[Coarseconditioning]
\lean{CommunicationComplexity.Functions.Disjointness.RandomizedLowerBound.coarseConditioning}
\label{CommunicationComplexity.Functions.Disjointness.RandomizedLowerBound.coarseConditioning}
\leanok
\uses{CommunicationComplexity.Functions.Disjointness.RandomizedLowerBound.HardSample, CommunicationComplexity.Functions.Disjointness.RandomizedLowerBound.specialCoordinate, CommunicationComplexity.Functions.Disjointness.RandomizedLowerBound.xBeforeSpecial, CommunicationComplexity.Functions.Disjointness.RandomizedLowerBound.yAfterSpecial}
The common coarse variable $T, X_{<T}, Y_{>T}$ used for special-coordinate fibers.
\end{definition}

\begin{definition}[Aliceclaimconditioning]
\lean{CommunicationComplexity.Functions.Disjointness.RandomizedLowerBound.aliceClaimConditioning}
\label{CommunicationComplexity.Functions.Disjointness.RandomizedLowerBound.aliceClaimConditioning}
\leanok
\uses{CommunicationComplexity.Functions.Disjointness.RandomizedLowerBound.HardSample, CommunicationComplexity.Functions.Disjointness.RandomizedLowerBound.specialCoordinate, CommunicationComplexity.Functions.Disjointness.RandomizedLowerBound.xBeforeSpecial, CommunicationComplexity.Functions.Disjointness.RandomizedLowerBound.yGeSpecial}
Alice's corrected special-coordinate conditioning variable: $T, X_{<T}, Y_{≥T}$.
\end{definition}

\begin{definition}[Alicedynamicconditioning]
\lean{CommunicationComplexity.Functions.Disjointness.RandomizedLowerBound.aliceDynamicConditioning}
\label{CommunicationComplexity.Functions.Disjointness.RandomizedLowerBound.aliceDynamicConditioning}
\leanok
\uses{CommunicationComplexity.Functions.Disjointness.RandomizedLowerBound.HardSample, CommunicationComplexity.Functions.Disjointness.RandomizedLowerBound.xBeforeSpecial, CommunicationComplexity.Functions.Disjointness.RandomizedLowerBound.yGeSpecial}
Alice's corrected conditioning data without the special coordinate: $X_{<T}, Y_{≥T}$.
\end{definition}

\begin{theorem}[Alice's conditioning as special coordinate paired with dynamic data]
\lean{CommunicationComplexity.Functions.Disjointness.RandomizedLowerBound.aliceClaimConditioning_eq_specialCoordinate_prod_dynamic}
\label{CommunicationComplexity.Functions.Disjointness.RandomizedLowerBound.aliceClaimConditioning_eq_specialCoordinate_prod_dynamic}
\leanok
\uses{CommunicationComplexity.Functions.Disjointness.RandomizedLowerBound.aliceClaimConditioning, CommunicationComplexity.Functions.Disjointness.RandomizedLowerBound.aliceDynamicConditioning, CommunicationComplexity.Functions.Disjointness.RandomizedLowerBound.specialCoordinate}
\difficulty{1}
Fix $n \ge 1$, and for a sample $\omega$ from the hard distribution write $T$ for its
special coordinate, $X_{<T}$ for Alice's input bits at the coordinates before $T$
(padded with $\mathrm{false}$ from $T$ onward), and $Y_{\ge T}$ for Bob's input bits at
the coordinates from $T$ onward (padded with $\mathrm{false}$ before $T$). Then Alice's
corrected special-coordinate conditioning variable, namely the triple $(T,\, X_{<T},\,
Y_{\ge T})$, coincides with the pair consisting of $T$ together with Alice's corrected
dynamic conditioning data $(X_{<T},\, Y_{\ge T})$; that is, the two agree as functions
of $\omega$.
\end{theorem}

\begin{definition}[Bobclaimconditioning]
\lean{CommunicationComplexity.Functions.Disjointness.RandomizedLowerBound.bobClaimConditioning}
\label{CommunicationComplexity.Functions.Disjointness.RandomizedLowerBound.bobClaimConditioning}
\leanok
\uses{CommunicationComplexity.Functions.Disjointness.RandomizedLowerBound.HardSample, CommunicationComplexity.Functions.Disjointness.RandomizedLowerBound.specialCoordinate, CommunicationComplexity.Functions.Disjointness.RandomizedLowerBound.xLeSpecial, CommunicationComplexity.Functions.Disjointness.RandomizedLowerBound.yAfterSpecial}
Bob's corrected special-coordinate conditioning variable: $T, X_{≤T}, Y_{>T}$.
\end{definition}

\begin{definition}[Fixedxbit]
\lean{CommunicationComplexity.Functions.Disjointness.RandomizedLowerBound.fixedXBit}
\label{CommunicationComplexity.Functions.Disjointness.RandomizedLowerBound.fixedXBit}
\leanok
\uses{CommunicationComplexity.Functions.Disjointness.RandomizedLowerBound.HardSample, CommunicationComplexity.Functions.Disjointness.RandomizedLowerBound.xBit}
Alice's bit at a fixed coordinate.
\end{definition}

\begin{definition}[Fixedxstrictprefix]
\lean{CommunicationComplexity.Functions.Disjointness.RandomizedLowerBound.fixedXStrictPrefix}
\label{CommunicationComplexity.Functions.Disjointness.RandomizedLowerBound.fixedXStrictPrefix}
\leanok
\uses{CommunicationComplexity.Functions.Disjointness.RandomizedLowerBound.HardSample, CommunicationComplexity.Functions.Disjointness.RandomizedLowerBound.xBit}
Alice's fixed-coordinate strict prefix $X_{<i}$, represented without padding.
\end{definition}

\begin{definition}[Fixedxbefore]
\lean{CommunicationComplexity.Functions.Disjointness.RandomizedLowerBound.fixedXBefore}
\label{CommunicationComplexity.Functions.Disjointness.RandomizedLowerBound.fixedXBefore}
\leanok
\uses{CommunicationComplexity.Functions.Disjointness.RandomizedLowerBound.HardSample, CommunicationComplexity.Functions.Disjointness.RandomizedLowerBound.xBit}
Alice's fixed-coordinate $X_{<i}$ prefix.
\end{definition}

\begin{definition}[Fixedybefore]
\lean{CommunicationComplexity.Functions.Disjointness.RandomizedLowerBound.fixedYBefore}
\label{CommunicationComplexity.Functions.Disjointness.RandomizedLowerBound.fixedYBefore}
\leanok
\uses{CommunicationComplexity.Functions.Disjointness.RandomizedLowerBound.HardSample, CommunicationComplexity.Functions.Disjointness.RandomizedLowerBound.yBit}
Bob's fixed-coordinate $Y_{<i}$ prefix.
\end{definition}

\begin{definition}[Fixedyge]
\lean{CommunicationComplexity.Functions.Disjointness.RandomizedLowerBound.fixedYGe}
\label{CommunicationComplexity.Functions.Disjointness.RandomizedLowerBound.fixedYGe}
\leanok
\uses{CommunicationComplexity.Functions.Disjointness.RandomizedLowerBound.HardSample, CommunicationComplexity.Functions.Disjointness.RandomizedLowerBound.yBit}
Bob's fixed-coordinate $Y_{≥i}$ suffix.
\end{definition}

\begin{definition}[Fixedaliceconditioning]
\lean{CommunicationComplexity.Functions.Disjointness.RandomizedLowerBound.fixedAliceConditioning}
\label{CommunicationComplexity.Functions.Disjointness.RandomizedLowerBound.fixedAliceConditioning}
\leanok
\uses{CommunicationComplexity.Functions.Disjointness.RandomizedLowerBound.HardSample, CommunicationComplexity.Functions.Disjointness.RandomizedLowerBound.fixedXBefore, CommunicationComplexity.Functions.Disjointness.RandomizedLowerBound.fixedYGe}
The fixed-coordinate Alice conditioning variable $X_{<i}, Y_{≥i}$.
\end{definition}

\begin{definition}[Fixedalicefullyconditioning]
\lean{CommunicationComplexity.Functions.Disjointness.RandomizedLowerBound.fixedAliceFullYConditioning}
\label{CommunicationComplexity.Functions.Disjointness.RandomizedLowerBound.fixedAliceFullYConditioning}
\leanok
\uses{CommunicationComplexity.Functions.Disjointness.RandomizedLowerBound.HardSample, CommunicationComplexity.Functions.Disjointness.RandomizedLowerBound.fixedXStrictPrefix, CommunicationComplexity.Functions.Disjointness.RandomizedLowerBound.yVector}
The fixed-coordinate conditioning variable $(X_{<i}, Y)$ used in the chain-rule sum. The $X_{<i}$ prefix is represented without padding so that recodings are injective.
\end{definition}

\begin{definition}[Fixedalicechainconditioningvalue]
\lean{CommunicationComplexity.Functions.Disjointness.RandomizedLowerBound.fixedAliceChainConditioningValue}
\label{CommunicationComplexity.Functions.Disjointness.RandomizedLowerBound.fixedAliceChainConditioningValue}
\leanok
Recode $(X_{<i}, Y)$ as $(Y_{<i}, X_{<i}, Y_{≥i})$, the conditioning variable produced by the chain-rule split.
\end{definition}

\begin{theorem}[Injectivity of the chain-rule conditioning recoding]
\lean{CommunicationComplexity.Functions.Disjointness.RandomizedLowerBound.fixedAliceChainConditioningValue_injective}
\label{CommunicationComplexity.Functions.Disjointness.RandomizedLowerBound.fixedAliceChainConditioningValue_injective}
\leanok
\uses{CommunicationComplexity.Functions.Disjointness.RandomizedLowerBound.fixedAliceChainConditioningValue}
\difficulty{4}
Fix $n \ge 1$ and an index $i \in \{0, 1, \dots, n-1\}$, and consider Boolean strings
valued in $\bbf_2 = \{0,1\}$. Define a map that sends a pair $(x, y)$, where $x \colon
\{0,\dots,i-1\} \to \bbf_2$ records the first $i$ coordinates of Alice's input and $y
\colon \{0,\dots,n-1\} \to \bbf_2$ is Bob's full input, to the triple $(u, v, w)$ of
Boolean functions on all of $\{0,\dots,n-1\}$ given by $u(j) = y(j)$ if $j < i$ and
$u(j) = 0$ otherwise, $v(j) = x(j)$ if $j < i$ and $v(j) = 0$ otherwise, and $w(j) =
y(j)$ if $j \ge i$ and $w(j) = 0$ otherwise; that is, $(x,y) \mapsto (y_{<i}, x_{<i},
y_{\ge i})$ with all out-of-range coordinates zero-padded. Then this map is injective.
\end{theorem}

\begin{theorem}[Chain-rule recoding of the full-Y conditioning variable]
\lean{CommunicationComplexity.Functions.Disjointness.RandomizedLowerBound.fixedAliceChainConditioningValue_fixedAliceFullYConditioning}
\label{CommunicationComplexity.Functions.Disjointness.RandomizedLowerBound.fixedAliceChainConditioningValue_fixedAliceFullYConditioning}
\leanok
\uses{CommunicationComplexity.Functions.Disjointness.RandomizedLowerBound.fixedAliceChainConditioningValue, CommunicationComplexity.Functions.Disjointness.RandomizedLowerBound.fixedAliceConditioning, CommunicationComplexity.Functions.Disjointness.RandomizedLowerBound.fixedAliceFullYConditioning, CommunicationComplexity.Functions.Disjointness.RandomizedLowerBound.fixedXBefore, CommunicationComplexity.Functions.Disjointness.RandomizedLowerBound.fixedXStrictPrefix, CommunicationComplexity.Functions.Disjointness.RandomizedLowerBound.fixedYBefore, CommunicationComplexity.Functions.Disjointness.RandomizedLowerBound.fixedYGe, CommunicationComplexity.Functions.Disjointness.RandomizedLowerBound.yVector}
\difficulty{3}
Fix a positive integer $n$ and a coordinate $i \in \{0,1,\dots,n-1\}$. Let $R$ be the
map that sends a pair $(a,b)$ — where $a$ is a Boolean vector indexed by
$\{0,\dots,i-1\}$ and $b$ is a Boolean vector of length $n$ — to
\[
R(a,b) = \Bigl(\ j\mapsto b_j\,[\,j<i\,],\ \ \bigl(\,j\mapsto a_j\,[\,j<i\,],\ \
j\mapsto b_j\,[\,i\le j\,]\,\bigr)\ \Bigr),
\]
whose three components are all Boolean vectors of length $n$; here $[\,\cdot\,]$ is the
Iverson bracket, so each component takes the value $0$ at every index falling outside
the range named in its own bracket. For every sample $\omega$ of the hard distribution,
let $b$ be Bob's full induced input vector of length $n$ and let $a$ be Alice's induced
input restricted to the coordinates $\{0,\dots,i-1\}$ (a Boolean vector on those
indices, carrying no padding). Then
\[
R(a,b) = \bigl(\,y_{<i},\ (\,x_{<i},\ y_{\ge i}\,)\,\bigr),
\]
where $x_{<i}$, $y_{<i}$, and $y_{\ge i}$ are the length-$n$ Boolean vectors that equal
Alice's induced input at each index below $i$, Bob's induced input at each index below
$i$, and Bob's induced input at each index at least $i$, respectively, and are $0$ at
every other index.
\end{theorem}

\begin{definition}[Aliceclaimconditioningyfalsevalues]
\lean{CommunicationComplexity.Functions.Disjointness.RandomizedLowerBound.aliceClaimConditioningYFalseValues}
\label{CommunicationComplexity.Functions.Disjointness.RandomizedLowerBound.aliceClaimConditioningYFalseValues}
\leanok
The values of Alice's corrected conditioning variable for which $Y_{T} = false$.
\end{definition}

\begin{theorem}[Value of Bob's truncated input at the special coordinate]
\lean{CommunicationComplexity.Functions.Disjointness.RandomizedLowerBound.yGeSpecial_specialCoordinate}
\label{CommunicationComplexity.Functions.Disjointness.RandomizedLowerBound.yGeSpecial_specialCoordinate}
\leanok
\uses{CommunicationComplexity.Functions.Disjointness.RandomizedLowerBound.HardSample, CommunicationComplexity.Functions.Disjointness.RandomizedLowerBound.specialCoordinate, CommunicationComplexity.Functions.Disjointness.RandomizedLowerBound.specialY, CommunicationComplexity.Functions.Disjointness.RandomizedLowerBound.yBit, CommunicationComplexity.Functions.Disjointness.RandomizedLowerBound.yGeSpecial}
\difficulty{1}
Fix a sample $\omega$ of the hard distribution, and write $T$ for its special
coordinate. Let $y_{\ge T}$ be Bob's truncated input vector, whose value at a coordinate
$i$ is Bob's bit at $i$ when $T \le i$ and is $\mathrm{false}$ otherwise. Then
evaluating this vector at the special coordinate itself gives Bob's bit at the special
coordinate: $y_{\ge T}(T)$ equals $y_T$.
\end{theorem}

\begin{theorem}[Event that Bob's special bit is false as a conditioning-variable preimage]
\lean{CommunicationComplexity.Functions.Disjointness.RandomizedLowerBound.specialY_false_eq_preimage_aliceClaimConditioningYFalseValues}
\label{CommunicationComplexity.Functions.Disjointness.RandomizedLowerBound.specialY_false_eq_preimage_aliceClaimConditioningYFalseValues}
\leanok
\uses{CommunicationComplexity.Functions.Disjointness.RandomizedLowerBound.HardSample, CommunicationComplexity.Functions.Disjointness.RandomizedLowerBound.aliceClaimConditioning, CommunicationComplexity.Functions.Disjointness.RandomizedLowerBound.aliceClaimConditioningYFalseValues, CommunicationComplexity.Functions.Disjointness.RandomizedLowerBound.specialY, CommunicationComplexity.Functions.Disjointness.RandomizedLowerBound.yGeSpecial_specialCoordinate}
\difficulty{2}
Fix a positive integer $n$ and work over the hard-distribution sample space, whose
samples $\omega$ record a special coordinate $T \in \{0,\dots,n-1\}$ together with
enough data to determine Alice's and Bob's bits at every coordinate. Let $Y_T(\omega)$
denote Bob's bit at the special coordinate, and let the conditioning variable be the
triple $\bigl(T,\,X_{<T},\,Y_{\ge T}\bigr)$ consisting of the special coordinate,
Alice's bits strictly before it (padded with $\mathrm{false}$ elsewhere), and Bob's bits
from it onward (padded with $\mathrm{false}$ elsewhere). Then the set of samples with
$Y_T = \mathrm{false}$ coincides exactly with the set of samples whose conditioning
variable $(t,x,y)$ satisfies $y(t) = \mathrm{false}$.
\end{theorem}

\begin{definition}[Dualconditioningvalue]
\lean{CommunicationComplexity.Functions.Disjointness.RandomizedLowerBound.dualConditioningValue}
\label{CommunicationComplexity.Functions.Disjointness.RandomizedLowerBound.dualConditioningValue}
\leanok
\uses{CommunicationComplexity.Functions.Disjointness.RandomizedLowerBound.reverseBoolVector}
The conditioning-value recoding induced by reversing coordinates and swapping Alice/Bob.
\end{definition}

\begin{theorem}[The dual conditioning-value recoding is an involution]
\lean{CommunicationComplexity.Functions.Disjointness.RandomizedLowerBound.dualConditioningValue_dualConditioningValue}
\label{CommunicationComplexity.Functions.Disjointness.RandomizedLowerBound.dualConditioningValue_dualConditioningValue}
\leanok
\uses{CommunicationComplexity.Functions.Disjointness.RandomizedLowerBound.dualConditioningValue, CommunicationComplexity.Functions.Disjointness.RandomizedLowerBound.reverseBoolVector_reverseBoolVector}
\difficulty{2}
Fix a positive integer $n$, and let a triple $c = (i, x, y)$ consist of a coordinate $i
\in \{0, 1, \dots, n-1\}$ together with two Boolean vectors $x, y \colon \{0, 1, \dots,
n-1\} \to \{\text{true}, \text{false}\}$. Applying the dual conditioning-value recoding
— which sends $c$ to $(n-1-i,\, y \circ \mathrm{rev},\, x \circ \mathrm{rev})$, where
$\mathrm{rev}$ reverses the coordinate order — twice in succession returns the original
triple $c$; that is, the recoding is an involution.
\end{theorem}

\begin{theorem}[Injectivity of the coordinate-reversal duality map]
\lean{CommunicationComplexity.Functions.Disjointness.RandomizedLowerBound.dualConditioningValue_injective}
\label{CommunicationComplexity.Functions.Disjointness.RandomizedLowerBound.dualConditioningValue_injective}
\leanok
\uses{CommunicationComplexity.Functions.Disjointness.RandomizedLowerBound.dualConditioningValue, CommunicationComplexity.Functions.Disjointness.RandomizedLowerBound.dualConditioningValue_dualConditioningValue}
\difficulty{3}
Fix an integer $n \ge 1$, and for a boolean vector $x \colon \{0,1,\dots,n-1\} \to
\{\text{true},\text{false}\}$ write $\bar x$ for its coordinate reversal, $\bar x(i) =
x(n-1-i)$. Consider the map on triples $(i, x, y)$, where $i \in \{0,1,\dots,n-1\}$ and
$x, y$ are boolean vectors of length $n$, defined by
\[
  (i, x, y) \;\longmapsto\; \bigl(n-1-i,\ \bar y,\ \bar x\bigr),
\]
which reverses the index, swaps the two vectors, and reverses the coordinate order of
each. This map is injective.
\end{theorem}

\begin{definition}[Protocoltype]
\lean{CommunicationComplexity.Functions.Disjointness.RandomizedLowerBound.ProtocolType}
\statementsource{roughgarden-cc-bootcamp}{pdf-d54d402b16e2-p002-b002}
\label{CommunicationComplexity.Functions.Disjointness.RandomizedLowerBound.ProtocolType}
\leanok
\uses{CommunicationComplexity.Deterministic.Protocol}
The deterministic protocol type for disjointness inputs of length $n$.
\end{definition}

\begin{definition}[Transcripttype]
\lean{CommunicationComplexity.Functions.Disjointness.RandomizedLowerBound.TranscriptType}
\statementsource{roughgarden-cc-bootcamp}{pdf-d54d402b16e2-p004-b005}
\label{CommunicationComplexity.Functions.Disjointness.RandomizedLowerBound.TranscriptType}
\leanok
\uses{CommunicationComplexity.Deterministic.Protocol.Transcript, CommunicationComplexity.Functions.Disjointness.RandomizedLowerBound.ProtocolType}
The syntactic transcript type for a disjointness protocol.
\end{definition}

\begin{definition}[Rawztype]
\lean{CommunicationComplexity.Functions.Disjointness.RandomizedLowerBound.RawZType}
\label{CommunicationComplexity.Functions.Disjointness.RandomizedLowerBound.RawZType}
\leanok
\uses{CommunicationComplexity.Functions.Disjointness.RandomizedLowerBound.ProtocolType, CommunicationComplexity.Functions.Disjointness.RandomizedLowerBound.TranscriptType}
The raw type of values of the $Z = (M, T, X_{<T}, Y_{>T})$ variable for a protocol.
\end{definition}

\begin{theorem}[Extensionality for raw $Z$-variables]
\lean{CommunicationComplexity.Functions.Disjointness.RandomizedLowerBound.RawZType.ext}
\label{CommunicationComplexity.Functions.Disjointness.RandomizedLowerBound.RawZType.ext}
\leanok
\uses{CommunicationComplexity.Functions.Disjointness.RandomizedLowerBound.RawZType}
\difficulty{3}
Fix a positive integer $n$ and a deterministic disjointness protocol $p$ on inputs of
length $n$, and let $z$ and $z'$ be two raw $Z$-values for $p$ — each consisting of a
syntactic transcript of $p$, a special coordinate in $\{0,1,\dots,n-1\}$, and two
Boolean vectors indexed by $\{0,1,\dots,n-1\}$ recording the bits before and after that
coordinate. If $z$ and $z'$ have the same transcript, the same special coordinate, the
same "before" vector, and the same "after" vector, then $z = z'$.
\end{theorem}

\begin{definition}[Equivprod]
\lean{CommunicationComplexity.Functions.Disjointness.RandomizedLowerBound.RawZType.equivProd}
\label{CommunicationComplexity.Functions.Disjointness.RandomizedLowerBound.RawZType.equivProd}
\leanok
\uses{CommunicationComplexity.Functions.Disjointness.RandomizedLowerBound.ProtocolType, CommunicationComplexity.Functions.Disjointness.RandomizedLowerBound.RawZType, CommunicationComplexity.Functions.Disjointness.RandomizedLowerBound.TranscriptType}
A raw $Z$ value is equivalent to the product of its named fields.
\end{definition}

\begin{definition}[Rawzvariable]
\lean{CommunicationComplexity.Functions.Disjointness.RandomizedLowerBound.rawZVariable}
\label{CommunicationComplexity.Functions.Disjointness.RandomizedLowerBound.rawZVariable}
\leanok
\uses{CommunicationComplexity.Deterministic.Protocol.transcript, CommunicationComplexity.Functions.Disjointness.RandomizedLowerBound.HardSample, CommunicationComplexity.Functions.Disjointness.RandomizedLowerBound.ProtocolType, CommunicationComplexity.Functions.Disjointness.RandomizedLowerBound.RawZType, CommunicationComplexity.Functions.Disjointness.RandomizedLowerBound.input, CommunicationComplexity.Functions.Disjointness.RandomizedLowerBound.specialCoordinate, CommunicationComplexity.Functions.Disjointness.RandomizedLowerBound.xBeforeSpecial, CommunicationComplexity.Functions.Disjointness.RandomizedLowerBound.yAfterSpecial}
The raw $Z = (M, T, X_{<T}, Y_{>T})$ variable used for coarse special-coordinate fibers.
\end{definition}

\begin{definition}[Ztype]
\lean{CommunicationComplexity.Functions.Disjointness.RandomizedLowerBound.ZType}
\label{CommunicationComplexity.Functions.Disjointness.RandomizedLowerBound.ZType}
\leanok
\uses{CommunicationComplexity.Functions.Disjointness.RandomizedLowerBound.HardSample, CommunicationComplexity.Functions.Disjointness.RandomizedLowerBound.ProtocolType, CommunicationComplexity.Functions.Disjointness.RandomizedLowerBound.RawZType, CommunicationComplexity.Functions.Disjointness.RandomizedLowerBound.rawZVariable}
The type of achievable values of the $Z = (M, T, X_{<T}, Y_{>T})$ variable for a protocol.
\end{definition}

\begin{definition}[Transcript]
\lean{CommunicationComplexity.Functions.Disjointness.RandomizedLowerBound.ZType.transcript}
\label{CommunicationComplexity.Functions.Disjointness.RandomizedLowerBound.ZType.transcript}
\leanok
\uses{CommunicationComplexity.Functions.Disjointness.RandomizedLowerBound.TranscriptType, CommunicationComplexity.Functions.Disjointness.RandomizedLowerBound.ZType}
The transcript component of an achievable $Z$ value.
\end{definition}

\begin{definition}[Specialcoordinate]
\lean{CommunicationComplexity.Functions.Disjointness.RandomizedLowerBound.ZType.specialCoordinate}
\label{CommunicationComplexity.Functions.Disjointness.RandomizedLowerBound.ZType.specialCoordinate}
\leanok
\uses{CommunicationComplexity.Functions.Disjointness.RandomizedLowerBound.ZType}
The special-coordinate component of an achievable $Z$ value.
\end{definition}

\begin{definition}[Xbefore]
\lean{CommunicationComplexity.Functions.Disjointness.RandomizedLowerBound.ZType.xBefore}
\label{CommunicationComplexity.Functions.Disjointness.RandomizedLowerBound.ZType.xBefore}
\leanok
\uses{CommunicationComplexity.Functions.Disjointness.RandomizedLowerBound.ZType}
The $X_{<T}$ component of an achievable $Z$ value.
\end{definition}

\begin{definition}[Yafter]
\lean{CommunicationComplexity.Functions.Disjointness.RandomizedLowerBound.ZType.yAfter}
\label{CommunicationComplexity.Functions.Disjointness.RandomizedLowerBound.ZType.yAfter}
\leanok
\uses{CommunicationComplexity.Functions.Disjointness.RandomizedLowerBound.ZType}
The $Y_{>T}$ component of an achievable $Z$ value.
\end{definition}

\begin{definition}[Raw]
\lean{CommunicationComplexity.Functions.Disjointness.RandomizedLowerBound.ZType.raw}
\label{CommunicationComplexity.Functions.Disjointness.RandomizedLowerBound.ZType.raw}
\leanok
\uses{CommunicationComplexity.Functions.Disjointness.RandomizedLowerBound.RawZType, CommunicationComplexity.Functions.Disjointness.RandomizedLowerBound.ZType}
The raw value underlying an achievable $Z$ value.
\end{definition}

\begin{theorem}[Every achievable Z-value comes from a hard sample]
\lean{CommunicationComplexity.Functions.Disjointness.RandomizedLowerBound.ZType.achievable}
\label{CommunicationComplexity.Functions.Disjointness.RandomizedLowerBound.ZType.achievable}
\leanok
\uses{CommunicationComplexity.Functions.Disjointness.RandomizedLowerBound.HardSample, CommunicationComplexity.Functions.Disjointness.RandomizedLowerBound.ZType, CommunicationComplexity.Functions.Disjointness.RandomizedLowerBound.ZType.raw, CommunicationComplexity.Functions.Disjointness.RandomizedLowerBound.rawZVariable}
\difficulty{1}
Fix a positive integer $n$ and a deterministic protocol $p$ for disjointness inputs of
length $n$, and recall that the raw $Z$-variable sends each hard-distribution sample
$\omega$ to the tuple $Z(\omega) = (M, T, X_{<T}, Y_{>T})$ consisting of the transcript
of $p$ on $\omega$'s input, the special coordinate, and the input bits before and after
that coordinate. An achievable $Z$-value is a raw value together with a witness that it
arises this way. Then for every achievable $Z$-value $z$, there exists a
hard-distribution sample $\omega$ such that $Z(\omega)$ equals the underlying raw value
of $z$.
\end{theorem}

\begin{theorem}[Componentwise equality of the Z variable]
\lean{CommunicationComplexity.Functions.Disjointness.RandomizedLowerBound.ZType.ext}
\label{CommunicationComplexity.Functions.Disjointness.RandomizedLowerBound.ZType.ext}
\leanok
\uses{CommunicationComplexity.Functions.Disjointness.RandomizedLowerBound.RawZType.ext, CommunicationComplexity.Functions.Disjointness.RandomizedLowerBound.ZType, CommunicationComplexity.Functions.Disjointness.RandomizedLowerBound.ZType.specialCoordinate, CommunicationComplexity.Functions.Disjointness.RandomizedLowerBound.ZType.transcript, CommunicationComplexity.Functions.Disjointness.RandomizedLowerBound.ZType.xBefore, CommunicationComplexity.Functions.Disjointness.RandomizedLowerBound.ZType.yAfter}
\difficulty{1}
Fix a deterministic communication protocol $p$ for disjointness inputs of length $n$,
and consider achievable values of the variable $Z = (M, T, X_{<T}, Y_{>T})$, each
carrying a transcript $M$, a special coordinate $T \in \{0, \dots, n-1\}$, a Boolean
string $X_{<T}$ recording Alice's bits before the special coordinate, and a Boolean
string $Y_{>T}$ recording Bob's bits after it. Then two such achievable values $z$ and
$z'$ coincide whenever they have the same transcript, the same special coordinate, the
same $X_{<T}$, and the same $Y_{>T}$.
\end{theorem}

\begin{definition}[Zvariable]
\lean{CommunicationComplexity.Functions.Disjointness.RandomizedLowerBound.zVariable}
\label{CommunicationComplexity.Functions.Disjointness.RandomizedLowerBound.zVariable}
\leanok
\uses{CommunicationComplexity.Functions.Disjointness.RandomizedLowerBound.HardSample, CommunicationComplexity.Functions.Disjointness.RandomizedLowerBound.ProtocolType, CommunicationComplexity.Functions.Disjointness.RandomizedLowerBound.ZType, CommunicationComplexity.Functions.Disjointness.RandomizedLowerBound.rawZVariable}
The $Z = (M, T, X_{<T}, Y_{>T})$ variable used for coarse special-coordinate fibers.
\end{definition}

\begin{definition}[Zfiber]
\lean{CommunicationComplexity.Functions.Disjointness.RandomizedLowerBound.zFiber}
\label{CommunicationComplexity.Functions.Disjointness.RandomizedLowerBound.zFiber}
\leanok
\uses{CommunicationComplexity.Functions.Disjointness.RandomizedLowerBound.HardSample, CommunicationComplexity.Functions.Disjointness.RandomizedLowerBound.ProtocolType, CommunicationComplexity.Functions.Disjointness.RandomizedLowerBound.ZType, CommunicationComplexity.Functions.Disjointness.RandomizedLowerBound.zVariable}
The fiber of the $Z = (M, T, X_{<T}, Y_{>T})$ variable above a value $z$.
\end{definition}

\begin{definition}[Zoutput]
\lean{CommunicationComplexity.Functions.Disjointness.RandomizedLowerBound.zOutput}
\label{CommunicationComplexity.Functions.Disjointness.RandomizedLowerBound.zOutput}
\leanok
\uses{CommunicationComplexity.Deterministic.Protocol.Transcript.output, CommunicationComplexity.Functions.Disjointness.RandomizedLowerBound.ProtocolType, CommunicationComplexity.Functions.Disjointness.RandomizedLowerBound.ZType, CommunicationComplexity.Functions.Disjointness.RandomizedLowerBound.ZType.transcript}
The protocol output determined by a $Z$ value, through its transcript leaf.
\end{definition}

\begin{theorem}[Protocol output agrees with the Z-value output]
\lean{CommunicationComplexity.Functions.Disjointness.RandomizedLowerBound.run_eq_zOutput_of_zVariable_eq}
\label{CommunicationComplexity.Functions.Disjointness.RandomizedLowerBound.run_eq_zOutput_of_zVariable_eq}
\leanok
\uses{CommunicationComplexity.Deterministic.Protocol.Transcript.output, CommunicationComplexity.Deterministic.Protocol.run, CommunicationComplexity.Deterministic.Protocol.run_eq_transcript_output, CommunicationComplexity.Functions.Disjointness.RandomizedLowerBound.HardSample, CommunicationComplexity.Functions.Disjointness.RandomizedLowerBound.ProtocolType, CommunicationComplexity.Functions.Disjointness.RandomizedLowerBound.X, CommunicationComplexity.Functions.Disjointness.RandomizedLowerBound.Y, CommunicationComplexity.Functions.Disjointness.RandomizedLowerBound.ZType, CommunicationComplexity.Functions.Disjointness.RandomizedLowerBound.ZType.transcript, CommunicationComplexity.Functions.Disjointness.RandomizedLowerBound.input, CommunicationComplexity.Functions.Disjointness.RandomizedLowerBound.zOutput, CommunicationComplexity.Functions.Disjointness.RandomizedLowerBound.zVariable}
\difficulty{3}
Fix $n \ge 1$ and a deterministic protocol $p$ for disjointness inputs of length $n$.
Let $\omega$ be a sample of the hard distribution, and let $z$ be an achievable value of
the $Z$-variable $Z = (M, T, X_{<T}, Y_{>T})$ for $p$. If $z$ is exactly the value of
the $Z$-variable at $\omega$, then running $p$ on Alice's input set $X(\omega)$ and
Bob's input set $Y(\omega)$ produces the same Boolean output as the output determined by
$z$ through its transcript leaf: \[ p\bigl(X(\omega),\, Y(\omega)\bigr) =
\mathrm{zOutput}(z). \]
\end{theorem}

\begin{theorem}[Achievable Z-fibers have nonzero volume]
\lean{CommunicationComplexity.Functions.Disjointness.RandomizedLowerBound.volume_zFiber_ne_zero}
\label{CommunicationComplexity.Functions.Disjointness.RandomizedLowerBound.volume_zFiber_ne_zero}
\leanok
\uses{CommunicationComplexity.Functions.Disjointness.RandomizedLowerBound.HardSample, CommunicationComplexity.Functions.Disjointness.RandomizedLowerBound.ProtocolType, CommunicationComplexity.Functions.Disjointness.RandomizedLowerBound.ZType, CommunicationComplexity.Functions.Disjointness.RandomizedLowerBound.ZType.achievable, CommunicationComplexity.Functions.Disjointness.RandomizedLowerBound.ZType.raw, CommunicationComplexity.Functions.Disjointness.RandomizedLowerBound.zFiber, CommunicationComplexity.Functions.Disjointness.RandomizedLowerBound.zVariable, CommunicationComplexity.uniformOn_univ_measureReal_eq_card_subtype}
\difficulty{4}
Fix a positive integer $n$ and a deterministic communication protocol $p$ for
disjointness inputs of length $n$. On the finite sample space of hard-distribution
samples $\omega$, let $Z(\omega) = (M, T, X_{<T}, Y_{>T})$ be the statistic recording
the transcript $M$ produced by running $p$ on the input pair generated by $\omega$, the
special coordinate $T$, Alice's input bits strictly before $T$, and Bob's input bits
strictly after $T$. Then for every achievable value $z$ of $Z$—that is, every $z$
realized by at least one sample—the fiber $Z^{-1}(\{z\})$, consisting of all samples
$\omega$ with $Z(\omega) = z$, has nonzero volume under the ambient measure on the
sample space.
\end{theorem}

\begin{definition}[Inputmeasure]
\lean{CommunicationComplexity.Functions.Disjointness.RandomizedLowerBound.inputMeasure}
\label{CommunicationComplexity.Functions.Disjointness.RandomizedLowerBound.inputMeasure}
\leanok
\uses{CommunicationComplexity.Functions.Disjointness.RandomizedLowerBound.input}
The hard input distribution as a measure on input pairs.
\end{definition}

\begin{definition}[Inputdist]
\lean{CommunicationComplexity.Functions.Disjointness.RandomizedLowerBound.inputDist}
\label{CommunicationComplexity.Functions.Disjointness.RandomizedLowerBound.inputDist}
\leanok
\uses{CommunicationComplexity.FiniteProbabilitySpace, CommunicationComplexity.FiniteProbabilitySpace.ofMeasure, CommunicationComplexity.Functions.Disjointness.RandomizedLowerBound.inputMeasure}
The hard input distribution packaged as a finite probability space for distributional error.
\end{definition}

\begin{definition}[Dualprotocol]
\lean{CommunicationComplexity.Functions.Disjointness.RandomizedLowerBound.dualProtocol}
\label{CommunicationComplexity.Functions.Disjointness.RandomizedLowerBound.dualProtocol}
\leanok
\uses{CommunicationComplexity.Deterministic.Protocol.comap, CommunicationComplexity.Deterministic.Protocol.swap, CommunicationComplexity.Functions.Disjointness.RandomizedLowerBound.ProtocolType, CommunicationComplexity.Functions.Disjointness.RandomizedLowerBound.reverseSet}
The protocol dual swaps Alice/Bob and reverses both coordinate sets.
\end{definition}

\begin{definition}[Dualprotocoltranscriptmap]
\lean{CommunicationComplexity.Functions.Disjointness.RandomizedLowerBound.dualProtocolTranscriptMap}
\label{CommunicationComplexity.Functions.Disjointness.RandomizedLowerBound.dualProtocolTranscriptMap}
\leanok
\uses{CommunicationComplexity.Deterministic.Protocol.swap, CommunicationComplexity.Deterministic.Protocol.transcriptComap, CommunicationComplexity.Deterministic.Protocol.transcriptSwap, CommunicationComplexity.Functions.Disjointness.RandomizedLowerBound.ProtocolType, CommunicationComplexity.Functions.Disjointness.RandomizedLowerBound.TranscriptType, CommunicationComplexity.Functions.Disjointness.RandomizedLowerBound.dualProtocol, CommunicationComplexity.Functions.Disjointness.RandomizedLowerBound.reverseSet}
The transcript-level map induced by protocol duality.
\end{definition}

\begin{theorem}[Injectivity of the dual-protocol transcript map]
\lean{CommunicationComplexity.Functions.Disjointness.RandomizedLowerBound.dualProtocolTranscriptMap_injective}
\label{CommunicationComplexity.Functions.Disjointness.RandomizedLowerBound.dualProtocolTranscriptMap_injective}
\leanok
\uses{CommunicationComplexity.Deterministic.Protocol.swap, CommunicationComplexity.Deterministic.Protocol.transcriptComap_injective, CommunicationComplexity.Deterministic.Protocol.transcriptSwap_injective, CommunicationComplexity.Functions.Disjointness.RandomizedLowerBound.ProtocolType, CommunicationComplexity.Functions.Disjointness.RandomizedLowerBound.dualProtocolTranscriptMap, CommunicationComplexity.Functions.Disjointness.RandomizedLowerBound.reverseSet}
\difficulty{1}
Let $n$ be a positive integer, and let $p$ be a deterministic two-party communication
protocol in which both Alice's and Bob's inputs are subsets of $\{0,1,\dots,n-1\}$ and
the output is a single Boolean. Protocol duality produces from $p$ a protocol over the
same input and output types by exchanging the roles of Alice and Bob and reversing the
coordinate order $i \mapsto n-1-i$ in both players' inputs, and it induces a map on
syntactic transcripts that sends each execution transcript of $p$ to the transcript of
the dual protocol carrying the very same sequence of Boolean messages. This induced
transcript map is injective.
\end{theorem}

\begin{theorem}[Alice's bit away from the special coordinate]
\lean{CommunicationComplexity.Functions.Disjointness.RandomizedLowerBound.xBit_of_ne_special}
\label{CommunicationComplexity.Functions.Disjointness.RandomizedLowerBound.xBit_of_ne_special}
\leanok
\uses{CommunicationComplexity.Functions.Disjointness.RandomizedLowerBound.DisjointCoordinate.xBit, CommunicationComplexity.Functions.Disjointness.RandomizedLowerBound.HardSample, CommunicationComplexity.Functions.Disjointness.RandomizedLowerBound.xBit}
\difficulty{1}
Let $\omega$ be a sample of the hard distribution, with special coordinate $T$,
per-coordinate assignment $\text{other}$, and let $\text{xBit}(\omega, i)$ denote
Alice's bit at coordinate $i$. Then for every coordinate $i \neq T$, Alice's bit at $i$
equals Alice's bit in the disjoint coordinate-pair $\text{other}(i)$ that $\omega$
assigns to $i$.
\end{theorem}

\begin{theorem}[Bob's bit away from the special coordinate]
\lean{CommunicationComplexity.Functions.Disjointness.RandomizedLowerBound.yBit_of_ne_special}
\label{CommunicationComplexity.Functions.Disjointness.RandomizedLowerBound.yBit_of_ne_special}
\leanok
\uses{CommunicationComplexity.Functions.Disjointness.RandomizedLowerBound.DisjointCoordinate.yBit, CommunicationComplexity.Functions.Disjointness.RandomizedLowerBound.HardSample, CommunicationComplexity.Functions.Disjointness.RandomizedLowerBound.yBit}
\difficulty{1}
Fix a positive integer $n$ and a sample $\omega$ of the hard distribution, consisting of
a distinguished coordinate $T \in \{0, 1, \dots, n-1\}$, two bits $x_T$ and $y_T$, and
an assignment $i \mapsto \omega_{\text{other}}(i)$ of a disjoint coordinate-type (one of
*neither*, *left only*, *right only*) to each coordinate. For every coordinate $i \ne
T$, Bob's bit at $i$ under this sample equals the bit $y(\omega_{\text{other}}(i))$ read
off from the coordinate-type assigned to $i$; that is, away from the distinguished
coordinate the distinguished bit $y_T$ plays no role.
\end{theorem}

\begin{theorem}[Disjointness away from the special coordinate]
\lean{CommunicationComplexity.Functions.Disjointness.RandomizedLowerBound.not_xBit_and_yBit_of_ne_special}
\label{CommunicationComplexity.Functions.Disjointness.RandomizedLowerBound.not_xBit_and_yBit_of_ne_special}
\leanok
\uses{CommunicationComplexity.Functions.Disjointness.RandomizedLowerBound.DisjointCoordinate.not_xBit_and_yBit, CommunicationComplexity.Functions.Disjointness.RandomizedLowerBound.HardSample, CommunicationComplexity.Functions.Disjointness.RandomizedLowerBound.xBit, CommunicationComplexity.Functions.Disjointness.RandomizedLowerBound.xBit_of_ne_special, CommunicationComplexity.Functions.Disjointness.RandomizedLowerBound.yBit, CommunicationComplexity.Functions.Disjointness.RandomizedLowerBound.yBit_of_ne_special}
\difficulty{2}
Let $\omega$ be a sample of the hard distribution on $n$ coordinates, with special
coordinate $T$, and let $i$ be any coordinate with $i \neq T$. Then Alice's bit and
Bob's bit at coordinate $i$ are never both equal to $\mathrm{true}$; that is, it is not
the case that both the Alice-bit at $i$ and the Bob-bit at $i$ hold.
\end{theorem}

\begin{definition}[Coordinateofbits]
\lean{CommunicationComplexity.Functions.Disjointness.RandomizedLowerBound.coordinateOfBits}
\label{CommunicationComplexity.Functions.Disjointness.RandomizedLowerBound.coordinateOfBits}
\leanok
The disjoint coordinate with the given bits. On the invalid $(true, true)$ input this chooses an arbitrary disjoint coordinate; the projection lemmas below assume that invalid case away.
\end{definition}

\begin{theorem}[Recovering Alice's bit from a disjoint coordinate]
\lean{CommunicationComplexity.Functions.Disjointness.RandomizedLowerBound.coordinateOfBits_xBit}
\label{CommunicationComplexity.Functions.Disjointness.RandomizedLowerBound.coordinateOfBits_xBit}
\leanok
\uses{CommunicationComplexity.Functions.Disjointness.RandomizedLowerBound.DisjointCoordinate.xBit, CommunicationComplexity.Functions.Disjointness.RandomizedLowerBound.coordinateOfBits}
\difficulty{2}
Let $x, y \in \{\mathrm{true}, \mathrm{false}\}$ be a pair of bits that are not both
$\mathrm{true}$. Then the disjoint coordinate built from $(x, y)$ has Alice's bit equal
to $x$; that is, extracting Alice's bit from this coordinate returns $x$ exactly.
\end{theorem}

\begin{theorem}[Recovering Bob's bit from a disjoint coordinate]
\lean{CommunicationComplexity.Functions.Disjointness.RandomizedLowerBound.coordinateOfBits_yBit}
\label{CommunicationComplexity.Functions.Disjointness.RandomizedLowerBound.coordinateOfBits_yBit}
\leanok
\uses{CommunicationComplexity.Functions.Disjointness.RandomizedLowerBound.DisjointCoordinate.yBit, CommunicationComplexity.Functions.Disjointness.RandomizedLowerBound.coordinateOfBits}
\difficulty{2}
Let $x, y \in \{\mathrm{true}, \mathrm{false}\}$ be Alice's and Bob's bits in a single
coordinate, and suppose they are not both $\mathrm{true}$. Then Bob's bit extracted from
the disjoint coordinate assembled from $x$ and $y$ equals $y$.
\end{theorem}

\begin{theorem}[Recovering a disjointness coordinate from its two bits]
\lean{CommunicationComplexity.Functions.Disjointness.RandomizedLowerBound.coordinateOfBits_xBit_yBit}
\label{CommunicationComplexity.Functions.Disjointness.RandomizedLowerBound.coordinateOfBits_xBit_yBit}
\leanok
\uses{CommunicationComplexity.Functions.Disjointness.RandomizedLowerBound.DisjointCoordinate.xBit, CommunicationComplexity.Functions.Disjointness.RandomizedLowerBound.DisjointCoordinate.yBit, CommunicationComplexity.Functions.Disjointness.RandomizedLowerBound.coordinateOfBits}
\difficulty{2}
Let a *disjointness coordinate* be one of three values, recording, for a single index,
whether that index belongs to neither of two subsets, to the first subset only
(Alice's), or to the second subset only (Bob's). To each coordinate $c$ attach two
Boolean bits: Alice's bit $x(c)$, which is set precisely when $c$ marks the index as
belonging to Alice's subset only, and Bob's bit $y(c)$, which is set precisely when $c$
marks it as belonging to Bob's subset only. Conversely, from a pair of bits $(x, y)$
form a coordinate $\kappa(x, y)$ by taking "Alice only" if $x$ is set, otherwise "Bob
only" if $y$ is set, and otherwise "neither." Then $\kappa\bigl(x(c), y(c)\bigr) = c$
for every coordinate $c$; that is, splitting a coordinate into its two bits and
reassembling recovers the original coordinate.
\end{theorem}

\begin{definition}[Mix]
\lean{CommunicationComplexity.Functions.Disjointness.RandomizedLowerBound.mix}
\label{CommunicationComplexity.Functions.Disjointness.RandomizedLowerBound.mix}
\leanok
\uses{CommunicationComplexity.Functions.Disjointness.RandomizedLowerBound.HardSample, CommunicationComplexity.Functions.Disjointness.RandomizedLowerBound.coordinateOfBits, CommunicationComplexity.Functions.Disjointness.RandomizedLowerBound.xBit, CommunicationComplexity.Functions.Disjointness.RandomizedLowerBound.yBit}
Mix Alice's input from $ωX$ with Bob's input from $ωY$. This is only intended to be used when the two samples agree on $T$, $X_{<T}$, and $Y_{>T}$; under those hypotheses the mixed coordinate requests are always disjoint.
\end{definition}

\begin{theorem}[Disjointness of the hard-sample input sets]
\lean{CommunicationComplexity.Functions.Disjointness.RandomizedLowerBound.disjoint_X_Y_iff}
\label{CommunicationComplexity.Functions.Disjointness.RandomizedLowerBound.disjoint_X_Y_iff}
\leanok
\uses{CommunicationComplexity.Functions.Disjointness.RandomizedLowerBound.HardSample, CommunicationComplexity.Functions.Disjointness.RandomizedLowerBound.X, CommunicationComplexity.Functions.Disjointness.RandomizedLowerBound.Y, CommunicationComplexity.Functions.Disjointness.RandomizedLowerBound.not_xBit_and_yBit_of_ne_special, CommunicationComplexity.Functions.Disjointness.RandomizedLowerBound.xBit, CommunicationComplexity.Functions.Disjointness.RandomizedLowerBound.yBit}
\difficulty{4}
Let $\omega$ be a hard-distribution sample, with distinguished coordinate $T$ carrying
Alice's bit $x_T$ and Bob's bit $y_T$, and let $X(\omega)$ and $Y(\omega)$ be Alice's
and Bob's input sets under $\omega$. Then $X(\omega)$ and $Y(\omega)$ are disjoint if
and only if it is not the case that $x_T$ and $y_T$ are both $\mathrm{true}$.
\end{theorem}

\begin{definition}[Disjointcoordinatevector]
\lean{CommunicationComplexity.Functions.Disjointness.RandomizedLowerBound.disjointCoordinateVector}
\label{CommunicationComplexity.Functions.Disjointness.RandomizedLowerBound.disjointCoordinateVector}
\leanok
\uses{CommunicationComplexity.Functions.Disjointness.RandomizedLowerBound.HardSample, CommunicationComplexity.Functions.Disjointness.RandomizedLowerBound.coordinateOfBits}
The generated disjoint coordinate-pair vector. On samples outside $D$, the special coordinate uses the arbitrary default built into \texttt{coordinateOfBits}.
\end{definition}

\begin{theorem}[Alice's bit is recovered on disjoint samples]
\lean{CommunicationComplexity.Functions.Disjointness.RandomizedLowerBound.disjointCoordinateVector_xBit_of_mem_disjointEvent}
\label{CommunicationComplexity.Functions.Disjointness.RandomizedLowerBound.disjointCoordinateVector_xBit_of_mem_disjointEvent}
\leanok
\uses{CommunicationComplexity.Functions.Disjointness.RandomizedLowerBound.DisjointCoordinate.xBit, CommunicationComplexity.Functions.Disjointness.RandomizedLowerBound.HardSample, CommunicationComplexity.Functions.Disjointness.RandomizedLowerBound.coordinateOfBits_xBit, CommunicationComplexity.Functions.Disjointness.RandomizedLowerBound.disjointCoordinateVector, CommunicationComplexity.Functions.Disjointness.RandomizedLowerBound.disjointEvent, CommunicationComplexity.Functions.Disjointness.RandomizedLowerBound.disjoint_X_Y_iff, CommunicationComplexity.Functions.Disjointness.RandomizedLowerBound.xBit}
\difficulty{4}
Let $\omega$ be a hard-distribution sample whose induced input pair, Alice's set
$X(\omega)$ and Bob's set $Y(\omega)$, are disjoint. Then for every coordinate $i$,
Alice's bit read off from the disjoint coordinate-pair generated at $i$ agrees with
Alice's bit at $i$ under $\omega$.
\end{theorem}

\begin{theorem}[Bob's bit recovered from the disjoint coordinate vector]
\lean{CommunicationComplexity.Functions.Disjointness.RandomizedLowerBound.disjointCoordinateVector_yBit_of_mem_disjointEvent}
\label{CommunicationComplexity.Functions.Disjointness.RandomizedLowerBound.disjointCoordinateVector_yBit_of_mem_disjointEvent}
\leanok
\uses{CommunicationComplexity.Functions.Disjointness.RandomizedLowerBound.DisjointCoordinate.yBit, CommunicationComplexity.Functions.Disjointness.RandomizedLowerBound.HardSample, CommunicationComplexity.Functions.Disjointness.RandomizedLowerBound.coordinateOfBits_yBit, CommunicationComplexity.Functions.Disjointness.RandomizedLowerBound.disjointCoordinateVector, CommunicationComplexity.Functions.Disjointness.RandomizedLowerBound.disjointEvent, CommunicationComplexity.Functions.Disjointness.RandomizedLowerBound.disjoint_X_Y_iff, CommunicationComplexity.Functions.Disjointness.RandomizedLowerBound.yBit}
\difficulty{4}
Let $\omega$ be a sample of the hard distribution lying in the disjoint event — that is,
one whose induced input sets $X(\omega)$ and $Y(\omega)$ are disjoint. Then for every
coordinate $i$, the Bob-bit of the disjoint coordinate assigned to $i$ by the generated
coordinate-pair vector agrees with Bob's bit $\mathrm{yBit}(\omega, i)$ at that
coordinate.
\end{theorem}

\begin{definition}[Ignoredcoordinate]
\lean{CommunicationComplexity.Functions.Disjointness.RandomizedLowerBound.ignoredCoordinate}
\label{CommunicationComplexity.Functions.Disjointness.RandomizedLowerBound.ignoredCoordinate}
\leanok
\uses{CommunicationComplexity.Functions.Disjointness.RandomizedLowerBound.HardSample}
The \texttt{other} value at the special coordinate is ignored by the generated input.
\end{definition}

\begin{definition}[Disjointmodel]
\lean{CommunicationComplexity.Functions.Disjointness.RandomizedLowerBound.disjointModel}
\label{CommunicationComplexity.Functions.Disjointness.RandomizedLowerBound.disjointModel}
\leanok
\uses{CommunicationComplexity.Functions.Disjointness.RandomizedLowerBound.HardSample, CommunicationComplexity.Functions.Disjointness.RandomizedLowerBound.disjointCoordinateVector, CommunicationComplexity.Functions.Disjointness.RandomizedLowerBound.ignoredCoordinate}
Coordinates for the conditioned-on-disjointness sample space: the special coordinate, the generated disjoint coordinate-pair vector, and the ignored \texttt{other} value at the special coordinate.
\end{definition}

\begin{definition}[Disjointeventequiv]
\lean{CommunicationComplexity.Functions.Disjointness.RandomizedLowerBound.disjointEventEquiv}
\label{CommunicationComplexity.Functions.Disjointness.RandomizedLowerBound.disjointEventEquiv}
\leanok
\uses{CommunicationComplexity.Functions.Disjointness.RandomizedLowerBound.DisjointCoordinate.not_xBit_and_yBit, CommunicationComplexity.Functions.Disjointness.RandomizedLowerBound.DisjointCoordinate.xBit, CommunicationComplexity.Functions.Disjointness.RandomizedLowerBound.DisjointCoordinate.yBit, CommunicationComplexity.Functions.Disjointness.RandomizedLowerBound.HardSample, CommunicationComplexity.Functions.Disjointness.RandomizedLowerBound.X, CommunicationComplexity.Functions.Disjointness.RandomizedLowerBound.Y, CommunicationComplexity.Functions.Disjointness.RandomizedLowerBound.coordinateOfBits_xBit, CommunicationComplexity.Functions.Disjointness.RandomizedLowerBound.coordinateOfBits_xBit_yBit, CommunicationComplexity.Functions.Disjointness.RandomizedLowerBound.coordinateOfBits_yBit, CommunicationComplexity.Functions.Disjointness.RandomizedLowerBound.disjointCoordinateVector, CommunicationComplexity.Functions.Disjointness.RandomizedLowerBound.disjointEvent, CommunicationComplexity.Functions.Disjointness.RandomizedLowerBound.disjointModel, CommunicationComplexity.Functions.Disjointness.RandomizedLowerBound.disjoint_X_Y_iff, CommunicationComplexity.Functions.Disjointness.RandomizedLowerBound.ignoredCoordinate}
Technical auxiliary result for \texttt{disjointEventEquiv}.
\end{definition}

\begin{definition}[Disjointeventinterdisjointmodelequiv]
\lean{CommunicationComplexity.Functions.Disjointness.RandomizedLowerBound.disjointEventInterDisjointModelEquiv}
\label{CommunicationComplexity.Functions.Disjointness.RandomizedLowerBound.disjointEventInterDisjointModelEquiv}
\leanok
\uses{CommunicationComplexity.Functions.Disjointness.RandomizedLowerBound.HardSample, CommunicationComplexity.Functions.Disjointness.RandomizedLowerBound.disjointEvent, CommunicationComplexity.Functions.Disjointness.RandomizedLowerBound.disjointEventEquiv, CommunicationComplexity.Functions.Disjointness.RandomizedLowerBound.disjointModel}
Technical auxiliary result for \texttt{disjointEventInterDisjointModelEquiv}.
\end{definition}

\begin{theorem}[Disjointness event fibers of the model map are singletons]
\lean{CommunicationComplexity.Functions.Disjointness.RandomizedLowerBound.card_disjointEvent_inter_disjointModel_fiber}
\label{CommunicationComplexity.Functions.Disjointness.RandomizedLowerBound.card_disjointEvent_inter_disjointModel_fiber}
\leanok
\uses{CommunicationComplexity.Functions.Disjointness.RandomizedLowerBound.HardSample, CommunicationComplexity.Functions.Disjointness.RandomizedLowerBound.disjointEvent, CommunicationComplexity.Functions.Disjointness.RandomizedLowerBound.disjointEventInterDisjointModelEquiv, CommunicationComplexity.Functions.Disjointness.RandomizedLowerBound.disjointModel}
\difficulty{2}
Fix $n \ge 1$, and consider the finite sample space whose points $\omega$ each consist
of a distinguished coordinate $T \in \{1,\dots,n\}$, a pair of bits $x_T, y_T$ at that
coordinate, and a labelling of every coordinate by one of the three disjoint-coordinate
types (neither, left-only, right-only). From such an $\omega$ one reads off the
generated coordinate-pair vector together with the pair of input sets $X(\omega),
Y(\omega)$, and one records the model data $\mathrm{model}(\omega) = (T,\
\text{coordinate-pair vector},\ \text{the ignored label at } T)$. Then for every triple
$z$ in the target of the model map, there is exactly one sample $\omega$ whose generated
inputs are disjoint, that is $X(\omega)$ and $Y(\omega)$ are disjoint, and for which
$\mathrm{model}(\omega) = z$.
\end{theorem}

\begin{theorem}[Uniform mass of each disjoint-model fiber]
\lean{CommunicationComplexity.Functions.Disjointness.RandomizedLowerBound.measureReal_disjointEvent_inter_disjointModel_fiber}
\label{CommunicationComplexity.Functions.Disjointness.RandomizedLowerBound.measureReal_disjointEvent_inter_disjointModel_fiber}
\leanok
\uses{CommunicationComplexity.Functions.Disjointness.RandomizedLowerBound.HardSample, CommunicationComplexity.Functions.Disjointness.RandomizedLowerBound.HardSample.card, CommunicationComplexity.Functions.Disjointness.RandomizedLowerBound.card_disjointEvent_inter_disjointModel_fiber, CommunicationComplexity.Functions.Disjointness.RandomizedLowerBound.disjointEvent, CommunicationComplexity.Functions.Disjointness.RandomizedLowerBound.disjointModel, CommunicationComplexity.uniformOn_univ_measureReal_eq_card_subtype}
\difficulty{3}
Fix an integer $n \ge 1$, and equip the finite space of hard samples with its uniform
probability measure: a hard sample records a special coordinate $T \in \{1,\dots,n\}$,
two bits $x_T, y_T$, and a disjoint-coordinate value at each of the $n$ coordinates, so
there are $n\cdot 4\cdot 3^{\,n}$ of them, each carrying mass $1/(n\cdot 4\cdot
3^{\,n})$. Then for every point $z$ of the space
$\{1,\dots,n\}\times\big(\{1,\dots,n\}\to \text{DisjointCoordinate}\big)\times
\text{DisjointCoordinate}$, the set of hard samples that lie in the disjointness event
and whose disjoint model equals $z$ has measure exactly
\[
\frac{1}{n\cdot 4\cdot 3^{\,n}},
\]
the mass of a single sample.
\end{theorem}

\begin{definition}[Dualhardsample]
\lean{CommunicationComplexity.Functions.Disjointness.RandomizedLowerBound.dualHardSample}
\label{CommunicationComplexity.Functions.Disjointness.RandomizedLowerBound.dualHardSample}
\leanok
\uses{CommunicationComplexity.Functions.Disjointness.RandomizedLowerBound.DisjointCoordinate.swap, CommunicationComplexity.Functions.Disjointness.RandomizedLowerBound.HardSample}
Dualize the hard sample space by swapping Alice/Bob and reversing coordinate order.
\end{definition}

\begin{theorem}[Dualization of hard samples is an involution]
\lean{CommunicationComplexity.Functions.Disjointness.RandomizedLowerBound.dualHardSample_dualHardSample}
\label{CommunicationComplexity.Functions.Disjointness.RandomizedLowerBound.dualHardSample_dualHardSample}
\leanok
\uses{CommunicationComplexity.Functions.Disjointness.RandomizedLowerBound.DisjointCoordinate.swap_swap, CommunicationComplexity.Functions.Disjointness.RandomizedLowerBound.HardSample, CommunicationComplexity.Functions.Disjointness.RandomizedLowerBound.dualHardSample}
\difficulty{2}
Fix $n \ge 1$, and let $\omega$ be any hard sample: a distinguished coordinate $T \in
\{0,\dots,n-1\}$, two bits recording Alice's and Bob's values there, and an assignment
of a disjoint coordinate (neither, left-only, or right-only) to every coordinate.
Applying the dualization map to $\omega$ twice returns $\omega$ itself; that is, the
dualization — which reverses the coordinate order, interchanges Alice's and Bob's bits
at the distinguished coordinate, and swaps the two sides of every coordinate assignment
— is an involution.
\end{theorem}

\begin{theorem}[Alice's dual bit equals Bob's reversed bit]
\lean{CommunicationComplexity.Functions.Disjointness.RandomizedLowerBound.xBit_dualHardSample}
\label{CommunicationComplexity.Functions.Disjointness.RandomizedLowerBound.xBit_dualHardSample}
\leanok
\uses{CommunicationComplexity.Functions.Disjointness.RandomizedLowerBound.DisjointCoordinate.xBit_swap, CommunicationComplexity.Functions.Disjointness.RandomizedLowerBound.HardSample, CommunicationComplexity.Functions.Disjointness.RandomizedLowerBound.dualHardSample, CommunicationComplexity.Functions.Disjointness.RandomizedLowerBound.xBit, CommunicationComplexity.Functions.Disjointness.RandomizedLowerBound.yBit}
\difficulty{4}
Fix a positive integer $n$, and let $\omega$ be a sample of the hard distribution on $n$
coordinates. For every coordinate $i \in \{0, 1, \dots, n-1\}$, Alice's bit at
coordinate $i$ in the dualized sample $\omega^{\ast}$ coincides with Bob's bit at the
reversed coordinate $\mathrm{rev}(i) = n-1-i$ in the original sample $\omega$; that is,
\[
\mathrm{xBit}(\omega^{\ast})(i) \;=\; \mathrm{yBit}(\omega)\bigl(\mathrm{rev}(i)\bigr).
\]
\end{theorem}

\begin{theorem}[Bob's bit of the dual sample equals Alice's reversed bit]
\lean{CommunicationComplexity.Functions.Disjointness.RandomizedLowerBound.yBit_dualHardSample}
\label{CommunicationComplexity.Functions.Disjointness.RandomizedLowerBound.yBit_dualHardSample}
\leanok
\uses{CommunicationComplexity.Functions.Disjointness.RandomizedLowerBound.HardSample, CommunicationComplexity.Functions.Disjointness.RandomizedLowerBound.dualHardSample, CommunicationComplexity.Functions.Disjointness.RandomizedLowerBound.dualHardSample_dualHardSample, CommunicationComplexity.Functions.Disjointness.RandomizedLowerBound.xBit, CommunicationComplexity.Functions.Disjointness.RandomizedLowerBound.xBit_dualHardSample, CommunicationComplexity.Functions.Disjointness.RandomizedLowerBound.yBit}
\difficulty{2}
Let $\omega$ be a sample of the hard distribution over $n$ coordinates, and write
$\mathrm{rev}(i) = n-1-i$ for the reversal of a coordinate $i \in \{0,1,\dots,n-1\}$.
Form the dual sample by swapping Alice's and Bob's roles and reversing the coordinate
order. Then Bob's bit at coordinate $i$ of the dual sample equals Alice's bit at the
reversed coordinate $\mathrm{rev}(i)$ of the original sample.
\end{theorem}

\begin{theorem}[Alice's dual set as the reversal of Bob's set]
\lean{CommunicationComplexity.Functions.Disjointness.RandomizedLowerBound.X_dualHardSample}
\label{CommunicationComplexity.Functions.Disjointness.RandomizedLowerBound.X_dualHardSample}
\leanok
\uses{CommunicationComplexity.Functions.Disjointness.RandomizedLowerBound.HardSample, CommunicationComplexity.Functions.Disjointness.RandomizedLowerBound.X, CommunicationComplexity.Functions.Disjointness.RandomizedLowerBound.Y, CommunicationComplexity.Functions.Disjointness.RandomizedLowerBound.dualHardSample, CommunicationComplexity.Functions.Disjointness.RandomizedLowerBound.reverseSet, CommunicationComplexity.Functions.Disjointness.RandomizedLowerBound.xBit_dualHardSample}
\difficulty{2}
Fix a positive integer $n$ of coordinates, and let $\omega$ be any sample of the hard
distribution. Then Alice's input set of the dualized sample coincides with the
coordinate-reversal of Bob's input set of the original sample; that is,
$X(\mathrm{dual}(\omega)) = \mathrm{reverse}(Y(\omega))$, where reversal sends each
coordinate $i$ to $n-1-i$.
\end{theorem}

\begin{theorem}[Bob's dual input set is the reversal of Alice's set]
\lean{CommunicationComplexity.Functions.Disjointness.RandomizedLowerBound.Y_dualHardSample}
\label{CommunicationComplexity.Functions.Disjointness.RandomizedLowerBound.Y_dualHardSample}
\leanok
\uses{CommunicationComplexity.Functions.Disjointness.RandomizedLowerBound.HardSample, CommunicationComplexity.Functions.Disjointness.RandomizedLowerBound.X, CommunicationComplexity.Functions.Disjointness.RandomizedLowerBound.X_dualHardSample, CommunicationComplexity.Functions.Disjointness.RandomizedLowerBound.Y, CommunicationComplexity.Functions.Disjointness.RandomizedLowerBound.dualHardSample, CommunicationComplexity.Functions.Disjointness.RandomizedLowerBound.dualHardSample_dualHardSample, CommunicationComplexity.Functions.Disjointness.RandomizedLowerBound.reverseSet, CommunicationComplexity.Functions.Disjointness.RandomizedLowerBound.reverseSet_reverseSet}
\difficulty{3}
Fix $n \ge 1$, and let $\omega$ be a sample of the hard distribution on $n$ coordinates,
with Alice's input set $X(\omega) \subseteq \{0,\dots,n-1\}$ and Bob's input set
$Y(\omega)$. Passing to the dual sample $\omega'$ — obtained by swapping Alice's and
Bob's roles and reversing the coordinate order — Bob's input set of $\omega'$ is exactly
the coordinate reversal of Alice's input set of $\omega$: \[ Y(\omega') = \{\, i :
\operatorname{rev}(i) \in X(\omega) \,\}, \] where $\operatorname{rev}(i) = n-1-i$.
Equivalently, a coordinate $i$ belongs to Bob's set for the dualized sample precisely
when the reversed coordinate $\operatorname{rev}(i)$ belongs to Alice's set for the
original sample.
\end{theorem}

\begin{theorem}[Dual sample reverses and swaps the input pair]
\lean{CommunicationComplexity.Functions.Disjointness.RandomizedLowerBound.input_dualHardSample}
\label{CommunicationComplexity.Functions.Disjointness.RandomizedLowerBound.input_dualHardSample}
\leanok
\uses{CommunicationComplexity.Functions.Disjointness.RandomizedLowerBound.HardSample, CommunicationComplexity.Functions.Disjointness.RandomizedLowerBound.X, CommunicationComplexity.Functions.Disjointness.RandomizedLowerBound.X_dualHardSample, CommunicationComplexity.Functions.Disjointness.RandomizedLowerBound.Y, CommunicationComplexity.Functions.Disjointness.RandomizedLowerBound.Y_dualHardSample, CommunicationComplexity.Functions.Disjointness.RandomizedLowerBound.dualHardSample, CommunicationComplexity.Functions.Disjointness.RandomizedLowerBound.input, CommunicationComplexity.Functions.Disjointness.RandomizedLowerBound.reverseSet}
\difficulty{1}
Fix $n \ge 1$, and let $\omega$ be a sample of the hard distribution, from which one
reads off Alice's input set $X(\omega)$ and Bob's input set $Y(\omega)$, packaged as the
input pair $(X(\omega), Y(\omega))$. Form the dual sample by swapping the roles of Alice
and Bob and reversing the coordinate order $i \mapsto \bar\imath$. Then the input pair
generated by the dual sample is obtained from $(X(\omega), Y(\omega))$ by exchanging the
two sets and reversing the coordinate order of each: Alice receives the
coordinate-reversal of $Y(\omega)$ and Bob receives the coordinate-reversal of
$X(\omega)$, i.e.
\[
\bigl(X, Y\bigr)\text{ of the dual sample} \;=\; \bigl(\{\, i : \bar\imath \in
Y(\omega)\,\},\; \{\, i : \bar\imath \in X(\omega)\,\}\bigr).
\]
\end{theorem}

\begin{theorem}[Duality preserves the disjointness event]
\lean{CommunicationComplexity.Functions.Disjointness.RandomizedLowerBound.disjointEvent_dualHardSample}
\label{CommunicationComplexity.Functions.Disjointness.RandomizedLowerBound.disjointEvent_dualHardSample}
\leanok
\uses{CommunicationComplexity.Functions.Disjointness.RandomizedLowerBound.HardSample, CommunicationComplexity.Functions.Disjointness.RandomizedLowerBound.X, CommunicationComplexity.Functions.Disjointness.RandomizedLowerBound.Y, CommunicationComplexity.Functions.Disjointness.RandomizedLowerBound.disjointEvent, CommunicationComplexity.Functions.Disjointness.RandomizedLowerBound.disjoint_X_Y_iff, CommunicationComplexity.Functions.Disjointness.RandomizedLowerBound.dualHardSample}
\difficulty{3}
Let $\omega$ be a sample of the hard distribution, and let $\omega^\ast$ be its dual,
obtained by exchanging Alice's and Bob's roles while reversing the order of the
coordinates. Then Alice's and Bob's input sets $X(\omega^\ast)$ and $Y(\omega^\ast)$ are
disjoint if and only if $X(\omega)$ and $Y(\omega)$ are disjoint; that is, $\omega^\ast$
lies in the disjointness event exactly when $\omega$ does.
\end{theorem}

\begin{theorem}[Alice's dual special bit equals Bob's original]
\lean{CommunicationComplexity.Functions.Disjointness.RandomizedLowerBound.specialX_dualHardSample}
\label{CommunicationComplexity.Functions.Disjointness.RandomizedLowerBound.specialX_dualHardSample}
\leanok
\uses{CommunicationComplexity.Functions.Disjointness.RandomizedLowerBound.HardSample, CommunicationComplexity.Functions.Disjointness.RandomizedLowerBound.dualHardSample, CommunicationComplexity.Functions.Disjointness.RandomizedLowerBound.specialX, CommunicationComplexity.Functions.Disjointness.RandomizedLowerBound.specialY}
\difficulty{1}
Fix a positive integer $n$, and let $\omega$ be a point of the hard sample space on $n$
coordinates. Passing to the dualized sample — which swaps the roles of Alice and Bob and
reverses the coordinate order — turns Alice's bit at the special coordinate into Bob's
bit at the special coordinate of $\omega$: Alice's special bit of the dual of $\omega$
equals Bob's special bit of $\omega$.
\end{theorem}

\begin{theorem}[Bob's dual special bit equals Alice's original]
\lean{CommunicationComplexity.Functions.Disjointness.RandomizedLowerBound.specialY_dualHardSample}
\label{CommunicationComplexity.Functions.Disjointness.RandomizedLowerBound.specialY_dualHardSample}
\leanok
\uses{CommunicationComplexity.Functions.Disjointness.RandomizedLowerBound.HardSample, CommunicationComplexity.Functions.Disjointness.RandomizedLowerBound.dualHardSample, CommunicationComplexity.Functions.Disjointness.RandomizedLowerBound.specialX, CommunicationComplexity.Functions.Disjointness.RandomizedLowerBound.specialY}
\difficulty{1}
Let $\omega$ be a sample of the hard distribution, and let $\omega'$ be its dual,
obtained by exchanging the roles of Alice and Bob and reversing the order of the
coordinates. Then Bob's bit at the special coordinate of $\omega'$ equals Alice's bit at
the special coordinate of $\omega$.
\end{theorem}

\begin{theorem}[Duality invariance of the all-zero special-bit event]
\lean{CommunicationComplexity.Functions.Disjointness.RandomizedLowerBound.specialZeroZero_dualHardSample}
\label{CommunicationComplexity.Functions.Disjointness.RandomizedLowerBound.specialZeroZero_dualHardSample}
\leanok
\uses{CommunicationComplexity.Functions.Disjointness.RandomizedLowerBound.HardSample, CommunicationComplexity.Functions.Disjointness.RandomizedLowerBound.dualHardSample, CommunicationComplexity.Functions.Disjointness.RandomizedLowerBound.specialBitsEvent, CommunicationComplexity.Functions.Disjointness.RandomizedLowerBound.specialZeroZero}
\difficulty{1}
Fix $n \ge 1$, and let a hard sample $\omega$ record a special coordinate $T \in
\{0,\dots,n-1\}$, Alice's and Bob's bits $x_T, y_T \in \{\mathrm{false},\mathrm{true}\}$
at that coordinate, and a labelling of every coordinate by one of the three disjointness
states. Let its dual be the hard sample obtained by reversing the coordinate order
(sending each index $i$ to $n-1-i$), swapping the sides of every coordinate label, and
interchanging Alice's and Bob's special bits, so that the dual's Alice-bit is $y_T$ and
its Bob-bit is $x_T$. Then the dual of $\omega$ has both of its special bits equal to
$\mathrm{false}$ if and only if $\omega$ itself has both special bits equal to
$\mathrm{false}$.
\end{theorem}

\begin{theorem}[Alice's dual input as Bob's reversed input]
\lean{CommunicationComplexity.Functions.Disjointness.RandomizedLowerBound.xVector_dualHardSample}
\label{CommunicationComplexity.Functions.Disjointness.RandomizedLowerBound.xVector_dualHardSample}
\leanok
\uses{CommunicationComplexity.Functions.Disjointness.RandomizedLowerBound.HardSample, CommunicationComplexity.Functions.Disjointness.RandomizedLowerBound.dualHardSample, CommunicationComplexity.Functions.Disjointness.RandomizedLowerBound.reverseBoolVector, CommunicationComplexity.Functions.Disjointness.RandomizedLowerBound.xBit_dualHardSample, CommunicationComplexity.Functions.Disjointness.RandomizedLowerBound.xVector, CommunicationComplexity.Functions.Disjointness.RandomizedLowerBound.yVector}
\difficulty{2}
Fix a positive dimension $n$ and let $\omega$ be a hard-distribution sample. Then
Alice's full input vector of the dual sample $\omega^{\vee}$ (obtained by swapping the
roles of Alice and Bob and reversing the coordinate order) coincides, coordinate by
coordinate, with the reversal of Bob's full input vector of $\omega$; that is, for every
coordinate the value Alice receives in the dual is the value Bob receives in the
original at the mirror-image coordinate $n-1-i$.
\end{theorem}

\begin{theorem}[Duality exchanges Bob's vector with Alice's reversed vector]
\lean{CommunicationComplexity.Functions.Disjointness.RandomizedLowerBound.yVector_dualHardSample}
\label{CommunicationComplexity.Functions.Disjointness.RandomizedLowerBound.yVector_dualHardSample}
\leanok
\uses{CommunicationComplexity.Functions.Disjointness.RandomizedLowerBound.HardSample, CommunicationComplexity.Functions.Disjointness.RandomizedLowerBound.dualHardSample, CommunicationComplexity.Functions.Disjointness.RandomizedLowerBound.reverseBoolVector, CommunicationComplexity.Functions.Disjointness.RandomizedLowerBound.xVector, CommunicationComplexity.Functions.Disjointness.RandomizedLowerBound.yBit_dualHardSample, CommunicationComplexity.Functions.Disjointness.RandomizedLowerBound.yVector}
\difficulty{2}
Fix a positive integer $n$ and let $\omega$ be a hard sample on the $n$ coordinates.
Then Bob's full input vector of the dualized sample equals the coordinate reversal of
Alice's full input vector of $\omega$: for every coordinate the value Bob receives after
dualizing agrees with the value Alice receives at the mirror-image coordinate of the
original sample.
\end{theorem}

\begin{theorem}[Duality of the before-special and after-special vectors]
\lean{CommunicationComplexity.Functions.Disjointness.RandomizedLowerBound.xBeforeSpecial_dualHardSample}
\label{CommunicationComplexity.Functions.Disjointness.RandomizedLowerBound.xBeforeSpecial_dualHardSample}
\leanok
\uses{CommunicationComplexity.Functions.Disjointness.RandomizedLowerBound.HardSample, CommunicationComplexity.Functions.Disjointness.RandomizedLowerBound.dualHardSample, CommunicationComplexity.Functions.Disjointness.RandomizedLowerBound.reverseBoolVector, CommunicationComplexity.Functions.Disjointness.RandomizedLowerBound.xBeforeSpecial, CommunicationComplexity.Functions.Disjointness.RandomizedLowerBound.xBit, CommunicationComplexity.Functions.Disjointness.RandomizedLowerBound.xBit_dualHardSample, CommunicationComplexity.Functions.Disjointness.RandomizedLowerBound.yAfterSpecial, CommunicationComplexity.Functions.Disjointness.RandomizedLowerBound.yBit}
\difficulty{4}
Fix a positive integer $n$, and let $\omega$ be a sample of the hard distribution on $n$
coordinates — a choice of special coordinate $T$, Alice's and Bob's bits there, and a
disjointness type at every remaining coordinate. Write $\omega^\ast$ for the dual sample
obtained by exchanging Alice's and Bob's roles and reversing the coordinate order. Then
Alice's before-special vector of $\omega^\ast$ equals the coordinate-reversal of Bob's
after-special vector of $\omega$; concretely, at each coordinate $i$ Alice's
before-special bit of $\omega^\ast$ at $i$ coincides with Bob's after-special bit of
$\omega$ at the mirror-image coordinate $n-1-i$.
\end{theorem}

\begin{theorem}[Duality exchanges Alice's prefix vector and Bob's suffix vector]
\lean{CommunicationComplexity.Functions.Disjointness.RandomizedLowerBound.yAfterSpecial_dualHardSample}
\label{CommunicationComplexity.Functions.Disjointness.RandomizedLowerBound.yAfterSpecial_dualHardSample}
\leanok
\uses{CommunicationComplexity.Functions.Disjointness.RandomizedLowerBound.HardSample, CommunicationComplexity.Functions.Disjointness.RandomizedLowerBound.dualHardSample, CommunicationComplexity.Functions.Disjointness.RandomizedLowerBound.dualHardSample_dualHardSample, CommunicationComplexity.Functions.Disjointness.RandomizedLowerBound.reverseBoolVector, CommunicationComplexity.Functions.Disjointness.RandomizedLowerBound.reverseBoolVector_reverseBoolVector, CommunicationComplexity.Functions.Disjointness.RandomizedLowerBound.xBeforeSpecial, CommunicationComplexity.Functions.Disjointness.RandomizedLowerBound.xBeforeSpecial_dualHardSample, CommunicationComplexity.Functions.Disjointness.RandomizedLowerBound.yAfterSpecial}
\difficulty{3}
Fix a coordinate count $n \ge 1$ and let $\omega$ be a sample of the hard distribution,
with special coordinate $T$. Form the dual sample $\omega^{*}$ by swapping the roles of
Alice and Bob and reversing the coordinate order $i \mapsto \bar\imath$. Then Bob's
after-special vector of $\omega^{*}$ — the vector whose entry at each coordinate $i$ is
Bob's bit of $\omega^{*}$ at $i$ when $i$ lies strictly after the special coordinate of
$\omega^{*}$, and is $\mathrm{false}$ otherwise — equals the coordinate reversal of
Alice's before-special vector of $\omega$, whose entry at each coordinate $i$ is Alice's
bit of $\omega$ at $i$ when $i$ lies strictly before $T$, and is $\mathrm{false}$
otherwise. Equivalently, for every coordinate $i$ the $i$-th entry of Bob's
after-special vector of $\omega^{*}$ agrees with the $\bar\imath$-th entry of Alice's
before-special vector of $\omega$.
\end{theorem}

\begin{theorem}[Coordinate-reversal duality of the truncated inputs]
\lean{CommunicationComplexity.Functions.Disjointness.RandomizedLowerBound.xLeSpecial_dualHardSample}
\label{CommunicationComplexity.Functions.Disjointness.RandomizedLowerBound.xLeSpecial_dualHardSample}
\leanok
\uses{CommunicationComplexity.Functions.Disjointness.RandomizedLowerBound.HardSample, CommunicationComplexity.Functions.Disjointness.RandomizedLowerBound.dualHardSample, CommunicationComplexity.Functions.Disjointness.RandomizedLowerBound.reverseBoolVector, CommunicationComplexity.Functions.Disjointness.RandomizedLowerBound.xBit, CommunicationComplexity.Functions.Disjointness.RandomizedLowerBound.xBit_dualHardSample, CommunicationComplexity.Functions.Disjointness.RandomizedLowerBound.xLeSpecial, CommunicationComplexity.Functions.Disjointness.RandomizedLowerBound.yBit, CommunicationComplexity.Functions.Disjointness.RandomizedLowerBound.yGeSpecial}
\difficulty{4}
Fix a positive integer $n$, and let $\omega$ be a sample of the hard distribution on $n$
coordinates. Form the dual sample by swapping Alice's and Bob's roles and reversing the
coordinate order. Then Alice's truncated input vector in the dual sample — her bits up
to and including the special coordinate, padded by $\mathrm{false}$ elsewhere — agrees
coordinate-for-coordinate with the coordinate-reversal of Bob's truncated input vector
in $\omega$, namely his bits from the special coordinate onward, padded by
$\mathrm{false}$ elsewhere.
\end{theorem}

\begin{theorem}[Dual suffix vector as the reversed prefix vector]
\lean{CommunicationComplexity.Functions.Disjointness.RandomizedLowerBound.yGeSpecial_dualHardSample}
\label{CommunicationComplexity.Functions.Disjointness.RandomizedLowerBound.yGeSpecial_dualHardSample}
\leanok
\uses{CommunicationComplexity.Functions.Disjointness.RandomizedLowerBound.HardSample, CommunicationComplexity.Functions.Disjointness.RandomizedLowerBound.dualHardSample, CommunicationComplexity.Functions.Disjointness.RandomizedLowerBound.dualHardSample_dualHardSample, CommunicationComplexity.Functions.Disjointness.RandomizedLowerBound.reverseBoolVector, CommunicationComplexity.Functions.Disjointness.RandomizedLowerBound.reverseBoolVector_reverseBoolVector, CommunicationComplexity.Functions.Disjointness.RandomizedLowerBound.xLeSpecial, CommunicationComplexity.Functions.Disjointness.RandomizedLowerBound.xLeSpecial_dualHardSample, CommunicationComplexity.Functions.Disjointness.RandomizedLowerBound.yGeSpecial}
\difficulty{3}
Let $n$ be a positive integer and let $\omega$ be a sample of the hard distribution,
with special coordinate $T$, and write $\mathrm{dual}(\omega)$ for its dual sample,
obtained by exchanging Alice's and Bob's bits and reversing the coordinate order. Then
Bob's suffix vector $y_{\ge T}$ of $\mathrm{dual}(\omega)$ equals the
coordinate-reversal of Alice's prefix vector $x_{\le T}$ of $\omega$: for every
coordinate $i$ with $0 \le i \le n-1$,
\[
    y_{\ge T}(\mathrm{dual}(\omega))_i \;=\; x_{\le T}(\omega)_{\,n-1-i}.
\]
\end{theorem}

\begin{theorem}[Dual of Bob's conditioning equals Alice's conditioning on the dual sample]
\lean{CommunicationComplexity.Functions.Disjointness.RandomizedLowerBound.aliceClaimConditioning_dualHardSample}
\label{CommunicationComplexity.Functions.Disjointness.RandomizedLowerBound.aliceClaimConditioning_dualHardSample}
\leanok
\uses{CommunicationComplexity.Functions.Disjointness.RandomizedLowerBound.HardSample, CommunicationComplexity.Functions.Disjointness.RandomizedLowerBound.aliceClaimConditioning, CommunicationComplexity.Functions.Disjointness.RandomizedLowerBound.bobClaimConditioning, CommunicationComplexity.Functions.Disjointness.RandomizedLowerBound.dualConditioningValue, CommunicationComplexity.Functions.Disjointness.RandomizedLowerBound.dualHardSample, CommunicationComplexity.Functions.Disjointness.RandomizedLowerBound.specialCoordinate, CommunicationComplexity.Functions.Disjointness.RandomizedLowerBound.xBeforeSpecial_dualHardSample, CommunicationComplexity.Functions.Disjointness.RandomizedLowerBound.yGeSpecial_dualHardSample}
\difficulty{2}
Let $\omega$ be a sample of the hard distribution over $n \ge 1$ coordinates, with
special coordinate $T$ and associated Alice/Bob bit vectors, and form its dual sample by
reversing the coordinate order and interchanging the roles of Alice and Bob. Then
Alice's corrected conditioning variable of the dual sample, namely the triple $(T,
X_{<T}, Y_{\ge T})$, equals the recoding of Bob's corrected conditioning variable $(T,
X_{\le T}, Y_{>T})$ of $\omega$ that reverses the special coordinate and interchanges
the two boolean vectors while reversing the coordinate order of each.
\end{theorem}

\begin{theorem}[Duality exchanges Alice's and Bob's conditioning variables]
\lean{CommunicationComplexity.Functions.Disjointness.RandomizedLowerBound.bobClaimConditioning_dualHardSample}
\label{CommunicationComplexity.Functions.Disjointness.RandomizedLowerBound.bobClaimConditioning_dualHardSample}
\leanok
\uses{CommunicationComplexity.Functions.Disjointness.RandomizedLowerBound.HardSample, CommunicationComplexity.Functions.Disjointness.RandomizedLowerBound.aliceClaimConditioning, CommunicationComplexity.Functions.Disjointness.RandomizedLowerBound.bobClaimConditioning, CommunicationComplexity.Functions.Disjointness.RandomizedLowerBound.dualConditioningValue, CommunicationComplexity.Functions.Disjointness.RandomizedLowerBound.dualHardSample, CommunicationComplexity.Functions.Disjointness.RandomizedLowerBound.specialCoordinate, CommunicationComplexity.Functions.Disjointness.RandomizedLowerBound.xLeSpecial_dualHardSample, CommunicationComplexity.Functions.Disjointness.RandomizedLowerBound.yAfterSpecial_dualHardSample}
\difficulty{2}
Fix a positive integer $n$ and a hard-distribution sample $\omega$ with special
coordinate $T \in \{0,\dots,n-1\}$. Write $\mathrm{rev}$ for the order-reversing map $i
\mapsto n-1-i$ on coordinates, and for a Boolean vector $u$ write $\bar u = u \circ
\mathrm{rev}$ for its coordinate reversal. Recall that Alice's conditioning variable of
$\omega$ is the triple $\bigl(T,\ X_{<T},\ Y_{\ge T}\bigr)$ — the special coordinate
together with Alice's bits strictly before $T$ (padded with $\mathrm{false}$ elsewhere)
and Bob's bits from $T$ onward — while Bob's conditioning variable is $\bigl(T,\ X_{\le
T},\ Y_{>T}\bigr)$. Let $\omega^\ast$ denote the dual sample, obtained from $\omega$ by
reversing the coordinate order and interchanging Alice's and Bob's roles. Then Bob's
conditioning variable of $\omega^\ast$ equals the dual recoding of Alice's conditioning
variable of $\omega$; that is, if $\bigl(T,\ X_{<T},\ Y_{\ge T}\bigr)$ is Alice's
conditioning variable of $\omega$, then Bob's conditioning variable of $\omega^\ast$ is
exactly $\bigl(\mathrm{rev}(T),\ \overline{Y_{\ge T}},\ \overline{X_{<T}}\bigr)$.
\end{theorem}

\begin{theorem}[Cross-sample disjointness away from the special coordinate]
\lean{CommunicationComplexity.Functions.Disjointness.RandomizedLowerBound.not_xBit_left_and_yBit_right_of_same_conditioning}
\label{CommunicationComplexity.Functions.Disjointness.RandomizedLowerBound.not_xBit_left_and_yBit_right_of_same_conditioning}
\leanok
\uses{CommunicationComplexity.Functions.Disjointness.RandomizedLowerBound.HardSample, CommunicationComplexity.Functions.Disjointness.RandomizedLowerBound.not_xBit_and_yBit_of_ne_special, CommunicationComplexity.Functions.Disjointness.RandomizedLowerBound.xBeforeSpecial, CommunicationComplexity.Functions.Disjointness.RandomizedLowerBound.xBit, CommunicationComplexity.Functions.Disjointness.RandomizedLowerBound.yAfterSpecial, CommunicationComplexity.Functions.Disjointness.RandomizedLowerBound.yBit}
\difficulty{5}
Let $\omega_X$ and $\omega_Y$ be two samples of the hard distribution that share the
same special coordinate $T$, whose Alice-before-special vectors agree and whose
Bob-after-special vectors agree. Then for every coordinate $i \neq T$, it is never the
case that Alice's bit at $i$ under $\omega_X$ and Bob's bit at $i$ under $\omega_Y$ are
both true.
\end{theorem}

\begin{theorem}[Mixing preserves Alice's input bits]
\lean{CommunicationComplexity.Functions.Disjointness.RandomizedLowerBound.xBit_mix}
\label{CommunicationComplexity.Functions.Disjointness.RandomizedLowerBound.xBit_mix}
\leanok
\uses{CommunicationComplexity.Functions.Disjointness.RandomizedLowerBound.DisjointCoordinate.xBit, CommunicationComplexity.Functions.Disjointness.RandomizedLowerBound.HardSample, CommunicationComplexity.Functions.Disjointness.RandomizedLowerBound.coordinateOfBits_xBit, CommunicationComplexity.Functions.Disjointness.RandomizedLowerBound.mix, CommunicationComplexity.Functions.Disjointness.RandomizedLowerBound.not_xBit_left_and_yBit_right_of_same_conditioning, CommunicationComplexity.Functions.Disjointness.RandomizedLowerBound.xBeforeSpecial, CommunicationComplexity.Functions.Disjointness.RandomizedLowerBound.xBit, CommunicationComplexity.Functions.Disjointness.RandomizedLowerBound.xBit_of_ne_special, CommunicationComplexity.Functions.Disjointness.RandomizedLowerBound.yAfterSpecial, CommunicationComplexity.Functions.Disjointness.RandomizedLowerBound.yBit}
\difficulty{4}
Let $\omega_X$ and $\omega_Y$ be two hard-distribution samples that share the same
special coordinate $T$, whose Alice-inputs on the coordinates before $T$ agree, and
whose Bob-inputs on the coordinates after $T$ agree. Then at every coordinate $i$,
Alice's bit under the mixed sample $\mathrm{mix}(\omega_X,\omega_Y)$ equals Alice's bit
under $\omega_X$ at that coordinate.
\end{theorem}

\begin{theorem}[Bob's bits are preserved by the mix]
\lean{CommunicationComplexity.Functions.Disjointness.RandomizedLowerBound.yBit_mix}
\label{CommunicationComplexity.Functions.Disjointness.RandomizedLowerBound.yBit_mix}
\leanok
\uses{CommunicationComplexity.Functions.Disjointness.RandomizedLowerBound.DisjointCoordinate.yBit, CommunicationComplexity.Functions.Disjointness.RandomizedLowerBound.HardSample, CommunicationComplexity.Functions.Disjointness.RandomizedLowerBound.X_mix, CommunicationComplexity.Functions.Disjointness.RandomizedLowerBound.coordinateOfBits_yBit, CommunicationComplexity.Functions.Disjointness.RandomizedLowerBound.mix, CommunicationComplexity.Functions.Disjointness.RandomizedLowerBound.not_xBit_left_and_yBit_right_of_same_conditioning, CommunicationComplexity.Functions.Disjointness.RandomizedLowerBound.xBeforeSpecial, CommunicationComplexity.Functions.Disjointness.RandomizedLowerBound.xBit, CommunicationComplexity.Functions.Disjointness.RandomizedLowerBound.yAfterSpecial, CommunicationComplexity.Functions.Disjointness.RandomizedLowerBound.yBit, CommunicationComplexity.Functions.Disjointness.RandomizedLowerBound.yBit_of_ne_special}
\difficulty{4}
Let $\omega_X$ and $\omega_Y$ be two samples of the hard distribution that share the
same special coordinate $T$ (so their $T$-fields agree), that agree on the profile of
Alice's bits at the coordinates strictly below $T$ (each padded by $\mathrm{false}$ at
the remaining coordinates), and that agree on the profile of Bob's bits at the
coordinates strictly above $T$ (again each padded by $\mathrm{false}$ elsewhere). Form
the mixed sample that takes its special coordinate, Alice's bit there, and Alice's input
from $\omega_X$ while taking Bob's bit at $T$ and Bob's input from $\omega_Y$. Then for
every coordinate $i$, Bob's bit at $i$ in this mixed sample equals Bob's bit at $i$ in
$\omega_Y$.
\end{theorem}

\begin{theorem}[Alice's set is preserved under mixing]
\lean{CommunicationComplexity.Functions.Disjointness.RandomizedLowerBound.X_mix}
\label{CommunicationComplexity.Functions.Disjointness.RandomizedLowerBound.X_mix}
\leanok
\uses{CommunicationComplexity.Functions.Disjointness.RandomizedLowerBound.HardSample, CommunicationComplexity.Functions.Disjointness.RandomizedLowerBound.X, CommunicationComplexity.Functions.Disjointness.RandomizedLowerBound.mix, CommunicationComplexity.Functions.Disjointness.RandomizedLowerBound.xBeforeSpecial, CommunicationComplexity.Functions.Disjointness.RandomizedLowerBound.xBit_mix, CommunicationComplexity.Functions.Disjointness.RandomizedLowerBound.yAfterSpecial}
\difficulty{2}
Let $\omega_X$ and $\omega_Y$ be two samples of the hard distribution that agree on the
special coordinate $T$, on Alice's padded input bits before $T$, and on Bob's padded
input bits after $T$. Then the mixed sample, which takes Alice's data from $\omega_X$
and Bob's data from $\omega_Y$, induces exactly Alice's input set of $\omega_X$; that
is, $X(\mathrm{mix}(\omega_X,\omega_Y)) = X(\omega_X)$.
\end{theorem}

\begin{theorem}[Bob's input set of the mixed sample]
\lean{CommunicationComplexity.Functions.Disjointness.RandomizedLowerBound.Y_mix}
\label{CommunicationComplexity.Functions.Disjointness.RandomizedLowerBound.Y_mix}
\leanok
\uses{CommunicationComplexity.Functions.Disjointness.RandomizedLowerBound.HardSample, CommunicationComplexity.Functions.Disjointness.RandomizedLowerBound.X_mix, CommunicationComplexity.Functions.Disjointness.RandomizedLowerBound.Y, CommunicationComplexity.Functions.Disjointness.RandomizedLowerBound.mix, CommunicationComplexity.Functions.Disjointness.RandomizedLowerBound.xBeforeSpecial, CommunicationComplexity.Functions.Disjointness.RandomizedLowerBound.yAfterSpecial, CommunicationComplexity.Functions.Disjointness.RandomizedLowerBound.yBit_mix}
\difficulty{2}
Fix a positive integer $n$ and consider two samples $\omega_X$ and $\omega_Y$ drawn from
the hard-distribution sample space on the coordinates $\{0,1,\dots,n-1\}$. Suppose the
two samples share the same special coordinate $T$, that they agree on Alice's input bits
before the special coordinate, and that they agree on Bob's input bits after the special
coordinate. Then Bob's input set of the mixed sample
$\operatorname{mix}(\omega_X,\omega_Y)$ — the sample that takes Alice's data from
$\omega_X$ and Bob's data from $\omega_Y$ — coincides with Bob's input set of
$\omega_Y$.
\end{theorem}

\begin{theorem}[Input of a mixed hard-distribution sample]
\lean{CommunicationComplexity.Functions.Disjointness.RandomizedLowerBound.input_mix}
\label{CommunicationComplexity.Functions.Disjointness.RandomizedLowerBound.input_mix}
\leanok
\uses{CommunicationComplexity.Functions.Disjointness.RandomizedLowerBound.HardSample, CommunicationComplexity.Functions.Disjointness.RandomizedLowerBound.X, CommunicationComplexity.Functions.Disjointness.RandomizedLowerBound.X_mix, CommunicationComplexity.Functions.Disjointness.RandomizedLowerBound.Y, CommunicationComplexity.Functions.Disjointness.RandomizedLowerBound.Y_mix, CommunicationComplexity.Functions.Disjointness.RandomizedLowerBound.input, CommunicationComplexity.Functions.Disjointness.RandomizedLowerBound.mix, CommunicationComplexity.Functions.Disjointness.RandomizedLowerBound.xBeforeSpecial, CommunicationComplexity.Functions.Disjointness.RandomizedLowerBound.yAfterSpecial}
\difficulty{1}
Let $\omega_X$ and $\omega_Y$ be two hard-distribution samples that agree on their
special coordinate $T$, whose associated Alice-bit vectors before the special coordinate
coincide, and whose associated Bob-bit vectors after the special coordinate coincide.
Then the input pair generated by mixing Alice's data from $\omega_X$ with Bob's data
from $\omega_Y$ is exactly $(X(\omega_X),\, Y(\omega_Y))$: the Alice set of the mixed
sample equals Alice's input set under $\omega_X$, and the Bob set of the mixed sample
equals Bob's input set under $\omega_Y$.
\end{theorem}

\begin{theorem}[Mixing preserves Alice's before-special vector]
\lean{CommunicationComplexity.Functions.Disjointness.RandomizedLowerBound.xBeforeSpecial_mix}
\label{CommunicationComplexity.Functions.Disjointness.RandomizedLowerBound.xBeforeSpecial_mix}
\leanok
\uses{CommunicationComplexity.Functions.Disjointness.RandomizedLowerBound.HardSample, CommunicationComplexity.Functions.Disjointness.RandomizedLowerBound.mix, CommunicationComplexity.Functions.Disjointness.RandomizedLowerBound.specialY_mix, CommunicationComplexity.Functions.Disjointness.RandomizedLowerBound.xBeforeSpecial, CommunicationComplexity.Functions.Disjointness.RandomizedLowerBound.xBit_mix, CommunicationComplexity.Functions.Disjointness.RandomizedLowerBound.yAfterSpecial}
\difficulty{3}
Let $\omega_X$ and $\omega_Y$ be two samples of the hard distribution that share the
same special coordinate $T$, that agree on Alice's before-special vector (the bits Alice
holds at the coordinates strictly below $T$, padded with $\mathrm{false}$ elsewhere),
and that agree on Bob's after-special vector (the bits Bob holds at the coordinates
strictly above $T$, padded with $\mathrm{false}$ elsewhere). Then the mixed sample
$\mathrm{mix}(\omega_X,\omega_Y)$, which takes Alice's data from $\omega_X$ and Bob's
data from $\omega_Y$, has the same before-special vector as $\omega_X$.
\end{theorem}

\begin{theorem}[Mixing preserves Bob's after-special profile]
\lean{CommunicationComplexity.Functions.Disjointness.RandomizedLowerBound.yAfterSpecial_mix}
\label{CommunicationComplexity.Functions.Disjointness.RandomizedLowerBound.yAfterSpecial_mix}
\leanok
\uses{CommunicationComplexity.Functions.Disjointness.RandomizedLowerBound.HardSample, CommunicationComplexity.Functions.Disjointness.RandomizedLowerBound.mix, CommunicationComplexity.Functions.Disjointness.RandomizedLowerBound.specialY_mix, CommunicationComplexity.Functions.Disjointness.RandomizedLowerBound.xBeforeSpecial, CommunicationComplexity.Functions.Disjointness.RandomizedLowerBound.yAfterSpecial, CommunicationComplexity.Functions.Disjointness.RandomizedLowerBound.yBit_mix}
\difficulty{3}
Let $\omega_X$ and $\omega_Y$ be two samples from the hard distribution on $n$
coordinates, and suppose they share the same special coordinate, $T(\omega_X) =
T(\omega_Y)$, that Alice's before-special profiles agree, $x_{<T}(\omega_X) =
x_{<T}(\omega_Y)$, and that Bob's after-special profiles agree, $y_{>T}(\omega_X) =
y_{>T}(\omega_Y)$. Then Bob's after-special profile of the mixed sample
$\mathrm{mix}(\omega_X,\omega_Y)$ coincides with that of $\omega_Y$: \[
y_{>T}\bigl(\mathrm{mix}(\omega_X,\omega_Y)\bigr) = y_{>T}(\omega_Y). \]
\end{theorem}

\begin{theorem}[Special bit is preserved by mixing]
\lean{CommunicationComplexity.Functions.Disjointness.RandomizedLowerBound.specialX_mix}
\label{CommunicationComplexity.Functions.Disjointness.RandomizedLowerBound.specialX_mix}
\leanok
\uses{CommunicationComplexity.Functions.Disjointness.RandomizedLowerBound.HardSample, CommunicationComplexity.Functions.Disjointness.RandomizedLowerBound.mix, CommunicationComplexity.Functions.Disjointness.RandomizedLowerBound.specialX}
\difficulty{1}
Fix a positive integer $n$, and let $\omega_X$ and $\omega_Y$ be two samples from the
hard distribution, each carrying a special coordinate together with Alice's and Bob's
bits there and a disjoint coordinate on every other index. Form the mixed sample that
takes Alice's data from $\omega_X$ and Bob's data from $\omega_Y$. Then Alice's bit at
the special coordinate of the mixed sample equals Alice's bit at the special coordinate
of $\omega_X$.
\end{theorem}

\begin{theorem}[Bob's special bit is preserved by mixing]
\lean{CommunicationComplexity.Functions.Disjointness.RandomizedLowerBound.specialY_mix}
\label{CommunicationComplexity.Functions.Disjointness.RandomizedLowerBound.specialY_mix}
\leanok
\uses{CommunicationComplexity.Functions.Disjointness.RandomizedLowerBound.HardSample, CommunicationComplexity.Functions.Disjointness.RandomizedLowerBound.mix, CommunicationComplexity.Functions.Disjointness.RandomizedLowerBound.specialY}
\difficulty{1}
Let $\omega_X$ and $\omega_Y$ be two samples of the hard distribution, and form the
mixed sample that takes Alice's input from $\omega_X$ and Bob's input from $\omega_Y$.
Then Bob's bit at the special coordinate of the mixed sample equals Bob's bit at the
special coordinate of $\omega_Y$.
\end{theorem}

\begin{theorem}[Double mixing recovers the first sample]
\lean{CommunicationComplexity.Functions.Disjointness.RandomizedLowerBound.mix_mix_swap}
\label{CommunicationComplexity.Functions.Disjointness.RandomizedLowerBound.mix_mix_swap}
\leanok
\uses{CommunicationComplexity.Functions.Disjointness.RandomizedLowerBound.HardSample, CommunicationComplexity.Functions.Disjointness.RandomizedLowerBound.coordinateOfBits, CommunicationComplexity.Functions.Disjointness.RandomizedLowerBound.coordinateOfBits_xBit_yBit, CommunicationComplexity.Functions.Disjointness.RandomizedLowerBound.mix, CommunicationComplexity.Functions.Disjointness.RandomizedLowerBound.xBeforeSpecial, CommunicationComplexity.Functions.Disjointness.RandomizedLowerBound.xBit, CommunicationComplexity.Functions.Disjointness.RandomizedLowerBound.xBit_mix, CommunicationComplexity.Functions.Disjointness.RandomizedLowerBound.xBit_of_ne_special, CommunicationComplexity.Functions.Disjointness.RandomizedLowerBound.yAfterSpecial, CommunicationComplexity.Functions.Disjointness.RandomizedLowerBound.yBit, CommunicationComplexity.Functions.Disjointness.RandomizedLowerBound.yBit_mix, CommunicationComplexity.Functions.Disjointness.RandomizedLowerBound.yBit_of_ne_special}
\difficulty{5}
For two samples $\omega,\omega'$ of the hard distribution, let $\omega\ast\omega'$
denote their *mixture*: the sample built by taking Alice's input from $\omega$ and Bob's
input from $\omega'$. Let $\omega_X$ and $\omega_Y$ be two samples that share the same
special coordinate $T$, whose Alice-bits before the special coordinate agree, and whose
Bob-bits after the special coordinate agree. Then
\[
(\omega_X\ast\omega_Y)\ast(\omega_Y\ast\omega_X)=\omega_X .
\]
\end{theorem}

\begin{theorem}[Rectangle events of the hard sample factor as a product]
\lean{CommunicationComplexity.Functions.Disjointness.RandomizedLowerBound.hardSample_measureReal_fieldProduct}
\label{CommunicationComplexity.Functions.Disjointness.RandomizedLowerBound.hardSample_measureReal_fieldProduct}
\leanok
\uses{CommunicationComplexity.Functions.Disjointness.RandomizedLowerBound.HardSample, CommunicationComplexity.Functions.Disjointness.RandomizedLowerBound.HardSample.card, CommunicationComplexity.Functions.Disjointness.RandomizedLowerBound.uniformBool, CommunicationComplexity.Functions.Disjointness.RandomizedLowerBound.uniformBool_real, CommunicationComplexity.Functions.Disjointness.RandomizedLowerBound.uniformDisjointCoordinateVector, CommunicationComplexity.Functions.Disjointness.RandomizedLowerBound.uniformDisjointCoordinateVector_real, CommunicationComplexity.Functions.Disjointness.RandomizedLowerBound.uniformFin, CommunicationComplexity.Functions.Disjointness.RandomizedLowerBound.uniformFin_real, CommunicationComplexity.uniformOn_univ_measureReal_eq_card_subtype}
\difficulty{5}
Fix a positive integer $n$, and let $\mathrm{HardSample}(n)$ be the finite sample space
whose points $\omega$ record a special coordinate $\omega.T \in \{0,\dots,n-1\}$, two
bits $\omega.x_T,\omega.y_T \in \{0,1\}$, and a coordinate vector
$\omega.\mathrm{other}\colon \{0,\dots,n-1\} \to \mathrm{DisjointCoordinate}$. For any
set $\mathcal T$ of special coordinates, any sets $\mathcal X,\mathcal Y$ of bits, and
any set $\mathcal O$ of coordinate vectors, the ambient measure of the rectangle event
\[
\bigl\{\omega : \omega.T \in \mathcal T,\ \omega.x_T \in \mathcal X,\ \omega.y_T \in
\mathcal Y,\ \omega.\mathrm{other} \in \mathcal O\bigr\},
\]
taken as a real number, equals the product
\[
\mu_{\mathrm{Fin}}(\mathcal T)\cdot \mu_{\mathrm{bit}}(\mathcal X)\cdot
\mu_{\mathrm{bit}}(\mathcal Y)\cdot \mu_{\mathrm{vec}}(\mathcal O),
\]
where $\mu_{\mathrm{Fin}}$ is the uniform law on the special coordinate,
$\mu_{\mathrm{bit}}$ is the uniform law on one bit, and $\mu_{\mathrm{vec}}$ is the
uniform law on coordinate vectors, each evaluated on the corresponding set as a real
number.
\end{theorem}

\begin{theorem}[Probability of a genuine intersection at the special coordinate]
\lean{CommunicationComplexity.Functions.Disjointness.RandomizedLowerBound.measureReal_specialIntersect}
\statementsource{roughgarden-cc-bootcamp}{pdf-d54d402b16e2-p016-b002}
\label{CommunicationComplexity.Functions.Disjointness.RandomizedLowerBound.measureReal_specialIntersect}
\leanok
\uses{CommunicationComplexity.Functions.Disjointness.RandomizedLowerBound.HardSample, CommunicationComplexity.Functions.Disjointness.RandomizedLowerBound.hardSample_measureReal_fieldProduct, CommunicationComplexity.Functions.Disjointness.RandomizedLowerBound.specialIntersect, CommunicationComplexity.Functions.Disjointness.RandomizedLowerBound.uniformBool_singleton, CommunicationComplexity.Functions.Disjointness.RandomizedLowerBound.uniformDisjointCoordinateVector_univ, CommunicationComplexity.Functions.Disjointness.RandomizedLowerBound.uniformFin_univ}
\difficulty{3}
Fix an integer $n \ge 1$, and let the sample space consist of all tuples $(T, x_T, y_T,
c)$, where $T \in \{1,\dots,n\}$ is a distinguished coordinate, $x_T, y_T \in
\{\mathrm{true}, \mathrm{false}\}$ are two Boolean values, and $c$ assigns to each
coordinate one of the three labels *neither*, *leftOnly*, or *rightOnly*; equip this
finite space with the uniform probability measure. Then the event that the distinguished
coordinate is a genuine intersection, namely $\{x_T = \mathrm{true} \text{ and } y_T =
\mathrm{true}\}$, has probability $\tfrac14$.
\end{theorem}

\begin{theorem}[Probability of prescribed special bits]
\lean{CommunicationComplexity.Functions.Disjointness.RandomizedLowerBound.measureReal_specialBitsEvent}
\label{CommunicationComplexity.Functions.Disjointness.RandomizedLowerBound.measureReal_specialBitsEvent}
\leanok
\uses{CommunicationComplexity.Functions.Disjointness.RandomizedLowerBound.HardSample, CommunicationComplexity.Functions.Disjointness.RandomizedLowerBound.hardSample_measureReal_fieldProduct, CommunicationComplexity.Functions.Disjointness.RandomizedLowerBound.specialBitsEvent, CommunicationComplexity.Functions.Disjointness.RandomizedLowerBound.uniformBool_singleton, CommunicationComplexity.Functions.Disjointness.RandomizedLowerBound.uniformDisjointCoordinateVector_univ, CommunicationComplexity.Functions.Disjointness.RandomizedLowerBound.uniformFin_univ}
\difficulty{3}
Consider the finite sample space whose elements record a distinguished coordinate $T$,
two special bits $x_T, y_T \in \{0,1\}$, and a labelling of every other coordinate by
one of three values, equipped with the uniform probability measure. For any prescribed
pair of bits $b_x, b_y \in \{0,1\}$, the event that $x_T = b_x$ and $y_T = b_y$ has
probability $\tfrac{1}{4}$.
\end{theorem}

\begin{theorem}[Uniform law of the special coordinate]
\lean{CommunicationComplexity.Functions.Disjointness.RandomizedLowerBound.measureReal_specialCoordinateEvent}
\label{CommunicationComplexity.Functions.Disjointness.RandomizedLowerBound.measureReal_specialCoordinateEvent}
\leanok
\uses{CommunicationComplexity.Functions.Disjointness.RandomizedLowerBound.HardSample, CommunicationComplexity.Functions.Disjointness.RandomizedLowerBound.hardSample_measureReal_fieldProduct, CommunicationComplexity.Functions.Disjointness.RandomizedLowerBound.specialCoordinateEvent, CommunicationComplexity.Functions.Disjointness.RandomizedLowerBound.uniformBool_univ, CommunicationComplexity.Functions.Disjointness.RandomizedLowerBound.uniformDisjointCoordinateVector_univ, CommunicationComplexity.Functions.Disjointness.RandomizedLowerBound.uniformFin_singleton}
\difficulty{3}
Fix a positive integer $n$ and an index $i \in \{0, 1, \dots, n-1\}$. Equip the finite
sample space of hard samples with its uniform probability measure. Then the probability
of the event that the special coordinate $T$ equals $i$ is
\[
\frac{1}{n}.
\]
\end{theorem}

\begin{theorem}[Probability the special coordinate is a fixed intersecting index]
\lean{CommunicationComplexity.Functions.Disjointness.RandomizedLowerBound.measureReal_specialCoordinateEvent_inter_specialIntersect}
\label{CommunicationComplexity.Functions.Disjointness.RandomizedLowerBound.measureReal_specialCoordinateEvent_inter_specialIntersect}
\leanok
\uses{CommunicationComplexity.Functions.Disjointness.RandomizedLowerBound.HardSample, CommunicationComplexity.Functions.Disjointness.RandomizedLowerBound.hardSample_measureReal_fieldProduct, CommunicationComplexity.Functions.Disjointness.RandomizedLowerBound.specialCoordinateEvent, CommunicationComplexity.Functions.Disjointness.RandomizedLowerBound.specialIntersect, CommunicationComplexity.Functions.Disjointness.RandomizedLowerBound.uniformBool_singleton, CommunicationComplexity.Functions.Disjointness.RandomizedLowerBound.uniformDisjointCoordinateVector_univ, CommunicationComplexity.Functions.Disjointness.RandomizedLowerBound.uniformFin_singleton}
\difficulty{3}
Equip the finite space of hard samples with the uniform probability measure, where each
sample consists of a special coordinate $T \in \{1,\dots,n\}$, two Boolean bits $x_T$
and $y_T$, and a value in $\{\text{neither},\text{leftOnly},\text{rightOnly}\}$ for
every one of the $n$ coordinates. Then for each fixed coordinate $i \in \{1,\dots,n\}$,
the probability that the special coordinate equals $i$ and that both special bits are
true, that is $T = i$ and $x_T = y_T = \text{true}$, is exactly $\dfrac{1}{4n}$.
\end{theorem}

\begin{theorem}[Probability that both special bits vanish]
\lean{CommunicationComplexity.Functions.Disjointness.RandomizedLowerBound.measureReal_specialZeroZero}
\label{CommunicationComplexity.Functions.Disjointness.RandomizedLowerBound.measureReal_specialZeroZero}
\leanok
\uses{CommunicationComplexity.Functions.Disjointness.RandomizedLowerBound.measureReal_specialBitsEvent, CommunicationComplexity.Functions.Disjointness.RandomizedLowerBound.specialZeroZero}
\difficulty{1}
Consider the hard distribution's sample space, whose points carry two special-coordinate
bits $x_T, y_T \in \{0,1\}$ (among other data), equipped with the uniform probability
measure. The event that both special bits vanish, $x_T = 0$ and $y_T = 0$, has
probability $\tfrac{1}{4}$.
\end{theorem}

\begin{theorem}[Bob's special bit is false with probability one half]
\lean{CommunicationComplexity.Functions.Disjointness.RandomizedLowerBound.measureReal_specialY_false}
\label{CommunicationComplexity.Functions.Disjointness.RandomizedLowerBound.measureReal_specialY_false}
\leanok
\uses{CommunicationComplexity.Functions.Disjointness.RandomizedLowerBound.HardSample, CommunicationComplexity.Functions.Disjointness.RandomizedLowerBound.hardSample_measureReal_fieldProduct, CommunicationComplexity.Functions.Disjointness.RandomizedLowerBound.specialY, CommunicationComplexity.Functions.Disjointness.RandomizedLowerBound.uniformBool_singleton, CommunicationComplexity.Functions.Disjointness.RandomizedLowerBound.uniformBool_univ, CommunicationComplexity.Functions.Disjointness.RandomizedLowerBound.uniformDisjointCoordinateVector_univ, CommunicationComplexity.Functions.Disjointness.RandomizedLowerBound.uniformFin_univ}
\difficulty{3}
Consider the finite sample space of hard samples for a fixed number of coordinates $n
\ge 1$, equipped with the uniform probability distribution, where each sample records in
particular Bob's bit $y_T \in \{\mathrm{true}, \mathrm{false}\}$ at the special
coordinate. Then the event that this bit takes the value $\mathrm{false}$ has
probability $\tfrac{1}{2}$.
\end{theorem}

\begin{theorem}[The disjoint event as the complement of a common special coordinate]
\lean{CommunicationComplexity.Functions.Disjointness.RandomizedLowerBound.disjointEvent_eq_compl_specialIntersect}
\label{CommunicationComplexity.Functions.Disjointness.RandomizedLowerBound.disjointEvent_eq_compl_specialIntersect}
\leanok
\uses{CommunicationComplexity.Functions.Disjointness.RandomizedLowerBound.disjointEvent, CommunicationComplexity.Functions.Disjointness.RandomizedLowerBound.disjoint_X_Y_iff, CommunicationComplexity.Functions.Disjointness.RandomizedLowerBound.specialIntersect}
\difficulty{2}
Fix an integer $n \ge 1$, and consider the finite space of hard-distribution samples on
$n$ coordinates, each sample consisting of a special coordinate $T \in \{1,\dots,n\}$,
two bits $x_T$ and $y_T$, and a disjoint-coordinate label (one of *neither*, *left
only*, *right only*) attached to every coordinate. From such a sample one forms Alice's
set $X$ and Bob's set $Y$. Then, as subsets of this sample space, the event that $X$ and
$Y$ are disjoint equals the complement of the special-intersection event; equivalently,
$X$ and $Y$ fail to be disjoint precisely on those samples for which $x_T$ and $y_T$ are
both true.
\end{theorem}

\begin{theorem}[Disjointness probability of the hard distribution]
\lean{CommunicationComplexity.Functions.Disjointness.RandomizedLowerBound.measureReal_disjointEvent}
\statementsource{roughgarden-cc-bootcamp}{pdf-d54d402b16e2-p016-b002}
\label{CommunicationComplexity.Functions.Disjointness.RandomizedLowerBound.measureReal_disjointEvent}
\leanok
\uses{CommunicationComplexity.Functions.Disjointness.RandomizedLowerBound.disjointEvent, CommunicationComplexity.Functions.Disjointness.RandomizedLowerBound.disjointEvent_eq_compl_specialIntersect, CommunicationComplexity.Functions.Disjointness.RandomizedLowerBound.measureReal_specialIntersect, CommunicationComplexity.Functions.Disjointness.RandomizedLowerBound.specialIntersect}
\difficulty{3}
Fix a positive integer $n$, and draw a hard sample $\omega$ uniformly at random from the
finite sample space of the hard distribution on $n$ coordinates. Then the disjointness
event — the event that Alice's input set $X(\omega)$ and Bob's input set $Y(\omega)$ are
disjoint — has probability exactly $3/4$.
\end{theorem}

\begin{theorem}[Probability of the special coordinate and disjointness]
\lean{CommunicationComplexity.Functions.Disjointness.RandomizedLowerBound.measureReal_specialCoordinateEvent_inter_disjointEvent}
\label{CommunicationComplexity.Functions.Disjointness.RandomizedLowerBound.measureReal_specialCoordinateEvent_inter_disjointEvent}
\leanok
\uses{CommunicationComplexity.Functions.Disjointness.RandomizedLowerBound.disjointEvent, CommunicationComplexity.Functions.Disjointness.RandomizedLowerBound.disjointEvent_eq_compl_specialIntersect, CommunicationComplexity.Functions.Disjointness.RandomizedLowerBound.measureReal_specialCoordinateEvent, CommunicationComplexity.Functions.Disjointness.RandomizedLowerBound.measureReal_specialCoordinateEvent_inter_specialIntersect, CommunicationComplexity.Functions.Disjointness.RandomizedLowerBound.specialCoordinateEvent, CommunicationComplexity.Functions.Disjointness.RandomizedLowerBound.specialIntersect}
\difficulty{4}
Fix $n \ge 1$ and equip the sample space of the hard distribution with the uniform
probability measure, so that a sample consists of a special coordinate $T \in
\{1,\dots,n\}$, Alice's and Bob's bits $x_T, y_T$ at that coordinate, and an independent
uniformly chosen label in $\{\text{neither},\text{leftOnly},\text{rightOnly}\}$ at each
of the $n$ coordinates. For any fixed coordinate $i \in \{1,\dots,n\}$, the probability
that the special coordinate satisfies $T = i$ and that the generated sets $X$ and $Y$
are disjoint equals \[ \frac{3}{4n}. \]
\end{theorem}

\begin{theorem}[Positive measure of the disjointness event]
\lean{CommunicationComplexity.Functions.Disjointness.RandomizedLowerBound.measure_disjointEvent_ne_zero}
\label{CommunicationComplexity.Functions.Disjointness.RandomizedLowerBound.measure_disjointEvent_ne_zero}
\leanok
\uses{CommunicationComplexity.Functions.Disjointness.RandomizedLowerBound.disjointEvent, CommunicationComplexity.Functions.Disjointness.RandomizedLowerBound.measureReal_disjointEvent}
\difficulty{2}
Fix a positive integer $n$, and consider the finite sample space of hard-distribution
samples, each of which determines Alice's input set $X \subseteq \{1,\dots,n\}$ and
Bob's input set $Y \subseteq \{1,\dots,n\}$. Under the ambient measure on this sample
space, the disjoint event — the set of samples for which $X$ and $Y$ are disjoint — has
nonzero measure.
\end{theorem}

\begin{definition}[Disjointcondmeasure]
\lean{CommunicationComplexity.Functions.Disjointness.RandomizedLowerBound.disjointCondMeasure}
\label{CommunicationComplexity.Functions.Disjointness.RandomizedLowerBound.disjointCondMeasure}
\leanok
\uses{CommunicationComplexity.Functions.Disjointness.RandomizedLowerBound.HardSample, CommunicationComplexity.Functions.Disjointness.RandomizedLowerBound.disjointEvent}
The hard distribution conditioned on the generated input being disjoint.
\end{definition}

\begin{theorem}[Uniformity of the special coordinate under the disjoint-conditioned distribution]
\lean{CommunicationComplexity.Functions.Disjointness.RandomizedLowerBound.disjointCondMeasure_measureReal_specialCoordinateEvent}
\label{CommunicationComplexity.Functions.Disjointness.RandomizedLowerBound.disjointCondMeasure_measureReal_specialCoordinateEvent}
\leanok
\uses{CommunicationComplexity.Functions.Disjointness.RandomizedLowerBound.disjointCondMeasure, CommunicationComplexity.Functions.Disjointness.RandomizedLowerBound.disjointEvent, CommunicationComplexity.Functions.Disjointness.RandomizedLowerBound.measureReal_disjointEvent, CommunicationComplexity.Functions.Disjointness.RandomizedLowerBound.measureReal_specialCoordinateEvent_inter_disjointEvent, CommunicationComplexity.Functions.Disjointness.RandomizedLowerBound.specialCoordinateEvent}
\difficulty{3}
Fix $n \ge 1$, and consider the hard distribution on $\mathrm{HardSample}(n)$
conditioned on the disjointness event (that Alice's set $X$ and Bob's set $Y$ are
disjoint). Then for every coordinate $i \in \{0,1,\dots,n-1\}$, the conditional
probability that the special coordinate $T$ equals $i$ is exactly
\[
\frac{1}{n}.
\]
In other words, conditioning on disjointness leaves the special coordinate $T$ uniformly
distributed over the $n$ coordinates.
\end{theorem}

\begin{theorem}[Preimage of a value under the special coordinate]
\lean{CommunicationComplexity.Functions.Disjointness.RandomizedLowerBound.specialCoordinate_preimage_singleton}
\label{CommunicationComplexity.Functions.Disjointness.RandomizedLowerBound.specialCoordinate_preimage_singleton}
\leanok
\uses{CommunicationComplexity.Functions.Disjointness.RandomizedLowerBound.specialCoordinate, CommunicationComplexity.Functions.Disjointness.RandomizedLowerBound.specialCoordinateEvent}
\difficulty{2}
Fix a positive integer $n$ and a coordinate $i \in \{0,1,\dots,n-1\}$. The set of hard
samples whose special coordinate equals $i$ — that is, the preimage under the
special-coordinate map of the singleton $\{i\}$ — coincides exactly with the event that
the special coordinate takes the value $i$.
\end{theorem}

\begin{theorem}[Uniformity of the special coordinate under the disjoint conditioning]
\lean{CommunicationComplexity.Functions.Disjointness.RandomizedLowerBound.disjointCondMeasure_measureReal_specialCoordinate_preimage_singleton}
\label{CommunicationComplexity.Functions.Disjointness.RandomizedLowerBound.disjointCondMeasure_measureReal_specialCoordinate_preimage_singleton}
\leanok
\uses{CommunicationComplexity.Functions.Disjointness.RandomizedLowerBound.disjointCondMeasure, CommunicationComplexity.Functions.Disjointness.RandomizedLowerBound.disjointCondMeasure_measureReal_specialCoordinateEvent, CommunicationComplexity.Functions.Disjointness.RandomizedLowerBound.specialCoordinate, CommunicationComplexity.Functions.Disjointness.RandomizedLowerBound.specialCoordinate_preimage_singleton}
\difficulty{1}
Fix a positive integer $n$, and consider the hard distribution on the sample space
conditioned on the event that the generated input pair $(X,Y)$ is disjoint. Writing $T
\in \{0,1,\dots,n-1\}$ for the special coordinate of a sample, then for every index $i
\in \{0,1,\dots,n-1\}$ the conditional probability that $T = i$ equals $\tfrac{1}{n}$;
that is, the special coordinate remains uniformly distributed after conditioning on
disjointness.
\end{theorem}

\begin{theorem}[Disjointness holds almost surely under the conditioned distribution]
\lean{CommunicationComplexity.Functions.Disjointness.RandomizedLowerBound.disjointCondMeasure_ae_disjointEvent}
\label{CommunicationComplexity.Functions.Disjointness.RandomizedLowerBound.disjointCondMeasure_ae_disjointEvent}
\leanok
\uses{CommunicationComplexity.Functions.Disjointness.RandomizedLowerBound.disjointCondMeasure, CommunicationComplexity.Functions.Disjointness.RandomizedLowerBound.disjointEvent}
\difficulty{1}
Consider the hard distribution on samples $\omega$, each of which determines Alice's
input set $X(\omega)$ and Bob's input set $Y(\omega)$, and let $D$ be the disjointness
event $\{\omega : X(\omega) \cap Y(\omega) = \varnothing\}$. Under the distribution
obtained by conditioning the hard distribution on $D$, almost every sample belongs to
$D$; that is, with respect to the conditioned measure, the sets $X(\omega)$ and
$Y(\omega)$ are disjoint almost surely.
\end{theorem}

\begin{theorem}[Disjointness when Bob's special bit is off]
\lean{CommunicationComplexity.Functions.Disjointness.RandomizedLowerBound.mem_disjointEvent_of_specialY_eq_false}
\label{CommunicationComplexity.Functions.Disjointness.RandomizedLowerBound.mem_disjointEvent_of_specialY_eq_false}
\leanok
\uses{CommunicationComplexity.Functions.Disjointness.RandomizedLowerBound.HardSample, CommunicationComplexity.Functions.Disjointness.RandomizedLowerBound.X, CommunicationComplexity.Functions.Disjointness.RandomizedLowerBound.Y, CommunicationComplexity.Functions.Disjointness.RandomizedLowerBound.disjointEvent, CommunicationComplexity.Functions.Disjointness.RandomizedLowerBound.disjoint_X_Y_iff, CommunicationComplexity.Functions.Disjointness.RandomizedLowerBound.specialY}
\difficulty{3}
Fix $n \ge 1$, and consider a sample from the hard distribution: a special coordinate $T
\in \{1,\dots,n\}$, bits $x_T, y_T \in \{0,1\}$ at that coordinate, and for every
coordinate a label drawn from $\{\text{neither},\text{leftOnly},\text{rightOnly}\}$ that
fixes the two players' bits away from $T$. This sample determines Alice's set $X = \{\,i
: \text{Alice's bit at } i = 1\,\}$ and Bob's set $Y = \{\,i : \text{Bob's bit at } i =
1\,\}$. If Bob's bit at the special coordinate satisfies $y_T = 0$, then $X$ and $Y$ are
disjoint.
\end{theorem}

\begin{theorem}[Special bits not both set force disjointness]
\lean{CommunicationComplexity.Functions.Disjointness.RandomizedLowerBound.specialBitsEvent_subset_disjointEvent}
\label{CommunicationComplexity.Functions.Disjointness.RandomizedLowerBound.specialBitsEvent_subset_disjointEvent}
\leanok
\uses{CommunicationComplexity.Functions.Disjointness.RandomizedLowerBound.X, CommunicationComplexity.Functions.Disjointness.RandomizedLowerBound.Y, CommunicationComplexity.Functions.Disjointness.RandomizedLowerBound.disjointEvent, CommunicationComplexity.Functions.Disjointness.RandomizedLowerBound.disjoint_X_Y_iff, CommunicationComplexity.Functions.Disjointness.RandomizedLowerBound.specialBitsEvent}
\difficulty{3}
Let $n \ge 1$, and let $b_x, b_y \in \{\mathrm{false}, \mathrm{true}\}$ be prescribed
bits that are not both $\mathrm{true}$. Then every hard-distribution sample whose
special coordinate carries Alice's bit $b_x$ and Bob's bit $b_y$ generates a disjoint
input pair: the special-bits event for $(b_x, b_y)$ is contained in the event that
Alice's set $X$ and Bob's set $Y$ are disjoint.
\end{theorem}

\begin{theorem}[Both special bits zero forces disjoint inputs]
\lean{CommunicationComplexity.Functions.Disjointness.RandomizedLowerBound.specialZeroZero_subset_disjointEvent}
\label{CommunicationComplexity.Functions.Disjointness.RandomizedLowerBound.specialZeroZero_subset_disjointEvent}
\leanok
\uses{CommunicationComplexity.Functions.Disjointness.RandomizedLowerBound.disjointEvent, CommunicationComplexity.Functions.Disjointness.RandomizedLowerBound.mem_disjointEvent_of_specialY_eq_false, CommunicationComplexity.Functions.Disjointness.RandomizedLowerBound.specialBitsEvent, CommunicationComplexity.Functions.Disjointness.RandomizedLowerBound.specialY, CommunicationComplexity.Functions.Disjointness.RandomizedLowerBound.specialZeroZero}
\difficulty{2}
Fix $n \ge 1$ and consider a hard-distribution sample $\omega$: a distinguished
coordinate $T \in \{1,\dots,n\}$, two special bits $x_T, y_T \in \{0,1\}$, and, at every
coordinate, a disjoint coordinate-type (neither, left-only, or right-only) that
determines Alice's and Bob's bits there. If both special bits vanish, $x_T = 0$ and $y_T
= 0$, then Alice's input set $X(\omega)$ and Bob's input set $Y(\omega)$ are disjoint.
Equivalently, the event that both special-coordinate bits are zero is contained in the
event that the generated input pair is disjoint.
\end{theorem}

\begin{theorem}[Special coordinate is doubly zero with probability one third]
\lean{CommunicationComplexity.Functions.Disjointness.RandomizedLowerBound.disjointCondMeasure_measureReal_specialZeroZero}
\label{CommunicationComplexity.Functions.Disjointness.RandomizedLowerBound.disjointCondMeasure_measureReal_specialZeroZero}
\leanok
\uses{CommunicationComplexity.Functions.Disjointness.RandomizedLowerBound.disjointCondMeasure, CommunicationComplexity.Functions.Disjointness.RandomizedLowerBound.disjointEvent, CommunicationComplexity.Functions.Disjointness.RandomizedLowerBound.measureReal_disjointEvent, CommunicationComplexity.Functions.Disjointness.RandomizedLowerBound.measureReal_specialZeroZero, CommunicationComplexity.Functions.Disjointness.RandomizedLowerBound.specialZeroZero, CommunicationComplexity.Functions.Disjointness.RandomizedLowerBound.specialZeroZero_subset_disjointEvent}
\difficulty{3}
Fix $n \ge 1$, and consider a sample drawn uniformly from the finite space consisting of
a distinguished coordinate $T \in \{0,\dots,n-1\}$, two special bits $x_T, y_T \in
\{0,1\}$, and, for every coordinate, an independent label from
$\{\text{neither},\text{leftOnly},\text{rightOnly}\}$; these data determine Alice's set
$X = \{i : x_i = 1\}$ and Bob's set $Y = \{i : y_i = 1\}$, where at the coordinate $T$
the two bits are $x_T$ and $y_T$ and at every other coordinate they are read off the
label (with $\text{neither}\mapsto(0,0)$, $\text{leftOnly}\mapsto(1,0)$,
$\text{rightOnly}\mapsto(0,1)$). Conditioning this uniform measure on the event that $X$
and $Y$ are disjoint, the conditional probability of the event $x_T = 0$ and $y_T = 0$
is exactly $\tfrac{1}{3}$.
\end{theorem}

\begin{theorem}[Uniform special-bit values under the disjoint-conditioned hard distribution]
\lean{CommunicationComplexity.Functions.Disjointness.RandomizedLowerBound.disjointCondMeasure_measureReal_specialBitsEvent}
\label{CommunicationComplexity.Functions.Disjointness.RandomizedLowerBound.disjointCondMeasure_measureReal_specialBitsEvent}
\leanok
\uses{CommunicationComplexity.Functions.Disjointness.RandomizedLowerBound.disjointCondMeasure, CommunicationComplexity.Functions.Disjointness.RandomizedLowerBound.disjointEvent, CommunicationComplexity.Functions.Disjointness.RandomizedLowerBound.measureReal_disjointEvent, CommunicationComplexity.Functions.Disjointness.RandomizedLowerBound.measureReal_specialBitsEvent, CommunicationComplexity.Functions.Disjointness.RandomizedLowerBound.specialBitsEvent, CommunicationComplexity.Functions.Disjointness.RandomizedLowerBound.specialBitsEvent_subset_disjointEvent}
\difficulty{3}
Consider the hard distribution on samples $\omega$, conditioned on the event that
Alice's set and Bob's set are disjoint. Let $b_x, b_y \in \{0,1\}$ be a pair of
prescribed values that are not both $1$. Then the probability, under this conditioned
distribution, that the special coordinate $T$ carries Alice's bit equal to $b_x$ and
Bob's bit equal to $b_y$ is exactly $\tfrac{1}{3}$.
\end{theorem}

\begin{theorem}[Special Bob-bit is false with probability two-thirds]
\lean{CommunicationComplexity.Functions.Disjointness.RandomizedLowerBound.disjointCondMeasure_measureReal_specialY_false}
\label{CommunicationComplexity.Functions.Disjointness.RandomizedLowerBound.disjointCondMeasure_measureReal_specialY_false}
\leanok
\uses{CommunicationComplexity.Functions.Disjointness.RandomizedLowerBound.HardSample, CommunicationComplexity.Functions.Disjointness.RandomizedLowerBound.disjointCondMeasure, CommunicationComplexity.Functions.Disjointness.RandomizedLowerBound.disjointEvent, CommunicationComplexity.Functions.Disjointness.RandomizedLowerBound.measureReal_disjointEvent, CommunicationComplexity.Functions.Disjointness.RandomizedLowerBound.measureReal_specialY_false, CommunicationComplexity.Functions.Disjointness.RandomizedLowerBound.mem_disjointEvent_of_specialY_eq_false, CommunicationComplexity.Functions.Disjointness.RandomizedLowerBound.specialY}
\difficulty{3}
Fix a positive integer $n$, and draw a point of the hard-distribution sample space
uniformly at random, then condition on the event that Alice's and Bob's generated input
sets are disjoint. Under this conditional distribution, the probability that Bob's bit
at the special coordinate $T$ equals $\mathrm{false}$ is exactly $2/3$.
\end{theorem}

\begin{definition}[Disjointspecialyfalsemeasure]
\lean{CommunicationComplexity.Functions.Disjointness.RandomizedLowerBound.disjointSpecialYFalseMeasure}
\label{CommunicationComplexity.Functions.Disjointness.RandomizedLowerBound.disjointSpecialYFalseMeasure}
\leanok
\uses{CommunicationComplexity.Functions.Disjointness.RandomizedLowerBound.HardSample, CommunicationComplexity.Functions.Disjointness.RandomizedLowerBound.disjointCondMeasure, CommunicationComplexity.Functions.Disjointness.RandomizedLowerBound.specialY}
The disjoint-conditioned hard distribution, further conditioned on Bob's special bit being \texttt{false}. This is the Alice-side conditioning used in the one-bit Pinsker step.
\end{definition}

\begin{theorem}[Total mass of the doubly-conditioned hard distribution]
\lean{CommunicationComplexity.Functions.Disjointness.RandomizedLowerBound.disjointSpecialYFalseMeasure_isProbabilityMeasure}
\label{CommunicationComplexity.Functions.Disjointness.RandomizedLowerBound.disjointSpecialYFalseMeasure_isProbabilityMeasure}
\leanok
\uses{CommunicationComplexity.Functions.Disjointness.RandomizedLowerBound.disjointCondMeasure, CommunicationComplexity.Functions.Disjointness.RandomizedLowerBound.disjointCondMeasure_measureReal_specialY_false, CommunicationComplexity.Functions.Disjointness.RandomizedLowerBound.disjointSpecialYFalseMeasure, CommunicationComplexity.Functions.Disjointness.RandomizedLowerBound.measure_disjointEvent_ne_zero}
\difficulty{3}
Fix a positive integer $n$, and work on the finite sample space of hard samples on $n$
coordinates, equipped with its ambient counting measure. Let $\mu$ be the measure
obtained by conditioning this counting measure first on the disjointness event—that
Alice's input set and Bob's input set are disjoint—and then on the event that Bob's bit
at the special coordinate equals $\mathrm{false}$. Then $\mu$ is a probability measure;
equivalently, $\mu$ assigns total mass $1$ to the whole sample space.
\end{theorem}

\begin{theorem}[Uniformity of Alice's special bit under disjoint, Bob-false conditioning]
\lean{CommunicationComplexity.Functions.Disjointness.RandomizedLowerBound.disjointSpecialYFalseMeasure_measureReal_specialX_singleton}
\label{CommunicationComplexity.Functions.Disjointness.RandomizedLowerBound.disjointSpecialYFalseMeasure_measureReal_specialX_singleton}
\leanok
\uses{CommunicationComplexity.Functions.Disjointness.RandomizedLowerBound.HardSample, CommunicationComplexity.Functions.Disjointness.RandomizedLowerBound.disjointCondMeasure, CommunicationComplexity.Functions.Disjointness.RandomizedLowerBound.disjointCondMeasure_measureReal_specialBitsEvent, CommunicationComplexity.Functions.Disjointness.RandomizedLowerBound.disjointCondMeasure_measureReal_specialY_false, CommunicationComplexity.Functions.Disjointness.RandomizedLowerBound.disjointCondMeasure_measureReal_specialZeroZero, CommunicationComplexity.Functions.Disjointness.RandomizedLowerBound.disjointSpecialYFalseMeasure, CommunicationComplexity.Functions.Disjointness.RandomizedLowerBound.specialBitsEvent, CommunicationComplexity.Functions.Disjointness.RandomizedLowerBound.specialX, CommunicationComplexity.Functions.Disjointness.RandomizedLowerBound.specialY, CommunicationComplexity.Functions.Disjointness.RandomizedLowerBound.specialZeroZero}
\difficulty{4}
Fix a positive integer $n$ and a Boolean value $b$, and work on the finite sample space
of hard-distribution samples over $n$ coordinates. Form the probability distribution
obtained by taking the uniform hard distribution, conditioning on the disjointness event
that Alice's set $X$ and Bob's set $Y$ are disjoint, and conditioning again on Bob's
special bit $y_T$ taking the value $\mathrm{false}$. Under this distribution, the
probability that Alice's special bit $x_T$ equals $b$ is exactly $\tfrac{1}{2}$.
\end{theorem}

\begin{theorem}[Uniform pushforward of the disjointness-conditioned hard distribution]
\lean{CommunicationComplexity.Functions.Disjointness.RandomizedLowerBound.disjointCondMeasure_measureReal_disjointModel_fiber}
\label{CommunicationComplexity.Functions.Disjointness.RandomizedLowerBound.disjointCondMeasure_measureReal_disjointModel_fiber}
\leanok
\uses{CommunicationComplexity.Functions.Disjointness.RandomizedLowerBound.disjointCondMeasure, CommunicationComplexity.Functions.Disjointness.RandomizedLowerBound.disjointEvent, CommunicationComplexity.Functions.Disjointness.RandomizedLowerBound.disjointModel, CommunicationComplexity.Functions.Disjointness.RandomizedLowerBound.measureReal_disjointEvent, CommunicationComplexity.Functions.Disjointness.RandomizedLowerBound.measureReal_disjointEvent_inter_disjointModel_fiber}
\difficulty{3}
Let $n \ge 1$, and write $\mathcal{C} = \{\text{neither}, \text{left-only},
\text{right-only}\}$ for the three-element type of disjoint coordinates. Consider the
hard distribution on samples conditioned on the event that the generated input pair of
sets is disjoint. For every point $z$ of the product space $\{1,\dots,n\} \times
\bigl(\{1,\dots,n\} \to \mathcal{C}\bigr) \times \mathcal{C}$, the conditional
probability that a sample's disjoint model — namely the triple consisting of its special
coordinate, its generated disjoint coordinate-pair vector, and the ignored value at the
special coordinate — equals $z$ is exactly
\[ \frac{1}{n \cdot 3^{n} \cdot 3}. \]
Since this value is independent of $z$ and equals the reciprocal of the cardinality of
the product space, the disjoint model pushes the conditioned distribution forward to the
uniform distribution on that space.
\end{theorem}

\begin{definition}[Disjointmodelfiberforspecialcoordinatecoordinatevectorequiv]
\lean{CommunicationComplexity.Functions.Disjointness.RandomizedLowerBound.disjointModelFiberForSpecialCoordinateCoordinateVectorEquiv}
\label{CommunicationComplexity.Functions.Disjointness.RandomizedLowerBound.disjointModelFiberForSpecialCoordinateCoordinateVectorEquiv}
\leanok
Technical auxiliary result for \texttt{disjointModelFiberForSpecialCoordinateCoordinateVectorEquiv}.
\end{definition}

\begin{theorem}[Counting triples with a fixed index and coordinate vector]
\lean{CommunicationComplexity.Functions.Disjointness.RandomizedLowerBound.card_disjointModel_fiber_for_specialCoordinate_coordinateVector}
\label{CommunicationComplexity.Functions.Disjointness.RandomizedLowerBound.card_disjointModel_fiber_for_specialCoordinate_coordinateVector}
\leanok
\uses{CommunicationComplexity.Functions.Disjointness.RandomizedLowerBound.DisjointCoordinate.card, CommunicationComplexity.Functions.Disjointness.RandomizedLowerBound.disjointModelFiberForSpecialCoordinateCoordinateVectorEquiv}
\difficulty{1}
Fix a positive integer $n$, an index $i \in \{0, 1, \dots, n-1\}$, and a function
$\mathrm{coords}$ assigning to each index in $\{0, 1, \dots, n-1\}$ one of the three
disjointness labels *neither*, *leftOnly*, or *rightOnly*. Then there are exactly three
triples $(a, f, c)$, where $a$ is an index, $f$ is such a labelling function, and $c$ is
a single disjointness label, subject to the constraints $a = i$ and $f =
\mathrm{coords}$.
\end{theorem}

\begin{theorem}[Uniform joint law of the special coordinate and the coordinate-pair vector]
\lean{CommunicationComplexity.Functions.Disjointness.RandomizedLowerBound.disjointCondMeasure_measureReal_specialCoordinate_disjointCoordinateVector_fiber}
\label{CommunicationComplexity.Functions.Disjointness.RandomizedLowerBound.disjointCondMeasure_measureReal_specialCoordinate_disjointCoordinateVector_fiber}
\leanok
\uses{CommunicationComplexity.FiniteMeasureSpace.measureReal_preimage_eq_sum_fibers, CommunicationComplexity.Functions.Disjointness.RandomizedLowerBound.HardSample, CommunicationComplexity.Functions.Disjointness.RandomizedLowerBound.card_disjointModel_fiber_for_specialCoordinate_coordinateVector, CommunicationComplexity.Functions.Disjointness.RandomizedLowerBound.disjointCondMeasure, CommunicationComplexity.Functions.Disjointness.RandomizedLowerBound.disjointCondMeasure_measureReal_disjointModel_fiber, CommunicationComplexity.Functions.Disjointness.RandomizedLowerBound.disjointCoordinateVector, CommunicationComplexity.Functions.Disjointness.RandomizedLowerBound.disjointModel, CommunicationComplexity.Functions.Disjointness.RandomizedLowerBound.specialCoordinate}
\difficulty{4}
Fix a positive integer $n$, and let $\mathcal{D}$ be the three-element set of disjoint
coordinate types, each coordinate belonging to neither player's set, to Alice's set
only, or to Bob's set only. Under the hard distribution conditioned on the generated
input pair being disjoint, the joint law of the special coordinate together with the
generated disjoint coordinate-pair vector is uniform: for every index $i \in \{1, \dots,
n\}$ and every assignment $c \colon \{1, \dots, n\} \to \mathcal{D}$, the probability
that the special coordinate equals $i$ and the generated coordinate-pair vector equals
$c$ is exactly
\[
\frac{1}{n \cdot 3^{\,n}}.
\]
\end{theorem}

\begin{theorem}[Uniform law on coordinate vectors as a product measure]
\lean{CommunicationComplexity.Functions.Disjointness.RandomizedLowerBound.uniformDisjointCoordinateVector_eq_pi}
\label{CommunicationComplexity.Functions.Disjointness.RandomizedLowerBound.uniformDisjointCoordinateVector_eq_pi}
\leanok
\uses{CommunicationComplexity.Functions.Disjointness.RandomizedLowerBound.uniformDisjointCoordinate, CommunicationComplexity.Functions.Disjointness.RandomizedLowerBound.uniformDisjointCoordinateVector, CommunicationComplexity.Functions.Disjointness.RandomizedLowerBound.uniformDisjointCoordinateVector_singleton, CommunicationComplexity.Functions.Disjointness.RandomizedLowerBound.uniformDisjointCoordinate_singleton}
\difficulty{4}
Fix a positive integer $n$, and consider length-$n$ vectors whose entries are disjoint
coordinates, each drawn from the three-element type $\{\text{neither}, \text{leftOnly},
\text{rightOnly}\}$. Then the uniform probability law on the set of all such vectors
coincides with the product measure of $n$ independent copies of the uniform law on a
single disjoint coordinate.
\end{theorem}

\begin{theorem}[Independence of coordinates under the uniform law]
\lean{CommunicationComplexity.Functions.Disjointness.RandomizedLowerBound.uniformDisjointCoordinateVector_iIndepFun}
\label{CommunicationComplexity.Functions.Disjointness.RandomizedLowerBound.uniformDisjointCoordinateVector_iIndepFun}
\leanok
\uses{CommunicationComplexity.Functions.Disjointness.RandomizedLowerBound.uniformDisjointCoordinateVector, CommunicationComplexity.Functions.Disjointness.RandomizedLowerBound.uniformDisjointCoordinateVector_eq_pi}
\difficulty{2}
Let $n$ be a positive integer, and equip the finite set of functions $x\colon
\{0,1,\dots,n-1\} \to \{\text{neither},\text{leftOnly},\text{rightOnly}\}$ with the
uniform probability law (each of the $3^n$ such functions is equally likely). Then the
coordinate maps $x \mapsto x(i)$, for $i = 0,1,\dots,n-1$, are mutually independent
random variables.
\end{theorem}

\begin{definition}[Coordinatexbit]
\lean{CommunicationComplexity.Functions.Disjointness.RandomizedLowerBound.coordinateXBit}
\label{CommunicationComplexity.Functions.Disjointness.RandomizedLowerBound.coordinateXBit}
\leanok
\uses{CommunicationComplexity.Functions.Disjointness.RandomizedLowerBound.DisjointCoordinate.xBit}
Alice's bit at a fixed coordinate of a generated disjoint coordinate vector.
\end{definition}

\begin{definition}[Coordinateybit]
\lean{CommunicationComplexity.Functions.Disjointness.RandomizedLowerBound.coordinateYBit}
\label{CommunicationComplexity.Functions.Disjointness.RandomizedLowerBound.coordinateYBit}
\leanok
\uses{CommunicationComplexity.Functions.Disjointness.RandomizedLowerBound.DisjointCoordinate.yBit}
Bob's bit at a fixed coordinate of a generated disjoint coordinate vector.
\end{definition}

\begin{definition}[Coordinatexset]
\lean{CommunicationComplexity.Functions.Disjointness.RandomizedLowerBound.coordinateXSet}
\label{CommunicationComplexity.Functions.Disjointness.RandomizedLowerBound.coordinateXSet}
\leanok
\uses{CommunicationComplexity.Functions.Disjointness.RandomizedLowerBound.coordinateXBit}
Alice's input set represented by a generated disjoint coordinate vector.
\end{definition}

\begin{definition}[Coordinateyset]
\lean{CommunicationComplexity.Functions.Disjointness.RandomizedLowerBound.coordinateYSet}
\label{CommunicationComplexity.Functions.Disjointness.RandomizedLowerBound.coordinateYSet}
\leanok
\uses{CommunicationComplexity.Functions.Disjointness.RandomizedLowerBound.coordinateYBit}
Bob's input set represented by a generated disjoint coordinate vector.
\end{definition}

\begin{definition}[Coordinateinput]
\lean{CommunicationComplexity.Functions.Disjointness.RandomizedLowerBound.coordinateInput}
\label{CommunicationComplexity.Functions.Disjointness.RandomizedLowerBound.coordinateInput}
\leanok
\uses{CommunicationComplexity.Functions.Disjointness.RandomizedLowerBound.coordinateXSet, CommunicationComplexity.Functions.Disjointness.RandomizedLowerBound.coordinateYSet}
The input pair represented by a generated disjoint coordinate vector.
\end{definition}

\begin{definition}[Coordinatemessage]
\lean{CommunicationComplexity.Functions.Disjointness.RandomizedLowerBound.coordinateMessage}
\label{CommunicationComplexity.Functions.Disjointness.RandomizedLowerBound.coordinateMessage}
\leanok
\uses{CommunicationComplexity.Deterministic.Protocol.transcript, CommunicationComplexity.Functions.Disjointness.RandomizedLowerBound.ProtocolType, CommunicationComplexity.Functions.Disjointness.RandomizedLowerBound.TranscriptType, CommunicationComplexity.Functions.Disjointness.RandomizedLowerBound.coordinateInput}
The protocol transcript as a function of the generated disjoint coordinate vector.
\end{definition}

\begin{definition}[Coordinatexbefore]
\lean{CommunicationComplexity.Functions.Disjointness.RandomizedLowerBound.coordinateXBefore}
\label{CommunicationComplexity.Functions.Disjointness.RandomizedLowerBound.coordinateXBefore}
\leanok
\uses{CommunicationComplexity.Functions.Disjointness.RandomizedLowerBound.coordinateXBit}
Alice's coordinate-vector bits before a fixed coordinate, padded by \texttt{false} elsewhere.
\end{definition}

\begin{definition}[Coordinateybefore]
\lean{CommunicationComplexity.Functions.Disjointness.RandomizedLowerBound.coordinateYBefore}
\label{CommunicationComplexity.Functions.Disjointness.RandomizedLowerBound.coordinateYBefore}
\leanok
\uses{CommunicationComplexity.Functions.Disjointness.RandomizedLowerBound.coordinateYBit}
Bob's coordinate-vector bits before a fixed coordinate, padded by \texttt{false} elsewhere.
\end{definition}

\begin{definition}[Coordinateyge]
\lean{CommunicationComplexity.Functions.Disjointness.RandomizedLowerBound.coordinateYGe}
\label{CommunicationComplexity.Functions.Disjointness.RandomizedLowerBound.coordinateYGe}
\leanok
\uses{CommunicationComplexity.Functions.Disjointness.RandomizedLowerBound.coordinateYBit}
Bob's coordinate-vector bits from a fixed coordinate onward, padded by \texttt{false} elsewhere.
\end{definition}

\begin{definition}[Coordinatealiceconditioning]
\lean{CommunicationComplexity.Functions.Disjointness.RandomizedLowerBound.coordinateAliceConditioning}
\label{CommunicationComplexity.Functions.Disjointness.RandomizedLowerBound.coordinateAliceConditioning}
\leanok
\uses{CommunicationComplexity.Functions.Disjointness.RandomizedLowerBound.coordinateXBefore, CommunicationComplexity.Functions.Disjointness.RandomizedLowerBound.coordinateYGe}
The coordinate-vector version of Alice's fixed-coordinate conditioning.
\end{definition}

\begin{definition}[Coordinatealicecondbit]
\lean{CommunicationComplexity.Functions.Disjointness.RandomizedLowerBound.coordinateAliceCondBit}
\label{CommunicationComplexity.Functions.Disjointness.RandomizedLowerBound.coordinateAliceCondBit}
\leanok
\uses{CommunicationComplexity.Functions.Disjointness.RandomizedLowerBound.coordinateXBit, CommunicationComplexity.Functions.Disjointness.RandomizedLowerBound.coordinateYBit}
The one-bit-per-coordinate version of Alice's fixed conditioning: use $X_{j}$ before $i$ and $Y_{j}$ from $i$ onward.
\end{definition}

\begin{definition}[Coordinatealicecondbits]
\lean{CommunicationComplexity.Functions.Disjointness.RandomizedLowerBound.coordinateAliceCondBits}
\label{CommunicationComplexity.Functions.Disjointness.RandomizedLowerBound.coordinateAliceCondBits}
\leanok
\uses{CommunicationComplexity.Functions.Disjointness.RandomizedLowerBound.coordinateAliceCondBit}
Alice's fixed conditioning as one bit from each coordinate.
\end{definition}

\begin{definition}[Coordinatealiceresidualbit]
\lean{CommunicationComplexity.Functions.Disjointness.RandomizedLowerBound.coordinateAliceResidualBit}
\label{CommunicationComplexity.Functions.Disjointness.RandomizedLowerBound.coordinateAliceResidualBit}
\leanok
\uses{CommunicationComplexity.Functions.Disjointness.RandomizedLowerBound.coordinateXBit, CommunicationComplexity.Functions.Disjointness.RandomizedLowerBound.coordinateYBit}
The residual coordinate bit left after conditioning on \texttt{coordinateAliceCondBit}. Before $i$ this is Bob's bit; at $i$ it is Alice's bit; after $i$ it is unused.
\end{definition}

\begin{definition}[Coordinatealiceconditioningofcondbits]
\lean{CommunicationComplexity.Functions.Disjointness.RandomizedLowerBound.coordinateAliceConditioningOfCondBits}
\label{CommunicationComplexity.Functions.Disjointness.RandomizedLowerBound.coordinateAliceConditioningOfCondBits}
\leanok
Recode the one-bit-per-coordinate Alice conditioning into the padded pair used in the information term.
\end{definition}

\begin{theorem}[Recoding recovers Alice's coordinate conditioning]
\lean{CommunicationComplexity.Functions.Disjointness.RandomizedLowerBound.coordinateAliceConditioning_eq_recode_condBits}
\label{CommunicationComplexity.Functions.Disjointness.RandomizedLowerBound.coordinateAliceConditioning_eq_recode_condBits}
\leanok
\uses{CommunicationComplexity.Functions.Disjointness.RandomizedLowerBound.coordinateAliceCondBit, CommunicationComplexity.Functions.Disjointness.RandomizedLowerBound.coordinateAliceCondBits, CommunicationComplexity.Functions.Disjointness.RandomizedLowerBound.coordinateAliceConditioning, CommunicationComplexity.Functions.Disjointness.RandomizedLowerBound.coordinateAliceConditioningOfCondBits, CommunicationComplexity.Functions.Disjointness.RandomizedLowerBound.coordinateXBefore, CommunicationComplexity.Functions.Disjointness.RandomizedLowerBound.coordinateYGe}
\difficulty{4}
Fix a positive integer $n$ and a coordinate index $i \in \{0, 1, \dots, n-1\}$, and
consider any assignment $c = (c_0, \dots, c_{n-1})$ of a disjointness value to each
coordinate, where each $c_k$ carries Alice's bit $x_k$ and Bob's bit $y_k$. Then Alice's
fixed-coordinate conditioning at $i$ — the pair whose first component holds $x_j$ for $j
< i$ and $\mathrm{false}$ elsewhere, and whose second component holds $y_j$ for $j \ge
i$ and $\mathrm{false}$ elsewhere — coincides with the result of first forming the
single bit vector $w$ with $w_k = x_k$ for $k < i$ and $w_k = y_k$ for $k \ge i$, and
then recoding $w$ into the padded pair $\bigl(\,j \mapsto w_j \text{ if } j < i \text{
else } \mathrm{false},\ \ j \mapsto w_j \text{ if } i \le j \text{ else }
\mathrm{false}\,\bigr)$. The two constructions agree as functions of the coordinate
assignment.
\end{theorem}

\begin{theorem}[Injectivity of the split-at-a-coordinate encoding]
\lean{CommunicationComplexity.Functions.Disjointness.RandomizedLowerBound.coordinateAliceConditioningOfCondBits_injective}
\label{CommunicationComplexity.Functions.Disjointness.RandomizedLowerBound.coordinateAliceConditioningOfCondBits_injective}
\leanok
\uses{CommunicationComplexity.Functions.Disjointness.RandomizedLowerBound.coordinateAliceConditioningOfCondBits}
\difficulty{4}
Fix $n \ge 1$ and an index $i \in \{0,\dots,n-1\}$, and consider the map that sends a
bit-vector $b \colon \{0,\dots,n-1\} \to \{0,1\}$ to the pair of bit-vectors $(u, v)$
obtained by splitting $b$ at coordinate $i$: namely $u(j) = b(j)$ for $j < i$ and $u(j)
= 0$ otherwise, while $v(j) = b(j)$ for $j \ge i$ and $v(j) = 0$ otherwise. This map is
injective.
\end{theorem}

\begin{theorem}[Alice's conditioning fiber as an intersection of coordinate fibers]
\lean{CommunicationComplexity.Functions.Disjointness.RandomizedLowerBound.coordinateAliceCondBits_fiber_eq_iInter}
\label{CommunicationComplexity.Functions.Disjointness.RandomizedLowerBound.coordinateAliceCondBits_fiber_eq_iInter}
\leanok
\uses{CommunicationComplexity.Functions.Disjointness.RandomizedLowerBound.coordinateAliceCondBit, CommunicationComplexity.Functions.Disjointness.RandomizedLowerBound.coordinateAliceCondBits}
\difficulty{2}
Fix $n \ge 1$, an index $i \in \{0, 1, \dots, n-1\}$, and a target bit-string $b \colon
\{0,1,\dots,n-1\} \to \{\mathrm{false},\mathrm{true}\}$. For a configuration $c$ that
assigns to each coordinate one of the three disjointness states, write $A_i(c) \in
\{\mathrm{false},\mathrm{true}\}^n$ for Alice's fixed conditioning string and
$A_{i,k}(c)$ for its value at coordinate $k$. Then the set of configurations whose
conditioning string equals $b$ exactly is the intersection over all coordinates of the
sets on which each individual conditioning bit is correct:
\[
A_i^{-1}(\{b\}) \;=\; \bigcap_{k=0}^{n-1} A_{i,k}^{-1}(\{b_k\}).
\]
\end{theorem}

\begin{definition}[Coordinatewithcondbit]
\lean{CommunicationComplexity.Functions.Disjointness.RandomizedLowerBound.coordinateWithCondBit}
\label{CommunicationComplexity.Functions.Disjointness.RandomizedLowerBound.coordinateWithCondBit}
\leanok
Technical auxiliary result for \texttt{coordinateWithCondBit}.
\end{definition}

\begin{theorem}[Conditional coordinate recovers its bit]
\lean{CommunicationComplexity.Functions.Disjointness.RandomizedLowerBound.coordinateWithCondBit_spec}
\label{CommunicationComplexity.Functions.Disjointness.RandomizedLowerBound.coordinateWithCondBit_spec}
\leanok
\uses{CommunicationComplexity.Functions.Disjointness.RandomizedLowerBound.DisjointCoordinate.xBit, CommunicationComplexity.Functions.Disjointness.RandomizedLowerBound.DisjointCoordinate.yBit, CommunicationComplexity.Functions.Disjointness.RandomizedLowerBound.coordinateWithCondBit}
\difficulty{2}
Each coordinate of a disjointness instance takes one of three values — belonging to
neither player, to the left player (Alice) only, or to the right player (Bob) only — and
one records Alice's bit, equal to $1$ exactly on the left-only value, and Bob's bit,
equal to $1$ exactly on the right-only value. Fix a positive integer $n$, indices $i, k
\in \{0, 1, \dots, n-1\}$, and a bit $b \in \{0,1\}$, and form the coordinate value that
is left-only if $k < i$ and $b = 1$, right-only if $k \ge i$ and $b = 1$, and neither in
every case where $b = 0$. Then reading Alice's bit of this value when $k < i$, and Bob's
bit of it when $k \ge i$, returns $b$.
\end{theorem}

\begin{theorem}[Positivity of each conditioning-bit event]
\lean{CommunicationComplexity.Functions.Disjointness.RandomizedLowerBound.uniformDisjointCoordinateVector_measure_condBit_ne_zero}
\label{CommunicationComplexity.Functions.Disjointness.RandomizedLowerBound.uniformDisjointCoordinateVector_measure_condBit_ne_zero}
\leanok
\uses{CommunicationComplexity.Functions.Disjointness.RandomizedLowerBound.coordinateAliceCondBit, CommunicationComplexity.Functions.Disjointness.RandomizedLowerBound.coordinateWithCondBit, CommunicationComplexity.Functions.Disjointness.RandomizedLowerBound.coordinateWithCondBit_spec, CommunicationComplexity.Functions.Disjointness.RandomizedLowerBound.coordinateXBit, CommunicationComplexity.Functions.Disjointness.RandomizedLowerBound.coordinateYBit, CommunicationComplexity.Functions.Disjointness.RandomizedLowerBound.uniformDisjointCoordinateVector, CommunicationComplexity.Functions.Disjointness.RandomizedLowerBound.uniformDisjointCoordinateVector_singleton}
\difficulty{4}
Fix a positive integer $n$, and draw a random vector of $n$ disjoint-coordinate symbols
by sampling each of its $n$ coordinates independently and uniformly from the three
values *neither*, *left-only*, and *right-only*. For a coordinate index $k$ write $X_k$
for Alice's bit at that coordinate (equal to $1$ exactly when the $k$-th symbol is
*left-only*) and $Y_k$ for Bob's bit (equal to $1$ exactly when it is *right-only*).
Then for every pair of indices $i,k \in \{0,1,\dots,n-1\}$ and every bit $b \in
\{0,1\}$, the event that Alice's conditioning bit at coordinate $k$ under cut-off $i$ —
namely $X_k$ when $k<i$ and $Y_k$ otherwise — takes the value $b$ has strictly positive
probability under this uniform law.
\end{theorem}

\begin{theorem}[Independence of the residual–conditioning coordinate pairs]
\lean{CommunicationComplexity.Functions.Disjointness.RandomizedLowerBound.uniformDisjointCoordinateVector_residual_cond_iIndepFun}
\label{CommunicationComplexity.Functions.Disjointness.RandomizedLowerBound.uniformDisjointCoordinateVector_residual_cond_iIndepFun}
\leanok
\uses{CommunicationComplexity.Functions.Disjointness.RandomizedLowerBound.DisjointCoordinate.xBit, CommunicationComplexity.Functions.Disjointness.RandomizedLowerBound.DisjointCoordinate.yBit, CommunicationComplexity.Functions.Disjointness.RandomizedLowerBound.coordinateAliceCondBit, CommunicationComplexity.Functions.Disjointness.RandomizedLowerBound.coordinateAliceResidualBit, CommunicationComplexity.Functions.Disjointness.RandomizedLowerBound.coordinateXBit, CommunicationComplexity.Functions.Disjointness.RandomizedLowerBound.coordinateYBit, CommunicationComplexity.Functions.Disjointness.RandomizedLowerBound.uniformDisjointCoordinateVector, CommunicationComplexity.Functions.Disjointness.RandomizedLowerBound.uniformDisjointCoordinateVector_iIndepFun}
\difficulty{3}
Fix a positive integer $n$ and an index $i \in \{0,\dots,n-1\}$, and let the coordinate
vector $c = (c_0,\dots,c_{n-1})$ be drawn from the uniform law on
$\{\text{neither},\text{leftOnly},\text{rightOnly}\}^{n}$. For each coordinate $k$ form
the pair $\bigl(R_i(k,c),\,A_i(k,c)\bigr)$, where $A_i(k,c)$ is Alice's conditioning bit
— the $X$-bit of $c_k$ when $k<i$ and the $Y$-bit of $c_k$ when $k\ge i$ — and
$R_i(k,c)$ is the corresponding residual bit — the $Y$-bit of $c_k$ when $k<i$, the
$X$-bit of $c_k$ when $k=i$, and $\mathrm{false}$ when $k>i$. Then the family of these
$\{0,1\}\times\{0,1\}$-valued pairs, indexed by $k \in \{0,\dots,n-1\}$, is mutually
independent.
\end{theorem}

\begin{theorem}[Conditional independence of Alice's coordinate bit and Bob's earlier bits]
\lean{CommunicationComplexity.Functions.Disjointness.RandomizedLowerBound.uniformDisjointCoordinateVector_indep_coordinateXBit_coordinateYBefore_condBits}
\label{CommunicationComplexity.Functions.Disjointness.RandomizedLowerBound.uniformDisjointCoordinateVector_indep_coordinateXBit_coordinateYBefore_condBits}
\leanok
\uses{CommunicationComplexity.Functions.Disjointness.RandomizedLowerBound.coordinateAliceCondBit, CommunicationComplexity.Functions.Disjointness.RandomizedLowerBound.coordinateAliceCondBits, CommunicationComplexity.Functions.Disjointness.RandomizedLowerBound.coordinateAliceCondBits_fiber_eq_iInter, CommunicationComplexity.Functions.Disjointness.RandomizedLowerBound.coordinateAliceResidualBit, CommunicationComplexity.Functions.Disjointness.RandomizedLowerBound.coordinateXBit, CommunicationComplexity.Functions.Disjointness.RandomizedLowerBound.coordinateYBefore, CommunicationComplexity.Functions.Disjointness.RandomizedLowerBound.coordinateYBit, CommunicationComplexity.Functions.Disjointness.RandomizedLowerBound.uniformDisjointCoordinateVector, CommunicationComplexity.Functions.Disjointness.RandomizedLowerBound.uniformDisjointCoordinateVector_measure_condBit_ne_zero, CommunicationComplexity.Functions.Disjointness.RandomizedLowerBound.uniformDisjointCoordinateVector_residual_cond_iIndepFun}
\difficulty{6}
Fix a positive integer $n$, a coordinate $i \in \{0,1,\dots,n-1\}$, and a bit vector $b
\colon \{0,\dots,n-1\} \to \{0,1\}$. Draw a coordinate vector uniformly from the
$n$-fold product of the three-element set whose members mark a coordinate as belonging
to neither party, to Alice only, or to Bob only; write $X_k = 1$ exactly when coordinate
$k$ belongs to Alice only and $Y_k = 1$ exactly when it belongs to Bob only. Condition
on the event that Alice's conditioning bit vector — whose $k$-th entry is $X_k$ when $k
< i$ and $Y_k$ when $k \ge i$ — equals $b$. Under this conditional law, Alice's bit
$X_i$ at coordinate $i$ and the vector of Bob's earlier bits, equal to $Y_j$ for each $j
< i$ and to $0$ at every $j \ge i$, are independent random variables.
\end{theorem}

\begin{theorem}[Conditional independence of Alice's bit from Bob's earlier bits]
\lean{CommunicationComplexity.Functions.Disjointness.RandomizedLowerBound.uniformDisjointCoordinateVector_crossInfo_condBits_eq_zero}
\label{CommunicationComplexity.Functions.Disjointness.RandomizedLowerBound.uniformDisjointCoordinateVector_crossInfo_condBits_eq_zero}
\leanok
\uses{CommunicationComplexity.Functions.Disjointness.RandomizedLowerBound.coordinateAliceCondBits, CommunicationComplexity.Functions.Disjointness.RandomizedLowerBound.coordinateXBit, CommunicationComplexity.Functions.Disjointness.RandomizedLowerBound.coordinateYBefore, CommunicationComplexity.Functions.Disjointness.RandomizedLowerBound.uniformDisjointCoordinateVector, CommunicationComplexity.Functions.Disjointness.RandomizedLowerBound.uniformDisjointCoordinateVector_indep_coordinateXBit_coordinateYBefore_condBits}
\difficulty{3}
Fix $n \ge 1$ and draw a random vector $C = (C_1, \dots, C_n)$ whose $n$ coordinates are
independent, each uniform over the three states of a disjoint coordinate-pair: one in
which both parties' bits are $0$, one carrying only Alice's bit, and one carrying only
Bob's bit. For each coordinate $k$ let $X_k$ be Alice's bit, equal to $1$ exactly in the
second state, and let $Y_k$ be Bob's bit, equal to $1$ exactly in the third; in
particular $X_k$ and $Y_k$ are never simultaneously $1$. Fix a coordinate $i$ and form
Alice's conditioning vector $A = (A_1, \dots, A_n)$ by setting $A_k = X_k$ for $k < i$
and $A_k = Y_k$ for $k \ge i$. Then the conditional mutual information between $X_i$ and
the tuple $(Y_k)_{k < i}$ of Bob's bits at the coordinates strictly before $i$, given
$A$, vanishes:
\[ I\big(X_i : (Y_k)_{k<i} \;\big|\; A\big) = 0. \]
\end{theorem}

\begin{theorem}[Conditional independence in the uniform disjoint-coordinate distribution]
\lean{CommunicationComplexity.Functions.Disjointness.RandomizedLowerBound.uniformDisjointCoordinateVector_crossInfo_eq_zero}
\label{CommunicationComplexity.Functions.Disjointness.RandomizedLowerBound.uniformDisjointCoordinateVector_crossInfo_eq_zero}
\leanok
\uses{CommunicationComplexity.Functions.Disjointness.RandomizedLowerBound.coordinateAliceCondBits, CommunicationComplexity.Functions.Disjointness.RandomizedLowerBound.coordinateAliceConditioning, CommunicationComplexity.Functions.Disjointness.RandomizedLowerBound.coordinateAliceConditioningOfCondBits, CommunicationComplexity.Functions.Disjointness.RandomizedLowerBound.coordinateAliceConditioningOfCondBits_injective, CommunicationComplexity.Functions.Disjointness.RandomizedLowerBound.coordinateAliceConditioning_eq_recode_condBits, CommunicationComplexity.Functions.Disjointness.RandomizedLowerBound.coordinateXBit, CommunicationComplexity.Functions.Disjointness.RandomizedLowerBound.coordinateYBefore, CommunicationComplexity.Functions.Disjointness.RandomizedLowerBound.uniformDisjointCoordinateVector, CommunicationComplexity.Functions.Disjointness.RandomizedLowerBound.uniformDisjointCoordinateVector_crossInfo_condBits_eq_zero}
\difficulty{2}
Fix a positive integer $n$ and draw $C = (C_1,\dots,C_n)$ uniformly at random from
$\{\text{neither},\text{leftOnly},\text{rightOnly}\}^n$, so that the $n$ coordinates are
independent and each is uniform on the three possibilities. For each index $j$ set
Alice's bit $X_j = \mathbf{1}[C_j = \text{leftOnly}]$ and Bob's bit $Y_j =
\mathbf{1}[C_j = \text{rightOnly}]$. Then for every coordinate $i$, the conditional
mutual information, taken under this uniform law, between Alice's bit at $i$ and Bob's
bits strictly before $i$, given the pair consisting of Alice's bits strictly before $i$
and Bob's bits from $i$ onward, vanishes:
\[
I\big(X_i \,;\, (Y_j)_{j<i} \,\big|\, (X_j)_{j<i},\,(Y_j)_{j\ge i}\big) = 0.
\]
\end{theorem}

\begin{theorem}[Fixed x-bit equals the coordinate-vector x-projection almost surely]
\lean{CommunicationComplexity.Functions.Disjointness.RandomizedLowerBound.fixedXBit_ae_eq_coordinateXBit}
\label{CommunicationComplexity.Functions.Disjointness.RandomizedLowerBound.fixedXBit_ae_eq_coordinateXBit}
\leanok
\uses{CommunicationComplexity.Functions.Disjointness.RandomizedLowerBound.coordinateXBit, CommunicationComplexity.Functions.Disjointness.RandomizedLowerBound.disjointCondMeasure, CommunicationComplexity.Functions.Disjointness.RandomizedLowerBound.disjointCondMeasure_ae_disjointEvent, CommunicationComplexity.Functions.Disjointness.RandomizedLowerBound.disjointCoordinateVector, CommunicationComplexity.Functions.Disjointness.RandomizedLowerBound.disjointCoordinateVector_xBit_of_mem_disjointEvent, CommunicationComplexity.Functions.Disjointness.RandomizedLowerBound.fixedXBit}
\difficulty{2}
Fix a positive integer $n$ and a coordinate $i \in \{1,\dots,n\}$. Under the hard
distribution conditioned on the event that the generated input pair is disjoint, Alice's
bit at coordinate $i$ agrees almost everywhere with the $x$-bit of the $i$-th entry of
the disjoint coordinate-vector generated from the sample; that is, the two functions
$\omega \mapsto \text{(Alice's bit at $i$)}$ and $\omega \mapsto \text{($x$-bit of the
$i$-th generated coordinate)}$ coincide up to a set of measure zero for the conditioned
distribution.
\end{theorem}

\begin{theorem}[Almost-sure agreement of Bob's fixed prefix and coordinate-vector prefix]
\lean{CommunicationComplexity.Functions.Disjointness.RandomizedLowerBound.fixedYBefore_ae_eq_coordinateYBefore}
\label{CommunicationComplexity.Functions.Disjointness.RandomizedLowerBound.fixedYBefore_ae_eq_coordinateYBefore}
\leanok
\uses{CommunicationComplexity.Functions.Disjointness.RandomizedLowerBound.coordinateYBefore, CommunicationComplexity.Functions.Disjointness.RandomizedLowerBound.coordinateYBit, CommunicationComplexity.Functions.Disjointness.RandomizedLowerBound.disjointCondMeasure, CommunicationComplexity.Functions.Disjointness.RandomizedLowerBound.disjointCondMeasure_ae_disjointEvent, CommunicationComplexity.Functions.Disjointness.RandomizedLowerBound.disjointCoordinateVector, CommunicationComplexity.Functions.Disjointness.RandomizedLowerBound.disjointCoordinateVector_yBit_of_mem_disjointEvent, CommunicationComplexity.Functions.Disjointness.RandomizedLowerBound.fixedYBefore}
\difficulty{3}
Fix a coordinate index $i$. Consider the hard distribution on samples $\omega$
conditioned on the event that the generated input pair is disjoint. Then, almost surely
with respect to this conditioned distribution, the two prefix functions on coordinates
agree: the map sending each $j$ to Bob's bit $\mathrm{yBit}(\omega, j)$ when $j < i$ and
to $\mathrm{false}$ otherwise coincides with the map sending each $j$ to Bob's bit at
coordinate $j$ of the generated disjoint coordinate vector of $\omega$ when $j < i$ and
to $\mathrm{false}$ otherwise.
\end{theorem}

\begin{theorem}[Two forms of Alice's conditioning agree almost surely]
\lean{CommunicationComplexity.Functions.Disjointness.RandomizedLowerBound.fixedAliceConditioning_ae_eq_coordinateAliceConditioning}
\label{CommunicationComplexity.Functions.Disjointness.RandomizedLowerBound.fixedAliceConditioning_ae_eq_coordinateAliceConditioning}
\leanok
\uses{CommunicationComplexity.Functions.Disjointness.RandomizedLowerBound.coordinateAliceConditioning, CommunicationComplexity.Functions.Disjointness.RandomizedLowerBound.coordinateXBefore, CommunicationComplexity.Functions.Disjointness.RandomizedLowerBound.coordinateXBit, CommunicationComplexity.Functions.Disjointness.RandomizedLowerBound.coordinateYBit, CommunicationComplexity.Functions.Disjointness.RandomizedLowerBound.coordinateYGe, CommunicationComplexity.Functions.Disjointness.RandomizedLowerBound.disjointCondMeasure, CommunicationComplexity.Functions.Disjointness.RandomizedLowerBound.disjointCondMeasure_ae_disjointEvent, CommunicationComplexity.Functions.Disjointness.RandomizedLowerBound.disjointCoordinateVector, CommunicationComplexity.Functions.Disjointness.RandomizedLowerBound.disjointCoordinateVector_xBit_of_mem_disjointEvent, CommunicationComplexity.Functions.Disjointness.RandomizedLowerBound.disjointCoordinateVector_yBit_of_mem_disjointEvent, CommunicationComplexity.Functions.Disjointness.RandomizedLowerBound.fixedAliceConditioning, CommunicationComplexity.Functions.Disjointness.RandomizedLowerBound.fixedXBefore, CommunicationComplexity.Functions.Disjointness.RandomizedLowerBound.fixedYGe}
\difficulty{3}
Fix a positive integer $n$ and a coordinate $i \in \{0, 1, \dots, n-1\}$, and equip the
sample space of the hard distribution with the measure obtained by conditioning on the
event that the generated pair of input sets is disjoint. There are two ways to record
Alice's conditioning information at coordinate $i$ from a sample $\omega$: directly, as
the pair $\big(X_{<i}(\omega),\, Y_{\ge i}(\omega)\big)$ formed from the sample's own
Alice-prefix and Bob-suffix bits; and indirectly, by first forming the generated
disjoint coordinate-pair vector of $\omega$ and then extracting the coordinate-vector
conditioning $\big(\text{$x$-bits before }i,\, \text{$y$-bits from }i\text{
onward}\big)$ from it. With respect to this conditioned measure, these two functions of
$\omega$ are equal almost everywhere; that is, they agree outside a set of measure zero.
\end{theorem}

\begin{theorem}[Special coordinate and coordinate vector are jointly uniform and independent]
\lean{CommunicationComplexity.Functions.Disjointness.RandomizedLowerBound.identDistrib_specialCoordinate_disjointCoordinateVector_uniform_prod}
\label{CommunicationComplexity.Functions.Disjointness.RandomizedLowerBound.identDistrib_specialCoordinate_disjointCoordinateVector_uniform_prod}
\leanok
\uses{CommunicationComplexity.Functions.Disjointness.RandomizedLowerBound.HardSample, CommunicationComplexity.Functions.Disjointness.RandomizedLowerBound.disjointCondMeasure, CommunicationComplexity.Functions.Disjointness.RandomizedLowerBound.disjointCondMeasure_measureReal_specialCoordinate_disjointCoordinateVector_fiber, CommunicationComplexity.Functions.Disjointness.RandomizedLowerBound.disjointCoordinateVector, CommunicationComplexity.Functions.Disjointness.RandomizedLowerBound.specialCoordinate, CommunicationComplexity.Functions.Disjointness.RandomizedLowerBound.uniformDisjointCoordinateVector, CommunicationComplexity.Functions.Disjointness.RandomizedLowerBound.uniformDisjointCoordinateVector_singleton, CommunicationComplexity.Functions.Disjointness.RandomizedLowerBound.uniformFin, CommunicationComplexity.Functions.Disjointness.RandomizedLowerBound.uniformFin_singleton}
\difficulty{5}
Fix an integer $n \ge 1$, and consider the hard distribution on samples conditioned on
the event that the generated input pair $(X,Y)$ is disjoint. Under this conditional law,
the pair consisting of the special coordinate $T \in \{1,\dots,n\}$ and the generated
disjoint coordinate vector $c \colon \{1,\dots,n\} \to \{\text{neither},
\text{leftOnly}, \text{rightOnly}\}$ has the same distribution as a pair of independent
random variables in which $T$ is uniform on $\{1,\dots,n\}$ and $c$ is uniform on all
$3^{n}$ coordinate vectors. In particular $T$ and $c$ are independent, and each is
uniformly distributed on its range.
\end{theorem}

\begin{theorem}[Second marginal of the product law is the uniform disjoint-coordinate-vector law]
\lean{CommunicationComplexity.Functions.Disjointness.RandomizedLowerBound.identDistrib_uniformProd_snd_uniformDisjointCoordinateVector}
\label{CommunicationComplexity.Functions.Disjointness.RandomizedLowerBound.identDistrib_uniformProd_snd_uniformDisjointCoordinateVector}
\leanok
\uses{CommunicationComplexity.Functions.Disjointness.RandomizedLowerBound.uniformDisjointCoordinateVector, CommunicationComplexity.Functions.Disjointness.RandomizedLowerBound.uniformFin}
\difficulty{2}
Fix a positive integer $n$, and consider pairs $(t, c)$ in which $t$ is an element of
$\{0, 1, \dots, n-1\}$ and $c$ is a disjoint-coordinate vector, that is, a function
assigning to each of the $n$ coordinates one of the three values "neither", "left only",
or "right only". Draw $(t, c)$ from the product of the uniform law on the special
coordinate and the uniform law on disjoint-coordinate vectors. Then the second component
$c$ of this pair is identically distributed to a disjoint-coordinate vector drawn
directly from the uniform law on disjoint-coordinate vectors; equivalently, the second
marginal of the product law coincides with that uniform law.
\end{theorem}

\begin{theorem}[Uniformity of the conditioned disjoint coordinate vector]
\lean{CommunicationComplexity.Functions.Disjointness.RandomizedLowerBound.identDistrib_disjointCoordinateVector_uniform}
\label{CommunicationComplexity.Functions.Disjointness.RandomizedLowerBound.identDistrib_disjointCoordinateVector_uniform}
\leanok
\uses{CommunicationComplexity.Functions.Disjointness.RandomizedLowerBound.disjointCondMeasure, CommunicationComplexity.Functions.Disjointness.RandomizedLowerBound.disjointCoordinateVector, CommunicationComplexity.Functions.Disjointness.RandomizedLowerBound.identDistrib_specialCoordinate_disjointCoordinateVector_uniform_prod, CommunicationComplexity.Functions.Disjointness.RandomizedLowerBound.identDistrib_uniformProd_snd_uniformDisjointCoordinateVector, CommunicationComplexity.Functions.Disjointness.RandomizedLowerBound.uniformDisjointCoordinateVector, CommunicationComplexity.Functions.Disjointness.RandomizedLowerBound.uniformFin}
\difficulty{3}
Fix a positive integer $n$, and consider the hard distribution on samples conditioned on
the disjointness event — the event that the generated input pair $(X,Y)$ is disjoint.
Under this conditioned distribution, the generated disjoint coordinate vector, viewed as
a function $\{1,\dots,n\} \to \{\text{neither}, \text{leftOnly}, \text{rightOnly}\}$, is
identically distributed to a coordinate vector drawn from the uniform law on all such
functions; that is, its law is exactly the uniform distribution over the $3^n$ possible
disjoint coordinate vectors.
\end{theorem}

\begin{theorem}[Independence of coordinate vector from the special coordinate]
\lean{CommunicationComplexity.Functions.Disjointness.RandomizedLowerBound.identDistrib_uniformProd_snd_uniformDisjointCoordinateVector_cond_fst}
\label{CommunicationComplexity.Functions.Disjointness.RandomizedLowerBound.identDistrib_uniformProd_snd_uniformDisjointCoordinateVector_cond_fst}
\leanok
\uses{CommunicationComplexity.Functions.Disjointness.RandomizedLowerBound.uniformDisjointCoordinateVector, CommunicationComplexity.Functions.Disjointness.RandomizedLowerBound.uniformDisjointCoordinateVector_singleton, CommunicationComplexity.Functions.Disjointness.RandomizedLowerBound.uniformDisjointCoordinateVector_univ, CommunicationComplexity.Functions.Disjointness.RandomizedLowerBound.uniformFin, CommunicationComplexity.Functions.Disjointness.RandomizedLowerBound.uniformFin_singleton}
\difficulty{5}
Fix a positive integer $n$ and an index $i \in \mathrm{Fin}\,n$, and let
$\mathrm{DisjointCoordinate} = \{\text{neither}, \text{leftOnly}, \text{rightOnly}\}$ be
the three-element type. On the product space $\mathrm{Fin}\,n \times (\mathrm{Fin}\,n
\to \mathrm{DisjointCoordinate})$ form the product of the uniform law on the special
coordinate and the uniform law on disjoint-coordinate vectors, and condition it on the
event that the first coordinate equals $i$. Then under this conditional law the second
coordinate is identically distributed to a uniformly random disjoint-coordinate vector.
\end{theorem}

\begin{theorem}[Uniform law of Alice's fixed-coordinate information triple]
\lean{CommunicationComplexity.Functions.Disjointness.RandomizedLowerBound.identDistrib_fixedAliceCrossInfoTriple_uniform}
\label{CommunicationComplexity.Functions.Disjointness.RandomizedLowerBound.identDistrib_fixedAliceCrossInfoTriple_uniform}
\leanok
\uses{CommunicationComplexity.Functions.Disjointness.RandomizedLowerBound.HardSample, CommunicationComplexity.Functions.Disjointness.RandomizedLowerBound.coordinateAliceConditioning, CommunicationComplexity.Functions.Disjointness.RandomizedLowerBound.coordinateXBit, CommunicationComplexity.Functions.Disjointness.RandomizedLowerBound.coordinateYBefore, CommunicationComplexity.Functions.Disjointness.RandomizedLowerBound.disjointCondMeasure, CommunicationComplexity.Functions.Disjointness.RandomizedLowerBound.disjointCoordinateVector, CommunicationComplexity.Functions.Disjointness.RandomizedLowerBound.fixedAliceConditioning, CommunicationComplexity.Functions.Disjointness.RandomizedLowerBound.fixedAliceConditioning_ae_eq_coordinateAliceConditioning, CommunicationComplexity.Functions.Disjointness.RandomizedLowerBound.fixedXBit, CommunicationComplexity.Functions.Disjointness.RandomizedLowerBound.fixedXBit_ae_eq_coordinateXBit, CommunicationComplexity.Functions.Disjointness.RandomizedLowerBound.fixedYBefore, CommunicationComplexity.Functions.Disjointness.RandomizedLowerBound.fixedYBefore_ae_eq_coordinateYBefore, CommunicationComplexity.Functions.Disjointness.RandomizedLowerBound.identDistrib_disjointCoordinateVector_uniform, CommunicationComplexity.Functions.Disjointness.RandomizedLowerBound.uniformDisjointCoordinateVector}
\difficulty{4}
Fix a positive integer $n$ and a coordinate $i \in \{0, 1, \dots, n-1\}$. Consider two
ways of producing a triple. First, draw a hard-distribution sample conditioned on the
event that Alice's input set $X$ and Bob's input set $Y$ are disjoint, and record the
triple $\bigl(x_i,\; Y_{<i},\; (X_{<i}, Y_{\ge i})\bigr)$, where $x_i$ is Alice's bit at
coordinate $i$, $Y_{<i}$ is Bob's bit vector on the coordinates $j < i$ (set to
$\mathrm{false}$ elsewhere), $X_{<i}$ is Alice's bit vector on the coordinates $j < i$,
and $Y_{\ge i}$ is Bob's bit vector on the coordinates $j \ge i$ (each padded with
$\mathrm{false}$ outside its range). Second, draw a vector of disjoint coordinates
uniformly at random and form the same triple from Alice's and Bob's bits read off that
vector. Then these two random triples are identically distributed.
\end{theorem}

\begin{theorem}[Uniform law of the conditioned disjoint-coordinate vector]
\lean{CommunicationComplexity.Functions.Disjointness.RandomizedLowerBound.identDistrib_disjointCoordinateVector_uniform_cond_specialCoordinate}
\label{CommunicationComplexity.Functions.Disjointness.RandomizedLowerBound.identDistrib_disjointCoordinateVector_uniform_cond_specialCoordinate}
\leanok
\uses{CommunicationComplexity.Functions.Disjointness.RandomizedLowerBound.HardSample, CommunicationComplexity.Functions.Disjointness.RandomizedLowerBound.disjointCondMeasure, CommunicationComplexity.Functions.Disjointness.RandomizedLowerBound.disjointCoordinateVector, CommunicationComplexity.Functions.Disjointness.RandomizedLowerBound.identDistrib_specialCoordinate_disjointCoordinateVector_uniform_prod, CommunicationComplexity.Functions.Disjointness.RandomizedLowerBound.identDistrib_uniformProd_snd_uniformDisjointCoordinateVector_cond_fst, CommunicationComplexity.Functions.Disjointness.RandomizedLowerBound.specialCoordinate, CommunicationComplexity.Functions.Disjointness.RandomizedLowerBound.uniformDisjointCoordinateVector, CommunicationComplexity.Functions.Disjointness.RandomizedLowerBound.uniformFin, ProbabilityTheory.IdentDistrib.cond_of_pair}
\difficulty{5}
Let $n \ge 1$ and fix a coordinate $i \in \{0, \dots, n-1\}$. Consider the hard
distribution on samples, conditioned both on the event that the generated pair of sets
is disjoint and on the special coordinate being equal to $i$. Then the generated vector
of disjoint coordinates, viewed as a map $\{0,\dots,n-1\} \to \{\text{neither},
\text{leftOnly}, \text{rightOnly}\}$, has the same law as the identity map applied to a
vector drawn uniformly at random from all such maps; that is, under this conditioning
the generated vector is uniformly distributed over the $3^n$ possible coordinate
assignments.
\end{theorem}

\begin{theorem}[Identical distribution under a measure-preserving map]
\lean{CommunicationComplexity.Functions.Disjointness.RandomizedLowerBound.identDistrib_self_comp_measurePreserving}
\label{CommunicationComplexity.Functions.Disjointness.RandomizedLowerBound.identDistrib_self_comp_measurePreserving}
\leanok
\difficulty{3}
Let $(\Omega, \mu)$ and $\alpha$ be measurable spaces, let $T \colon \Omega \to \Omega$
be a measure-preserving map (so that $\mu(T^{-1}(A)) = \mu(A)$ for every measurable set
$A$), and let $F \colon \Omega \to \alpha$ be measurable. Then the random variables $F$
and $F \circ T$ are identically distributed under $\mu$; that is, they induce the same
pushforward measure on $\alpha$.
\end{theorem}

\begin{theorem}[Preimage of a singleton under hard-sample dualization]
\lean{CommunicationComplexity.Functions.Disjointness.RandomizedLowerBound.dualHardSample_preimage_singleton}
\label{CommunicationComplexity.Functions.Disjointness.RandomizedLowerBound.dualHardSample_preimage_singleton}
\leanok
\uses{CommunicationComplexity.Functions.Disjointness.RandomizedLowerBound.HardSample, CommunicationComplexity.Functions.Disjointness.RandomizedLowerBound.dualHardSample, CommunicationComplexity.Functions.Disjointness.RandomizedLowerBound.dualHardSample_dualHardSample}
\difficulty{3}
Fix a positive integer $n$, and let $D$ denote the dualization map on the finite space
of hard samples, which sends a hard sample to the one obtained by exchanging Alice's and
Bob's bits and reversing the order of the coordinates. Then for every hard sample
$\omega$, the preimage of the singleton $\{\omega\}$ under $D$ is exactly the singleton
$\{D(\omega)\}$; that is, $D^{-1}(\{\omega\}) = \{D(\omega)\}$.
\end{theorem}

\begin{theorem}[Dualization preserves singleton volume mass]
\lean{CommunicationComplexity.Functions.Disjointness.RandomizedLowerBound.volume_measureReal_singleton_dualHardSample}
\label{CommunicationComplexity.Functions.Disjointness.RandomizedLowerBound.volume_measureReal_singleton_dualHardSample}
\leanok
\uses{CommunicationComplexity.Functions.Disjointness.RandomizedLowerBound.HardSample, CommunicationComplexity.Functions.Disjointness.RandomizedLowerBound.dualHardSample, CommunicationComplexity.uniformOn_univ_measureReal_eq_card_subtype}
\difficulty{2}
Fix a positive integer $n$, and let $\omega$ be a hard sample for the disjointness lower
bound. Writing $\mathrm{vol}$ for the ambient volume measure on the finite sample space
of hard samples, and $\mathrm{vol}(A)$ for its real-valued mass on a set $A$, the
singleton consisting of the dual sample $\mathrm{dual}(\omega)$ — obtained by swapping
Alice's and Bob's roles and reversing the coordinate order — carries the same mass as
the singleton $\{\omega\}$; that is, $\mathrm{vol}(\{\mathrm{dual}(\omega)\}) =
\mathrm{vol}(\{\omega\})$.
\end{theorem}

\begin{theorem}[Dualization preserves the uniform hard-sample measure]
\lean{CommunicationComplexity.Functions.Disjointness.RandomizedLowerBound.volume_measurePreserving_dualHardSample}
\label{CommunicationComplexity.Functions.Disjointness.RandomizedLowerBound.volume_measurePreserving_dualHardSample}
\leanok
\uses{CommunicationComplexity.Functions.Disjointness.RandomizedLowerBound.dualHardSample, CommunicationComplexity.Functions.Disjointness.RandomizedLowerBound.dualHardSample_preimage_singleton, CommunicationComplexity.Functions.Disjointness.RandomizedLowerBound.volume_measureReal_singleton_dualHardSample}
\difficulty{3}
Fix a positive integer $n$, and consider the finite space of hard samples, each
consisting of a distinguished coordinate $T \in \{0, 1, \dots, n-1\}$, two bits $x_T,
y_T$, and a function assigning to every coordinate one of the three disjointness states
(neither, left-only, or right-only); equip this space with its uniform counting measure
$\mathrm{vol}$. The dualization map — which sends $T$ to its reversed index $n-1-T$,
exchanges the bits $x_T$ and $y_T$, and replaces the state function by the one that, at
each reversed coordinate, swaps the left-only and right-only states while fixing neither
— is measure-preserving for $\mathrm{vol}$: it is measurable, and the pushforward of
$\mathrm{vol}$ along it is again $\mathrm{vol}$.
\end{theorem}

\begin{theorem}[Duality invariance of the disjointness-conditioned hard distribution]
\lean{CommunicationComplexity.Functions.Disjointness.RandomizedLowerBound.disjointCondMeasure_measureReal_singleton_dualHardSample}
\label{CommunicationComplexity.Functions.Disjointness.RandomizedLowerBound.disjointCondMeasure_measureReal_singleton_dualHardSample}
\leanok
\uses{CommunicationComplexity.Functions.Disjointness.RandomizedLowerBound.HardSample, CommunicationComplexity.Functions.Disjointness.RandomizedLowerBound.disjointCondMeasure, CommunicationComplexity.Functions.Disjointness.RandomizedLowerBound.disjointEvent, CommunicationComplexity.Functions.Disjointness.RandomizedLowerBound.disjointEvent_dualHardSample, CommunicationComplexity.Functions.Disjointness.RandomizedLowerBound.dualHardSample, CommunicationComplexity.Functions.Disjointness.RandomizedLowerBound.volume_measureReal_singleton_dualHardSample}
\difficulty{4}
Fix a positive integer $n$, and consider the hard distribution on hard samples
conditioned on the event that the generated pair $(X,Y)$ is disjoint. For every hard
sample $\omega$, this conditioned distribution assigns the same mass to the singleton
$\{\omega^{*}\}$ as to the singleton $\{\omega\}$, where $\omega^{*}$ is the dual sample
obtained by swapping Alice's and Bob's roles and reversing the order of the coordinates.
That is, the conditioned probability of $\omega^{*}$ equals that of $\omega$.
\end{theorem}

\begin{theorem}[Invariance of the disjoint-conditioned hard distribution under dualization]
\lean{CommunicationComplexity.Functions.Disjointness.RandomizedLowerBound.disjointCondMeasure_measurePreserving_dualHardSample}
\label{CommunicationComplexity.Functions.Disjointness.RandomizedLowerBound.disjointCondMeasure_measurePreserving_dualHardSample}
\leanok
\uses{CommunicationComplexity.Functions.Disjointness.RandomizedLowerBound.disjointCondMeasure, CommunicationComplexity.Functions.Disjointness.RandomizedLowerBound.disjointCondMeasure_measureReal_singleton_dualHardSample, CommunicationComplexity.Functions.Disjointness.RandomizedLowerBound.dualHardSample, CommunicationComplexity.Functions.Disjointness.RandomizedLowerBound.dualHardSample_preimage_singleton}
\difficulty{3}
Fix a positive integer $n$, and consider the finite sample space of hard samples $\omega
= (T, x_T, y_T, \text{other})$, where $T \in \{0, \dots, n-1\}$, the bits $x_T, y_T$ lie
in $\{\text{false}, \text{true}\}$, and $\text{other}$ assigns to each coordinate one of
$\{\text{neither}, \text{leftOnly}, \text{rightOnly}\}$. Let the dual map send $\omega$
to the sample whose distinguished coordinate is the reversal $\mathrm{rev}(T) = n-1-T$,
whose distinguished bits are $x_T$ and $y_T$ interchanged, and whose remaining
coordinate is $i \mapsto \mathrm{swap}\big(\text{other}(\mathrm{rev}(i))\big)$; and let
the disjoint-conditioned measure be the uniform distribution on this sample space
conditioned on the event that Alice's and Bob's induced input sets are disjoint. Then
the dual map is measure-preserving for the disjoint-conditioned measure: it is
measurable, and the pushforward of the disjoint-conditioned measure along it is again
the disjoint-conditioned measure.
\end{theorem}

\begin{theorem}[Conditioning commutes with conditioning on disjointness]
\lean{CommunicationComplexity.Functions.Disjointness.RandomizedLowerBound.volume_cond_eq_disjointCondMeasure_cond_of_subset_disjointEvent}
\label{CommunicationComplexity.Functions.Disjointness.RandomizedLowerBound.volume_cond_eq_disjointCondMeasure_cond_of_subset_disjointEvent}
\leanok
\uses{CommunicationComplexity.Functions.Disjointness.RandomizedLowerBound.HardSample, CommunicationComplexity.Functions.Disjointness.RandomizedLowerBound.disjointCondMeasure, CommunicationComplexity.Functions.Disjointness.RandomizedLowerBound.disjointEvent}
\difficulty{3}
Fix $n \ge 1$ coordinates, and write $\mathbb{P}$ for the uniform probability measure on
the finite space of hard-distribution samples. Let $E$ denote the disjointness event —
the set of samples for which Alice's input set and Bob's input set are disjoint — and
let $\mathbb{P}(\,\cdot \mid E)$ be the hard distribution conditioned on disjointness.
Then for every event $A \subseteq E$, conditioning on $A$ produces the same measure
whether one starts from $\mathbb{P}$ or from $\mathbb{P}(\,\cdot \mid E)$; that is,
$\mathbb{P}(\,\cdot \mid A) = \mathbb{P}(\,\cdot \mid E)(\,\cdot \mid A)$.
\end{theorem}

\begin{theorem}[A factor-two bound between two conditionings of the hard distribution]
\lean{CommunicationComplexity.Functions.Disjointness.RandomizedLowerBound.volume_cond_specialZeroZero_measureReal_le_two_mul_disjointSpecialYFalseMeasure}
\label{CommunicationComplexity.Functions.Disjointness.RandomizedLowerBound.volume_cond_specialZeroZero_measureReal_le_two_mul_disjointSpecialYFalseMeasure}
\leanok
\uses{CommunicationComplexity.Functions.Disjointness.RandomizedLowerBound.HardSample, CommunicationComplexity.Functions.Disjointness.RandomizedLowerBound.disjointCondMeasure, CommunicationComplexity.Functions.Disjointness.RandomizedLowerBound.disjointCondMeasure_measureReal_specialY_false, CommunicationComplexity.Functions.Disjointness.RandomizedLowerBound.disjointCondMeasure_measureReal_specialZeroZero, CommunicationComplexity.Functions.Disjointness.RandomizedLowerBound.disjointSpecialYFalseMeasure, CommunicationComplexity.Functions.Disjointness.RandomizedLowerBound.specialBitsEvent, CommunicationComplexity.Functions.Disjointness.RandomizedLowerBound.specialY, CommunicationComplexity.Functions.Disjointness.RandomizedLowerBound.specialZeroZero, CommunicationComplexity.Functions.Disjointness.RandomizedLowerBound.specialZeroZero_subset_disjointEvent, CommunicationComplexity.Functions.Disjointness.RandomizedLowerBound.volume_cond_eq_disjointCondMeasure_cond_of_subset_disjointEvent}
\difficulty{5}
Fix $n\ge 1$ and let $\mu$ be the uniform distribution on the hard-sample space, where
each sample selects a special coordinate $T\in\{1,\dots,n\}$, Alice's and Bob's bits
$x_T,y_T$ at that coordinate, and a disjoint-coordinate type at every other coordinate;
these data determine Alice's set $X$ and Bob's set $Y$. Write $\mu[\,\cdot\mid E\,]$ for
$\mu$ conditioned on an event $E$, and let $\nu$ denote $\mu$ conditioned first on the
pair $(X,Y)$ being disjoint and then further on Bob's special bit satisfying
$y_T=\mathrm{false}$. Then for every set $S$ of samples,
\[
\mu\big(S \mid x_T=\mathrm{false}\ \text{and}\ y_T=\mathrm{false}\big)\ \le\ 2\,\nu(S).
\]
\end{theorem}

\begin{definition}[Message]
\lean{CommunicationComplexity.Functions.Disjointness.RandomizedLowerBound.message}
\label{CommunicationComplexity.Functions.Disjointness.RandomizedLowerBound.message}
\leanok
\uses{CommunicationComplexity.Deterministic.Protocol.transcript, CommunicationComplexity.Functions.Disjointness.RandomizedLowerBound.HardSample, CommunicationComplexity.Functions.Disjointness.RandomizedLowerBound.ProtocolType, CommunicationComplexity.Functions.Disjointness.RandomizedLowerBound.TranscriptType, CommunicationComplexity.Functions.Disjointness.RandomizedLowerBound.input}
The transcript random variable induced by a deterministic protocol on the hard sample space.
\end{definition}

\begin{theorem}[Transcript agrees with its coordinate reconstruction almost surely]
\lean{CommunicationComplexity.Functions.Disjointness.RandomizedLowerBound.message_ae_eq_coordinateMessage}
\label{CommunicationComplexity.Functions.Disjointness.RandomizedLowerBound.message_ae_eq_coordinateMessage}
\leanok
\uses{CommunicationComplexity.Functions.Disjointness.RandomizedLowerBound.ProtocolType, CommunicationComplexity.Functions.Disjointness.RandomizedLowerBound.X, CommunicationComplexity.Functions.Disjointness.RandomizedLowerBound.Y, CommunicationComplexity.Functions.Disjointness.RandomizedLowerBound.coordinateInput, CommunicationComplexity.Functions.Disjointness.RandomizedLowerBound.coordinateMessage, CommunicationComplexity.Functions.Disjointness.RandomizedLowerBound.coordinateXBit, CommunicationComplexity.Functions.Disjointness.RandomizedLowerBound.coordinateXSet, CommunicationComplexity.Functions.Disjointness.RandomizedLowerBound.coordinateYBit, CommunicationComplexity.Functions.Disjointness.RandomizedLowerBound.coordinateYSet, CommunicationComplexity.Functions.Disjointness.RandomizedLowerBound.disjointCondMeasure, CommunicationComplexity.Functions.Disjointness.RandomizedLowerBound.disjointCondMeasure_ae_disjointEvent, CommunicationComplexity.Functions.Disjointness.RandomizedLowerBound.disjointCoordinateVector, CommunicationComplexity.Functions.Disjointness.RandomizedLowerBound.disjointCoordinateVector_xBit_of_mem_disjointEvent, CommunicationComplexity.Functions.Disjointness.RandomizedLowerBound.disjointCoordinateVector_yBit_of_mem_disjointEvent, CommunicationComplexity.Functions.Disjointness.RandomizedLowerBound.input, CommunicationComplexity.Functions.Disjointness.RandomizedLowerBound.message}
\difficulty{3}
Let $n$ be a positive integer, and let $p$ be a deterministic communication protocol
whose inputs are pairs $(A,B)$ of subsets of $\{0,1,\dots,n-1\}$. Then, for almost every
sample $\omega$ under the hard distribution conditioned on the event that its generated
input pair is disjoint, the transcript that $p$ produces when run on the input pair
$(X(\omega),Y(\omega))$ generated directly by $\omega$ equals the transcript that $p$
produces when run on the input pair reconstructed from the disjoint coordinate vector
generated by $\omega$.
\end{theorem}

\begin{theorem}[Transcript entropy bounded by communication cost]
\lean{CommunicationComplexity.Functions.Disjointness.RandomizedLowerBound.entropy_message_le_complexity_mul_log_two_of_measure}
\label{CommunicationComplexity.Functions.Disjointness.RandomizedLowerBound.entropy_message_le_complexity_mul_log_two_of_measure}
\leanok
\uses{CommunicationComplexity.Deterministic.Protocol.card_transcript_le_two_pow_complexity, CommunicationComplexity.Deterministic.Protocol.complexity, CommunicationComplexity.Deterministic.Protocol.transcript, CommunicationComplexity.Functions.Disjointness.RandomizedLowerBound.HardSample, CommunicationComplexity.Functions.Disjointness.RandomizedLowerBound.ProtocolType, CommunicationComplexity.Functions.Disjointness.RandomizedLowerBound.TranscriptType, CommunicationComplexity.Functions.Disjointness.RandomizedLowerBound.message, ProbabilityTheory.entropy_le_nat_mul_log_two_of_card_le_two_pow}
\difficulty{2}
Fix $n \ge 1$, let $p$ be a deterministic two-party protocol for the disjointness
problem on inputs of length $n$, and let $\mu$ be any measure on the hard sample space.
Executing $p$ on the input pair generated by a sample produces a syntactic transcript,
and we regard this transcript as a random variable on the sample space. Then its Shannon
entropy, computed with respect to $\mu$, satisfies
\[
H[\text{transcript}] \;\le\; c(p)\,\log 2,
\]
where $c(p)$ is the communication complexity of $p$, i.e. the worst-case number of bits
exchanged.
\end{theorem}

\begin{theorem}[Duality commutes with protocol execution]
\lean{CommunicationComplexity.Functions.Disjointness.RandomizedLowerBound.message_dualProtocol_dualHardSample}
\label{CommunicationComplexity.Functions.Disjointness.RandomizedLowerBound.message_dualProtocol_dualHardSample}
\leanok
\uses{CommunicationComplexity.Deterministic.Protocol.comap, CommunicationComplexity.Deterministic.Protocol.swap, CommunicationComplexity.Deterministic.Protocol.transcript, CommunicationComplexity.Deterministic.Protocol.transcriptComap_transcript, CommunicationComplexity.Deterministic.Protocol.transcriptSwap_transcript, CommunicationComplexity.Functions.Disjointness.RandomizedLowerBound.HardSample, CommunicationComplexity.Functions.Disjointness.RandomizedLowerBound.ProtocolType, CommunicationComplexity.Functions.Disjointness.RandomizedLowerBound.X, CommunicationComplexity.Functions.Disjointness.RandomizedLowerBound.Y, CommunicationComplexity.Functions.Disjointness.RandomizedLowerBound.dualHardSample, CommunicationComplexity.Functions.Disjointness.RandomizedLowerBound.dualProtocol, CommunicationComplexity.Functions.Disjointness.RandomizedLowerBound.dualProtocolTranscriptMap, CommunicationComplexity.Functions.Disjointness.RandomizedLowerBound.input, CommunicationComplexity.Functions.Disjointness.RandomizedLowerBound.input_dualHardSample, CommunicationComplexity.Functions.Disjointness.RandomizedLowerBound.message, CommunicationComplexity.Functions.Disjointness.RandomizedLowerBound.reverseSet, CommunicationComplexity.Functions.Disjointness.RandomizedLowerBound.reverseSet_reverseSet}
\difficulty{3}
Fix a length $n \ge 1$, let $p$ be a deterministic communication protocol for
disjointness inputs on the coordinate set $\{0,\dots,n-1\}$, and let $\omega$ be a point
of the hard sample space. Running the dual protocol of $p$ — the protocol obtained by
exchanging the roles of Alice and Bob and reversing both coordinate sets — on the dual
sample of $\omega$ — the sample obtained by swapping Alice's and Bob's bits and
reversing coordinate order — produces exactly the transcript obtained from the
transcript of $p$ on the input pair generated by $\omega$ by applying the
transcript-level recoding induced by protocol duality (that is, swapping the two
players' nodes and then re-routing through the coordinate-reversal maps).
\end{theorem}

\begin{definition}[Dualrawzvalue]
\lean{CommunicationComplexity.Functions.Disjointness.RandomizedLowerBound.dualRawZValue}
\label{CommunicationComplexity.Functions.Disjointness.RandomizedLowerBound.dualRawZValue}
\leanok
\uses{CommunicationComplexity.Functions.Disjointness.RandomizedLowerBound.ProtocolType, CommunicationComplexity.Functions.Disjointness.RandomizedLowerBound.RawZType, CommunicationComplexity.Functions.Disjointness.RandomizedLowerBound.dualProtocol, CommunicationComplexity.Functions.Disjointness.RandomizedLowerBound.dualProtocolTranscriptMap, CommunicationComplexity.Functions.Disjointness.RandomizedLowerBound.reverseBoolVector}
Dualize a raw $Z = (M,T,X_{<T},Y_{>T})$ value by recoding the transcript and coarse conditioning data.
\end{definition}

\begin{theorem}[Dualization commutes with the raw Z-variable]
\lean{CommunicationComplexity.Functions.Disjointness.RandomizedLowerBound.rawZVariable_dualProtocol_dualHardSample}
\label{CommunicationComplexity.Functions.Disjointness.RandomizedLowerBound.rawZVariable_dualProtocol_dualHardSample}
\leanok
\uses{CommunicationComplexity.Functions.Disjointness.RandomizedLowerBound.HardSample, CommunicationComplexity.Functions.Disjointness.RandomizedLowerBound.ProtocolType, CommunicationComplexity.Functions.Disjointness.RandomizedLowerBound.RawZType.ext, CommunicationComplexity.Functions.Disjointness.RandomizedLowerBound.dualHardSample, CommunicationComplexity.Functions.Disjointness.RandomizedLowerBound.dualProtocol, CommunicationComplexity.Functions.Disjointness.RandomizedLowerBound.dualRawZValue, CommunicationComplexity.Functions.Disjointness.RandomizedLowerBound.message_dualProtocol_dualHardSample, CommunicationComplexity.Functions.Disjointness.RandomizedLowerBound.rawZVariable, CommunicationComplexity.Functions.Disjointness.RandomizedLowerBound.specialCoordinate, CommunicationComplexity.Functions.Disjointness.RandomizedLowerBound.xBeforeSpecial_dualHardSample, CommunicationComplexity.Functions.Disjointness.RandomizedLowerBound.yAfterSpecial_dualHardSample}
\difficulty{3}
Fix an integer $n \ge 1$, let $p$ be a deterministic communication protocol for
disjointness inputs of length $n$, and let $\omega$ be a sample of the hard
distribution. Then forming the raw $Z$-variable of the dual protocol at the dual sample
yields exactly the dual of the raw $Z$-variable of $p$ at $\omega$; that is,
\[
Z\bigl(\mathrm{dual}(p),\,\mathrm{dual}(\omega)\bigr) \;=\;
\mathrm{dual}\bigl(Z(p,\omega)\bigr),
\]
where $Z(p,\omega)$ is the raw $Z$-variable — recording the transcript of $p$ on the
input pair generated by $\omega$, the special coordinate of $\omega$, Alice's bits
strictly before the special coordinate, and Bob's bits strictly after it — and where
$\mathrm{dual}(\cdot)$ denotes, respectively, the dual protocol, the dual hard sample,
and the dual raw $Z$-value, each obtained by swapping the roles of Alice and Bob and
reversing the coordinate order.
\end{theorem}

\begin{definition}[Dualzvalue]
\lean{CommunicationComplexity.Functions.Disjointness.RandomizedLowerBound.dualZValue}
\label{CommunicationComplexity.Functions.Disjointness.RandomizedLowerBound.dualZValue}
\leanok
\uses{CommunicationComplexity.Functions.Disjointness.RandomizedLowerBound.ProtocolType, CommunicationComplexity.Functions.Disjointness.RandomizedLowerBound.ZType, CommunicationComplexity.Functions.Disjointness.RandomizedLowerBound.ZType.achievable, CommunicationComplexity.Functions.Disjointness.RandomizedLowerBound.ZType.raw, CommunicationComplexity.Functions.Disjointness.RandomizedLowerBound.dualHardSample, CommunicationComplexity.Functions.Disjointness.RandomizedLowerBound.dualProtocol, CommunicationComplexity.Functions.Disjointness.RandomizedLowerBound.dualRawZValue, CommunicationComplexity.Functions.Disjointness.RandomizedLowerBound.rawZVariable_dualProtocol_dualHardSample}
Dualize an achievable $Z = (M,T,X_{<T},Y_{>T})$ value by recoding the transcript and coarse conditioning data.
\end{definition}

\begin{theorem}[Injectivity of raw-$Z$ dualization]
\lean{CommunicationComplexity.Functions.Disjointness.RandomizedLowerBound.dualRawZValue_injective}
\label{CommunicationComplexity.Functions.Disjointness.RandomizedLowerBound.dualRawZValue_injective}
\leanok
\uses{CommunicationComplexity.Functions.Disjointness.RandomizedLowerBound.ProtocolType, CommunicationComplexity.Functions.Disjointness.RandomizedLowerBound.RawZType.ext, CommunicationComplexity.Functions.Disjointness.RandomizedLowerBound.dualProtocolTranscriptMap_injective, CommunicationComplexity.Functions.Disjointness.RandomizedLowerBound.dualRawZValue, CommunicationComplexity.Functions.Disjointness.RandomizedLowerBound.reverseBoolVector_injective}
\difficulty{3}
Fix a positive integer $n$ and a deterministic communication protocol $p$ for
disjointness inputs of length $n$. The dualization map sends a raw $Z = (M, T, X_{<T},
Y_{>T})$ value for $p$ to the corresponding value for the dual protocol — the protocol
obtained from $p$ by exchanging the roles of Alice and Bob and reversing both coordinate
sets — as follows: the transcript $M$ is carried through the transcript-level map
induced by this duality; the special coordinate $T \in \{0, \dots, n-1\}$ is replaced by
its reversal $n - 1 - T$; the new ``before'' vector is the coordinate-reversal of
$Y_{>T}$; and the new ``after'' vector is the coordinate-reversal of $X_{<T}$. This map
is injective.
\end{theorem}

\begin{theorem}[Injectivity of the dual Z-value map]
\lean{CommunicationComplexity.Functions.Disjointness.RandomizedLowerBound.dualZValue_injective}
\label{CommunicationComplexity.Functions.Disjointness.RandomizedLowerBound.dualZValue_injective}
\leanok
\uses{CommunicationComplexity.Functions.Disjointness.RandomizedLowerBound.ProtocolType, CommunicationComplexity.Functions.Disjointness.RandomizedLowerBound.dualRawZValue_injective, CommunicationComplexity.Functions.Disjointness.RandomizedLowerBound.dualZValue, CommunicationComplexity.Functions.Disjointness.RandomizedLowerBound.zVariable_dualProtocol_dualHardSample}
\difficulty{1}
Fix a positive integer $n$ and a deterministic communication protocol $p$ over
disjointness inputs of length $n$. An (achievable) $Z$-value of $p$ is a tuple recording
a transcript of $p$, a special coordinate $t \in \{0,\dots,n-1\}$, and two Boolean
vectors capturing Alice's bits before $t$ and Bob's bits after $t$, subject to the
constraint that the tuple arises from some hard-distribution sample. Consider the
dualization map that sends each such $Z$-value of $p$ to the corresponding $Z$-value of
the dual protocol, by dualizing the transcript (swapping the roles of Alice and Bob and
reversing both coordinate sets), replacing the special coordinate $t$ by its reverse
$n-1-t$, and interchanging the two Boolean vectors while reversing the coordinate order
of each. This map is injective.
\end{theorem}

\begin{theorem}[Z-variable commutes with dualization]
\lean{CommunicationComplexity.Functions.Disjointness.RandomizedLowerBound.zVariable_dualProtocol_dualHardSample}
\label{CommunicationComplexity.Functions.Disjointness.RandomizedLowerBound.zVariable_dualProtocol_dualHardSample}
\leanok
\uses{CommunicationComplexity.Functions.Disjointness.RandomizedLowerBound.HardSample, CommunicationComplexity.Functions.Disjointness.RandomizedLowerBound.ProtocolType, CommunicationComplexity.Functions.Disjointness.RandomizedLowerBound.dualHardSample, CommunicationComplexity.Functions.Disjointness.RandomizedLowerBound.dualProtocol, CommunicationComplexity.Functions.Disjointness.RandomizedLowerBound.dualZValue, CommunicationComplexity.Functions.Disjointness.RandomizedLowerBound.rawZVariable_dualProtocol_dualHardSample, CommunicationComplexity.Functions.Disjointness.RandomizedLowerBound.zVariable}
\difficulty{1}
Let $n$ be a positive integer, let $p$ be a deterministic disjointness protocol on
inputs of length $n$, and let $\omega$ be a sample from the hard distribution. Forming
the $Z$-variable — the packaged record consisting of the transcript of $p$ on the input
generated by $\omega$, together with the special coordinate, Alice's bits strictly
before it, and Bob's bits strictly after it — commutes with dualization: the
$Z$-variable of the dual protocol evaluated at the dualized sample coincides with the
dualization of the $Z$-variable of $p$ at $\omega$. That is,
\[
Z\bigl(p^\ast,\ \omega^\ast\bigr) \;=\; \bigl(Z(p,\omega)\bigr)^\ast,
\]
where $p^\ast$ is the protocol obtained by interchanging Alice and Bob and reversing
both coordinate sets, $\omega^\ast$ is the sample obtained by swapping Alice's and Bob's
bits and reversing the coordinate order, and $(\,\cdot\,)^\ast$ on $Z$-values is the
corresponding recoding.
\end{theorem}

\begin{theorem}[Pullback of the dual $Z$-fiber under sample duality]
\lean{CommunicationComplexity.Functions.Disjointness.RandomizedLowerBound.zVariable_dualProtocol_preimage_dualZValue}
\label{CommunicationComplexity.Functions.Disjointness.RandomizedLowerBound.zVariable_dualProtocol_preimage_dualZValue}
\leanok
\uses{CommunicationComplexity.Functions.Disjointness.RandomizedLowerBound.ProtocolType, CommunicationComplexity.Functions.Disjointness.RandomizedLowerBound.ZType, CommunicationComplexity.Functions.Disjointness.RandomizedLowerBound.dualHardSample, CommunicationComplexity.Functions.Disjointness.RandomizedLowerBound.dualProtocol, CommunicationComplexity.Functions.Disjointness.RandomizedLowerBound.dualZValue, CommunicationComplexity.Functions.Disjointness.RandomizedLowerBound.dualZValue_injective, CommunicationComplexity.Functions.Disjointness.RandomizedLowerBound.zFiber, CommunicationComplexity.Functions.Disjointness.RandomizedLowerBound.zVariable, CommunicationComplexity.Functions.Disjointness.RandomizedLowerBound.zVariable_dualProtocol_dualHardSample}
\difficulty{3}
Fix a positive integer $n$ and a deterministic communication protocol $p$ for
disjointness inputs of length $n$. For a sample $\omega$ of the hard distribution, let
$Z_p(\omega)$ denote its $Z$-statistic: the tuple consisting of the transcript of $p$ on
the input pair generated by $\omega$, the special coordinate of $\omega$, Alice's input
bits strictly before that coordinate, and Bob's input bits strictly after it (each bit
vector padded with $\mathrm{false}$ off its range); and for an achievable value $z$ of
this statistic, write its fiber as $\{\omega : Z_p(\omega) = z\}$. Let $\omega \mapsto
\bar\omega$ be the sample duality, which swaps Alice's and Bob's roles and reverses the
coordinate order, let $p^\dagger$ be the dual protocol, and let $z^\dagger$ be the dual
value obtained from $z$ by dualizing its transcript, reversing its special coordinate,
and interchanging (with coordinate reversal) its before- and after-special bit vectors.
Then a sample $\omega$ satisfies $Z_{p^\dagger}(\bar\omega) = z^\dagger$ if and only if
$Z_p(\omega) = z$; equivalently, the preimage under sample duality of the fiber of
$Z_{p^\dagger}$ over $z^\dagger$ is exactly the fiber of $Z_p$ over $z$.
\end{theorem}

\begin{theorem}[Duality preserves the mass of a Z-fiber]
\lean{CommunicationComplexity.Functions.Disjointness.RandomizedLowerBound.volume_zVariable_dualProtocol_dualZValue}
\label{CommunicationComplexity.Functions.Disjointness.RandomizedLowerBound.volume_zVariable_dualProtocol_dualZValue}
\leanok
\uses{CommunicationComplexity.Functions.Disjointness.RandomizedLowerBound.HardSample, CommunicationComplexity.Functions.Disjointness.RandomizedLowerBound.ProtocolType, CommunicationComplexity.Functions.Disjointness.RandomizedLowerBound.ZType, CommunicationComplexity.Functions.Disjointness.RandomizedLowerBound.dualHardSample, CommunicationComplexity.Functions.Disjointness.RandomizedLowerBound.dualProtocol, CommunicationComplexity.Functions.Disjointness.RandomizedLowerBound.dualZValue, CommunicationComplexity.Functions.Disjointness.RandomizedLowerBound.volume_measurePreserving_dualHardSample, CommunicationComplexity.Functions.Disjointness.RandomizedLowerBound.zFiber, CommunicationComplexity.Functions.Disjointness.RandomizedLowerBound.zVariable_dualProtocol_preimage_dualZValue}
\difficulty{3}
Fix a positive integer $n$, let $p$ be a deterministic communication protocol for
disjointness inputs of length $n$, and let $z$ be an achievable value of the
$Z$-variable of $p$ — that is, a $Z$-value realized by at least one sample of the hard
distribution. Form the dual protocol, obtained by exchanging the roles of Alice and Bob
and reversing the coordinate order of both input sets, together with the dual $Z$-value
induced by $z$. Then, under the ambient measure on the hard-sample space, the fiber of
the dual $Z$-variable over the dual value of $z$ (the set of hard samples whose dual
$Z$-variable equals it) has the same measure as the fiber of the $Z$-variable of $p$
over $z$.
\end{theorem}

\begin{theorem}[Dualization preserves the volume of Z-value fibers]
\lean{CommunicationComplexity.Functions.Disjointness.RandomizedLowerBound.volume_measureReal_zVariable_dualProtocol_dualZValue}
\label{CommunicationComplexity.Functions.Disjointness.RandomizedLowerBound.volume_measureReal_zVariable_dualProtocol_dualZValue}
\leanok
\uses{CommunicationComplexity.Functions.Disjointness.RandomizedLowerBound.ProtocolType, CommunicationComplexity.Functions.Disjointness.RandomizedLowerBound.ZType, CommunicationComplexity.Functions.Disjointness.RandomizedLowerBound.dualProtocol, CommunicationComplexity.Functions.Disjointness.RandomizedLowerBound.dualZValue, CommunicationComplexity.Functions.Disjointness.RandomizedLowerBound.volume_zVariable_dualProtocol_dualZValue, CommunicationComplexity.Functions.Disjointness.RandomizedLowerBound.zFiber}
\difficulty{2}
Fix a positive integer $n$, and let $p$ be a deterministic communication protocol whose
two inputs are subsets of $\{0,1,\dots,n-1\}$ — one held by Alice, one by Bob — and
whose output is a single bit. Each sample $\omega$ of the hard distribution determines
an input pair, and running $p$ on that pair records a transcript; together with the
sample's special coordinate $T$, Alice's input bits at the coordinates strictly below
$T$, and Bob's input bits at the coordinates strictly above $T$, this transcript forms
the value $Z_p(\omega)$. For an achievable value $z$ of this record, let $Z_p^{-1}(z)$
be the set of samples $\omega$ with $Z_p(\omega)=z$. Let $p^{\vee}$ be the protocol
obtained from $p$ by exchanging the roles of Alice and Bob and by reversing the
coordinate order $i \mapsto n-1-i$ in both input subsets. Let $z^{\vee}$ be the value
whose special coordinate is $n-1-T$, whose Alice-bits-before-the-special-coordinate
vector is the coordinate reversal of the Bob-bits-after vector recorded by $z$, whose
Bob-bits-after vector is the coordinate reversal of the Alice-bits-before vector
recorded by $z$, and whose transcript is the transcript recorded by $z$ reinterpreted as
a transcript of $p^{\vee}$ by swapping the two players and re-routing each message
through the coordinate reversal $i \mapsto n-1-i$. Then the volume, taken as a real
number, of $Z_{p^{\vee}}^{-1}(z^{\vee})$ equals the volume, taken as a real number, of
$Z_p^{-1}(z)$.
\end{theorem}

\begin{theorem}[Hard input probability as a preimage measure]
\lean{CommunicationComplexity.Functions.Disjointness.RandomizedLowerBound.inputDist_measureReal_eq_preimage}
\label{CommunicationComplexity.Functions.Disjointness.RandomizedLowerBound.inputDist_measureReal_eq_preimage}
\leanok
\uses{CommunicationComplexity.Functions.Disjointness.RandomizedLowerBound.HardSample, CommunicationComplexity.Functions.Disjointness.RandomizedLowerBound.input, CommunicationComplexity.Functions.Disjointness.RandomizedLowerBound.inputDist, CommunicationComplexity.Functions.Disjointness.RandomizedLowerBound.inputMeasure}
\difficulty{2}
Fix a positive integer $n$, and let $S$ be a set of pairs $(A,B)$ of subsets of
$\{0,1,\dots,n-1\}$. Then the real-valued probability that the pair of input sets
$(X,Y)$ generated by the hard input distribution lies in $S$ equals the real-valued
measure, under the hard-distribution measure on the sample space, of the preimage
$\{\omega : (X(\omega),Y(\omega)) \in S\}$ — that is, the probability assigned to $S$ by
the hard input distribution is computed by pulling $S$ back through the input map to the
explicit hard sample space.
\end{theorem}

\begin{definition}[Protocolerrorevent]
\lean{CommunicationComplexity.Functions.Disjointness.RandomizedLowerBound.protocolErrorEvent}
\label{CommunicationComplexity.Functions.Disjointness.RandomizedLowerBound.protocolErrorEvent}
\leanok
\uses{CommunicationComplexity.Deterministic.Protocol.run, CommunicationComplexity.Functions.Disjointness.RandomizedLowerBound.HardSample, CommunicationComplexity.Functions.Disjointness.RandomizedLowerBound.ProtocolType, CommunicationComplexity.Functions.Disjointness.RandomizedLowerBound.X, CommunicationComplexity.Functions.Disjointness.RandomizedLowerBound.Y, CommunicationComplexity.Functions.Disjointness.disjointness}
The sample-space event where a deterministic protocol errs on disjointness.
\end{definition}

\begin{theorem}[Error rate over inputs equals error probability over samples]
\lean{CommunicationComplexity.Functions.Disjointness.RandomizedLowerBound.distributionalError_inputDist_eq_protocolErrorEvent}
\label{CommunicationComplexity.Functions.Disjointness.RandomizedLowerBound.distributionalError_inputDist_eq_protocolErrorEvent}
\leanok
\uses{CommunicationComplexity.Deterministic.Protocol.distributionalError, CommunicationComplexity.Deterministic.Protocol.run, CommunicationComplexity.Functions.Disjointness.RandomizedLowerBound.ProtocolType, CommunicationComplexity.Functions.Disjointness.RandomizedLowerBound.inputDist, CommunicationComplexity.Functions.Disjointness.RandomizedLowerBound.inputDist_measureReal_eq_preimage, CommunicationComplexity.Functions.Disjointness.RandomizedLowerBound.protocolErrorEvent, CommunicationComplexity.Functions.Disjointness.disjointness}
\difficulty{3}
Let $n$ be a positive integer and let $p$ be a deterministic two-party communication
protocol taking a pair of subsets of $[n]$ to a Boolean value. Then the distributional
error of $p$ against the disjointness function, measured under the hard input
distribution on pairs $(X,Y)$ of subsets of $[n]$, equals the probability, taken over
the hard sample space, of the event that $p$ errs — that is, the total mass of those
samples $\omega$ for which $p$, run on the input pair $(X(\omega), Y(\omega))$ generated
by $\omega$, disagrees with $\mathrm{disjointness}(n, X(\omega), Y(\omega))$.
\end{theorem}

\begin{definition}[Specialpair]
\lean{CommunicationComplexity.Functions.Disjointness.RandomizedLowerBound.specialPair}
\label{CommunicationComplexity.Functions.Disjointness.RandomizedLowerBound.specialPair}
\leanok
\uses{CommunicationComplexity.Functions.Disjointness.RandomizedLowerBound.HardSample, CommunicationComplexity.Functions.Disjointness.RandomizedLowerBound.specialX, CommunicationComplexity.Functions.Disjointness.RandomizedLowerBound.specialY}
The special-coordinate bit-pair.
\end{definition}

\begin{theorem}[Uniformity of Alice's special bit under the disjoint, $y_T$-false conditioning]
\lean{CommunicationComplexity.Functions.Disjointness.RandomizedLowerBound.disjointSpecialYFalseMeasure_specialX_law_eq_uniformBool}
\label{CommunicationComplexity.Functions.Disjointness.RandomizedLowerBound.disjointSpecialYFalseMeasure_specialX_law_eq_uniformBool}
\leanok
\uses{CommunicationComplexity.Functions.Disjointness.RandomizedLowerBound.disjointSpecialYFalseMeasure, CommunicationComplexity.Functions.Disjointness.RandomizedLowerBound.disjointSpecialYFalseMeasure_isProbabilityMeasure, CommunicationComplexity.Functions.Disjointness.RandomizedLowerBound.disjointSpecialYFalseMeasure_measureReal_specialX_singleton, CommunicationComplexity.Functions.Disjointness.RandomizedLowerBound.specialX, CommunicationComplexity.Functions.Disjointness.RandomizedLowerBound.uniformBool, CommunicationComplexity.Functions.Disjointness.RandomizedLowerBound.uniformBool_singleton}
\difficulty{3}
Fix an integer $n \ge 1$, and let the hard distribution be the uniform law on samples
that each specify a special coordinate $T \in \{1,\dots,n\}$, Alice's and Bob's bits
$x_T, y_T$ at that coordinate, and, at every coordinate, a disjoint coordinate-type
recording which of the two players (if either) occupies it. Conditioning the hard
distribution on the event that the induced input sets of Alice and Bob are disjoint and
on Bob's special bit satisfying $y_T = \mathrm{false}$, Alice's special bit $x_T$
becomes uniformly distributed on $\{\mathrm{true}, \mathrm{false}\}$.
\end{theorem}

\begin{theorem}[Full support of the uniform bit distribution]
\lean{CommunicationComplexity.Functions.Disjointness.RandomizedLowerBound.uniformBool_toPMF_ne_zero}
\label{CommunicationComplexity.Functions.Disjointness.RandomizedLowerBound.uniformBool_toPMF_ne_zero}
\leanok
\uses{CommunicationComplexity.Functions.Disjointness.RandomizedLowerBound.uniformBool, CommunicationComplexity.Functions.Disjointness.RandomizedLowerBound.uniformBool_singleton}
\difficulty{2}
Let the uniform law on one bit assign probability $\tfrac{1}{2}$ to each of the two
values in $\{\text{false}, \text{true}\}$. Then for every bit $b$, the probability mass
assigned to $b$ is nonzero.
\end{theorem}

\begin{definition}[Uniformboolpair]
\lean{CommunicationComplexity.Functions.Disjointness.RandomizedLowerBound.uniformBoolPair}
\label{CommunicationComplexity.Functions.Disjointness.RandomizedLowerBound.uniformBoolPair}
\leanok
Uniform law on a bit-pair.
\end{definition}

\begin{theorem}[Uniform measure of a singleton bit-pair]
\lean{CommunicationComplexity.Functions.Disjointness.RandomizedLowerBound.uniformBoolPair_singleton}
\label{CommunicationComplexity.Functions.Disjointness.RandomizedLowerBound.uniformBoolPair_singleton}
\leanok
\uses{CommunicationComplexity.Functions.Disjointness.RandomizedLowerBound.uniformBoolPair}
\difficulty{2}
Let $\mu$ be the uniform probability law on the four-element set $\{0,1\}\times\{0,1\}$
of bit-pairs. Then for every bit-pair $b$, the measure of the singleton $\{b\}$ is
$\mu(\{b\}) = \tfrac{1}{4}$.
\end{theorem}

\begin{theorem}[Uniform bit-pair law as a product]
\lean{CommunicationComplexity.Functions.Disjointness.RandomizedLowerBound.uniformBoolPair_eq_prod}
\label{CommunicationComplexity.Functions.Disjointness.RandomizedLowerBound.uniformBoolPair_eq_prod}
\leanok
\uses{CommunicationComplexity.Functions.Disjointness.RandomizedLowerBound.uniformBool, CommunicationComplexity.Functions.Disjointness.RandomizedLowerBound.uniformBoolPair, CommunicationComplexity.Functions.Disjointness.RandomizedLowerBound.uniformBoolPair_singleton, CommunicationComplexity.Functions.Disjointness.RandomizedLowerBound.uniformBool_singleton}
\difficulty{4}
Write the uniform probability distribution on a single bit $\{0,1\}$, which assigns mass
$\tfrac12$ to each value. Then the uniform probability distribution on the four ordered
bit-pairs $\{0,1\}\times\{0,1\}$, which assigns mass $\tfrac14$ to each pair, is equal
to the product of two independent copies of the single-bit uniform distribution.
\end{theorem}

\begin{definition}[Zfibermeasure]
\lean{CommunicationComplexity.Functions.Disjointness.RandomizedLowerBound.zFiberMeasure}
\label{CommunicationComplexity.Functions.Disjointness.RandomizedLowerBound.zFiberMeasure}
\leanok
\uses{CommunicationComplexity.Functions.Disjointness.RandomizedLowerBound.HardSample, CommunicationComplexity.Functions.Disjointness.RandomizedLowerBound.ProtocolType, CommunicationComplexity.Functions.Disjointness.RandomizedLowerBound.ZType, CommunicationComplexity.Functions.Disjointness.RandomizedLowerBound.zFiber}
The hard-distribution measure conditioned on a $Z=z$ fiber.
\end{definition}

\begin{theorem}[Conditional measure on a $Z$-fiber as a normalized intersection]
\lean{CommunicationComplexity.Functions.Disjointness.RandomizedLowerBound.zFiberMeasure_real_apply}
\label{CommunicationComplexity.Functions.Disjointness.RandomizedLowerBound.zFiberMeasure_real_apply}
\leanok
\uses{CommunicationComplexity.Functions.Disjointness.RandomizedLowerBound.HardSample, CommunicationComplexity.Functions.Disjointness.RandomizedLowerBound.ProtocolType, CommunicationComplexity.Functions.Disjointness.RandomizedLowerBound.ZType, CommunicationComplexity.Functions.Disjointness.RandomizedLowerBound.zFiber, CommunicationComplexity.Functions.Disjointness.RandomizedLowerBound.zFiberMeasure}
\difficulty{2}
Fix a positive integer $n$, a deterministic disjointness protocol $p$ on inputs of
length $n$, and an achievable value $z$ of the variable $Z = (M, T, X_{<T}, Y_{>T})$;
let $F$ denote the fiber of $Z$ above $z$, that is, the set of hard-distribution samples
$\omega$ with $Z(\omega) = z$. Write $\mathrm{vol}$ for the ambient volume measure on
the (finite) sample space of hard-distribution samples, and let $\mathrm{vol}(\cdot \mid
F)$ be the associated measure conditioned on $F$. Then for every set $S$ of
hard-distribution samples, the real number assigned to $S$ by the conditioned measure
satisfies
\[
\mathrm{vol}(S \mid F) \;=\; \mathrm{vol}(F)^{-1}\,\cdot\,\mathrm{vol}(F \cap S),
\]
where the factor $\mathrm{vol}(F)^{-1}$ is the real reciprocal, so that both sides equal
$0$ when $F$ has measure zero.
\end{theorem}

\begin{theorem}[Special-bit conditioning rescales a $Z$-fiber's mass by $3/2$]
\lean{CommunicationComplexity.Functions.Disjointness.RandomizedLowerBound.disjointSpecialYFalseMeasure_measureReal_zVariable}
\label{CommunicationComplexity.Functions.Disjointness.RandomizedLowerBound.disjointSpecialYFalseMeasure_measureReal_zVariable}
\leanok
\uses{CommunicationComplexity.Functions.Disjointness.RandomizedLowerBound.ProtocolType, CommunicationComplexity.Functions.Disjointness.RandomizedLowerBound.ZType, CommunicationComplexity.Functions.Disjointness.RandomizedLowerBound.disjointCondMeasure, CommunicationComplexity.Functions.Disjointness.RandomizedLowerBound.disjointCondMeasure_measureReal_specialY_false, CommunicationComplexity.Functions.Disjointness.RandomizedLowerBound.disjointSpecialYFalseMeasure, CommunicationComplexity.Functions.Disjointness.RandomizedLowerBound.specialY, CommunicationComplexity.Functions.Disjointness.RandomizedLowerBound.zFiber}
\difficulty{2}
Fix a positive length $n$ and a deterministic communication protocol $p$ for
disjointness inputs of length $n$, and let $z$ be an achievable value of the variable
$Z=(M,T,X_{<T},Y_{>T})$ for $p$, with fiber $Z^{-1}(z)$. Write $D$ for the hard
distribution conditioned on the generated input pair being disjoint, and let $D'$ be $D$
conditioned further on the event $\{Y_T=\mathrm{false}\}$ that Bob's bit at the special
coordinate $T$ is $\mathrm{false}$. Then
\[
D'\big(Z^{-1}(z)\big) \;=\; \frac{3}{2}\, D\big(\{Y_T=\mathrm{false}\}\cap
Z^{-1}(z)\big).
\]
\end{theorem}

\begin{theorem}[Iterated conditioning collapses to the joint event]
\lean{CommunicationComplexity.Functions.Disjointness.RandomizedLowerBound.disjointSpecialYFalseMeasure_cond_zVariable_eq_cond_inter}
\label{CommunicationComplexity.Functions.Disjointness.RandomizedLowerBound.disjointSpecialYFalseMeasure_cond_zVariable_eq_cond_inter}
\leanok
\uses{CommunicationComplexity.Functions.Disjointness.RandomizedLowerBound.ProtocolType, CommunicationComplexity.Functions.Disjointness.RandomizedLowerBound.ZType, CommunicationComplexity.Functions.Disjointness.RandomizedLowerBound.disjointCondMeasure, CommunicationComplexity.Functions.Disjointness.RandomizedLowerBound.disjointSpecialYFalseMeasure, CommunicationComplexity.Functions.Disjointness.RandomizedLowerBound.specialY, CommunicationComplexity.Functions.Disjointness.RandomizedLowerBound.zFiber, CommunicationComplexity.Functions.Disjointness.RandomizedLowerBound.zVariable}
\difficulty{2}
Fix an input length $n \ge 1$, a deterministic disjointness protocol $p$ on inputs of
length $n$, and an achievable value $z$ of the summary variable $Z = (M, T, X_{<T},
Y_{>T})$. Let $D$ denote the hard distribution conditioned on the event that the two
generated input sets are disjoint, and write $Y_T$ for Bob's bit at the special
coordinate. Then conditioning $D$ first on the event $\{Y_T = \mathrm{false}\}$ and then
on the fiber $\{Z = z\}$ produces exactly the same measure as conditioning $D$ in a
single step on the intersection $\{Y_T = \mathrm{false}\} \cap \{Z = z\}$.
\end{theorem}

\begin{theorem}[Iterated conditioning on a fiber and Bob's special bit]
\lean{CommunicationComplexity.Functions.Disjointness.RandomizedLowerBound.zFiberMeasure_cond_specialY_eq_volume_cond_inter}
\label{CommunicationComplexity.Functions.Disjointness.RandomizedLowerBound.zFiberMeasure_cond_specialY_eq_volume_cond_inter}
\leanok
\uses{CommunicationComplexity.Functions.Disjointness.RandomizedLowerBound.ProtocolType, CommunicationComplexity.Functions.Disjointness.RandomizedLowerBound.ZType, CommunicationComplexity.Functions.Disjointness.RandomizedLowerBound.specialY, CommunicationComplexity.Functions.Disjointness.RandomizedLowerBound.zFiber, CommunicationComplexity.Functions.Disjointness.RandomizedLowerBound.zFiberMeasure}
\difficulty{2}
Fix a positive integer $n$, a deterministic protocol $p$ for the disjointness problem on
inputs of length $n$, an achievable value $z$ of the coarse variable $Z = (M, T, X_{<T},
Y_{>T})$, and a Boolean $b$. On the sample space of hard samples, write $\mu$ for the
ambient measure, let $F = \{Z = z\}$ be the fiber above $z$, and let $E$ be the event
that Bob's bit at the special coordinate equals $b$. Then conditioning the fiber measure
$\mu[\,\cdot \mid F]$ further on $E$ yields the same measure as conditioning $\mu$
directly on the intersection $F \cap E$:
\[
\bigl(\mu[\,\cdot \mid F]\bigr)[\,\cdot \mid E] \;=\; \mu[\,\cdot \mid F \cap E].
\]
\end{theorem}

\begin{theorem}[Commuting the $Z$-fiber and disjointness conditionings]
\lean{CommunicationComplexity.Functions.Disjointness.RandomizedLowerBound.zFiberMeasure_cond_specialYFalse_eq_disjointSpecialYFalseMeasure_cond_zVariable}
\label{CommunicationComplexity.Functions.Disjointness.RandomizedLowerBound.zFiberMeasure_cond_specialYFalse_eq_disjointSpecialYFalseMeasure_cond_zVariable}
\leanok
\uses{CommunicationComplexity.Functions.Disjointness.RandomizedLowerBound.ProtocolType, CommunicationComplexity.Functions.Disjointness.RandomizedLowerBound.ZType, CommunicationComplexity.Functions.Disjointness.RandomizedLowerBound.disjointSpecialYFalseMeasure, CommunicationComplexity.Functions.Disjointness.RandomizedLowerBound.disjointSpecialYFalseMeasure_cond_zVariable_eq_cond_inter, CommunicationComplexity.Functions.Disjointness.RandomizedLowerBound.mem_disjointEvent_of_specialY_eq_false, CommunicationComplexity.Functions.Disjointness.RandomizedLowerBound.specialY, CommunicationComplexity.Functions.Disjointness.RandomizedLowerBound.volume_cond_eq_disjointCondMeasure_cond_of_subset_disjointEvent, CommunicationComplexity.Functions.Disjointness.RandomizedLowerBound.zFiber, CommunicationComplexity.Functions.Disjointness.RandomizedLowerBound.zFiberMeasure, CommunicationComplexity.Functions.Disjointness.RandomizedLowerBound.zFiberMeasure_cond_specialY_eq_volume_cond_inter, CommunicationComplexity.Functions.Disjointness.RandomizedLowerBound.zVariable}
\difficulty{3}
Fix a length $n \ge 1$ and a deterministic disjointness protocol $p$ on inputs of length
$n$. Let $\mu$ be the uniform distribution on the hard sample space, and let $Z = (M, T,
X_{<T}, Y_{>T})$ be the random variable recording the transcript $M$ of $p$ on the
generated input pair, the special coordinate $T$, Alice's input bits strictly before
$T$, and Bob's input bits strictly after $T$; write $Y_T$ for Bob's bit at the special
coordinate, and call a sample disjoint when Alice's and Bob's generated sets are
disjoint. Then for every achievable value $z$ of $Z$, the measure obtained from $\mu$ by
first conditioning on the fiber $\{Z = z\}$ and then further conditioning on $\{Y_T =
\mathrm{false}\}$ coincides with the measure obtained from $\mu$ by conditioning on
disjointness and on $\{Y_T = \mathrm{false}\}$ and then conditioning on $\{Z = z\}$.
\end{theorem}

\begin{theorem}[Positivity transfers to the uniform fiber measure]
\lean{CommunicationComplexity.Functions.Disjointness.RandomizedLowerBound.zFiberMeasure_specialYFalse_ne_zero_of_disjointSpecialYFalseMeasure_ne_zero}
\label{CommunicationComplexity.Functions.Disjointness.RandomizedLowerBound.zFiberMeasure_specialYFalse_ne_zero_of_disjointSpecialYFalseMeasure_ne_zero}
\leanok
\uses{CommunicationComplexity.Functions.Disjointness.RandomizedLowerBound.HardSample, CommunicationComplexity.Functions.Disjointness.RandomizedLowerBound.ProtocolType, CommunicationComplexity.Functions.Disjointness.RandomizedLowerBound.ZType, CommunicationComplexity.Functions.Disjointness.RandomizedLowerBound.disjointCondMeasure, CommunicationComplexity.Functions.Disjointness.RandomizedLowerBound.disjointSpecialYFalseMeasure, CommunicationComplexity.Functions.Disjointness.RandomizedLowerBound.disjointSpecialYFalseMeasure_measureReal_zVariable, CommunicationComplexity.Functions.Disjointness.RandomizedLowerBound.specialY, CommunicationComplexity.Functions.Disjointness.RandomizedLowerBound.volume_zFiber_ne_zero, CommunicationComplexity.Functions.Disjointness.RandomizedLowerBound.zFiber, CommunicationComplexity.Functions.Disjointness.RandomizedLowerBound.zFiberMeasure}
\difficulty{4}
Fix a positive integer $n$, a deterministic protocol $p$ for disjointness on inputs of
length $n$, and an achievable value $z$ of the variable $Z = (M, T, X_{<T}, Y_{>T})$
recorded during an execution, so that the fiber $\{Z = z\}$ is a subset of the sample
space of hard instances. Let $\mu$ denote the hard distribution conditioned both on the
generated input pair being disjoint and on Bob's bit $Y_T$ at the special coordinate
being $\mathrm{false}$, and let $\nu$ denote the uniform measure on the sample space
conditioned on the fiber $\{Z = z\}$. If the $\mu$-measure of the fiber $\{Z = z\}$ is
nonzero, then the $\nu$-measure of the event $\{Y_T = \mathrm{false}\}$ is nonzero.
\end{theorem}

\begin{definition}[Conditionalspecialpairlaw]
\lean{CommunicationComplexity.Functions.Disjointness.RandomizedLowerBound.conditionalSpecialPairLaw}
\label{CommunicationComplexity.Functions.Disjointness.RandomizedLowerBound.conditionalSpecialPairLaw}
\leanok
\uses{CommunicationComplexity.Functions.Disjointness.RandomizedLowerBound.ProtocolType, CommunicationComplexity.Functions.Disjointness.RandomizedLowerBound.ZType, CommunicationComplexity.Functions.Disjointness.RandomizedLowerBound.specialPair, CommunicationComplexity.Functions.Disjointness.RandomizedLowerBound.zFiberMeasure}
The conditional law of the special bit-pair on a positive-mass $Z=z$ fiber.
\end{definition}

\begin{theorem}[Singleton mass of the conditional special-pair law]
\lean{CommunicationComplexity.Functions.Disjointness.RandomizedLowerBound.conditionalSpecialPairLaw_singleton}
\label{CommunicationComplexity.Functions.Disjointness.RandomizedLowerBound.conditionalSpecialPairLaw_singleton}
\leanok
\uses{CommunicationComplexity.Functions.Disjointness.RandomizedLowerBound.ProtocolType, CommunicationComplexity.Functions.Disjointness.RandomizedLowerBound.ZType, CommunicationComplexity.Functions.Disjointness.RandomizedLowerBound.conditionalSpecialPairLaw, CommunicationComplexity.Functions.Disjointness.RandomizedLowerBound.specialPair, CommunicationComplexity.Functions.Disjointness.RandomizedLowerBound.zFiberMeasure}
\difficulty{2}
Fix a positive length $n$, a deterministic disjointness protocol $p$ on inputs of length
$n$, and an achievable value $z$ of the variable $Z = (M, T, X_{<T}, Y_{>T})$. Let
$\mu_z$ be the conditional distribution of a hard-distribution sample given $Z = z$, and
for a sample $\omega$ write its special pair as $(x_T(\omega), y_T(\omega))$, the two
players' bits at the special coordinate $T = T(\omega)$. Then for every bit-pair $b \in
\{0,1\} \times \{0,1\}$, the conditional law of the special pair — the pushforward of
$\mu_z$ along the special-pair map — assigns to the singleton $\{b\}$ exactly the
$\mu_z$-probability of the event that the special pair equals $b$; that is, its
real-valued mass at $\{b\}$ equals the $\mu_z$-measure of the set of samples $\omega$
with $(x_T(\omega), y_T(\omega)) = b$.
\end{theorem}

\begin{definition}[Conditionalspecialxlaw]
\lean{CommunicationComplexity.Functions.Disjointness.RandomizedLowerBound.conditionalSpecialXLaw}
\label{CommunicationComplexity.Functions.Disjointness.RandomizedLowerBound.conditionalSpecialXLaw}
\leanok
\uses{CommunicationComplexity.Functions.Disjointness.RandomizedLowerBound.ProtocolType, CommunicationComplexity.Functions.Disjointness.RandomizedLowerBound.ZType, CommunicationComplexity.Functions.Disjointness.RandomizedLowerBound.specialX, CommunicationComplexity.Functions.Disjointness.RandomizedLowerBound.zFiberMeasure}
The conditional law of Alice's special bit on a positive-mass $Z=z$ fiber.
\end{definition}

\begin{definition}[Conditionalspecialylaw]
\lean{CommunicationComplexity.Functions.Disjointness.RandomizedLowerBound.conditionalSpecialYLaw}
\label{CommunicationComplexity.Functions.Disjointness.RandomizedLowerBound.conditionalSpecialYLaw}
\leanok
\uses{CommunicationComplexity.Functions.Disjointness.RandomizedLowerBound.ProtocolType, CommunicationComplexity.Functions.Disjointness.RandomizedLowerBound.ZType, CommunicationComplexity.Functions.Disjointness.RandomizedLowerBound.specialY, CommunicationComplexity.Functions.Disjointness.RandomizedLowerBound.zFiberMeasure}
The conditional law of Bob's special bit on a positive-mass $Z=z$ fiber.
\end{definition}

\begin{theorem}[Singleton mass of the conditional special-bit law]
\lean{CommunicationComplexity.Functions.Disjointness.RandomizedLowerBound.conditionalSpecialXLaw_singleton}
\label{CommunicationComplexity.Functions.Disjointness.RandomizedLowerBound.conditionalSpecialXLaw_singleton}
\leanok
\uses{CommunicationComplexity.Functions.Disjointness.RandomizedLowerBound.ProtocolType, CommunicationComplexity.Functions.Disjointness.RandomizedLowerBound.ZType, CommunicationComplexity.Functions.Disjointness.RandomizedLowerBound.conditionalSpecialXLaw, CommunicationComplexity.Functions.Disjointness.RandomizedLowerBound.specialX, CommunicationComplexity.Functions.Disjointness.RandomizedLowerBound.zFiberMeasure}
\difficulty{2}
Fix a positive integer $n$, a deterministic protocol $p$ for length-$n$ set
disjointness, a value $z$ in the range of the variable $Z = (M, T, X_{<T}, Y_{>T})$, and
a bit $b$. Let $\mu_z$ denote the hard input distribution conditioned on the fiber $\{Z
= z\}$, and let $X_T$ denote Alice's bit at the special coordinate $T$. Then the
conditional law of $X_T$ on this fiber — the pushforward of $\mu_z$ along the map $X_T$
— assigns to the singleton $\{b\}$ exactly the $\mu_z$-probability of the event that
Alice's special bit equals $b$; that is, its real-valued mass at $b$ equals
$\mu_z\big(\{\,X_T = b\,\}\big)$.
\end{theorem}

\begin{theorem}[Singleton mass of the conditional special-bit law]
\lean{CommunicationComplexity.Functions.Disjointness.RandomizedLowerBound.conditionalSpecialYLaw_singleton}
\label{CommunicationComplexity.Functions.Disjointness.RandomizedLowerBound.conditionalSpecialYLaw_singleton}
\leanok
\uses{CommunicationComplexity.Functions.Disjointness.RandomizedLowerBound.ProtocolType, CommunicationComplexity.Functions.Disjointness.RandomizedLowerBound.ZType, CommunicationComplexity.Functions.Disjointness.RandomizedLowerBound.conditionalSpecialYLaw, CommunicationComplexity.Functions.Disjointness.RandomizedLowerBound.specialY, CommunicationComplexity.Functions.Disjointness.RandomizedLowerBound.zFiberMeasure}
\difficulty{2}
Let $n$ be a positive integer, let $p$ be a deterministic protocol for disjointness
inputs of length $n$, let $z$ be an achievable value of the variable
$Z=(M,T,X_{<T},Y_{>T})$ for $p$, and let $b$ be a Boolean value. Write $\mu_z$ for the
hard-distribution measure conditioned on the fiber $\{Z=z\}$, and let $S$ be the map
sending a sample to Bob's bit at its special coordinate. Taking the conditional law of
Bob's special bit to be the pushforward $S_*\mu_z$, its real-valued mass on the
singleton $\{b\}$ equals the real-valued mass that $\mu_z$ assigns to the preimage
$S^{-1}(\{b\})$.
\end{theorem}

\begin{definition}[Zdistance]
\lean{CommunicationComplexity.Functions.Disjointness.RandomizedLowerBound.zDistance}
\label{CommunicationComplexity.Functions.Disjointness.RandomizedLowerBound.zDistance}
\leanok
\uses{CommunicationComplexity.Functions.Disjointness.RandomizedLowerBound.ProtocolType, CommunicationComplexity.Functions.Disjointness.RandomizedLowerBound.ZType, CommunicationComplexity.Functions.Disjointness.RandomizedLowerBound.conditionalSpecialPairLaw, CommunicationComplexity.Functions.Disjointness.RandomizedLowerBound.uniformBoolPair, CommunicationComplexity.tvDistance}
The TV distance between the conditional special-pair law on a $Z=z$ fiber and the uniform bit-pair law.
\end{definition}

\begin{definition}[Xdistance]
\lean{CommunicationComplexity.Functions.Disjointness.RandomizedLowerBound.xDistance}
\label{CommunicationComplexity.Functions.Disjointness.RandomizedLowerBound.xDistance}
\leanok
\uses{CommunicationComplexity.Functions.Disjointness.RandomizedLowerBound.ProtocolType, CommunicationComplexity.Functions.Disjointness.RandomizedLowerBound.ZType, CommunicationComplexity.Functions.Disjointness.RandomizedLowerBound.conditionalSpecialXLaw, CommunicationComplexity.Functions.Disjointness.RandomizedLowerBound.uniformBool, CommunicationComplexity.tvDistance}
The TV distance between Alice's conditional special-bit law and a uniform bit.
\end{definition}

\begin{definition}[Ydistance]
\lean{CommunicationComplexity.Functions.Disjointness.RandomizedLowerBound.yDistance}
\label{CommunicationComplexity.Functions.Disjointness.RandomizedLowerBound.yDistance}
\leanok
\uses{CommunicationComplexity.Functions.Disjointness.RandomizedLowerBound.ProtocolType, CommunicationComplexity.Functions.Disjointness.RandomizedLowerBound.ZType, CommunicationComplexity.Functions.Disjointness.RandomizedLowerBound.conditionalSpecialYLaw, CommunicationComplexity.Functions.Disjointness.RandomizedLowerBound.uniformBool, CommunicationComplexity.tvDistance}
The TV distance between Bob's conditional special-bit law and a uniform bit.
\end{definition}

\begin{theorem}[Nonnegativity of the conditional bit distance]
\lean{CommunicationComplexity.Functions.Disjointness.RandomizedLowerBound.xDistance_nonneg}
\label{CommunicationComplexity.Functions.Disjointness.RandomizedLowerBound.xDistance_nonneg}
\leanok
\uses{CommunicationComplexity.Functions.Disjointness.RandomizedLowerBound.ProtocolType, CommunicationComplexity.Functions.Disjointness.RandomizedLowerBound.ZType, CommunicationComplexity.Functions.Disjointness.RandomizedLowerBound.xDistance, CommunicationComplexity.TVDistance.tvDistance_nonneg}
\difficulty{1}
Fix a positive integer $n$, a deterministic communication protocol $p$ for disjointness
inputs of length $n$, and an achievable value $z$ of the random variable $Z = (M, T,
X_{<T}, Y_{>T})$ for $p$. Writing $\mathrm{TV}$ for the total variation distance, the
quantity
\[
\mathrm{TV}\bigl(\mathrm{Law}(X_T \mid Z = z),\; \mathrm{Unif}\{0,1\}\bigr),
\]
the total variation distance between the conditional law of Alice's bit at the special
coordinate given $Z = z$ and the uniform distribution on one bit, is nonnegative.
\end{theorem}

\begin{theorem}[Sample-space duality on Z-fibers and special bits]
\lean{CommunicationComplexity.Functions.Disjointness.RandomizedLowerBound.dualHardSample_preimage_zVariable_dualZValue_inter_specialX}
\label{CommunicationComplexity.Functions.Disjointness.RandomizedLowerBound.dualHardSample_preimage_zVariable_dualZValue_inter_specialX}
\leanok
\uses{CommunicationComplexity.Functions.Disjointness.RandomizedLowerBound.ProtocolType, CommunicationComplexity.Functions.Disjointness.RandomizedLowerBound.ZType, CommunicationComplexity.Functions.Disjointness.RandomizedLowerBound.dualHardSample, CommunicationComplexity.Functions.Disjointness.RandomizedLowerBound.dualProtocol, CommunicationComplexity.Functions.Disjointness.RandomizedLowerBound.dualZValue, CommunicationComplexity.Functions.Disjointness.RandomizedLowerBound.dualZValue_injective, CommunicationComplexity.Functions.Disjointness.RandomizedLowerBound.specialX, CommunicationComplexity.Functions.Disjointness.RandomizedLowerBound.specialX_dualHardSample, CommunicationComplexity.Functions.Disjointness.RandomizedLowerBound.specialY, CommunicationComplexity.Functions.Disjointness.RandomizedLowerBound.zFiber, CommunicationComplexity.Functions.Disjointness.RandomizedLowerBound.zVariable, CommunicationComplexity.Functions.Disjointness.RandomizedLowerBound.zVariable_dualProtocol_dualHardSample}
\difficulty{3}
Fix a positive integer $n$ and a deterministic communication protocol $p$ for the
disjointness problem on inputs of length $n$. For a hard-distribution sample $\omega$,
let $Z_p(\omega)$ be its $Z$-value — the transcript of $p$ on the input pair induced by
$\omega$, together with the special coordinate and the input bits lying strictly before
and strictly after it — and let $\alpha(\omega)$ and $\beta(\omega)$ denote,
respectively, Alice's and Bob's bit at that special coordinate. Let $D$ be the
sample-space duality that swaps the roles of Alice and Bob and reverses the coordinate
order, let $p^{\ast}$ be the dual protocol, and let $z^{\ast}$ be the dual of an
achievable $Z$-value $z$. Then for every achievable $Z$-value $z$ of $p$ and every bit
$b$,
\[
D^{-1}\!\Big(\, Z_{p^{\ast}}^{-1}(\{z^{\ast}\}) \cap \alpha^{-1}(\{b\}) \,\Big)
\;=\;
Z_{p}^{-1}(\{z\}) \cap \beta^{-1}(\{b\}).
\]
Equivalently, a sample $\omega$ has $Z_p(\omega) = z$ with Bob's special bit equal to
$b$ if and only if its dual $D(\omega)$ has dual $Z$-value $z^{\ast}$ with Alice's
special bit equal to $b$.
\end{theorem}

\begin{theorem}[Mass of dual and original special-bit Z-fibers]
\lean{CommunicationComplexity.Functions.Disjointness.RandomizedLowerBound.volume_zVariable_dualProtocol_dualZValue_inter_specialX}
\label{CommunicationComplexity.Functions.Disjointness.RandomizedLowerBound.volume_zVariable_dualProtocol_dualZValue_inter_specialX}
\leanok
\uses{CommunicationComplexity.Functions.Disjointness.RandomizedLowerBound.HardSample, CommunicationComplexity.Functions.Disjointness.RandomizedLowerBound.ProtocolType, CommunicationComplexity.Functions.Disjointness.RandomizedLowerBound.ZType, CommunicationComplexity.Functions.Disjointness.RandomizedLowerBound.dualHardSample, CommunicationComplexity.Functions.Disjointness.RandomizedLowerBound.dualHardSample_preimage_zVariable_dualZValue_inter_specialX, CommunicationComplexity.Functions.Disjointness.RandomizedLowerBound.dualProtocol, CommunicationComplexity.Functions.Disjointness.RandomizedLowerBound.dualZValue, CommunicationComplexity.Functions.Disjointness.RandomizedLowerBound.specialX, CommunicationComplexity.Functions.Disjointness.RandomizedLowerBound.specialY, CommunicationComplexity.Functions.Disjointness.RandomizedLowerBound.volume_measurePreserving_dualHardSample, CommunicationComplexity.Functions.Disjointness.RandomizedLowerBound.zFiber}
\difficulty{3}
Let $n$ be a positive integer and let $p$ be a deterministic communication protocol for
the disjointness inputs of length $n$. For a sample $\omega$ of the hard distribution,
let $Z_p(\omega)$ denote the statistic recording the transcript of $p$ on the input pair
generated by $\omega$, together with the special coordinate $T$, Alice's input bits
strictly before $T$, and Bob's input bits strictly after $T$ (each padded with
$\mathrm{false}$ off its range); write $Z_p^{-1}(z)$ for the set of samples on which
this statistic takes a fixed achievable value $z$. Let $p^{*}$ be the dual protocol,
obtained by exchanging the roles of Alice and Bob and reversing both coordinate sets,
and let $z^{*}$ be the correspondingly dualized value of the statistic for $p^{*}$. Then
for each bit $b$, the set of samples with $Z_{p^{*}}(\omega) = z^{*}$ whose special
coordinate carries Alice's bit equal to $b$ has the same mass, under the ambient measure
on the finite space of hard samples, as the set of samples with $Z_p(\omega) = z$ whose
special coordinate carries Bob's bit equal to $b$.
\end{theorem}

\begin{theorem}[Duality exchanges the special bits on $Z$-fibers]
\lean{CommunicationComplexity.Functions.Disjointness.RandomizedLowerBound.volume_measureReal_zVariable_dualProtocol_dualZValue_inter_specialX}
\label{CommunicationComplexity.Functions.Disjointness.RandomizedLowerBound.volume_measureReal_zVariable_dualProtocol_dualZValue_inter_specialX}
\leanok
\uses{CommunicationComplexity.Functions.Disjointness.RandomizedLowerBound.ProtocolType, CommunicationComplexity.Functions.Disjointness.RandomizedLowerBound.ZType, CommunicationComplexity.Functions.Disjointness.RandomizedLowerBound.dualProtocol, CommunicationComplexity.Functions.Disjointness.RandomizedLowerBound.dualZValue, CommunicationComplexity.Functions.Disjointness.RandomizedLowerBound.specialX, CommunicationComplexity.Functions.Disjointness.RandomizedLowerBound.specialY, CommunicationComplexity.Functions.Disjointness.RandomizedLowerBound.volume_zVariable_dualProtocol_dualZValue_inter_specialX, CommunicationComplexity.Functions.Disjointness.RandomizedLowerBound.zFiber}
\difficulty{2}
Fix a positive integer $n$ and let $p$ be a deterministic two-party protocol for the
length-$n$ disjointness problem, run on the input pair produced by a sample $\omega$
drawn from the hard sample space $\Omega$ (a finite type carrying the volume measure).
To each sample associate its $Z$-value $Z_p(\omega)$, which records the transcript of
$p$ on the input of $\omega$, the special coordinate $T = T(\omega)$, Alice's input bits
strictly before $T$, and Bob's input bits strictly after $T$; for an achievable value
$z$ write $Z_p^{-1}(z) \subseteq \Omega$ for the fiber of samples with $Z_p(\omega) =
z$. Let $p^{\dagger}$ be the dual protocol, obtained by exchanging the roles of Alice
and Bob and reversing the coordinate order, and let $z^{\dagger}$ be the dual $Z$-value
of $z$. Writing $x_T(\omega)$ and $y_T(\omega)$ for Alice's and Bob's bits at the
special coordinate, for every achievable $Z$-value $z$ of $p$ and every $b \in \{0,1\}$
the real-valued volume measure satisfies
\[
\operatorname{vol}\!\big(Z_{p^{\dagger}}^{-1}(z^{\dagger}) \cap \{\omega : x_T(\omega) =
b\}\big)
= \operatorname{vol}\!\big(Z_p^{-1}(z) \cap \{\omega : y_T(\omega) = b\}\big).
\]
\end{theorem}

\begin{theorem}[Special-bit conditional law under protocol duality]
\lean{CommunicationComplexity.Functions.Disjointness.RandomizedLowerBound.conditionalSpecialXLaw_dualProtocol_dualZValue}
\label{CommunicationComplexity.Functions.Disjointness.RandomizedLowerBound.conditionalSpecialXLaw_dualProtocol_dualZValue}
\leanok
\uses{CommunicationComplexity.Functions.Disjointness.RandomizedLowerBound.ProtocolType, CommunicationComplexity.Functions.Disjointness.RandomizedLowerBound.ZType, CommunicationComplexity.Functions.Disjointness.RandomizedLowerBound.conditionalSpecialXLaw, CommunicationComplexity.Functions.Disjointness.RandomizedLowerBound.conditionalSpecialXLaw_singleton, CommunicationComplexity.Functions.Disjointness.RandomizedLowerBound.conditionalSpecialYLaw, CommunicationComplexity.Functions.Disjointness.RandomizedLowerBound.conditionalSpecialYLaw_singleton, CommunicationComplexity.Functions.Disjointness.RandomizedLowerBound.dualProtocol, CommunicationComplexity.Functions.Disjointness.RandomizedLowerBound.dualZValue, CommunicationComplexity.Functions.Disjointness.RandomizedLowerBound.volume_measureReal_zVariable_dualProtocol_dualZValue, CommunicationComplexity.Functions.Disjointness.RandomizedLowerBound.volume_measureReal_zVariable_dualProtocol_dualZValue_inter_specialX, CommunicationComplexity.Functions.Disjointness.RandomizedLowerBound.zFiberMeasure_real_apply}
\difficulty{3}
Fix a positive integer $n$ and consider deterministic communication protocols whose two
players, Alice and Bob, each receive a subset of an $n$-element coordinate set. The hard
distribution is the uniform law on a finite sample space in which a sample fixes a
distinguished coordinate $T$ together with both players' bits at every coordinate, and
hence an input pair. Given a protocol $q$, define the summary of a sample to be the
four-tuple $(\tau, t, a, b)$, where $\tau$ is the transcript $q$ produces on the induced
input pair, $t = T$, the vector $a$ gives Alice's bits at the coordinates strictly below
$T$ and $\mathrm{false}$ at every other coordinate, and the vector $b$ gives Bob's bits
at the coordinates strictly above $T$ and $\mathrm{false}$ at every other coordinate.

Let $p$ be such a protocol and let $z = (\tau, t, a, b)$ be a summary realised by at
least one sample of $p$. Let $p^{\ast}$ be the dual protocol, formed by exchanging the
roles of Alice and Bob and reversing the coordinate order of both input sets via $i
\mapsto n+1-i$. Let $z^{\ast}$ be the summary of $p^{\ast}$ whose special coordinate is
the reversal $n+1-t$, whose Alice-vector is the coordinate reversal of $b$, whose
Bob-vector is the coordinate reversal of $a$, and whose transcript is the image of
$\tau$ under the map on transcripts induced by passing to $p^{\ast}$.

Then the law of Alice's bit at the special coordinate, for the hard distribution
conditioned on the event that the summary under $p^{\ast}$ equals $z^{\ast}$, coincides
with the law of Bob's bit at the special coordinate, for the hard distribution
conditioned on the event that the summary under $p$ equals $z$.
\end{theorem}

\begin{theorem}[Duality exchanges Alice's and Bob's special-bit distances]
\lean{CommunicationComplexity.Functions.Disjointness.RandomizedLowerBound.xDistance_dualProtocol_dualZValue_eq_yDistance}
\label{CommunicationComplexity.Functions.Disjointness.RandomizedLowerBound.xDistance_dualProtocol_dualZValue_eq_yDistance}
\leanok
\uses{CommunicationComplexity.Functions.Disjointness.RandomizedLowerBound.ProtocolType, CommunicationComplexity.Functions.Disjointness.RandomizedLowerBound.ZType, CommunicationComplexity.Functions.Disjointness.RandomizedLowerBound.conditionalSpecialXLaw_dualProtocol_dualZValue, CommunicationComplexity.Functions.Disjointness.RandomizedLowerBound.dualProtocol, CommunicationComplexity.Functions.Disjointness.RandomizedLowerBound.dualZValue, CommunicationComplexity.Functions.Disjointness.RandomizedLowerBound.xDistance, CommunicationComplexity.Functions.Disjointness.RandomizedLowerBound.yDistance}
\difficulty{1}
Fix a positive integer $n$, a deterministic communication protocol $p$ for disjointness
inputs of length $n$, and an achievable value $z$ of the summary variable $Z$ associated
with $p$. For any such pair, let the *special-bit distance from uniform* be the total
variation distance between the conditional law of the relevant party's bit at the
special coordinate — given that the summary variable equals $z$ — and the uniform law on
a single bit; write it $d_X$ when the party is Alice and $d_Y$ when the party is Bob.
Passing to the dual protocol $p^{*}$, obtained by interchanging the roles of Alice and
Bob and reversing the coordinate order of both input sets, and to the correspondingly
dualized value $z^{*}$ of $Z$, one has \[ d_X\bigl(p^{*},\,z^{*}\bigr) =
d_Y\bigl(p,\,z\bigr). \] That is, Alice's special-bit distance from uniform in the dual
system equals Bob's special-bit distance from uniform in the original.
\end{theorem}

\begin{theorem}[Duality exchanges the Alice and Bob bias events on the zero–zero slice]
\lean{CommunicationComplexity.Functions.Disjointness.RandomizedLowerBound.volume_specialZeroZero_inter_xDistance_dualProtocol_eq_yDistance}
\label{CommunicationComplexity.Functions.Disjointness.RandomizedLowerBound.volume_specialZeroZero_inter_xDistance_dualProtocol_eq_yDistance}
\leanok
\uses{CommunicationComplexity.Functions.Disjointness.RandomizedLowerBound.HardSample, CommunicationComplexity.Functions.Disjointness.RandomizedLowerBound.ProtocolType, CommunicationComplexity.Functions.Disjointness.RandomizedLowerBound.dualHardSample, CommunicationComplexity.Functions.Disjointness.RandomizedLowerBound.dualProtocol, CommunicationComplexity.Functions.Disjointness.RandomizedLowerBound.specialZeroZero, CommunicationComplexity.Functions.Disjointness.RandomizedLowerBound.specialZeroZero_dualHardSample, CommunicationComplexity.Functions.Disjointness.RandomizedLowerBound.volume_measurePreserving_dualHardSample, CommunicationComplexity.Functions.Disjointness.RandomizedLowerBound.xDistance, CommunicationComplexity.Functions.Disjointness.RandomizedLowerBound.xDistance_dualProtocol_dualZValue_eq_yDistance, CommunicationComplexity.Functions.Disjointness.RandomizedLowerBound.yDistance, CommunicationComplexity.Functions.Disjointness.RandomizedLowerBound.zVariable, CommunicationComplexity.Functions.Disjointness.RandomizedLowerBound.zVariable_dualProtocol_dualHardSample}
\difficulty{4}
Fix a positive integer $n$, a deterministic communication protocol $p$ for the
disjointness problem on inputs of length $n$, and a real number $\gamma$. For a protocol
$q$ and a hard sample $\omega$, let $Z_q(\omega)$ denote the associated statistic (the
transcript of $q$ on the induced input together with the special coordinate and the
padded bit vectors lying before and after it), write $d^X_q(z)$ for the total variation
distance between the conditional law of Alice's special-coordinate bit given $Z_q = z$
and the uniform law on one bit, and write $d^Y_q(z)$ for the same distance computed from
Bob's special-coordinate bit. Let $p^\ast$ be the dual protocol, obtained by exchanging
the roles of Alice and Bob and reversing both coordinate sets. Then, on the event that
both special-coordinate bits are $0$, the volume of the set of samples $\omega$ with
$\gamma < d^X_{p^\ast}\big(Z_{p^\ast}(\omega)\big)$ equals the volume of the set of
samples $\omega$ with $\gamma < d^Y_{p}\big(Z_{p}(\omega)\big)$, where the volume is
taken in the ambient measure on the finite hard sample space and read as a real number.
\end{theorem}

\begin{theorem}[Pinsker bound for the conditional special-bit law]
\lean{CommunicationComplexity.Functions.Disjointness.RandomizedLowerBound.two_mul_xDistance_sq_le_toReal_klDiv_uniformBool}
\label{CommunicationComplexity.Functions.Disjointness.RandomizedLowerBound.two_mul_xDistance_sq_le_toReal_klDiv_uniformBool}
\leanok
\uses{CommunicationComplexity.FiniteMeasureSpace.klDiv_ne_top_of_forall_toPMF_ne_zero, CommunicationComplexity.Functions.Disjointness.RandomizedLowerBound.ProtocolType, CommunicationComplexity.Functions.Disjointness.RandomizedLowerBound.ZType, CommunicationComplexity.Functions.Disjointness.RandomizedLowerBound.conditionalSpecialXLaw, CommunicationComplexity.Functions.Disjointness.RandomizedLowerBound.uniformBool, CommunicationComplexity.Functions.Disjointness.RandomizedLowerBound.uniformBool_toPMF_ne_zero, CommunicationComplexity.Functions.Disjointness.RandomizedLowerBound.xDistance, CommunicationComplexity.two_mul_tvDistance_sq_le_toReal_klDiv}
\difficulty{2}
Let $n$ be a positive integer, let $p$ be a deterministic communication protocol for the
disjointness problem on inputs of length $n$, and let $z$ be an achievable value of the
auxiliary variable $Z = (M, T, X_{<T}, Y_{>T})$ for $p$. Write $\mu_z$ for the
conditional law of Alice's bit at the special coordinate given $Z = z$, and let
$\upsilon$ denote the uniform distribution on a single bit. Then, taking the total
variation distance $\mathrm{TV}(\mu_z, \upsilon)$ and the Kullback–Leibler divergence
$D_{\mathrm{KL}}(\mu_z \,\|\, \upsilon)$ (as an extended real, read off as a real
number), the two satisfy
\[
2\,\mathrm{TV}(\mu_z, \upsilon)^2 \;\le\; D_{\mathrm{KL}}(\mu_z \,\|\, \upsilon).
\]
\end{theorem}

\begin{definition}[Xfiberkl]
\lean{CommunicationComplexity.Functions.Disjointness.RandomizedLowerBound.xFiberKL}
\label{CommunicationComplexity.Functions.Disjointness.RandomizedLowerBound.xFiberKL}
\leanok
\uses{CommunicationComplexity.Functions.Disjointness.RandomizedLowerBound.ProtocolType, CommunicationComplexity.Functions.Disjointness.RandomizedLowerBound.ZType, CommunicationComplexity.Functions.Disjointness.RandomizedLowerBound.conditionalSpecialXLaw, CommunicationComplexity.Functions.Disjointness.RandomizedLowerBound.uniformBool}
Alice's one-bit fiber KL cost on a $Z=z$ fiber.
\end{definition}

\begin{theorem}[Pinsker's inequality for the conditional law of Alice's special bit]
\lean{CommunicationComplexity.Functions.Disjointness.RandomizedLowerBound.two_mul_xDistance_sq_le_xFiberKL}
\label{CommunicationComplexity.Functions.Disjointness.RandomizedLowerBound.two_mul_xDistance_sq_le_xFiberKL}
\leanok
\uses{CommunicationComplexity.Functions.Disjointness.RandomizedLowerBound.ProtocolType, CommunicationComplexity.Functions.Disjointness.RandomizedLowerBound.ZType, CommunicationComplexity.Functions.Disjointness.RandomizedLowerBound.two_mul_xDistance_sq_le_toReal_klDiv_uniformBool, CommunicationComplexity.Functions.Disjointness.RandomizedLowerBound.xDistance, CommunicationComplexity.Functions.Disjointness.RandomizedLowerBound.xFiberKL}
\difficulty{1}
Fix $n \ge 1$ and a deterministic two-party communication protocol $p$ for disjointness
inputs of length $n$, and let $z$ be an achievable value of the
transcript-and-side-information variable $Z = (M, T, X_{<T}, Y_{>T})$. Write $\mu_z$ for
the conditional law of Alice's bit $X_T$ at the special coordinate given that $Z = z$ —
that is, the pushforward under the map $\omega \mapsto X_T(\omega)$ of the uniform
measure on the sample space conditioned on the fiber $\{Z = z\}$ — and write $\nu$ for
the uniform law on a single bit. Then twice the squared total variation distance between
$\mu_z$ and $\nu$ is at most their Kullback–Leibler divergence,
\[
2\,\mathrm{TV}(\mu_z, \nu)^2 \le D_{\mathrm{KL}}(\mu_z \,\|\, \nu).
\]
\end{theorem}

\begin{theorem}[Integrated Pinsker bound for Alice's special bit]
\lean{CommunicationComplexity.Functions.Disjointness.RandomizedLowerBound.two_mul_integral_xDistance_sq_le_integral_xFiberKL_disjointSpecialYFalse}
\label{CommunicationComplexity.Functions.Disjointness.RandomizedLowerBound.two_mul_integral_xDistance_sq_le_integral_xFiberKL_disjointSpecialYFalse}
\leanok
\uses{CommunicationComplexity.Functions.Disjointness.RandomizedLowerBound.HardSample, CommunicationComplexity.Functions.Disjointness.RandomizedLowerBound.ProtocolType, CommunicationComplexity.Functions.Disjointness.RandomizedLowerBound.disjointSpecialYFalseMeasure, CommunicationComplexity.Functions.Disjointness.RandomizedLowerBound.disjointSpecialYFalseMeasure_isProbabilityMeasure, CommunicationComplexity.Functions.Disjointness.RandomizedLowerBound.two_mul_xDistance_sq_le_xFiberKL, CommunicationComplexity.Functions.Disjointness.RandomizedLowerBound.xDistance, CommunicationComplexity.Functions.Disjointness.RandomizedLowerBound.xFiberKL, CommunicationComplexity.Functions.Disjointness.RandomizedLowerBound.zVariable}
\difficulty{3}
Fix a positive integer $n$ and a deterministic communication protocol $p$ for the
disjointness problem on inputs that are subsets of an $n$-element set. For a hard sample
$\omega$, let $Z(\omega)$ be the record consisting of the transcript produced by running
$p$ on the induced input pair, the special coordinate $T$, Alice's input bits at the
coordinates strictly below $T$ (padded by $\mathrm{false}$ elsewhere), and Bob's input
bits at the coordinates strictly above $T$. Write $\nu_{Z(\omega)}$ for the conditional
law of Alice's special bit $X_T$ given the value of $Z(\omega)$ — the pushforward, under
the ambient uniform law on hard samples restricted to the fibre $\{Z = Z(\omega)\}$, of
the map $\omega \mapsto X_T$ — and let $U$ denote the uniform distribution on a single
bit. Let $\mu$ be the hard distribution conditioned both on the induced input pair being
disjoint and on Bob's bit at the special coordinate being $\mathrm{false}$. Then
\[
2\int \mathrm{TV}\!\bigl(\nu_{Z(\omega)},\, U\bigr)^{2}\, d\mu(\omega)\ \le\ \int
D_{\mathrm{KL}}\!\bigl(\nu_{Z(\omega)}\,\big\|\,U\bigr)\, d\mu(\omega),
\]
where $\mathrm{TV}$ is the total variation distance and $D_{\mathrm{KL}}$ the
Kullback–Leibler divergence.
\end{theorem}

\begin{theorem}[Linear average-distance bound from a squared bound]
\lean{CommunicationComplexity.Functions.Disjointness.RandomizedLowerBound.disjointSpecialYFalseMeasure_integral_xDistance_le_of_integral_sq_le}
\label{CommunicationComplexity.Functions.Disjointness.RandomizedLowerBound.disjointSpecialYFalseMeasure_integral_xDistance_le_of_integral_sq_le}
\leanok
\uses{CommunicationComplexity.Functions.Disjointness.RandomizedLowerBound.HardSample, CommunicationComplexity.Functions.Disjointness.RandomizedLowerBound.ProtocolType, CommunicationComplexity.Functions.Disjointness.RandomizedLowerBound.disjointSpecialYFalseMeasure, CommunicationComplexity.Functions.Disjointness.RandomizedLowerBound.disjointSpecialYFalseMeasure_isProbabilityMeasure, CommunicationComplexity.Functions.Disjointness.RandomizedLowerBound.fixedAliceFullYInfoTerm, CommunicationComplexity.Functions.Disjointness.RandomizedLowerBound.xDistance, CommunicationComplexity.Functions.Disjointness.RandomizedLowerBound.xDistance_nonneg, CommunicationComplexity.Functions.Disjointness.RandomizedLowerBound.zVariable}
\difficulty{5}
Fix $n \ge 1$, a deterministic communication protocol $p$ for disjointness inputs on $n$
coordinates, and a real number $\gamma$. Let $\mu$ be the uniform distribution on hard
samples, conditioned on the event that the generated sets $X(\omega)$ and $Y(\omega)$
are disjoint and that Bob's bit at the special coordinate is $\mathrm{false}$. For a
sample $\omega$, let $Z(\omega)$ record the transcript of $p$ on the input pair together
with the special coordinate, Alice's input bits before it, and Bob's input bits after
it, and write
\[
\Delta(\omega) = \mathrm{TV}\!\left(\mathcal{L}\!\left(\text{Alice's special-coordinate
bit}\mid Z(\omega)\right),\, U\right),
\]
the total variation distance between the conditional law of Alice's bit at the special
coordinate given $Z(\omega)$ and the uniform law $U$ on one bit. If $\int
\Delta(\omega)^2 \, d\mu(\omega) \le \gamma^4$, then $\int \Delta(\omega) \,
d\mu(\omega) \le \gamma^2$.
\end{theorem}

\begin{theorem}[Markov tail bound for Alice's conditional bit distance]
\lean{CommunicationComplexity.Functions.Disjointness.RandomizedLowerBound.disjointSpecialYFalseMeasure_xDistance_bad_le_of_integral_le}
\label{CommunicationComplexity.Functions.Disjointness.RandomizedLowerBound.disjointSpecialYFalseMeasure_xDistance_bad_le_of_integral_le}
\leanok
\uses{CommunicationComplexity.Functions.Disjointness.RandomizedLowerBound.HardSample, CommunicationComplexity.Functions.Disjointness.RandomizedLowerBound.ProtocolType, CommunicationComplexity.Functions.Disjointness.RandomizedLowerBound.disjointSpecialYFalseMeasure, CommunicationComplexity.Functions.Disjointness.RandomizedLowerBound.disjointSpecialYFalseMeasure_isProbabilityMeasure, CommunicationComplexity.Functions.Disjointness.RandomizedLowerBound.xDistance, CommunicationComplexity.Functions.Disjointness.RandomizedLowerBound.xDistance_nonneg, CommunicationComplexity.Functions.Disjointness.RandomizedLowerBound.zVariable}
\difficulty{5}
Fix a positive integer $n$, a deterministic protocol $p$ for the disjointness function
on inputs of length $n$, and a real number $\gamma > 0$. Let $\mu$ denote the hard
distribution on samples, conditioned on the generated input pair being disjoint and on
Bob's bit at the special coordinate being $\mathrm{false}$. For a sample $\omega$, let
$Z(\omega)$ record the transcript produced by running $p$ on the input pair generated by
$\omega$, together with the special coordinate $T$, Alice's input bits strictly before
$T$, and Bob's input bits strictly after $T$; and let
\[
D(\omega) \;=\; \mathrm{TV}\!\big(\,\mathrm{Law}(x_T \mid Z = Z(\omega)),\
\mathrm{Unif}\{0,1\}\,\big)
\]
be the total variation distance between the conditional law of Alice's bit $x_T$ at the
special coordinate, given the value $Z(\omega)$, and the uniform distribution on a
single bit. If the expectation of $D$ under $\mu$ satisfies $\int D \, d\mu \le
\gamma^2$, then the $\mu$-probability that $D(\omega) > \gamma$ is at most $\gamma$.
\end{theorem}

\begin{definition}[Goodz]
\lean{CommunicationComplexity.Functions.Disjointness.RandomizedLowerBound.goodZ}
\label{CommunicationComplexity.Functions.Disjointness.RandomizedLowerBound.goodZ}
\leanok
\uses{CommunicationComplexity.Functions.Disjointness.RandomizedLowerBound.ProtocolType, CommunicationComplexity.Functions.Disjointness.RandomizedLowerBound.ZType, CommunicationComplexity.Functions.Disjointness.RandomizedLowerBound.zDistance}
A $Z$ value is good when the conditional special-pair law on its achievable fiber is within $2γ$ of uniform.
\end{definition}

\begin{definition}[Goodzevent]
\lean{CommunicationComplexity.Functions.Disjointness.RandomizedLowerBound.goodZEvent}
\label{CommunicationComplexity.Functions.Disjointness.RandomizedLowerBound.goodZEvent}
\leanok
\uses{CommunicationComplexity.Functions.Disjointness.RandomizedLowerBound.HardSample, CommunicationComplexity.Functions.Disjointness.RandomizedLowerBound.ProtocolType, CommunicationComplexity.Functions.Disjointness.RandomizedLowerBound.goodZ, CommunicationComplexity.Functions.Disjointness.RandomizedLowerBound.zVariable}
The good-fiber event, pulled back to the hard sample space through $Z$.
\end{definition}

\begin{theorem}[An input lies in the input set of its own transcript]
\lean{CommunicationComplexity.Functions.Disjointness.RandomizedLowerBound.input_mem_transcript}
\label{CommunicationComplexity.Functions.Disjointness.RandomizedLowerBound.input_mem_transcript}
\leanok
\uses{CommunicationComplexity.Deterministic.Protocol.Transcript.inputSet, CommunicationComplexity.Deterministic.Protocol.mem_transcript, CommunicationComplexity.Functions.Disjointness.RandomizedLowerBound.HardSample, CommunicationComplexity.Functions.Disjointness.RandomizedLowerBound.ProtocolType, CommunicationComplexity.Functions.Disjointness.RandomizedLowerBound.input, CommunicationComplexity.Functions.Disjointness.RandomizedLowerBound.message}
\difficulty{1}
Fix a positive integer $n$, and let $p$ be a deterministic two-party protocol whose
parties each receive a subset of $\{0, 1, \dots, n-1\}$ and which outputs a Boolean
value. For any sample $\omega$ of the hard distribution, write $\mathrm{input}(\omega) =
(X(\omega), Y(\omega))$ for the pair of input sets it produces, and let $t$ be the
transcript obtained by executing $p$ on $\mathrm{input}(\omega)$. Then
$\mathrm{input}(\omega)$ belongs to the input set of $t$; that is, the input pair
generated by $\omega$ lies among those input pairs whose execution of $p$ follows
exactly the path recorded by the transcript it produces.
\end{theorem}

\begin{theorem}[Sample input lies in its transcript's input set]
\lean{CommunicationComplexity.Functions.Disjointness.RandomizedLowerBound.input_mem_transcript_of_zVariable_eq}
\label{CommunicationComplexity.Functions.Disjointness.RandomizedLowerBound.input_mem_transcript_of_zVariable_eq}
\leanok
\uses{CommunicationComplexity.Deterministic.Protocol.Transcript.inputSet, CommunicationComplexity.Functions.Disjointness.RandomizedLowerBound.HardSample, CommunicationComplexity.Functions.Disjointness.RandomizedLowerBound.ProtocolType, CommunicationComplexity.Functions.Disjointness.RandomizedLowerBound.ZType, CommunicationComplexity.Functions.Disjointness.RandomizedLowerBound.ZType.transcript, CommunicationComplexity.Functions.Disjointness.RandomizedLowerBound.input, CommunicationComplexity.Functions.Disjointness.RandomizedLowerBound.input_mem_transcript, CommunicationComplexity.Functions.Disjointness.RandomizedLowerBound.message, CommunicationComplexity.Functions.Disjointness.RandomizedLowerBound.rawZVariable, CommunicationComplexity.Functions.Disjointness.RandomizedLowerBound.zVariable}
\difficulty{2}
Fix a positive integer $n$ and a deterministic protocol $p$ for disjointness inputs of
length $n$. Let $z$ be an achievable value of the variable $Z = (M, T, X_{<T}, Y_{>T})$
for $p$, and let $\omega$ be a hard-distribution sample whose $Z$-variable equals $z$.
Then the input pair generated by $\omega$ belongs to the input set of the transcript
component of $z$; that is, executing $p$ on this input pair follows exactly the path
recorded by that transcript.
\end{theorem}

\begin{theorem}[Reading the special coordinate off the Z value]
\lean{CommunicationComplexity.Functions.Disjointness.RandomizedLowerBound.specialCoordinate_eq_of_zVariable_eq}
\label{CommunicationComplexity.Functions.Disjointness.RandomizedLowerBound.specialCoordinate_eq_of_zVariable_eq}
\leanok
\uses{CommunicationComplexity.Functions.Disjointness.RandomizedLowerBound.HardSample, CommunicationComplexity.Functions.Disjointness.RandomizedLowerBound.ProtocolType, CommunicationComplexity.Functions.Disjointness.RandomizedLowerBound.ZType, CommunicationComplexity.Functions.Disjointness.RandomizedLowerBound.ZType.specialCoordinate, CommunicationComplexity.Functions.Disjointness.RandomizedLowerBound.rawZVariable, CommunicationComplexity.Functions.Disjointness.RandomizedLowerBound.specialCoordinate, CommunicationComplexity.Functions.Disjointness.RandomizedLowerBound.zVariable}
\difficulty{2}
Fix a disjointness protocol $p$ on inputs of length $n$, and recall that the special
coordinate of a hard-distribution sample $\omega$ is the index $T$ selected by that
sample. Let $z$ be an achievable value of the variable $Z = (M, T, X_{<T}, Y_{>T})$,
whose special-coordinate component records the index $T$. If the $Z$ value produced by
$\omega$ under $p$ equals $z$, then the special coordinate of $\omega$ coincides with
the special-coordinate component of $z$.
\end{theorem}

\begin{theorem}[Recovering Alice's early bits from the Z variable]
\lean{CommunicationComplexity.Functions.Disjointness.RandomizedLowerBound.xBeforeSpecial_eq_of_zVariable_eq}
\label{CommunicationComplexity.Functions.Disjointness.RandomizedLowerBound.xBeforeSpecial_eq_of_zVariable_eq}
\leanok
\uses{CommunicationComplexity.Functions.Disjointness.RandomizedLowerBound.HardSample, CommunicationComplexity.Functions.Disjointness.RandomizedLowerBound.ProtocolType, CommunicationComplexity.Functions.Disjointness.RandomizedLowerBound.ZType, CommunicationComplexity.Functions.Disjointness.RandomizedLowerBound.ZType.xBefore, CommunicationComplexity.Functions.Disjointness.RandomizedLowerBound.rawZVariable, CommunicationComplexity.Functions.Disjointness.RandomizedLowerBound.xBeforeSpecial, CommunicationComplexity.Functions.Disjointness.RandomizedLowerBound.zVariable}
\difficulty{2}
Fix a positive integer $n$ and a deterministic protocol $p$ for the length-$n$
disjointness problem, whose inputs are pairs of subsets of $\{1,\dots,n\}$. Recall that
each hard-distribution sample $\omega$ determines a value $Z(\omega) = (M, T, X_{<T},
Y_{>T})$ of the associated $Z$ variable, consisting of the transcript, the special
coordinate, Alice's bits strictly before the special coordinate (padded by
$\mathrm{false}$ elsewhere), and Bob's bits strictly after it. Let $z$ be any achievable
value of this $Z$ variable. Then for every hard-distribution sample $\omega$ with
$Z(\omega) = z$, the function recording Alice's bits before the special coordinate of
$\omega$ agrees with the $X_{<T}$ component of $z$.
\end{theorem}

\begin{theorem}[Recovering Bob's after-special bits from the Z-value]
\lean{CommunicationComplexity.Functions.Disjointness.RandomizedLowerBound.yAfterSpecial_eq_of_zVariable_eq}
\label{CommunicationComplexity.Functions.Disjointness.RandomizedLowerBound.yAfterSpecial_eq_of_zVariable_eq}
\leanok
\uses{CommunicationComplexity.Functions.Disjointness.RandomizedLowerBound.HardSample, CommunicationComplexity.Functions.Disjointness.RandomizedLowerBound.ProtocolType, CommunicationComplexity.Functions.Disjointness.RandomizedLowerBound.ZType, CommunicationComplexity.Functions.Disjointness.RandomizedLowerBound.ZType.yAfter, CommunicationComplexity.Functions.Disjointness.RandomizedLowerBound.rawZVariable, CommunicationComplexity.Functions.Disjointness.RandomizedLowerBound.yAfterSpecial, CommunicationComplexity.Functions.Disjointness.RandomizedLowerBound.zVariable}
\difficulty{2}
Fix a positive integer $n$ and a deterministic communication protocol $p$ for
disjointness inputs of length $n$. Each hard-distribution sample $\omega$ determines a
value $Z(\omega) = (M, T, X_{<T}, Y_{>T})$, whose four components are the transcript of
$p$ on the input pair induced by $\omega$, the special coordinate $T$, Alice's bits at
the coordinates strictly below $T$ (and false elsewhere), and Bob's bits at the
coordinates strictly above $T$ (and false elsewhere). If $z$ is an achievable value of
$Z$ and $\omega$ is any sample with $Z(\omega) = z$, then the bit-string assigning to
each coordinate $i$ Bob's bit at $i$ when $T < i$ and false otherwise coincides with the
$Y_{>T}$-component of $z$.
\end{theorem}

\begin{theorem}[Mixed inputs lie in the common transcript's rectangle]
\lean{CommunicationComplexity.Functions.Disjointness.RandomizedLowerBound.mixed_inputs_mem_transcript_of_zVariable_eq}
\proofsource{roughgarden-cc-bootcamp}{
  pdf-d54d402b16e2-p007-b002,
  pdf-d54d402b16e2-p007-b003,
  pdf-d54d402b16e2-p007-b004,
  pdf-d54d402b16e2-p008-b001
}
\label{CommunicationComplexity.Functions.Disjointness.RandomizedLowerBound.mixed_inputs_mem_transcript_of_zVariable_eq}
\leanok
\uses{CommunicationComplexity.Deterministic.Protocol.Transcript.inputSet, CommunicationComplexity.Deterministic.Protocol.Transcript.inputSet_isRectangle, CommunicationComplexity.Functions.Disjointness.RandomizedLowerBound.HardSample, CommunicationComplexity.Functions.Disjointness.RandomizedLowerBound.ProtocolType, CommunicationComplexity.Functions.Disjointness.RandomizedLowerBound.X, CommunicationComplexity.Functions.Disjointness.RandomizedLowerBound.Y, CommunicationComplexity.Functions.Disjointness.RandomizedLowerBound.ZType, CommunicationComplexity.Functions.Disjointness.RandomizedLowerBound.ZType.transcript, CommunicationComplexity.Functions.Disjointness.RandomizedLowerBound.input, CommunicationComplexity.Functions.Disjointness.RandomizedLowerBound.input_mem_transcript_of_zVariable_eq, CommunicationComplexity.Functions.Disjointness.RandomizedLowerBound.zVariable, CommunicationComplexity.Rectangle.IsRectangle, CommunicationComplexity.Rectangle.IsRectangle_iff}
\difficulty{3}
Fix $n \ge 1$ and a deterministic disjointness protocol $p$ on the coordinate set
$\{0,\dots,n-1\}$. For a hard-distribution sample $\omega$, let $X(\omega), Y(\omega)
\subseteq \{0,\dots,n-1\}$ denote Alice's and Bob's input sets, and recall that the
$Z$-variable $Z = (M, T, X_{<T}, Y_{>T})$ records, among other data, the transcript $M$
that $p$ produces on the input pair $(X(\omega), Y(\omega))$. Suppose two
hard-distribution samples $\omega$ and $\omega'$ yield the same value $z$ of this
$Z$-variable under $p$, and let $M$ be the transcript component of $z$. Then both mixed
input pairs $(X(\omega'), Y(\omega))$ and $(X(\omega), Y(\omega'))$ belong to the input
set of $M$, that is, the set of all input pairs whose execution of $p$ follows the
transcript $M$.
\end{theorem}

\begin{theorem}[Special coordinate is determined by the Z-variable]
\lean{CommunicationComplexity.Functions.Disjointness.RandomizedLowerBound.specialCoordinate_eq_of_same_zVariable}
\label{CommunicationComplexity.Functions.Disjointness.RandomizedLowerBound.specialCoordinate_eq_of_same_zVariable}
\leanok
\uses{CommunicationComplexity.Functions.Disjointness.RandomizedLowerBound.HardSample, CommunicationComplexity.Functions.Disjointness.RandomizedLowerBound.ProtocolType, CommunicationComplexity.Functions.Disjointness.RandomizedLowerBound.ZType, CommunicationComplexity.Functions.Disjointness.RandomizedLowerBound.specialCoordinate, CommunicationComplexity.Functions.Disjointness.RandomizedLowerBound.specialCoordinate_eq_of_zVariable_eq, CommunicationComplexity.Functions.Disjointness.RandomizedLowerBound.zVariable}
\difficulty{2}
Fix a deterministic protocol $p$ for disjointness inputs of length $n$, and recall that
the $Z$-variable of a hard-distribution sample $\omega$ records the tuple $(M, T,
X_{<T}, Y_{>T})$ consisting of the protocol's transcript on the input generated by
$\omega$, the special coordinate $T$, and the padded bit-vectors of Alice's input before
$T$ and Bob's input after $T$. If two hard-distribution samples $\omega$ and $\omega'$
yield the same value $z$ of this $Z$-variable, then they select the same special
coordinate; that is, $\omega.T = \omega'.T$.
\end{theorem}

\begin{theorem}[Agreement of the pre-special bits on a Z-fibre]
\lean{CommunicationComplexity.Functions.Disjointness.RandomizedLowerBound.xBeforeSpecial_eq_of_same_zVariable}
\label{CommunicationComplexity.Functions.Disjointness.RandomizedLowerBound.xBeforeSpecial_eq_of_same_zVariable}
\leanok
\uses{CommunicationComplexity.Functions.Disjointness.RandomizedLowerBound.HardSample, CommunicationComplexity.Functions.Disjointness.RandomizedLowerBound.ProtocolType, CommunicationComplexity.Functions.Disjointness.RandomizedLowerBound.ZType, CommunicationComplexity.Functions.Disjointness.RandomizedLowerBound.xBeforeSpecial, CommunicationComplexity.Functions.Disjointness.RandomizedLowerBound.xBeforeSpecial_eq_of_zVariable_eq, CommunicationComplexity.Functions.Disjointness.RandomizedLowerBound.zVariable}
\difficulty{2}
Fix an integer $n \ge 1$ and a deterministic communication protocol $p$ for disjointness
on $n$ coordinates, and consider the coarse variable $Z = (M, T, X_{<T}, Y_{>T})$ — the
transcript $M$ produced by running $p$ on the input pair, the special coordinate $T$,
the vector $X_{<T}$ of Alice's bits strictly below $T$ (padded with $\mathrm{false}$
elsewhere), and the vector $Y_{>T}$ of Bob's bits strictly above $T$ (padded with
$\mathrm{false}$ elsewhere). If two hard-distribution samples $\omega$ and $\omega'$
yield the same value $z$ of $Z$, then they determine the same pre-special vector:
$X_{<T}(\omega) = X_{<T}(\omega')$.
\end{theorem}

\begin{theorem}[Post-special bits determined by the coarse variable]
\lean{CommunicationComplexity.Functions.Disjointness.RandomizedLowerBound.yAfterSpecial_eq_of_same_zVariable}
\label{CommunicationComplexity.Functions.Disjointness.RandomizedLowerBound.yAfterSpecial_eq_of_same_zVariable}
\leanok
\uses{CommunicationComplexity.Functions.Disjointness.RandomizedLowerBound.HardSample, CommunicationComplexity.Functions.Disjointness.RandomizedLowerBound.ProtocolType, CommunicationComplexity.Functions.Disjointness.RandomizedLowerBound.ZType, CommunicationComplexity.Functions.Disjointness.RandomizedLowerBound.yAfterSpecial, CommunicationComplexity.Functions.Disjointness.RandomizedLowerBound.yAfterSpecial_eq_of_zVariable_eq, CommunicationComplexity.Functions.Disjointness.RandomizedLowerBound.zVariable}
\difficulty{2}
Let $n$ be a positive integer and let $p$ be a deterministic communication protocol for
disjointness inputs of length $n$. To each sample $\omega$ of the hard distribution one
associates the value $Z(\omega) = (M, T, X_{<T}, Y_{>T})$, recording the transcript $M$
of $p$ run on the input generated by $\omega$, the special coordinate $T$, Alice's input
bits strictly before $T$, and Bob's input bits strictly after $T$ (each of the last two
padded with $\mathrm{false}$ outside its range). If two samples $\omega$ and $\omega'$
yield the same value $z$ of this variable, then the post-special bit vectors they
determine agree: $Y_{>T}(\omega) = Y_{>T}(\omega')$.
\end{theorem}

\begin{theorem}[Invariance of the Z-variable under mixing]
\lean{CommunicationComplexity.Functions.Disjointness.RandomizedLowerBound.zVariable_mix_eq_of_same_zVariable}
\label{CommunicationComplexity.Functions.Disjointness.RandomizedLowerBound.zVariable_mix_eq_of_same_zVariable}
\leanok
\uses{CommunicationComplexity.Deterministic.Protocol.Transcript.inputSet, CommunicationComplexity.Deterministic.Protocol.transcript, CommunicationComplexity.Deterministic.Protocol.transcript_eq_of_mem, CommunicationComplexity.Functions.Disjointness.RandomizedLowerBound.HardSample, CommunicationComplexity.Functions.Disjointness.RandomizedLowerBound.ProtocolType, CommunicationComplexity.Functions.Disjointness.RandomizedLowerBound.RawZType.ext, CommunicationComplexity.Functions.Disjointness.RandomizedLowerBound.ZType, CommunicationComplexity.Functions.Disjointness.RandomizedLowerBound.ZType.specialCoordinate, CommunicationComplexity.Functions.Disjointness.RandomizedLowerBound.ZType.transcript, CommunicationComplexity.Functions.Disjointness.RandomizedLowerBound.ZType.xBefore, CommunicationComplexity.Functions.Disjointness.RandomizedLowerBound.ZType.yAfter, CommunicationComplexity.Functions.Disjointness.RandomizedLowerBound.input, CommunicationComplexity.Functions.Disjointness.RandomizedLowerBound.input_mix, CommunicationComplexity.Functions.Disjointness.RandomizedLowerBound.mix, CommunicationComplexity.Functions.Disjointness.RandomizedLowerBound.mixed_inputs_mem_transcript_of_zVariable_eq, CommunicationComplexity.Functions.Disjointness.RandomizedLowerBound.rawZVariable, CommunicationComplexity.Functions.Disjointness.RandomizedLowerBound.specialCoordinate, CommunicationComplexity.Functions.Disjointness.RandomizedLowerBound.specialCoordinate_eq_of_same_zVariable, CommunicationComplexity.Functions.Disjointness.RandomizedLowerBound.specialCoordinate_eq_of_zVariable_eq, CommunicationComplexity.Functions.Disjointness.RandomizedLowerBound.xBeforeSpecial, CommunicationComplexity.Functions.Disjointness.RandomizedLowerBound.xBeforeSpecial_eq_of_same_zVariable, CommunicationComplexity.Functions.Disjointness.RandomizedLowerBound.xBeforeSpecial_eq_of_zVariable_eq, CommunicationComplexity.Functions.Disjointness.RandomizedLowerBound.xBeforeSpecial_mix, CommunicationComplexity.Functions.Disjointness.RandomizedLowerBound.yAfterSpecial, CommunicationComplexity.Functions.Disjointness.RandomizedLowerBound.yAfterSpecial_eq_of_same_zVariable, CommunicationComplexity.Functions.Disjointness.RandomizedLowerBound.yAfterSpecial_eq_of_zVariable_eq, CommunicationComplexity.Functions.Disjointness.RandomizedLowerBound.yAfterSpecial_mix, CommunicationComplexity.Functions.Disjointness.RandomizedLowerBound.zVariable}
\difficulty{5}
Fix a deterministic protocol $p$ for set disjointness on $n$ coordinates, and for a
sample $\omega$ of the hard distribution write $Z(\omega) = (M, T, X_{<T}, Y_{>T})$ for
the associated variable, where $M$ is the transcript that $p$ produces on the input pair
generated by $\omega$, $T$ is the special coordinate, $X_{<T}$ is Alice's input
restricted to the coordinates before $T$, and $Y_{>T}$ is Bob's input restricted to the
coordinates after $T$. Let $\omega_X$ and $\omega_Y$ be two samples, and let $z$ be an
achievable value of this variable with $Z(\omega_X) = z$ and $Z(\omega_Y) = z$. Then the
mixed sample that takes Alice's data from $\omega_X$ and Bob's data from $\omega_Y$
again satisfies $Z(\mathrm{mix}(\omega_X, \omega_Y)) = z$.
\end{theorem}

\begin{theorem}[Mixing preserves the fiber and sets the special pair]
\lean{CommunicationComplexity.Functions.Disjointness.RandomizedLowerBound.mix_mem_fiber_inter_specialPair_of_mem_specialX_specialY}
\label{CommunicationComplexity.Functions.Disjointness.RandomizedLowerBound.mix_mem_fiber_inter_specialPair_of_mem_specialX_specialY}
\leanok
\uses{CommunicationComplexity.Functions.Disjointness.RandomizedLowerBound.HardSample, CommunicationComplexity.Functions.Disjointness.RandomizedLowerBound.ProtocolType, CommunicationComplexity.Functions.Disjointness.RandomizedLowerBound.ZType, CommunicationComplexity.Functions.Disjointness.RandomizedLowerBound.mix, CommunicationComplexity.Functions.Disjointness.RandomizedLowerBound.specialPair, CommunicationComplexity.Functions.Disjointness.RandomizedLowerBound.specialX, CommunicationComplexity.Functions.Disjointness.RandomizedLowerBound.specialX_mix, CommunicationComplexity.Functions.Disjointness.RandomizedLowerBound.specialY, CommunicationComplexity.Functions.Disjointness.RandomizedLowerBound.specialY_mix, CommunicationComplexity.Functions.Disjointness.RandomizedLowerBound.zFiber, CommunicationComplexity.Functions.Disjointness.RandomizedLowerBound.zVariable, CommunicationComplexity.Functions.Disjointness.RandomizedLowerBound.zVariable_mix_eq_of_same_zVariable}
\difficulty{3}
Fix a deterministic communication protocol $p$ for disjointness on $n$ coordinates, a
value $z$ of the statistic $Z = (M, T, X_{<T}, Y_{>T})$ recorded by the protocol (its
transcript, the special coordinate, Alice's input bits before the special coordinate,
and Bob's input bits after it), and two bits $b_X, b_Y$. Let $\omega_X$ and $\omega_Y$
be hard-distribution samples that both lie in the fiber $Z = z$, and suppose that
Alice's bit at the special coordinate of $\omega_X$ equals $b_X$ while Bob's bit at the
special coordinate of $\omega_Y$ equals $b_Y$. Then the mixed sample — taking Alice's
input from $\omega_X$ and Bob's input from $\omega_Y$ — again lies in the fiber $Z = z$,
and its pair of special-coordinate bits is $(b_X, b_Y)$.
\end{theorem}

\begin{theorem}[Mixed sample remains in the Z-variable fiber]
\lean{CommunicationComplexity.Functions.Disjointness.RandomizedLowerBound.mix_swap_mem_fiber_of_mem_specialX_specialY}
\label{CommunicationComplexity.Functions.Disjointness.RandomizedLowerBound.mix_swap_mem_fiber_of_mem_specialX_specialY}
\leanok
\uses{CommunicationComplexity.Functions.Disjointness.RandomizedLowerBound.HardSample, CommunicationComplexity.Functions.Disjointness.RandomizedLowerBound.ProtocolType, CommunicationComplexity.Functions.Disjointness.RandomizedLowerBound.ZType, CommunicationComplexity.Functions.Disjointness.RandomizedLowerBound.mix, CommunicationComplexity.Functions.Disjointness.RandomizedLowerBound.specialX, CommunicationComplexity.Functions.Disjointness.RandomizedLowerBound.specialY, CommunicationComplexity.Functions.Disjointness.RandomizedLowerBound.zFiber, CommunicationComplexity.Functions.Disjointness.RandomizedLowerBound.zVariable, CommunicationComplexity.Functions.Disjointness.RandomizedLowerBound.zVariable_mix_eq_of_same_zVariable}
\difficulty{2}
Fix a deterministic protocol $p$ for disjointness inputs of length $n$ and a value $z$
of the variable $Z = (M, T, X_{<T}, Y_{>T})$, recording the transcript, the special
coordinate, Alice's input bits before it and Bob's input bits after it. Let $\omega_X$
and $\omega_Y$ be two hard-distribution samples that both lie in the fiber above $z$, so
that each realizes $Z = z$, and suppose moreover that Alice's bit at the special
coordinate of $\omega_X$ equals $b_X$ and Bob's bit at the special coordinate of
$\omega_Y$ equals $b_Y$. Then the mixed sample obtained by taking Alice's coordinates
from $\omega_Y$ together with Bob's coordinates from $\omega_X$ again lies in the fiber
above $z$.
\end{theorem}

\begin{theorem}[Mixing preserves the Z-fiber and the special Alice bit]
\lean{CommunicationComplexity.Functions.Disjointness.RandomizedLowerBound.mix_mem_fiber_inter_specialX_of_mem_specialPair_fiber}
\label{CommunicationComplexity.Functions.Disjointness.RandomizedLowerBound.mix_mem_fiber_inter_specialX_of_mem_specialPair_fiber}
\leanok
\uses{CommunicationComplexity.Functions.Disjointness.RandomizedLowerBound.HardSample, CommunicationComplexity.Functions.Disjointness.RandomizedLowerBound.ProtocolType, CommunicationComplexity.Functions.Disjointness.RandomizedLowerBound.ZType, CommunicationComplexity.Functions.Disjointness.RandomizedLowerBound.mix, CommunicationComplexity.Functions.Disjointness.RandomizedLowerBound.specialPair, CommunicationComplexity.Functions.Disjointness.RandomizedLowerBound.specialX, CommunicationComplexity.Functions.Disjointness.RandomizedLowerBound.specialX_mix, CommunicationComplexity.Functions.Disjointness.RandomizedLowerBound.zFiber, CommunicationComplexity.Functions.Disjointness.RandomizedLowerBound.zVariable, CommunicationComplexity.Functions.Disjointness.RandomizedLowerBound.zVariable_mix_eq_of_same_zVariable}
\difficulty{3}
Fix a deterministic disjointness protocol $p$ on inputs of length $n$, and let $z$ be an
achievable value of the variable $Z = (M, T, X_{<T}, Y_{>T})$, whose fiber $Z^{-1}(z)$
is the set of hard-distribution samples mapping to $z$. Let $\omega_{\mathrm{pair}}$ and
$\omega$ be two samples of this fiber, and suppose that the pair of bits at the special
coordinate of $\omega_{\mathrm{pair}}$, namely Alice's bit and Bob's bit there, equals
$b = (b_1, b_2)$. Form the mixed sample by taking Alice's input from
$\omega_{\mathrm{pair}}$ and Bob's input from $\omega$ (keeping their common special
coordinate $T$, the bits of $X$ before $T$, and the bits of $Y$ after $T$). Then this
mixed sample again lies in $Z^{-1}(z)$, and Alice's bit at its special coordinate equals
$b_1$.
\end{theorem}

\begin{theorem}[Mixing two fiber samples preserves the fiber and Bob's special bit]
\lean{CommunicationComplexity.Functions.Disjointness.RandomizedLowerBound.mix_mem_fiber_inter_specialY_of_mem_fiber_specialPair}
\label{CommunicationComplexity.Functions.Disjointness.RandomizedLowerBound.mix_mem_fiber_inter_specialY_of_mem_fiber_specialPair}
\leanok
\uses{CommunicationComplexity.Functions.Disjointness.RandomizedLowerBound.HardSample, CommunicationComplexity.Functions.Disjointness.RandomizedLowerBound.ProtocolType, CommunicationComplexity.Functions.Disjointness.RandomizedLowerBound.ZType, CommunicationComplexity.Functions.Disjointness.RandomizedLowerBound.mix, CommunicationComplexity.Functions.Disjointness.RandomizedLowerBound.specialPair, CommunicationComplexity.Functions.Disjointness.RandomizedLowerBound.specialY, CommunicationComplexity.Functions.Disjointness.RandomizedLowerBound.specialY_mix, CommunicationComplexity.Functions.Disjointness.RandomizedLowerBound.zFiber, CommunicationComplexity.Functions.Disjointness.RandomizedLowerBound.zVariable, CommunicationComplexity.Functions.Disjointness.RandomizedLowerBound.zVariable_mix_eq_of_same_zVariable}
\difficulty{3}
Fix a positive integer $n$ and a deterministic communication protocol $p$ for the
length-$n$ disjointness problem, and consider the hard-distribution samples together
with the derived variable $Z = (M, T, X_{<T}, Y_{>T})$ recording the transcript, the
special coordinate, Alice's bits before it, and Bob's bits after it. Let $z$ be an
achievable value of $Z$, let $\omega$ and $\omega'$ be two samples, and let $b = (b_1,
b_2)$ be a pair of Booleans. Suppose that both $\omega$ and $\omega'$ produce the value
$z$ (that is, each lies in the fiber $Z^{-1}(z)$), and that the special-coordinate
bit-pair of $\omega'$ equals $b$, so that Alice's and Bob's bits at the special
coordinate of $\omega'$ are $b_1$ and $b_2$ respectively. Then the mixed sample formed
by taking Alice's data from $\omega$ and Bob's data from $\omega'$ again produces the
value $z$, and its bit for Bob at the special coordinate equals $b_2$.
\end{theorem}

\begin{theorem}[Fiberwise independence of the two special bits]
\lean{CommunicationComplexity.Functions.Disjointness.RandomizedLowerBound.card_fiber_inter_specialX_mul_card_fiber_inter_specialY}
\label{CommunicationComplexity.Functions.Disjointness.RandomizedLowerBound.card_fiber_inter_specialX_mul_card_fiber_inter_specialY}
\leanok
\uses{CommunicationComplexity.Functions.Disjointness.RandomizedLowerBound.HardSample, CommunicationComplexity.Functions.Disjointness.RandomizedLowerBound.ProtocolType, CommunicationComplexity.Functions.Disjointness.RandomizedLowerBound.ZType, CommunicationComplexity.Functions.Disjointness.RandomizedLowerBound.conditionalSpecialPairLaw_eq_prod_of_singleton_factorization, CommunicationComplexity.Functions.Disjointness.RandomizedLowerBound.mix, CommunicationComplexity.Functions.Disjointness.RandomizedLowerBound.mix_mem_fiber_inter_specialPair_of_mem_specialX_specialY, CommunicationComplexity.Functions.Disjointness.RandomizedLowerBound.mix_mem_fiber_inter_specialX_of_mem_specialPair_fiber, CommunicationComplexity.Functions.Disjointness.RandomizedLowerBound.mix_mem_fiber_inter_specialY_of_mem_fiber_specialPair, CommunicationComplexity.Functions.Disjointness.RandomizedLowerBound.mix_mix_swap, CommunicationComplexity.Functions.Disjointness.RandomizedLowerBound.mix_swap_mem_fiber_of_mem_specialX_specialY, CommunicationComplexity.Functions.Disjointness.RandomizedLowerBound.specialCoordinate_eq_of_same_zVariable, CommunicationComplexity.Functions.Disjointness.RandomizedLowerBound.specialPair, CommunicationComplexity.Functions.Disjointness.RandomizedLowerBound.specialX, CommunicationComplexity.Functions.Disjointness.RandomizedLowerBound.specialY, CommunicationComplexity.Functions.Disjointness.RandomizedLowerBound.xBeforeSpecial_eq_of_same_zVariable, CommunicationComplexity.Functions.Disjointness.RandomizedLowerBound.yAfterSpecial_eq_of_same_zVariable, CommunicationComplexity.Functions.Disjointness.RandomizedLowerBound.zFiber, CommunicationComplexity.Functions.Disjointness.RandomizedLowerBound.zVariable}
\difficulty{6}
Fix a positive integer $n$, a deterministic communication protocol $p$ for disjointness
inputs of length $n$, an achievable value $z$ of the associated $Z$-variable, and a pair
of bits $b=(b_1,b_2)$. Let $F$ be the fiber consisting of those hard-distribution
samples $\omega$ whose $Z$-variable equals $z$, and for such a sample write
$x_T(\omega)$ for Alice's bit at the special coordinate and $y_T(\omega)$ for Bob's bit
at the special coordinate. Then
\[
\abs{\{\omega\in F : x_T(\omega)=b_1\}}\cdot\abs{\{\omega\in F : y_T(\omega)=b_2\}}
=\abs{\{\omega\in F : (x_T(\omega),y_T(\omega))=b\}}\cdot\abs{F}.
\]
\end{theorem}

\begin{theorem}[Uniform measure as a cardinality ratio]
\lean{CommunicationComplexity.Functions.Disjointness.RandomizedLowerBound.measureReal_eq_card_subtype_div}
\label{CommunicationComplexity.Functions.Disjointness.RandomizedLowerBound.measureReal_eq_card_subtype_div}
\leanok
\uses{CommunicationComplexity.Functions.Disjointness.RandomizedLowerBound.HardSample, CommunicationComplexity.uniformOn_univ_measureReal_eq_card_filter}
\difficulty{2}
Fix $n \ge 1$, and equip the finite sample space of hard samples with its uniform
probability measure. Then for every subset $S$ of this space, the measure of $S$ equals
the ratio
\[
\frac{\abs{\{\omega \in S\}}}{\abs{\text{HardSample}}},
\]
that is, the number of hard samples lying in $S$ divided by the total number of hard
samples.
\end{theorem}

\begin{theorem}[Conditional independence of the special-coordinate bits]
\lean{CommunicationComplexity.Functions.Disjointness.RandomizedLowerBound.fiber_volume_factorization}
\proofsource{roughgarden-cc-bootcamp}{
  pdf-d54d402b16e2-p007-b002,
  pdf-d54d402b16e2-p007-b003
}
\label{CommunicationComplexity.Functions.Disjointness.RandomizedLowerBound.fiber_volume_factorization}
\leanok
\proofstep
  {CommunicationComplexity.Deterministic.Protocol}
  {direct}
  {roughgarden-cc-bootcamp}
  {pdf-d54d402b16e2-p002-b002,
    pdf-d54d402b16e2-p004-b005}
\proofstep
  {CommunicationComplexity.Deterministic.Protocol.run}
  {context}
  {roughgarden-cc-bootcamp}
  {pdf-d54d402b16e2-p002-b002}
\proofstep
  {CommunicationComplexity.Deterministic.Protocol.complexity}
  {direct}
  {roughgarden-cc-bootcamp}
  {pdf-d54d402b16e2-p002-b002}
\proofstep
  {CommunicationComplexity.Deterministic.Protocol.swap}
  {context}
  {roughgarden-cc-bootcamp}
  {pdf-d54d402b16e2-p002-b002}
\proofstep
  {CommunicationComplexity.Deterministic.Protocol.swap_run}
  {context}
  {roughgarden-cc-bootcamp}
  {pdf-d54d402b16e2-p002-b002}
\proofstep
  {CommunicationComplexity.Deterministic.Protocol.swap_complexity}
  {context}
  {roughgarden-cc-bootcamp}
  {pdf-d54d402b16e2-p002-b002}
\proofstep
  {CommunicationComplexity.Deterministic.Protocol.swap_swap}
  {context}
  {roughgarden-cc-bootcamp}
  {pdf-d54d402b16e2-p002-b002}
\proofstep
  {CommunicationComplexity.Deterministic.Protocol.comap}
  {context}
  {roughgarden-cc-bootcamp}
  {pdf-d54d402b16e2-p002-b002}
\proofstep
  {CommunicationComplexity.Deterministic.Protocol.comap_run}
  {context}
  {roughgarden-cc-bootcamp}
  {pdf-d54d402b16e2-p002-b002}
\proofstep
  {CommunicationComplexity.Deterministic.Protocol.comap_complexity}
  {context}
  {roughgarden-cc-bootcamp}
  {pdf-d54d402b16e2-p002-b002}
\proofstep
  {CommunicationComplexity.Deterministic.Protocol.IsSubprotocol}
  {context}
  {roughgarden-cc-bootcamp}
  {pdf-d54d402b16e2-p004-b005}
\proofstep
  {CommunicationComplexity.Deterministic.Protocol.SubprotocolPath}
  {context}
  {roughgarden-cc-bootcamp}
  {pdf-d54d402b16e2-p004-b005}
\proofstep
  {CommunicationComplexity.Deterministic.Protocol.path_exists_of_isSubprotocol}
  {context}
  {roughgarden-cc-bootcamp}
  {pdf-d54d402b16e2-p004-b005}
\proofstep
  {CommunicationComplexity.Deterministic.Protocol.choosePath}
  {context}
  {roughgarden-cc-bootcamp}
  {pdf-d54d402b16e2-p004-b005}
\proofstep
  {CommunicationComplexity.Deterministic.Protocol.reachXPath}
  {context}
  {roughgarden-cc-bootcamp}
  {pdf-d54d402b16e2-p007-b004}
\proofstep
  {CommunicationComplexity.Deterministic.Protocol.reachYPath}
  {context}
  {roughgarden-cc-bootcamp}
  {pdf-d54d402b16e2-p007-b004}
\proofstep
  {CommunicationComplexity.Deterministic.Protocol.reachX}
  {context}
  {roughgarden-cc-bootcamp}
  {pdf-d54d402b16e2-p007-b002}
\proofstep
  {CommunicationComplexity.Deterministic.Protocol.reachY}
  {context}
  {roughgarden-cc-bootcamp}
  {pdf-d54d402b16e2-p007-b002}
\proofstep
  {CommunicationComplexity.Deterministic.Protocol.testSubprotocol}
  {context}
  {roughgarden-cc-bootcamp}
  {pdf-d54d402b16e2-p007-b003}
\proofstep
  {CommunicationComplexity.Deterministic.Protocol.testSubprotocol_complexity}
  {context}
  {roughgarden-cc-bootcamp}
  {pdf-d54d402b16e2-p007-b003}
\proofstep
  {CommunicationComplexity.Deterministic.FiniteMessage.Protocol}
  {context}
  {roughgarden-cc-bootcamp}
  {pdf-d54d402b16e2-p002-b002}
\proofstep
  {CommunicationComplexity.Deterministic.FiniteMessage.Protocol.run}
  {context}
  {roughgarden-cc-bootcamp}
  {pdf-d54d402b16e2-p002-b002}
\proofstep
  {CommunicationComplexity.Deterministic.FiniteMessage.Protocol.complexity}
  {context}
  {roughgarden-cc-bootcamp}
  {pdf-d54d402b16e2-p002-b002}
\proofstep
  {CommunicationComplexity.Deterministic.FiniteMessage.Protocol.completeTreeAlice}
  {context}
  {roughgarden-cc-bootcamp}
  {pdf-d54d402b16e2-p004-b003}
\proofstep
  {CommunicationComplexity.Deterministic.FiniteMessage.Protocol.completeTreeAlice_run}
  {context}
  {roughgarden-cc-bootcamp}
  {pdf-d54d402b16e2-p004-b003}
\proofstep
  {CommunicationComplexity.Deterministic.FiniteMessage.Protocol.completeTreeAlice_complexity}
  {context}
  {roughgarden-cc-bootcamp}
  {pdf-d54d402b16e2-p002-b002}
\proofstep
  {CommunicationComplexity.Deterministic.FiniteMessage.Protocol.encode_alice}
  {context}
  {roughgarden-cc-bootcamp}
  {pdf-d54d402b16e2-p002-b002}
\proofstep
  {CommunicationComplexity.Deterministic.FiniteMessage.Protocol.toProtocol_exists}
  {context}
  {roughgarden-cc-bootcamp}
  {pdf-d54d402b16e2-p002-b002}
\proofstep
  {CommunicationComplexity.Deterministic.FiniteMessage.Protocol.toProtocol}
  {context}
  {roughgarden-cc-bootcamp}
  {pdf-d54d402b16e2-p002-b002}
\proofstep
  {CommunicationComplexity.Deterministic.FiniteMessage.Protocol.toProtocol_run}
  {context}
  {roughgarden-cc-bootcamp}
  {pdf-d54d402b16e2-p002-b002}
\proofstep
  {CommunicationComplexity.Deterministic.FiniteMessage.Protocol.toProtocol_complexity}
  {context}
  {roughgarden-cc-bootcamp}
  {pdf-d54d402b16e2-p002-b002}
\proofstep
  {CommunicationComplexity.Deterministic.FiniteMessage.Protocol.ofProtocol}
  {context}
  {roughgarden-cc-bootcamp}
  {pdf-d54d402b16e2-p002-b002}
\proofstep
  {CommunicationComplexity.Deterministic.FiniteMessage.Protocol.ofProtocol_run}
  {context}
  {roughgarden-cc-bootcamp}
  {pdf-d54d402b16e2-p002-b002}
\proofstep
  {CommunicationComplexity.Deterministic.FiniteMessage.Protocol.ofProtocol_complexity}
  {context}
  {roughgarden-cc-bootcamp}
  {pdf-d54d402b16e2-p002-b002}
\proofstep
  {CommunicationComplexity.Deterministic.FiniteMessage.Protocol.comap}
  {context}
  {roughgarden-cc-bootcamp}
  {pdf-d54d402b16e2-p002-b002}
\proofstep
  {CommunicationComplexity.Deterministic.FiniteMessage.Protocol.comap_run}
  {context}
  {roughgarden-cc-bootcamp}
  {pdf-d54d402b16e2-p002-b002}
\proofstep
  {CommunicationComplexity.Deterministic.FiniteMessage.Protocol.comap_complexity}
  {context}
  {roughgarden-cc-bootcamp}
  {pdf-d54d402b16e2-p002-b002}
\proofstep
  {CommunicationComplexity.Internal.enat_iInf_le_coe_iff}
  {context}
  {roughgarden-cc-bootcamp}
  {pdf-d54d402b16e2-p002-b002}
\proofstep
  {CommunicationComplexity.Rectangle.IsRectangle}
  {direct}
  {roughgarden-cc-bootcamp}
  {pdf-d54d402b16e2-p007-b003}
\proofstep
  {CommunicationComplexity.Rectangle.IsRectangle_iff}
  {direct}
  {roughgarden-cc-bootcamp}
  {pdf-d54d402b16e2-p007-b003}
\proofstep
  {CommunicationComplexity.Deterministic.Protocol.swapInputSet}
  {context}
  {roughgarden-cc-bootcamp}
  {pdf-d54d402b16e2-p007-b003}
\proofstep
  {CommunicationComplexity.Deterministic.Protocol.mem_swapInputSet}
  {context}
  {roughgarden-cc-bootcamp}
  {pdf-d54d402b16e2-p007-b003}
\proofstep
  {CommunicationComplexity.Deterministic.Protocol.swapInputSet_swapInputSet}
  {context}
  {roughgarden-cc-bootcamp}
  {pdf-d54d402b16e2-p007-b003}
\proofstep
  {CommunicationComplexity.Deterministic.Protocol.preimageInputSet}
  {context}
  {roughgarden-cc-bootcamp}
  {pdf-d54d402b16e2-p006-b005}
\proofstep
  {CommunicationComplexity.Deterministic.Protocol.mem_preimageInputSet}
  {context}
  {roughgarden-cc-bootcamp}
  {pdf-d54d402b16e2-p006-b005}
\proofstep
  {CommunicationComplexity.FiniteMeasureSpace}
  {context}
  {roughgarden-cc-bootcamp}
  {pdf-d54d402b16e2-p011-b003}
\proofstep
  {CommunicationComplexity.FiniteMeasureSpace.of}
  {context}
  {roughgarden-cc-bootcamp}
  {pdf-d54d402b16e2-p011-b003}
\proofstep
  {CommunicationComplexity.FiniteProbabilitySpace}
  {context}
  {roughgarden-cc-bootcamp}
  {pdf-d54d402b16e2-p011-b003}
\proofstep
  {CommunicationComplexity.FiniteProbabilitySpace.of}
  {context}
  {roughgarden-cc-bootcamp}
  {pdf-d54d402b16e2-p011-b003}
\proofstep
  {CommunicationComplexity.FiniteProbabilitySpace.ofMeasure}
  {context}
  {roughgarden-cc-bootcamp}
  {pdf-d54d402b16e2-p011-b003}
\proofstep
  {CommunicationComplexity.FiniteMeasureSpace.measureReal_eq_sum_singletons}
  {context}
  {roughgarden-cc-bootcamp}
  {pdf-d54d402b16e2-p012-b002}
\proofstep
  {CommunicationComplexity.FiniteMeasureSpace.measureReal_preimage_eq_sum_fibers}
  {context}
  {roughgarden-cc-bootcamp}
  {pdf-d54d402b16e2-p012-b002}
\proofstep
  {CommunicationComplexity.FiniteMeasureSpace.absolutelyContinuous_iff_forall_singletons}
  {context}
  {roughgarden-cc-bootcamp}
  {pdf-d54d402b16e2-p012-b002}
\proofstep
  {CommunicationComplexity.FiniteMeasureSpace.sq_integral_le_integral_sq}
  {context}
  {roughgarden-cc-bootcamp}
  {pdf-d54d402b16e2-p012-b002}
\proofstep
  {CommunicationComplexity.FiniteMeasureSpace.integral_comp_eq_sum_measureReal_fibers}
  {context}
  {roughgarden-cc-bootcamp}
  {pdf-d54d402b16e2-p012-b002}
\proofstep
  {CommunicationComplexity.FiniteMeasureSpace.measureReal_eq_sum_cond_fiber_real}
  {context}
  {roughgarden-cc-bootcamp}
  {pdf-d54d402b16e2-p011-b003}
\proofstep
  {CommunicationComplexity.finiteMeasureSpaceBool}
  {context}
  {roughgarden-cc-bootcamp}
  {pdf-d54d402b16e2-p011-b003}
\proofstep
  {CommunicationComplexity.finiteMeasureSpaceProd}
  {context}
  {roughgarden-cc-bootcamp}
  {pdf-d54d402b16e2-p012-b002}
\proofstep
  {CommunicationComplexity.finiteMeasureSpacePi}
  {context}
  {roughgarden-cc-bootcamp}
  {pdf-d54d402b16e2-p011-b003}
\proofstep
  {CommunicationComplexity.instProd}
  {context}
  {roughgarden-cc-bootcamp}
  {pdf-d54d402b16e2-p012-b002}
\proofstep
  {CommunicationComplexity.instPi}
  {context}
  {roughgarden-cc-bootcamp}
  {pdf-d54d402b16e2-p012-b002}
\proofstep
  {CommunicationComplexity.uniformOn_univ_measureReal_eq_card_filter}
  {context}
  {roughgarden-cc-bootcamp}
  {pdf-d54d402b16e2-p011-b003}
\proofstep
  {CommunicationComplexity.uniformOn_univ_measureReal_eq_card_subtype}
  {context}
  {roughgarden-cc-bootcamp}
  {pdf-d54d402b16e2-p011-b003}
\proofstep
  {CommunicationComplexity.FiniteProbabilitySpace.nonempty}
  {context}
  {roughgarden-cc-bootcamp}
  {pdf-d54d402b16e2-p012-b002}
\proofstep
  {CommunicationComplexity.FiniteProbabilitySpace.instNonempty}
  {context}
  {roughgarden-cc-bootcamp}
  {pdf-d54d402b16e2-p012-b002}
\proofstep
  {CommunicationComplexity.FiniteProbabilitySpace.toPMF}
  {context}
  {roughgarden-cc-bootcamp}
  {pdf-d54d402b16e2-p011-b003}
\proofstep
  {CommunicationComplexity.FiniteProbabilitySpace.measure_eq}
  {context}
  {roughgarden-cc-bootcamp}
  {pdf-d54d402b16e2-p012-b002}
\proofstep
  {CommunicationComplexity.FiniteProbabilitySpace.hasSum_measure_singletons}
  {context}
  {roughgarden-cc-bootcamp}
  {pdf-d54d402b16e2-p012-b002}
\proofstep
  {CommunicationComplexity.FiniteProbabilitySpace.pmf_prod}
  {context}
  {roughgarden-cc-bootcamp}
  {pdf-d54d402b16e2-p012-b002}
\proofstep
  {CommunicationComplexity.FiniteProbabilitySpace.measureReal_prod}
  {context}
  {roughgarden-cc-bootcamp}
  {pdf-d54d402b16e2-p012-b002}
\proofstep
  {CommunicationComplexity.FiniteProbabilitySpace.measureReal_iUnion_fintype}
  {context}
  {roughgarden-cc-bootcamp}
  {pdf-d54d402b16e2-p012-b002}
\proofstep
  {CommunicationComplexity.FiniteProbabilitySpace.measureReal_preimage_finset}
  {context}
  {roughgarden-cc-bootcamp}
  {pdf-d54d402b16e2-p012-b002}
\proofstep
  {CommunicationComplexity.FiniteProbabilitySpace.measureReal_finset}
  {context}
  {roughgarden-cc-bootcamp}
  {pdf-d54d402b16e2-p012-b002}
\proofstep
  {CommunicationComplexity.FiniteProbabilitySpace.integral_eq_pmf_sum}
  {context}
  {roughgarden-cc-bootcamp}
  {pdf-d54d402b16e2-p012-b002}
\proofstep
  {CommunicationComplexity.FiniteProbabilitySpace.sq_integral_le_integral_sq}
  {context}
  {roughgarden-cc-bootcamp}
  {pdf-d54d402b16e2-p012-b002}
\proofstep
  {CommunicationComplexity.FiniteProbabilitySpace.integral_comp_eval}
  {context}
  {roughgarden-cc-bootcamp}
  {pdf-d54d402b16e2-p012-b002}
\proofstep
  {CommunicationComplexity.FiniteProbabilitySpace.measureReal_pi_univ}
  {context}
  {roughgarden-cc-bootcamp}
  {pdf-d54d402b16e2-p012-b002}
\proofstep
  {CommunicationComplexity.FiniteProbabilitySpace.measureReal_eq_integral_indicator_one}
  {context}
  {roughgarden-cc-bootcamp}
  {pdf-d54d402b16e2-p012-b002}
\proofstep
  {CommunicationComplexity.FiniteProbabilitySpace.measureReal_eq_integral_indicator}
  {context}
  {roughgarden-cc-bootcamp}
  {pdf-d54d402b16e2-p012-b002}
\proofstep
  {CommunicationComplexity.FiniteProbabilitySpace.pmf_toReal_sum_eq_one}
  {context}
  {roughgarden-cc-bootcamp}
  {pdf-d54d402b16e2-p011-b003}
\proofstep
  {CommunicationComplexity.FiniteProbabilitySpace.pmf_toReal_nonneg}
  {context}
  {roughgarden-cc-bootcamp}
  {pdf-d54d402b16e2-p011-b003}
\proofstep
  {CommunicationComplexity.FiniteProbabilitySpace.exists_pmf_toReal_pos}
  {context}
  {roughgarden-cc-bootcamp}
  {pdf-d54d402b16e2-p012-b002}
\proofstep
  {CommunicationComplexity.FiniteProbabilitySpace.integral_le_of_le}
  {context}
  {roughgarden-cc-bootcamp}
  {pdf-d54d402b16e2-p012-b002}
\proofstep
  {CommunicationComplexity.FiniteProbabilitySpace.measureReal_ge_le_integral_div}
  {context}
  {roughgarden-cc-bootcamp}
  {pdf-d54d402b16e2-p012-b002}
\proofstep
  {CommunicationComplexity.FiniteProbabilitySpace.lt_integral_of_lt}
  {context}
  {roughgarden-cc-bootcamp}
  {pdf-d54d402b16e2-p012-b002}
\proofstep
  {CommunicationComplexity.CoinTape}
  {context}
  {roughgarden-cc-bootcamp}
  {pdf-d54d402b16e2-p011-b003,
    pdf-d54d402b16e2-p012-b002}
\proofstep
  {CommunicationComplexity.coinTapeMeasure}
  {context}
  {roughgarden-cc-bootcamp}
  {pdf-d54d402b16e2-p011-b003}
\proofstep
  {CommunicationComplexity.coinTapeIsProbabilityMeasure}
  {context}
  {roughgarden-cc-bootcamp}
  {pdf-d54d402b16e2-p011-b003}
\proofstep
  {CommunicationComplexity.coinTapeFiniteProbabilitySpace}
  {context}
  {roughgarden-cc-bootcamp}
  {pdf-d54d402b16e2-p011-b003}
\proofstep
  {CommunicationComplexity.PublicCoin.Protocol}
  {context}
  {roughgarden-cc-bootcamp}
  {pdf-d54d402b16e2-p011-b003,
    pdf-d54d402b16e2-p011-b005}
\proofstep
  {CommunicationComplexity.PublicCoin.Protocol.rrun}
  {context}
  {roughgarden-cc-bootcamp}
  {pdf-d54d402b16e2-p011-b003}
\proofstep
  {CommunicationComplexity.PublicCoin.Protocol.rrun_eq}
  {context}
  {roughgarden-cc-bootcamp}
  {pdf-d54d402b16e2-p011-b003}
\proofstep
  {CommunicationComplexity.PublicCoin.FiniteMessage.Protocol}
  {context}
  {roughgarden-cc-bootcamp}
  {pdf-d54d402b16e2-p011-b003}
\proofstep
  {CommunicationComplexity.PublicCoin.FiniteMessage.Protocol.rrun}
  {context}
  {roughgarden-cc-bootcamp}
  {pdf-d54d402b16e2-p011-b003}
\proofstep
  {CommunicationComplexity.PublicCoin.FiniteMessage.Protocol.rrun_eq}
  {context}
  {roughgarden-cc-bootcamp}
  {pdf-d54d402b16e2-p011-b003}
\proofstep
  {CommunicationComplexity.PublicCoin.FiniteMessage.Protocol.toProtocol}
  {context}
  {roughgarden-cc-bootcamp}
  {pdf-d54d402b16e2-p011-b003}
\proofstep
  {CommunicationComplexity.PublicCoin.FiniteMessage.Protocol.toProtocol_rrun}
  {context}
  {roughgarden-cc-bootcamp}
  {pdf-d54d402b16e2-p011-b003}
\proofstep
  {CommunicationComplexity.PublicCoin.FiniteMessage.Protocol.toProtocol_complexity}
  {context}
  {roughgarden-cc-bootcamp}
  {pdf-d54d402b16e2-p002-b002}
\proofstep
  {CommunicationComplexity.PublicCoin.FiniteMessage.Protocol.ofProtocol}
  {context}
  {roughgarden-cc-bootcamp}
  {pdf-d54d402b16e2-p011-b003}
\proofstep
  {CommunicationComplexity.PublicCoin.FiniteMessage.Protocol.ofProtocol_rrun}
  {context}
  {roughgarden-cc-bootcamp}
  {pdf-d54d402b16e2-p011-b003}
\proofstep
  {CommunicationComplexity.PublicCoin.FiniteMessage.Protocol.ofProtocol_complexity}
  {context}
  {roughgarden-cc-bootcamp}
  {pdf-d54d402b16e2-p002-b002}
\proofstep
  {CommunicationComplexity.PrivateCoin.Protocol}
  {context}
  {roughgarden-cc-bootcamp}
  {pdf-d54d402b16e2-p011-b003,
    pdf-d54d402b16e2-p011-b005}
\proofstep
  {CommunicationComplexity.PrivateCoin.Protocol.output}
  {context}
  {roughgarden-cc-bootcamp}
  {pdf-d54d402b16e2-p004-b005}
\proofstep
  {CommunicationComplexity.PrivateCoin.Protocol.alice}
  {context}
  {roughgarden-cc-bootcamp}
  {pdf-d54d402b16e2-p002-b002}
\proofstep
  {CommunicationComplexity.PrivateCoin.Protocol.bob}
  {context}
  {roughgarden-cc-bootcamp}
  {pdf-d54d402b16e2-p002-b002}
\proofstep
  {CommunicationComplexity.PrivateCoin.Protocol.rrun}
  {context}
  {roughgarden-cc-bootcamp}
  {pdf-d54d402b16e2-p011-b003}
\proofstep
  {CommunicationComplexity.PrivateCoin.Protocol.rrun_eq}
  {context}
  {roughgarden-cc-bootcamp}
  {pdf-d54d402b16e2-p011-b003}
\proofstep
  {CommunicationComplexity.PrivateCoin.Protocol.ApproxSatisfies}
  {context}
  {roughgarden-cc-bootcamp}
  {pdf-d54d402b16e2-p011-b003}
\proofstep
  {CommunicationComplexity.PrivateCoin.Protocol.ApproxComputes}
  {context}
  {roughgarden-cc-bootcamp}
  {pdf-d54d402b16e2-p011-b005}
\proofstep
  {CommunicationComplexity.PrivateCoin.Protocol.ApproxComputes_eq_ApproxSatisfies}
  {context}
  {roughgarden-cc-bootcamp}
  {pdf-d54d402b16e2-p011-b005}
\proofstep
  {CommunicationComplexity.PrivateCoin.FiniteMessage.Protocol}
  {context}
  {roughgarden-cc-bootcamp}
  {pdf-d54d402b16e2-p011-b003}
\proofstep
  {CommunicationComplexity.PrivateCoin.FiniteMessage.Protocol.output}
  {context}
  {roughgarden-cc-bootcamp}
  {pdf-d54d402b16e2-p004-b005}
\proofstep
  {CommunicationComplexity.PrivateCoin.FiniteMessage.Protocol.alice}
  {context}
  {roughgarden-cc-bootcamp}
  {pdf-d54d402b16e2-p002-b002}
\proofstep
  {CommunicationComplexity.PrivateCoin.FiniteMessage.Protocol.bob}
  {context}
  {roughgarden-cc-bootcamp}
  {pdf-d54d402b16e2-p002-b002}
\proofstep
  {CommunicationComplexity.PrivateCoin.FiniteMessage.Protocol.rrun}
  {context}
  {roughgarden-cc-bootcamp}
  {pdf-d54d402b16e2-p011-b003}
\proofstep
  {CommunicationComplexity.PrivateCoin.FiniteMessage.Protocol.rrun_eq}
  {context}
  {roughgarden-cc-bootcamp}
  {pdf-d54d402b16e2-p011-b003}
\proofstep
  {CommunicationComplexity.PrivateCoin.FiniteMessage.Protocol.ApproxSatisfies}
  {context}
  {roughgarden-cc-bootcamp}
  {pdf-d54d402b16e2-p011-b003}
\proofstep
  {CommunicationComplexity.PrivateCoin.FiniteMessage.Protocol.ApproxComputes}
  {context}
  {roughgarden-cc-bootcamp}
  {pdf-d54d402b16e2-p011-b005}
\proofstep
  {CommunicationComplexity.PrivateCoin.FiniteMessage.Protocol.ApproxComputes_eq_ApproxSatisfies}
  {context}
  {roughgarden-cc-bootcamp}
  {pdf-d54d402b16e2-p011-b005}
\proofstep
  {CommunicationComplexity.PrivateCoin.FiniteMessage.Protocol.toProtocol}
  {context}
  {roughgarden-cc-bootcamp}
  {pdf-d54d402b16e2-p011-b003}
\proofstep
  {CommunicationComplexity.PrivateCoin.FiniteMessage.Protocol.toProtocol_rrun}
  {context}
  {roughgarden-cc-bootcamp}
  {pdf-d54d402b16e2-p011-b003}
\proofstep
  {CommunicationComplexity.PrivateCoin.FiniteMessage.Protocol.toProtocol_complexity}
  {context}
  {roughgarden-cc-bootcamp}
  {pdf-d54d402b16e2-p002-b002}
\proofstep
  {CommunicationComplexity.PrivateCoin.FiniteMessage.Protocol.ofProtocol}
  {context}
  {roughgarden-cc-bootcamp}
  {pdf-d54d402b16e2-p011-b003}
\proofstep
  {CommunicationComplexity.PrivateCoin.FiniteMessage.Protocol.ofProtocol_rrun}
  {context}
  {roughgarden-cc-bootcamp}
  {pdf-d54d402b16e2-p011-b003}
\proofstep
  {CommunicationComplexity.PrivateCoin.FiniteMessage.Protocol.ofProtocol_complexity}
  {context}
  {roughgarden-cc-bootcamp}
  {pdf-d54d402b16e2-p002-b002}
\proofstep
  {CommunicationComplexity.PrivateCoin.FiniteMessage.Protocol.ofProtocol_equiv}
  {context}
  {roughgarden-cc-bootcamp}
  {pdf-d54d402b16e2-p002-b002}
\proofstep
  {CommunicationComplexity.Internal.cdf}
  {context}
  {roughgarden-cc-bootcamp}
  {pdf-d54d402b16e2-p012-b002}
\proofstep
  {CommunicationComplexity.Internal.cdf_zero}
  {context}
  {roughgarden-cc-bootcamp}
  {pdf-d54d402b16e2-p012-b002}
\proofstep
  {CommunicationComplexity.Internal.cdf_succ}
  {context}
  {roughgarden-cc-bootcamp}
  {pdf-d54d402b16e2-p012-b002}
\proofstep
  {CommunicationComplexity.Internal.cdf_one}
  {context}
  {roughgarden-cc-bootcamp}
  {pdf-d54d402b16e2-p012-b002}
\proofstep
  {CommunicationComplexity.Internal.cdf_mono}
  {context}
  {roughgarden-cc-bootcamp}
  {pdf-d54d402b16e2-p012-b002}
\proofstep
  {CommunicationComplexity.Internal.invCdf}
  {context}
  {roughgarden-cc-bootcamp}
  {pdf-d54d402b16e2-p012-b002}
\proofstep
  {CommunicationComplexity.Internal.invCdf_eq_iff}
  {context}
  {roughgarden-cc-bootcamp}
  {pdf-d54d402b16e2-p012-b002}
\proofstep
  {CommunicationComplexity.Internal.uniformApprox}
  {context}
  {roughgarden-cc-bootcamp}
  {pdf-d54d402b16e2-p012-b002}
\proofstep
  {CommunicationComplexity.Internal.card_nat_in_Ico}
  {context}
  {roughgarden-cc-bootcamp}
  {pdf-d54d402b16e2-p011-b003}
\proofstep
  {CommunicationComplexity.Internal.uniformApprox_approx}
  {context}
  {roughgarden-cc-bootcamp}
  {pdf-d54d402b16e2-p011-b003}
\proofstep
  {CommunicationComplexity.Internal.single_coin_approx}
  {context}
  {roughgarden-cc-bootcamp}
  {pdf-d54d402b16e2-p011-b003}
\proofstep
  {CommunicationComplexity.Internal.weighted_sum_approx}
  {context}
  {roughgarden-cc-bootcamp}
  {pdf-d54d402b16e2-p011-b003}
\proofstep
  {CommunicationComplexity.Internal.product_coin_approx}
  {context}
  {roughgarden-cc-bootcamp}
  {pdf-d54d402b16e2-p011-b003}
\proofstep
  {CommunicationComplexity.PrivateCoin.FiniteMessage.Protocol.toCoinTape}
  {context}
  {roughgarden-cc-bootcamp}
  {pdf-d54d402b16e2-p011-b003}
\proofstep
  {CommunicationComplexity.PrivateCoin.FiniteMessage.Protocol.toCoinTape_complexity}
  {context}
  {roughgarden-cc-bootcamp}
  {pdf-d54d402b16e2-p011-b003}
\proofstep
  {CommunicationComplexity.PublicCoin.FiniteMessage.Protocol.toCoinTape}
  {context}
  {roughgarden-cc-bootcamp}
  {pdf-d54d402b16e2-p011-b003}
\proofstep
  {CommunicationComplexity.PublicCoin.FiniteMessage.Protocol.toCoinTape_complexity}
  {context}
  {roughgarden-cc-bootcamp}
  {pdf-d54d402b16e2-p011-b003}
\proofstep
  {CommunicationComplexity.Deterministic.Protocol.Transcript}
  {context}
  {roughgarden-cc-bootcamp}
  {pdf-d54d402b16e2-p004-b005}
\proofstep
  {CommunicationComplexity.Deterministic.Protocol.Transcript.inputSet}
  {context}
  {roughgarden-cc-bootcamp}
  {pdf-d54d402b16e2-p004-b005}
\proofstep
  {CommunicationComplexity.Deterministic.Protocol.Transcript.inputSet_isRectangle}
  {direct}
  {roughgarden-cc-bootcamp}
  {pdf-d54d402b16e2-p007-b002,
    pdf-d54d402b16e2-p007-b004}
\proofstep
  {CommunicationComplexity.Deterministic.Protocol.transcriptFintype}
  {context}
  {roughgarden-cc-bootcamp}
  {pdf-d54d402b16e2-p008-b005}
\proofstep
  {CommunicationComplexity.Deterministic.Protocol.transcriptMeasurableSpace}
  {context}
  {roughgarden-cc-bootcamp}
  {pdf-d54d402b16e2-p004-b005}
\proofstep
  {CommunicationComplexity.Deterministic.Protocol.transcript}
  {context}
  {roughgarden-cc-bootcamp}
  {pdf-d54d402b16e2-p004-b005}
\proofstep
  {CommunicationComplexity.Deterministic.Protocol.mem_transcript}
  {context}
  {roughgarden-cc-bootcamp}
  {pdf-d54d402b16e2-p004-b005}
\proofstep
  {CommunicationComplexity.Deterministic.Protocol.transcript_eq_of_mem}
  {context}
  {roughgarden-cc-bootcamp}
  {pdf-d54d402b16e2-p004-b005}
\proofstep
  {CommunicationComplexity.Deterministic.Protocol.transcriptSwap}
  {context}
  {roughgarden-cc-bootcamp}
  {pdf-d54d402b16e2-p004-b005}
\proofstep
  {CommunicationComplexity.PublicCoin.Protocol.toDeterministic}
  {context}
  {roughgarden-cc-bootcamp}
  {pdf-d54d402b16e2-p011-b003}
\proofstep
  {CommunicationComplexity.PublicCoin.Protocol.toDeterministic_run}
  {context}
  {roughgarden-cc-bootcamp}
  {pdf-d54d402b16e2-p011-b003}
\proofstep
  {CommunicationComplexity.PublicCoin.Protocol.toDeterministic_complexity}
  {context}
  {roughgarden-cc-bootcamp}
  {pdf-d54d402b16e2-p011-b003}
\proofstep
  {CommunicationComplexity.Deterministic.FiniteMessage.Protocol.toPrivateCoin}
  {context}
  {roughgarden-cc-bootcamp}
  {pdf-d54d402b16e2-p011-b003}
\proofstep
  {CommunicationComplexity.Deterministic.FiniteMessage.Protocol.toPrivateCoin_rrun}
  {context}
  {roughgarden-cc-bootcamp}
  {pdf-d54d402b16e2-p011-b003}
\proofstep
  {CommunicationComplexity.Deterministic.FiniteMessage.Protocol.toPrivateCoin_complexity}
  {context}
  {roughgarden-cc-bootcamp}
  {pdf-d54d402b16e2-p011-b003}
\proofstep
  {CommunicationComplexity.PublicCoin.FiniteMessage.Protocol.toDeterministic}
  {context}
  {roughgarden-cc-bootcamp}
  {pdf-d54d402b16e2-p011-b003}
\proofstep
  {CommunicationComplexity.PublicCoin.FiniteMessage.Protocol.toDeterministic_run}
  {context}
  {roughgarden-cc-bootcamp}
  {pdf-d54d402b16e2-p011-b003}
\proofstep
  {CommunicationComplexity.PublicCoin.FiniteMessage.Protocol.toDeterministic_complexity}
  {context}
  {roughgarden-cc-bootcamp}
  {pdf-d54d402b16e2-p011-b003}
\proofstep
  {CommunicationComplexity.Functions.Disjointness.RandomizedLowerBound.setFinFiniteMeasureSpace}
  {context}
  {roughgarden-cc-bootcamp}
  {pdf-d54d402b16e2-p016-b002}
\proofstep
  {CommunicationComplexity.Functions.Disjointness.RandomizedLowerBound.DisjointCoordinate}
  {context}
  {roughgarden-cc-bootcamp}
  {pdf-d54d402b16e2-p016-b002}
\proofstep
  {CommunicationComplexity.Functions.Disjointness.RandomizedLowerBound.disjointCoordinateNonempty}
  {context}
  {roughgarden-cc-bootcamp}
  {pdf-d54d402b16e2-p016-b002}
\proofstep
  {CommunicationComplexity.Functions.Disjointness.RandomizedLowerBound.disjointCoordinateMeasurableSpace}
  {context}
  {roughgarden-cc-bootcamp}
  {pdf-d54d402b16e2-p016-b002}
\proofstep
  {CommunicationComplexity.Functions.Disjointness.RandomizedLowerBound.disjointCoordinateDiscreteMeasurableSpace}
  {context}
  {roughgarden-cc-bootcamp}
  {pdf-d54d402b16e2-p016-b002}
\proofstep
  {CommunicationComplexity.Functions.Disjointness.RandomizedLowerBound.disjointCoordinateFiniteMeasureSpace}
  {context}
  {roughgarden-cc-bootcamp}
  {pdf-d54d402b16e2-p016-b002}
\proofstep
  {CommunicationComplexity.Functions.Disjointness.RandomizedLowerBound.DisjointCoordinate.xBit}
  {context}
  {roughgarden-cc-bootcamp}
  {pdf-d54d402b16e2-p016-b002}
\proofstep
  {CommunicationComplexity.Functions.Disjointness.RandomizedLowerBound.DisjointCoordinate.yBit}
  {context}
  {roughgarden-cc-bootcamp}
  {pdf-d54d402b16e2-p016-b002}
\proofstep
  {CommunicationComplexity.Functions.Disjointness.RandomizedLowerBound.DisjointCoordinate.card}
  {context}
  {roughgarden-cc-bootcamp}
  {pdf-d54d402b16e2-p016-b002}
\proofstep
  {CommunicationComplexity.Functions.Disjointness.RandomizedLowerBound.DisjointCoordinate.not_xBit_and_yBit}
  {context}
  {roughgarden-cc-bootcamp}
  {pdf-d54d402b16e2-p016-b002}
\proofstep
  {CommunicationComplexity.Functions.Disjointness.RandomizedLowerBound.HardSample}
  {context}
  {roughgarden-cc-bootcamp}
  {pdf-d54d402b16e2-p016-b002}
\proofstep
  {CommunicationComplexity.Functions.Disjointness.RandomizedLowerBound.HardSample.instNonempty}
  {context}
  {roughgarden-cc-bootcamp}
  {pdf-d54d402b16e2-p016-b002}
\proofstep
  {CommunicationComplexity.Functions.Disjointness.RandomizedLowerBound.HardSample.instMeasurableSpace}
  {context}
  {roughgarden-cc-bootcamp}
  {pdf-d54d402b16e2-p016-b002}
\proofstep
  {CommunicationComplexity.Functions.Disjointness.RandomizedLowerBound.HardSample.measureSpace}
  {context}
  {roughgarden-cc-bootcamp}
  {pdf-d54d402b16e2-p016-b002}
\proofstep
  {CommunicationComplexity.Functions.Disjointness.RandomizedLowerBound.HardSample.isProbabilityMeasure}
  {context}
  {roughgarden-cc-bootcamp}
  {pdf-d54d402b16e2-p016-b002}
\proofstep
  {CommunicationComplexity.Functions.Disjointness.RandomizedLowerBound.HardSample.finiteProbabilitySpace}
  {context}
  {roughgarden-cc-bootcamp}
  {pdf-d54d402b16e2-p016-b002}
\proofstep
  {CommunicationComplexity.Functions.Disjointness.RandomizedLowerBound.HardSample.equivProd}
  {context}
  {roughgarden-cc-bootcamp}
  {pdf-d54d402b16e2-p016-b002}
\proofstep
  {CommunicationComplexity.Functions.Disjointness.RandomizedLowerBound.HardSample.card}
  {context}
  {roughgarden-cc-bootcamp}
  {pdf-d54d402b16e2-p016-b002}
\proofstep
  {CommunicationComplexity.Functions.Disjointness.RandomizedLowerBound.uniformFin}
  {context}
  {roughgarden-cc-bootcamp}
  {pdf-d54d402b16e2-p016-b002}
\proofstep
  {CommunicationComplexity.Functions.Disjointness.RandomizedLowerBound.uniformFin_isProbabilityMeasure}
  {context}
  {roughgarden-cc-bootcamp}
  {pdf-d54d402b16e2-p016-b002}
\proofstep
  {CommunicationComplexity.Functions.Disjointness.RandomizedLowerBound.uniformFin_real}
  {context}
  {roughgarden-cc-bootcamp}
  {pdf-d54d402b16e2-p016-b002}
\proofstep
  {CommunicationComplexity.Functions.Disjointness.RandomizedLowerBound.uniformFin_univ}
  {context}
  {roughgarden-cc-bootcamp}
  {pdf-d54d402b16e2-p016-b002}
\proofstep
  {CommunicationComplexity.Functions.Disjointness.RandomizedLowerBound.uniformBool}
  {context}
  {roughgarden-cc-bootcamp}
  {pdf-d54d402b16e2-p016-b002}
\proofstep
  {CommunicationComplexity.Functions.Disjointness.RandomizedLowerBound.uniformBool_isProbabilityMeasure}
  {context}
  {roughgarden-cc-bootcamp}
  {pdf-d54d402b16e2-p016-b002}
\proofstep
  {CommunicationComplexity.Functions.Disjointness.RandomizedLowerBound.uniformBool_singleton}
  {context}
  {roughgarden-cc-bootcamp}
  {pdf-d54d402b16e2-p016-b002}
\proofstep
  {CommunicationComplexity.Functions.Disjointness.RandomizedLowerBound.uniformBool_real}
  {context}
  {roughgarden-cc-bootcamp}
  {pdf-d54d402b16e2-p016-b002}
\proofstep
  {CommunicationComplexity.Functions.Disjointness.RandomizedLowerBound.uniformDisjointCoordinateVector}
  {context}
  {roughgarden-cc-bootcamp}
  {pdf-d54d402b16e2-p016-b002}
\proofstep
  {CommunicationComplexity.Functions.Disjointness.RandomizedLowerBound.uniformDisjointCoordinateVector_isProbabilityMeasure}
  {context}
  {roughgarden-cc-bootcamp}
  {pdf-d54d402b16e2-p016-b002}
\proofstep
  {CommunicationComplexity.Functions.Disjointness.RandomizedLowerBound.uniformDisjointCoordinate}
  {context}
  {roughgarden-cc-bootcamp}
  {pdf-d54d402b16e2-p016-b002}
\proofstep
  {CommunicationComplexity.Functions.Disjointness.RandomizedLowerBound.uniformDisjointCoordinate_isProbabilityMeasure}
  {context}
  {roughgarden-cc-bootcamp}
  {pdf-d54d402b16e2-p016-b002}
\proofstep
  {CommunicationComplexity.Functions.Disjointness.RandomizedLowerBound.uniformDisjointCoordinateVector_real}
  {context}
  {roughgarden-cc-bootcamp}
  {pdf-d54d402b16e2-p016-b002}
\proofstep
  {CommunicationComplexity.Functions.Disjointness.RandomizedLowerBound.uniformDisjointCoordinateVector_univ}
  {context}
  {roughgarden-cc-bootcamp}
  {pdf-d54d402b16e2-p016-b002}
\proofstep
  {CommunicationComplexity.Functions.Disjointness.RandomizedLowerBound.xBit}
  {context}
  {roughgarden-cc-bootcamp}
  {pdf-d54d402b16e2-p016-b002}
\proofstep
  {CommunicationComplexity.Functions.Disjointness.RandomizedLowerBound.yBit}
  {context}
  {roughgarden-cc-bootcamp}
  {pdf-d54d402b16e2-p016-b002}
\proofstep
  {CommunicationComplexity.Functions.Disjointness.RandomizedLowerBound.X}
  {context}
  {roughgarden-cc-bootcamp}
  {pdf-d54d402b16e2-p016-b002}
\proofstep
  {CommunicationComplexity.Functions.Disjointness.RandomizedLowerBound.Y}
  {context}
  {roughgarden-cc-bootcamp}
  {pdf-d54d402b16e2-p016-b002}
\proofstep
  {CommunicationComplexity.Functions.Disjointness.RandomizedLowerBound.input}
  {context}
  {roughgarden-cc-bootcamp}
  {pdf-d54d402b16e2-p016-b002}
\proofstep
  {CommunicationComplexity.Functions.Disjointness.RandomizedLowerBound.specialIntersect}
  {context}
  {roughgarden-cc-bootcamp}
  {pdf-d54d402b16e2-p016-b002}
\proofstep
  {CommunicationComplexity.Functions.Disjointness.RandomizedLowerBound.specialBitsEvent}
  {context}
  {roughgarden-cc-bootcamp}
  {pdf-d54d402b16e2-p016-b002}
\proofstep
  {CommunicationComplexity.Functions.Disjointness.RandomizedLowerBound.specialCoordinateEvent}
  {context}
  {roughgarden-cc-bootcamp}
  {pdf-d54d402b16e2-p016-b002}
\proofstep
  {CommunicationComplexity.Functions.Disjointness.RandomizedLowerBound.disjointEvent}
  {context}
  {roughgarden-cc-bootcamp}
  {pdf-d54d402b16e2-p016-b002}
\proofstep
  {CommunicationComplexity.Functions.Disjointness.RandomizedLowerBound.specialIntersectFintype}
  {context}
  {roughgarden-cc-bootcamp}
  {pdf-d54d402b16e2-p016-b002}
\proofstep
  {CommunicationComplexity.Functions.Disjointness.RandomizedLowerBound.specialBitsEventFintype}
  {context}
  {roughgarden-cc-bootcamp}
  {pdf-d54d402b16e2-p016-b002}
\proofstep
  {CommunicationComplexity.Functions.Disjointness.RandomizedLowerBound.specialCoordinateEventFintype}
  {context}
  {roughgarden-cc-bootcamp}
  {pdf-d54d402b16e2-p016-b002}
\proofstep
  {CommunicationComplexity.Functions.Disjointness.RandomizedLowerBound.specialCoordinate}
  {context}
  {roughgarden-cc-bootcamp}
  {pdf-d54d402b16e2-p016-b002}
\proofstep
  {CommunicationComplexity.Functions.Disjointness.RandomizedLowerBound.specialX}
  {context}
  {roughgarden-cc-bootcamp}
  {pdf-d54d402b16e2-p016-b002}
\proofstep
  {CommunicationComplexity.Functions.Disjointness.RandomizedLowerBound.specialY}
  {context}
  {roughgarden-cc-bootcamp}
  {pdf-d54d402b16e2-p016-b002}
\proofstep
  {CommunicationComplexity.Functions.Disjointness.RandomizedLowerBound.xBeforeSpecial}
  {context}
  {roughgarden-cc-bootcamp}
  {pdf-d54d402b16e2-p017-b001}
\proofstep
  {CommunicationComplexity.Functions.Disjointness.RandomizedLowerBound.yAfterSpecial}
  {context}
  {roughgarden-cc-bootcamp}
  {pdf-d54d402b16e2-p017-b001}
\proofstep
  {CommunicationComplexity.Functions.Disjointness.RandomizedLowerBound.ProtocolType}
  {context}
  {roughgarden-cc-bootcamp}
  {pdf-d54d402b16e2-p014-b006}
\proofstep
  {CommunicationComplexity.Functions.Disjointness.RandomizedLowerBound.TranscriptType}
  {context}
  {roughgarden-cc-bootcamp}
  {pdf-d54d402b16e2-p007-b002}
\proofstep
  {CommunicationComplexity.Functions.Disjointness.RandomizedLowerBound.RawZType}
  {context}
  {roughgarden-cc-bootcamp}
  {pdf-d54d402b16e2-p017-b001}
\proofstep
  {CommunicationComplexity.Functions.Disjointness.RandomizedLowerBound.RawZType.ext}
  {context}
  {roughgarden-cc-bootcamp}
  {pdf-d54d402b16e2-p017-b001}
\proofstep
  {CommunicationComplexity.Functions.Disjointness.RandomizedLowerBound.RawZType.equivProd}
  {context}
  {roughgarden-cc-bootcamp}
  {pdf-d54d402b16e2-p017-b001}
\proofstep
  {CommunicationComplexity.Functions.Disjointness.RandomizedLowerBound.rawZTypeFintype}
  {context}
  {roughgarden-cc-bootcamp}
  {pdf-d54d402b16e2-p017-b001}
\proofstep
  {CommunicationComplexity.Functions.Disjointness.RandomizedLowerBound.rawZTypeDecidableEq}
  {context}
  {roughgarden-cc-bootcamp}
  {pdf-d54d402b16e2-p017-b001}
\proofstep
  {CommunicationComplexity.Functions.Disjointness.RandomizedLowerBound.rawZVariable}
  {context}
  {roughgarden-cc-bootcamp}
  {pdf-d54d402b16e2-p017-b001}
\proofstep
  {CommunicationComplexity.Functions.Disjointness.RandomizedLowerBound.ZType}
  {context}
  {roughgarden-cc-bootcamp}
  {pdf-d54d402b16e2-p017-b001}
\proofstep
  {CommunicationComplexity.Functions.Disjointness.RandomizedLowerBound.ZType.transcript}
  {context}
  {roughgarden-cc-bootcamp}
  {pdf-d54d402b16e2-p007-b002}
\proofstep
  {CommunicationComplexity.Functions.Disjointness.RandomizedLowerBound.ZType.specialCoordinate}
  {context}
  {roughgarden-cc-bootcamp}
  {pdf-d54d402b16e2-p017-b001}
\proofstep
  {CommunicationComplexity.Functions.Disjointness.RandomizedLowerBound.ZType.xBefore}
  {context}
  {roughgarden-cc-bootcamp}
  {pdf-d54d402b16e2-p017-b001}
\proofstep
  {CommunicationComplexity.Functions.Disjointness.RandomizedLowerBound.ZType.yAfter}
  {context}
  {roughgarden-cc-bootcamp}
  {pdf-d54d402b16e2-p017-b001}
\proofstep
  {CommunicationComplexity.Functions.Disjointness.RandomizedLowerBound.ZType.raw}
  {context}
  {roughgarden-cc-bootcamp}
  {pdf-d54d402b16e2-p017-b001}
\proofstep
  {CommunicationComplexity.Functions.Disjointness.RandomizedLowerBound.ZType.achievable}
  {context}
  {roughgarden-cc-bootcamp}
  {pdf-d54d402b16e2-p017-b001}
\proofstep
  {CommunicationComplexity.Functions.Disjointness.RandomizedLowerBound.ZType.ext}
  {context}
  {roughgarden-cc-bootcamp}
  {pdf-d54d402b16e2-p017-b001}
\proofstep
  {CommunicationComplexity.Functions.Disjointness.RandomizedLowerBound.zTypeFintype}
  {context}
  {roughgarden-cc-bootcamp}
  {pdf-d54d402b16e2-p017-b001}
\proofstep
  {CommunicationComplexity.Functions.Disjointness.RandomizedLowerBound.zTypeMeasurableSpace}
  {context}
  {roughgarden-cc-bootcamp}
  {pdf-d54d402b16e2-p017-b001}
\proofstep
  {CommunicationComplexity.Functions.Disjointness.RandomizedLowerBound.zVariable}
  {context}
  {roughgarden-cc-bootcamp}
  {pdf-d54d402b16e2-p017-b001}
\proofstep
  {CommunicationComplexity.Functions.Disjointness.RandomizedLowerBound.zFiber}
  {context}
  {roughgarden-cc-bootcamp}
  {pdf-d54d402b16e2-p017-b001}
\proofstep
  {CommunicationComplexity.Functions.Disjointness.RandomizedLowerBound.volume_zFiber_ne_zero}
  {context}
  {roughgarden-cc-bootcamp}
  {pdf-d54d402b16e2-p017-b001}
\proofstep
  {CommunicationComplexity.Functions.Disjointness.RandomizedLowerBound.inputMeasure}
  {context}
  {roughgarden-cc-bootcamp}
  {pdf-d54d402b16e2-p016-b002}
\proofstep
  {CommunicationComplexity.Functions.Disjointness.RandomizedLowerBound.inputMeasure_isProbabilityMeasure}
  {context}
  {roughgarden-cc-bootcamp}
  {pdf-d54d402b16e2-p016-b002}
\proofstep
  {CommunicationComplexity.Functions.Disjointness.RandomizedLowerBound.xBit_of_ne_special}
  {context}
  {roughgarden-cc-bootcamp}
  {pdf-d54d402b16e2-p016-b002}
\proofstep
  {CommunicationComplexity.Functions.Disjointness.RandomizedLowerBound.yBit_of_ne_special}
  {context}
  {roughgarden-cc-bootcamp}
  {pdf-d54d402b16e2-p016-b002}
\proofstep
  {CommunicationComplexity.Functions.Disjointness.RandomizedLowerBound.not_xBit_and_yBit_of_ne_special}
  {context}
  {roughgarden-cc-bootcamp}
  {pdf-d54d402b16e2-p016-b002}
\proofstep
  {CommunicationComplexity.Functions.Disjointness.RandomizedLowerBound.coordinateOfBits}
  {context}
  {roughgarden-cc-bootcamp}
  {pdf-d54d402b16e2-p016-b002}
\proofstep
  {CommunicationComplexity.Functions.Disjointness.RandomizedLowerBound.coordinateOfBits_xBit}
  {context}
  {roughgarden-cc-bootcamp}
  {pdf-d54d402b16e2-p016-b002}
\proofstep
  {CommunicationComplexity.Functions.Disjointness.RandomizedLowerBound.coordinateOfBits_yBit}
  {context}
  {roughgarden-cc-bootcamp}
  {pdf-d54d402b16e2-p016-b002}
\proofstep
  {CommunicationComplexity.Functions.Disjointness.RandomizedLowerBound.coordinateOfBits_xBit_yBit}
  {context}
  {roughgarden-cc-bootcamp}
  {pdf-d54d402b16e2-p016-b002}
\proofstep
  {CommunicationComplexity.Functions.Disjointness.RandomizedLowerBound.mix}
  {context}
  {roughgarden-cc-bootcamp}
  {pdf-d54d402b16e2-p007-b003}
\proofstep
  {CommunicationComplexity.Functions.Disjointness.RandomizedLowerBound.disjoint_X_Y_iff}
  {context}
  {roughgarden-cc-bootcamp}
  {pdf-d54d402b16e2-p016-b002}
\proofstep
  {CommunicationComplexity.Functions.Disjointness.RandomizedLowerBound.disjointCoordinateVector}
  {context}
  {roughgarden-cc-bootcamp}
  {pdf-d54d402b16e2-p016-b002}
\proofstep
  {CommunicationComplexity.Functions.Disjointness.RandomizedLowerBound.ignoredCoordinate}
  {context}
  {roughgarden-cc-bootcamp}
  {pdf-d54d402b16e2-p016-b002}
\proofstep
  {CommunicationComplexity.Functions.Disjointness.RandomizedLowerBound.disjointModel}
  {context}
  {roughgarden-cc-bootcamp}
  {pdf-d54d402b16e2-p016-b002}
\proofstep
  {CommunicationComplexity.Functions.Disjointness.RandomizedLowerBound.disjointEventEquiv}
  {context}
  {roughgarden-cc-bootcamp}
  {pdf-d54d402b16e2-p016-b002}
\proofstep
  {CommunicationComplexity.Functions.Disjointness.RandomizedLowerBound.disjointEventFintype}
  {context}
  {roughgarden-cc-bootcamp}
  {pdf-d54d402b16e2-p016-b002}
\proofstep
  {CommunicationComplexity.Functions.Disjointness.RandomizedLowerBound.disjointEventInterDisjointModelEquiv}
  {context}
  {roughgarden-cc-bootcamp}
  {pdf-d54d402b16e2-p016-b002}
\proofstep
  {CommunicationComplexity.Functions.Disjointness.RandomizedLowerBound.disjointEventInterDisjointModelFintype}
  {context}
  {roughgarden-cc-bootcamp}
  {pdf-d54d402b16e2-p016-b002}
\proofstep
  {CommunicationComplexity.Functions.Disjointness.RandomizedLowerBound.not_xBit_left_and_yBit_right_of_same_conditioning}
  {context}
  {roughgarden-cc-bootcamp}
  {pdf-d54d402b16e2-p016-b002,
    pdf-d54d402b16e2-p007-b003}
\proofstep
  {CommunicationComplexity.Functions.Disjointness.RandomizedLowerBound.xBit_mix}
  {context}
  {roughgarden-cc-bootcamp}
  {pdf-d54d402b16e2-p007-b003}
\proofstep
  {CommunicationComplexity.Functions.Disjointness.RandomizedLowerBound.yBit_mix}
  {context}
  {roughgarden-cc-bootcamp}
  {pdf-d54d402b16e2-p007-b003}
\proofstep
  {CommunicationComplexity.Functions.Disjointness.RandomizedLowerBound.X_mix}
  {context}
  {roughgarden-cc-bootcamp}
  {pdf-d54d402b16e2-p007-b003}
\proofstep
  {CommunicationComplexity.Functions.Disjointness.RandomizedLowerBound.Y_mix}
  {context}
  {roughgarden-cc-bootcamp}
  {pdf-d54d402b16e2-p007-b003}
\proofstep
  {CommunicationComplexity.Functions.Disjointness.RandomizedLowerBound.input_mix}
  {context}
  {roughgarden-cc-bootcamp}
  {pdf-d54d402b16e2-p007-b003}
\proofstep
  {CommunicationComplexity.Functions.Disjointness.RandomizedLowerBound.xBeforeSpecial_mix}
  {context}
  {roughgarden-cc-bootcamp}
  {pdf-d54d402b16e2-p007-b003}
\proofstep
  {CommunicationComplexity.Functions.Disjointness.RandomizedLowerBound.yAfterSpecial_mix}
  {context}
  {roughgarden-cc-bootcamp}
  {pdf-d54d402b16e2-p007-b003}
\proofstep
  {CommunicationComplexity.Functions.Disjointness.RandomizedLowerBound.specialX_mix}
  {context}
  {roughgarden-cc-bootcamp}
  {pdf-d54d402b16e2-p007-b003}
\proofstep
  {CommunicationComplexity.Functions.Disjointness.RandomizedLowerBound.specialY_mix}
  {context}
  {roughgarden-cc-bootcamp}
  {pdf-d54d402b16e2-p007-b003}
\proofstep
  {CommunicationComplexity.Functions.Disjointness.RandomizedLowerBound.mix_mix_swap}
  {context}
  {roughgarden-cc-bootcamp}
  {pdf-d54d402b16e2-p007-b003}
\proofstep
  {CommunicationComplexity.Functions.Disjointness.RandomizedLowerBound.hardSample_measureReal_fieldProduct}
  {context}
  {roughgarden-cc-bootcamp}
  {pdf-d54d402b16e2-p016-b002}
\proofstep
  {CommunicationComplexity.Functions.Disjointness.RandomizedLowerBound.measureReal_specialIntersect}
  {direct}
  {roughgarden-cc-bootcamp}
  {pdf-d54d402b16e2-p016-b002,
    pdf-d54d402b16e2-p015-b002}
\proofstep
  {CommunicationComplexity.Functions.Disjointness.RandomizedLowerBound.disjointEvent_eq_compl_specialIntersect}
  {context}
  {roughgarden-cc-bootcamp}
  {pdf-d54d402b16e2-p016-b002}
\proofstep
  {CommunicationComplexity.Functions.Disjointness.RandomizedLowerBound.measureReal_disjointEvent}
  {direct}
  {roughgarden-cc-bootcamp}
  {pdf-d54d402b16e2-p016-b002}
\proofstep
  {CommunicationComplexity.Functions.Disjointness.RandomizedLowerBound.measure_disjointEvent_ne_zero}
  {context}
  {roughgarden-cc-bootcamp}
  {pdf-d54d402b16e2-p016-b002}
\proofstep
  {CommunicationComplexity.Functions.Disjointness.RandomizedLowerBound.disjointCondMeasure}
  {context}
  {roughgarden-cc-bootcamp}
  {pdf-d54d402b16e2-p016-b002}
\proofstep
  {CommunicationComplexity.Functions.Disjointness.RandomizedLowerBound.disjointCondMeasure_isProbabilityMeasure}
  {context}
  {roughgarden-cc-bootcamp}
  {pdf-d54d402b16e2-p016-b002}
\proofstep
  {CommunicationComplexity.Functions.Disjointness.RandomizedLowerBound.message}
  {context}
  {roughgarden-cc-bootcamp}
  {pdf-d54d402b16e2-p004-b005}
\proofstep
  {CommunicationComplexity.Functions.Disjointness.RandomizedLowerBound.specialPair}
  {context}
  {roughgarden-cc-bootcamp}
  {pdf-d54d402b16e2-p016-b002}
\proofstep
  {CommunicationComplexity.Functions.Disjointness.RandomizedLowerBound.uniformBoolPair}
  {context}
  {roughgarden-cc-bootcamp}
  {pdf-d54d402b16e2-p017-b001}
\proofstep
  {CommunicationComplexity.Functions.Disjointness.RandomizedLowerBound.uniformBoolPair_isProbabilityMeasure}
  {context}
  {roughgarden-cc-bootcamp}
  {pdf-d54d402b16e2-p017-b001}
\proofstep
  {CommunicationComplexity.Functions.Disjointness.RandomizedLowerBound.zFiberMeasure}
  {context}
  {roughgarden-cc-bootcamp}
  {pdf-d54d402b16e2-p017-b002}
\proofstep
  {CommunicationComplexity.Functions.Disjointness.RandomizedLowerBound.zFiberMeasure_isProbabilityMeasure}
  {context}
  {roughgarden-cc-bootcamp}
  {pdf-d54d402b16e2-p017-b002}
\proofstep
  {CommunicationComplexity.Functions.Disjointness.RandomizedLowerBound.conditionalSpecialPairLaw}
  {context}
  {roughgarden-cc-bootcamp}
  {pdf-d54d402b16e2-p017-b001,
    pdf-d54d402b16e2-p017-b005}
\proofstep
  {CommunicationComplexity.Functions.Disjointness.RandomizedLowerBound.conditionalSpecialPairLaw_isProbabilityMeasure}
  {context}
  {roughgarden-cc-bootcamp}
  {pdf-d54d402b16e2-p017-b001}
\proofstep
  {CommunicationComplexity.Functions.Disjointness.RandomizedLowerBound.conditionalSpecialXLaw}
  {context}
  {roughgarden-cc-bootcamp}
  {pdf-d54d402b16e2-p017-b001}
\proofstep
  {CommunicationComplexity.Functions.Disjointness.RandomizedLowerBound.conditionalSpecialXLaw_isProbabilityMeasure}
  {context}
  {roughgarden-cc-bootcamp}
  {pdf-d54d402b16e2-p017-b001}
\proofstep
  {CommunicationComplexity.Functions.Disjointness.RandomizedLowerBound.conditionalSpecialYLaw}
  {context}
  {roughgarden-cc-bootcamp}
  {pdf-d54d402b16e2-p017-b001}
\proofstep
  {CommunicationComplexity.Functions.Disjointness.RandomizedLowerBound.conditionalSpecialYLaw_isProbabilityMeasure}
  {context}
  {roughgarden-cc-bootcamp}
  {pdf-d54d402b16e2-p017-b001}
\proofstep
  {CommunicationComplexity.Functions.Disjointness.RandomizedLowerBound.input_mem_transcript}
  {context}
  {roughgarden-cc-bootcamp}
  {pdf-d54d402b16e2-p004-b005}
\proofstep
  {CommunicationComplexity.Functions.Disjointness.RandomizedLowerBound.input_mem_transcript_of_zVariable_eq}
  {context}
  {roughgarden-cc-bootcamp}
  {pdf-d54d402b16e2-p007-b002}
\proofstep
  {CommunicationComplexity.Functions.Disjointness.RandomizedLowerBound.specialCoordinate_eq_of_zVariable_eq}
  {context}
  {roughgarden-cc-bootcamp}
  {pdf-d54d402b16e2-p017-b002}
\proofstep
  {CommunicationComplexity.Functions.Disjointness.RandomizedLowerBound.xBeforeSpecial_eq_of_zVariable_eq}
  {context}
  {roughgarden-cc-bootcamp}
  {pdf-d54d402b16e2-p017-b002}
\proofstep
  {CommunicationComplexity.Functions.Disjointness.RandomizedLowerBound.yAfterSpecial_eq_of_zVariable_eq}
  {context}
  {roughgarden-cc-bootcamp}
  {pdf-d54d402b16e2-p017-b002}
\proofstep
  {CommunicationComplexity.Functions.Disjointness.RandomizedLowerBound.mixed_inputs_mem_transcript_of_zVariable_eq}
  {direct}
  {roughgarden-cc-bootcamp}
  {pdf-d54d402b16e2-p007-b003,
    pdf-d54d402b16e2-p007-b002}
\proofstep
  {CommunicationComplexity.Functions.Disjointness.RandomizedLowerBound.specialCoordinate_eq_of_same_zVariable}
  {context}
  {roughgarden-cc-bootcamp}
  {pdf-d54d402b16e2-p007-b003}
\proofstep
  {CommunicationComplexity.Functions.Disjointness.RandomizedLowerBound.xBeforeSpecial_eq_of_same_zVariable}
  {context}
  {roughgarden-cc-bootcamp}
  {pdf-d54d402b16e2-p007-b003}
\proofstep
  {CommunicationComplexity.Functions.Disjointness.RandomizedLowerBound.yAfterSpecial_eq_of_same_zVariable}
  {context}
  {roughgarden-cc-bootcamp}
  {pdf-d54d402b16e2-p007-b003}
\proofstep
  {CommunicationComplexity.Functions.Disjointness.RandomizedLowerBound.zVariable_mix_eq_of_same_zVariable}
  {context}
  {roughgarden-cc-bootcamp}
  {pdf-d54d402b16e2-p007-b003,
    pdf-d54d402b16e2-p007-b002,
    pdf-d54d402b16e2-p007-b004}
\proofstep
  {CommunicationComplexity.Functions.Disjointness.RandomizedLowerBound.mix_mem_fiber_inter_specialPair_of_mem_specialX_specialY}
  {context}
  {roughgarden-cc-bootcamp}
  {pdf-d54d402b16e2-p007-b003,
    pdf-d54d402b16e2-p007-b002}
\proofstep
  {CommunicationComplexity.Functions.Disjointness.RandomizedLowerBound.mix_swap_mem_fiber_of_mem_specialX_specialY}
  {context}
  {roughgarden-cc-bootcamp}
  {pdf-d54d402b16e2-p007-b003,
    pdf-d54d402b16e2-p007-b002}
\proofstep
  {CommunicationComplexity.Functions.Disjointness.RandomizedLowerBound.mix_mem_fiber_inter_specialX_of_mem_specialPair_fiber}
  {context}
  {roughgarden-cc-bootcamp}
  {pdf-d54d402b16e2-p007-b003,
    pdf-d54d402b16e2-p007-b002}
\proofstep
  {CommunicationComplexity.Functions.Disjointness.RandomizedLowerBound.mix_mem_fiber_inter_specialY_of_mem_fiber_specialPair}
  {context}
  {roughgarden-cc-bootcamp}
  {pdf-d54d402b16e2-p007-b003,
    pdf-d54d402b16e2-p007-b002}
\proofstep
  {CommunicationComplexity.Functions.Disjointness.RandomizedLowerBound.card_fiber_inter_specialX_mul_card_fiber_inter_specialY}
  {context}
  {roughgarden-cc-bootcamp}
  {pdf-d54d402b16e2-p007-b002,
    pdf-d54d402b16e2-p007-b003,
    pdf-d54d402b16e2-p017-b002}
\proofstep
  {CommunicationComplexity.Functions.Disjointness.RandomizedLowerBound.measureReal_eq_card_subtype_div}
  {context}
  {roughgarden-cc-bootcamp}
  {pdf-d54d402b16e2-p016-b005,
    pdf-d54d402b16e2-p015-b002}
\proofstep
  {CommunicationComplexity.Functions.Disjointness.RandomizedLowerBound.conditionalSpecialPairLaw_eq_prod_of_singleton_factorization}
  {context}
  {roughgarden-cc-bootcamp}
  {pdf-d54d402b16e2-p007-b002,
    pdf-d54d402b16e2-p015-b003}
\uses{CommunicationComplexity.Functions.Disjointness.RandomizedLowerBound.HardSample, CommunicationComplexity.Functions.Disjointness.RandomizedLowerBound.ProtocolType, CommunicationComplexity.Functions.Disjointness.RandomizedLowerBound.ZType, CommunicationComplexity.Functions.Disjointness.RandomizedLowerBound.card_fiber_inter_specialX_mul_card_fiber_inter_specialY, CommunicationComplexity.Functions.Disjointness.RandomizedLowerBound.measureReal_eq_card_subtype_div, CommunicationComplexity.Functions.Disjointness.RandomizedLowerBound.specialPair, CommunicationComplexity.Functions.Disjointness.RandomizedLowerBound.specialX, CommunicationComplexity.Functions.Disjointness.RandomizedLowerBound.specialY, CommunicationComplexity.Functions.Disjointness.RandomizedLowerBound.zFiber}
\difficulty{4}
Fix a deterministic disjointness protocol $p$ on inputs of length $n$, an achievable
value $z$ of the variable $Z = (M, T, X_{<T}, Y_{>T})$, and a pair of bits $b = (b_1,
b_2)$. Write $\mu$ for the uniform measure on the sample space of the hard distribution,
and let $F$ be the fiber of $Z$ above $z$, that is, the set of samples whose $Z$-value
equals $z$. Denote by $X_T$ and $Y_T$ Alice's and Bob's bits at the special coordinate
$T$. Then
\[
\mu(F)\cdot\mu\!\big(F \cap \{(X_T, Y_T) = (b_1, b_2)\}\big) \;=\; \mu\!\big(F \cap
\{X_T = b_1\}\big)\cdot\mu\!\big(F \cap \{Y_T = b_2\}\big).
\]
\end{theorem}

\begin{theorem}[Product-measure criterion for the conditional special-pair law]
\lean{CommunicationComplexity.Functions.Disjointness.RandomizedLowerBound.conditionalSpecialPairLaw_eq_prod_of_singleton_factorization}
\label{CommunicationComplexity.Functions.Disjointness.RandomizedLowerBound.conditionalSpecialPairLaw_eq_prod_of_singleton_factorization}
\leanok
\uses{CommunicationComplexity.Functions.Disjointness.RandomizedLowerBound.ProtocolType, CommunicationComplexity.Functions.Disjointness.RandomizedLowerBound.ZType, CommunicationComplexity.Functions.Disjointness.RandomizedLowerBound.conditionalSpecialPairLaw, CommunicationComplexity.Functions.Disjointness.RandomizedLowerBound.conditionalSpecialXLaw, CommunicationComplexity.Functions.Disjointness.RandomizedLowerBound.conditionalSpecialYLaw}
\difficulty{3}
Fix a positive integer $n$ and a deterministic disjointness protocol $p$ on inputs of
length $n$, and let $z$ be an achievable value of the variable $Z=(M,T,X_{<T},Y_{>T})$.
Write $\mu$ for the conditional law, on the fiber $\{Z=z\}$, of the special bit-pair
$(X_T,Y_T)\in\{0,1\}^2$, and write $\mu_X$ and $\mu_Y$ for the conditional laws on that
same fiber of the individual special bits $X_T$ and $Y_T$. Suppose that for every pair
$b=(b_1,b_2)\in\{0,1\}^2$ the mass assigned by $\mu$ to the singleton $\{b\}$ equals the
product $\mu_X(\{b_1\})\,\mu_Y(\{b_2\})$. Then $\mu$ is the product measure
$\mu_X\otimes\mu_Y$.
\end{theorem}

\begin{theorem}[Product law from singleton factorization of the special pair]
\lean{CommunicationComplexity.Functions.Disjointness.RandomizedLowerBound.conditionalSpecialPairLaw_eq_prod_of_zFiberMeasure_factorization}
\label{CommunicationComplexity.Functions.Disjointness.RandomizedLowerBound.conditionalSpecialPairLaw_eq_prod_of_zFiberMeasure_factorization}
\leanok
\uses{CommunicationComplexity.Functions.Disjointness.RandomizedLowerBound.ProtocolType, CommunicationComplexity.Functions.Disjointness.RandomizedLowerBound.ZType, CommunicationComplexity.Functions.Disjointness.RandomizedLowerBound.conditionalSpecialPairLaw, CommunicationComplexity.Functions.Disjointness.RandomizedLowerBound.conditionalSpecialPairLaw_eq_prod_of_singleton_factorization, CommunicationComplexity.Functions.Disjointness.RandomizedLowerBound.conditionalSpecialPairLaw_singleton, CommunicationComplexity.Functions.Disjointness.RandomizedLowerBound.conditionalSpecialXLaw, CommunicationComplexity.Functions.Disjointness.RandomizedLowerBound.conditionalSpecialXLaw_singleton, CommunicationComplexity.Functions.Disjointness.RandomizedLowerBound.conditionalSpecialYLaw, CommunicationComplexity.Functions.Disjointness.RandomizedLowerBound.conditionalSpecialYLaw_singleton, CommunicationComplexity.Functions.Disjointness.RandomizedLowerBound.specialPair, CommunicationComplexity.Functions.Disjointness.RandomizedLowerBound.specialX, CommunicationComplexity.Functions.Disjointness.RandomizedLowerBound.specialY, CommunicationComplexity.Functions.Disjointness.RandomizedLowerBound.zFiberMeasure}
\difficulty{2}
Fix a deterministic disjointness protocol $p$ on inputs of length $n$ and an achievable
value $z$ of the variable $Z = (M, T, X_{<T}, Y_{>T})$, and let $\mu_z$ denote the law
of the hard-distribution sample conditioned on the fiber $\{Z = z\}$. Suppose that for
every pair of bits $(b_1, b_2) \in \mathrm{Bool} \times \mathrm{Bool}$ the
$\mu_z$-probability that the special-coordinate bits of Alice and Bob equal $(b_1, b_2)$
is the product of the $\mu_z$-probability that Alice's special bit is $b_1$ and the
$\mu_z$-probability that Bob's special bit is $b_2$. Then the conditional law of the
special bit-pair under $\mu_z$ equals the product of the conditional law of Alice's
special bit and the conditional law of Bob's special bit.
\end{theorem}

\begin{theorem}[Product form of the conditional special-pair law from fiber factorization]
\lean{CommunicationComplexity.Functions.Disjointness.RandomizedLowerBound.conditionalSpecialPairLaw_eq_prod_of_fiber_volume_factorization}
\label{CommunicationComplexity.Functions.Disjointness.RandomizedLowerBound.conditionalSpecialPairLaw_eq_prod_of_fiber_volume_factorization}
\leanok
\uses{CommunicationComplexity.Functions.Disjointness.RandomizedLowerBound.ProtocolType, CommunicationComplexity.Functions.Disjointness.RandomizedLowerBound.ZType, CommunicationComplexity.Functions.Disjointness.RandomizedLowerBound.conditionalSpecialPairLaw, CommunicationComplexity.Functions.Disjointness.RandomizedLowerBound.conditionalSpecialPairLaw_eq_prod_of_zFiberMeasure_factorization, CommunicationComplexity.Functions.Disjointness.RandomizedLowerBound.conditionalSpecialXLaw, CommunicationComplexity.Functions.Disjointness.RandomizedLowerBound.conditionalSpecialYLaw, CommunicationComplexity.Functions.Disjointness.RandomizedLowerBound.specialPair, CommunicationComplexity.Functions.Disjointness.RandomizedLowerBound.specialX, CommunicationComplexity.Functions.Disjointness.RandomizedLowerBound.specialY, CommunicationComplexity.Functions.Disjointness.RandomizedLowerBound.volume_zFiber_ne_zero, CommunicationComplexity.Functions.Disjointness.RandomizedLowerBound.zFiber, CommunicationComplexity.Functions.Disjointness.RandomizedLowerBound.zFiberMeasure_real_apply}
\difficulty{3}
Fix a positive integer $n$, a deterministic disjointness protocol $p$ on inputs of
length $n$, and an achievable value $z$ of the variable $Z = (M, T, X_{<T}, Y_{>T})$.
Write $F$ for the fiber of $Z$ above $z$, and for a subset $S$ of the sample space let
$\mathrm{vol}(S)$ denote its real-valued volume. Suppose that for every pair of bits
$(b_1, b_2) \in \mathrm{Bool} \times \mathrm{Bool}$,
\[
\mathrm{vol}(F)\cdot\mathrm{vol}\!\left(F \cap \{\text{special bit-pair} = (b_1,
b_2)\}\right)
= \mathrm{vol}\!\left(F \cap \{\text{Alice's special bit} =
b_1\}\right)\cdot\mathrm{vol}\!\left(F \cap \{\text{Bob's special bit} = b_2\}\right).
\]
Then the conditional law of the special bit-pair on the fiber $F$ equals the product of
the conditional law of Alice's special bit and the conditional law of Bob's special bit.
\end{theorem}

\begin{theorem}[Product form of the conditional special-pair law]
\lean{CommunicationComplexity.Functions.Disjointness.RandomizedLowerBound.conditionalSpecialPairLaw_eq_prod}
\label{CommunicationComplexity.Functions.Disjointness.RandomizedLowerBound.conditionalSpecialPairLaw_eq_prod}
\leanok
\uses{CommunicationComplexity.Functions.Disjointness.RandomizedLowerBound.ProtocolType, CommunicationComplexity.Functions.Disjointness.RandomizedLowerBound.ZType, CommunicationComplexity.Functions.Disjointness.RandomizedLowerBound.conditionalSpecialPairLaw, CommunicationComplexity.Functions.Disjointness.RandomizedLowerBound.conditionalSpecialPairLaw_eq_prod_of_fiber_volume_factorization, CommunicationComplexity.Functions.Disjointness.RandomizedLowerBound.conditionalSpecialXLaw, CommunicationComplexity.Functions.Disjointness.RandomizedLowerBound.conditionalSpecialYLaw, CommunicationComplexity.Functions.Disjointness.RandomizedLowerBound.fiber_volume_factorization}
\difficulty{1}
Fix a deterministic communication protocol $p$ for the disjointness problem on $n$
coordinates and an achievable value $z$ of the summary variable $Z = (M, T, X_{<T},
Y_{>T})$, whose components are the transcript $M$, the special coordinate $T$, the
vector $X_{<T}$ of Alice's bits at coordinates below $T$, and the vector $Y_{>T}$ of
Bob's bits at coordinates above $T$. Let $\nu$ be the uniform distribution on hard
samples conditioned on the fiber $\{Z = z\}$, let $\mu$ be the law of the special
bit-pair $(X_T, Y_T)$ under $\nu$, and let $\mu_X$ and $\mu_Y$ be the laws of the single
bits $X_T$ and $Y_T$ under $\nu$. Then
\[
\mu = \mu_X \otimes \mu_Y.
\]
\end{theorem}

\begin{theorem}[Pointwise factorization of the conditional special-pair law]
\lean{CommunicationComplexity.Functions.Disjointness.RandomizedLowerBound.conditionalSpecialPairLaw_singleton_factorization_of_eq_prod}
\label{CommunicationComplexity.Functions.Disjointness.RandomizedLowerBound.conditionalSpecialPairLaw_singleton_factorization_of_eq_prod}
\leanok
\uses{CommunicationComplexity.Functions.Disjointness.RandomizedLowerBound.ProtocolType, CommunicationComplexity.Functions.Disjointness.RandomizedLowerBound.ZType, CommunicationComplexity.Functions.Disjointness.RandomizedLowerBound.conditionalSpecialPairLaw, CommunicationComplexity.Functions.Disjointness.RandomizedLowerBound.conditionalSpecialXLaw, CommunicationComplexity.Functions.Disjointness.RandomizedLowerBound.conditionalSpecialYLaw}
\difficulty{3}
Fix a deterministic disjointness protocol $p$ on inputs of length $n$ and an achievable
value $z$ of the variable $Z = (M, T, X_{<T}, Y_{>T})$. Write $\mu$ for the conditional
law of the special-coordinate bit-pair on the fiber $\{Z = z\}$, and $\mu_X$, $\mu_Y$
for the conditional laws on that same fiber of Alice's bit and of Bob's bit at the
special coordinate, all measures on Boolean values. Suppose the joint law factors as the
product of its marginals, $\mu = \mu_X \otimes \mu_Y$. Then for every pair of bits $b =
(b_1, b_2) \in \mathrm{Bool} \times \mathrm{Bool}$, the real-valued mass assigned to the
singleton is the product of the marginal masses, \[ \mu(\{b\}) = \mu_X(\{b_1\}) \cdot
\mu_Y(\{b_2\}). \]
\end{theorem}

\begin{theorem}[Product law implies conditioning on Bob's special bit leaves Alice's law unchanged]
\lean{CommunicationComplexity.Functions.Disjointness.RandomizedLowerBound.conditionalSpecialXLaw_eq_cond_specialY_of_prod}
\label{CommunicationComplexity.Functions.Disjointness.RandomizedLowerBound.conditionalSpecialXLaw_eq_cond_specialY_of_prod}
\leanok
\uses{CommunicationComplexity.Functions.Disjointness.RandomizedLowerBound.ProtocolType, CommunicationComplexity.Functions.Disjointness.RandomizedLowerBound.ZType, CommunicationComplexity.Functions.Disjointness.RandomizedLowerBound.conditionalSpecialPairLaw, CommunicationComplexity.Functions.Disjointness.RandomizedLowerBound.conditionalSpecialPairLaw_singleton, CommunicationComplexity.Functions.Disjointness.RandomizedLowerBound.conditionalSpecialPairLaw_singleton_factorization_of_eq_prod, CommunicationComplexity.Functions.Disjointness.RandomizedLowerBound.conditionalSpecialXLaw, CommunicationComplexity.Functions.Disjointness.RandomizedLowerBound.conditionalSpecialXLaw_singleton, CommunicationComplexity.Functions.Disjointness.RandomizedLowerBound.conditionalSpecialYLaw, CommunicationComplexity.Functions.Disjointness.RandomizedLowerBound.conditionalSpecialYLaw_singleton, CommunicationComplexity.Functions.Disjointness.RandomizedLowerBound.specialPair, CommunicationComplexity.Functions.Disjointness.RandomizedLowerBound.specialX, CommunicationComplexity.Functions.Disjointness.RandomizedLowerBound.specialY, CommunicationComplexity.Functions.Disjointness.RandomizedLowerBound.zFiberMeasure}
\difficulty{5}
Fix a deterministic disjointness protocol $p$ on inputs of length $n$ and an achievable
value $z$ of the record $Z=(M,T,X_{<T},Y_{>T})$, and write $\mu_z$ for the hard
distribution conditioned on the fiber $Z=z$. Let $X_T$ and $Y_T$ denote Alice's and
Bob's bits at the special coordinate $T$, and suppose that under $\mu_z$ the pair
$(X_T,Y_T)$ has product law, i.e. the law of $(X_T,Y_T)$ equals the product of the law
of $X_T$ and the law of $Y_T$. Then for every $b_Y\in\{\text{false},\text{true}\}$ such
that the event $\{Y_T=b_Y\}$ has positive $\mu_z$-mass, the law of $X_T$ under $\mu_z$
coincides with the law of $X_T$ under $\mu_z$ conditioned further on $\{Y_T=b_Y\}$.
\end{theorem}

\begin{theorem}[Conditioning on Bob's false special bit leaves Alice's special-bit law unchanged]
\lean{CommunicationComplexity.Functions.Disjointness.RandomizedLowerBound.conditionalSpecialXLaw_eq_cond_specialYFalse}
\label{CommunicationComplexity.Functions.Disjointness.RandomizedLowerBound.conditionalSpecialXLaw_eq_cond_specialYFalse}
\leanok
\uses{CommunicationComplexity.Functions.Disjointness.RandomizedLowerBound.ProtocolType, CommunicationComplexity.Functions.Disjointness.RandomizedLowerBound.ZType, CommunicationComplexity.Functions.Disjointness.RandomizedLowerBound.conditionalSpecialPairLaw_eq_prod, CommunicationComplexity.Functions.Disjointness.RandomizedLowerBound.conditionalSpecialXLaw, CommunicationComplexity.Functions.Disjointness.RandomizedLowerBound.conditionalSpecialXLaw_eq_cond_specialY_of_prod, CommunicationComplexity.Functions.Disjointness.RandomizedLowerBound.specialX, CommunicationComplexity.Functions.Disjointness.RandomizedLowerBound.specialY, CommunicationComplexity.Functions.Disjointness.RandomizedLowerBound.zFiberMeasure}
\difficulty{1}
Fix a positive integer $n$, a deterministic disjointness protocol $p$ on inputs of
length $n$, and an achievable value $z$ of the coarse variable $Z = (M, T, X_{<T},
Y_{>T})$. Let $\mu_z$ be the hard distribution conditioned on the fiber $\{Z = z\}$, and
suppose that under $\mu_z$ the event that Bob's bit at the special coordinate is
$\mathrm{false}$ (that is, $Y_T = \mathrm{false}$) has positive probability. Then the
law of Alice's special bit $X_T$ under $\mu_z$ coincides with the law of $X_T$ under
$\mu_z$ conditioned further on $\{Y_T = \mathrm{false}\}$.
\end{theorem}

\begin{theorem}[Fiber KL divergence unchanged by conditioning on the special Y-bit]
\lean{CommunicationComplexity.Functions.Disjointness.RandomizedLowerBound.xFiberKL_eq_cond_specialYFalse_klDiv}
\label{CommunicationComplexity.Functions.Disjointness.RandomizedLowerBound.xFiberKL_eq_cond_specialYFalse_klDiv}
\leanok
\uses{CommunicationComplexity.Functions.Disjointness.RandomizedLowerBound.ProtocolType, CommunicationComplexity.Functions.Disjointness.RandomizedLowerBound.ZType, CommunicationComplexity.Functions.Disjointness.RandomizedLowerBound.conditionalSpecialXLaw_eq_cond_specialYFalse, CommunicationComplexity.Functions.Disjointness.RandomizedLowerBound.specialX, CommunicationComplexity.Functions.Disjointness.RandomizedLowerBound.specialY, CommunicationComplexity.Functions.Disjointness.RandomizedLowerBound.uniformBool, CommunicationComplexity.Functions.Disjointness.RandomizedLowerBound.xFiberKL, CommunicationComplexity.Functions.Disjointness.RandomizedLowerBound.zFiberMeasure}
\difficulty{1}
Fix a positive integer $n$ and a deterministic communication protocol $p$ for the
disjointness function on inputs of length $n$, and let $z$ be an achievable value of the
variable $Z = (M, T, X_{<T}, Y_{>T})$, recording the transcript, the special coordinate
$T$, Alice's bits before $T$, and Bob's bits after $T$. Write $\mu_z$ for the hard
distribution conditioned on the fiber $\{Z = z\}$, and let $U$ be the uniform law on one
bit. Suppose the event that Bob's bit $Y_T$ at the special coordinate equals
$\mathrm{false}$ has nonzero probability under $\mu_z$. Then the Kullback–Leibler
divergence from $U$ of the law of Alice's special bit $X_T$ under $\mu_z$ equals the
Kullback–Leibler divergence from $U$ of the law of $X_T$ under $\mu_z$ additionally
conditioned on the event $\{Y_T = \mathrm{false}\}$.
\end{theorem}

\begin{theorem}[Fiberwise KL of Alice's special bit under disjointness conditioning]
\lean{CommunicationComplexity.Functions.Disjointness.RandomizedLowerBound.xFiberKL_eq_disjointSpecialYFalseMeasure_cond_zVariable_klDiv}
\label{CommunicationComplexity.Functions.Disjointness.RandomizedLowerBound.xFiberKL_eq_disjointSpecialYFalseMeasure_cond_zVariable_klDiv}
\leanok
\uses{CommunicationComplexity.Functions.Disjointness.RandomizedLowerBound.ProtocolType, CommunicationComplexity.Functions.Disjointness.RandomizedLowerBound.ZType, CommunicationComplexity.Functions.Disjointness.RandomizedLowerBound.disjointSpecialYFalseMeasure, CommunicationComplexity.Functions.Disjointness.RandomizedLowerBound.specialX, CommunicationComplexity.Functions.Disjointness.RandomizedLowerBound.specialY, CommunicationComplexity.Functions.Disjointness.RandomizedLowerBound.uniformBool, CommunicationComplexity.Functions.Disjointness.RandomizedLowerBound.xFiberKL, CommunicationComplexity.Functions.Disjointness.RandomizedLowerBound.xFiberKL_eq_cond_specialYFalse_klDiv, CommunicationComplexity.Functions.Disjointness.RandomizedLowerBound.zFiberMeasure, CommunicationComplexity.Functions.Disjointness.RandomizedLowerBound.zFiberMeasure_cond_specialYFalse_eq_disjointSpecialYFalseMeasure_cond_zVariable, CommunicationComplexity.Functions.Disjointness.RandomizedLowerBound.zVariable}
\difficulty{2}
Let $n$ be a positive integer and let $p$ be a deterministic communication protocol for
the length-$n$ set-disjointness problem, whose inputs are pairs of subsets of
$\{0,\dots,n-1\}$. Write $\mathbb{P}$ for the uniform hard distribution on samples
$\omega$ (each specifying a special coordinate together with the coordinatewise data
generating Alice's set $X(\omega)$ and Bob's set $Y(\omega)$), let $X_T(\omega)$ and
$Y_T(\omega)$ be Alice's and Bob's bits at the special coordinate, and let $D$ be the
event that the generated pair $(X,Y)$ is disjoint. Let $Z = Z_p$ be the statistic
recording the transcript produced by running $p$ on the generated input, together with
the special coordinate, Alice's input bits strictly before it, and Bob's input bits
strictly after it. Write $U$ for the uniform law on one bit and $D_{\mathrm{KL}}$ for
Kullback–Leibler divergence. Fix a value $z$ of $Z$ and suppose the conditional
probability $\mathbb{P}\big(Y_T = \mathrm{false} \mid Z = z\big)$ is nonzero. Then
\[
D_{\mathrm{KL}}\!\big(\mathrm{Law}_{\mathbb{P}}(X_T \mid Z = z)\,\big\|\,U\big)
=
D_{\mathrm{KL}}\!\big(\mathrm{Law}_{\mathbb{P}}(X_T \mid D,\; Y_T = \mathrm{false},\; Z
= z)\,\big\|\,U\big),
\]
both sides taken as real numbers.
\end{theorem}

\begin{theorem}[KL divergence of Alice's special bit under two conditionings]
\lean{CommunicationComplexity.Functions.Disjointness.RandomizedLowerBound.xFiberKL_eq_disjointSpecialYFalseMeasure_cond_zVariable_klDiv_of_ne_zero}
\label{CommunicationComplexity.Functions.Disjointness.RandomizedLowerBound.xFiberKL_eq_disjointSpecialYFalseMeasure_cond_zVariable_klDiv_of_ne_zero}
\leanok
\uses{CommunicationComplexity.Functions.Disjointness.RandomizedLowerBound.ProtocolType, CommunicationComplexity.Functions.Disjointness.RandomizedLowerBound.ZType, CommunicationComplexity.Functions.Disjointness.RandomizedLowerBound.disjointSpecialYFalseMeasure, CommunicationComplexity.Functions.Disjointness.RandomizedLowerBound.specialX, CommunicationComplexity.Functions.Disjointness.RandomizedLowerBound.uniformBool, CommunicationComplexity.Functions.Disjointness.RandomizedLowerBound.xFiberKL, CommunicationComplexity.Functions.Disjointness.RandomizedLowerBound.xFiberKL_eq_disjointSpecialYFalseMeasure_cond_zVariable_klDiv, CommunicationComplexity.Functions.Disjointness.RandomizedLowerBound.zFiber, CommunicationComplexity.Functions.Disjointness.RandomizedLowerBound.zFiberMeasure_specialYFalse_ne_zero_of_disjointSpecialYFalseMeasure_ne_zero, CommunicationComplexity.Functions.Disjointness.RandomizedLowerBound.zVariable}
\difficulty{1}
Fix $n \ge 1$ and a deterministic two-party protocol $p$ for length-$n$ set-disjointness
inputs. To each hard sample $\omega$ associate the record
\[
Z(\omega) = \bigl(\text{transcript of } p \text{ on the input pair generated by }
\omega,\ T(\omega),\ x_{<T}(\omega),\ y_{>T}(\omega)\bigr),
\]
consisting of the protocol's transcript on $\omega$'s input, the special coordinate $T$,
Alice's input bits strictly before $T$, and Bob's input bits strictly after $T$; for a
value $z$ write $\{Z = z\}$ for its fiber. Let $U$ be the uniform law on a single bit,
let $x_T$ denote Alice's bit at the special coordinate, and let $D(\mu\,\|\,\nu)$ denote
the Kullback–Leibler divergence, read as a nonnegative real number. Let $\lambda$ be the
hard distribution conditioned on the generated input pair being disjoint and on Bob's
special bit $y_T$ being false. If the fiber $\{Z = z\}$ has nonzero mass under
$\lambda$, then the divergence from $U$ of the law of $x_T$ under the uniform measure on
hard samples conditioned on $\{Z = z\}$ equals the divergence from $U$ of the law of
$x_T$ under $\lambda$ conditioned on $\{Z = z\}$; that is,
\[
D\!\left(\mu_z \,\middle\|\, U\right) = D\!\left(\lambda_z \,\middle\|\, U\right),
\]
where $\mu_z$ is the law of $x_T$ under the uniform measure on hard samples restricted
to $\{Z = z\}$, and $\lambda_z$ is the law of $x_T$ under $\lambda$ restricted to $\{Z =
z\}$.
\end{theorem}

\begin{theorem}[Total variation of the special bit-pair versus its marginals]
\lean{CommunicationComplexity.Functions.Disjointness.RandomizedLowerBound.zDistance_le_xDistance_add_yDistance}
\label{CommunicationComplexity.Functions.Disjointness.RandomizedLowerBound.zDistance_le_xDistance_add_yDistance}
\leanok
\uses{CommunicationComplexity.Functions.Disjointness.RandomizedLowerBound.ProtocolType, CommunicationComplexity.Functions.Disjointness.RandomizedLowerBound.ZType, CommunicationComplexity.Functions.Disjointness.RandomizedLowerBound.conditionalSpecialPairLaw_eq_prod, CommunicationComplexity.Functions.Disjointness.RandomizedLowerBound.conditionalSpecialXLaw, CommunicationComplexity.Functions.Disjointness.RandomizedLowerBound.conditionalSpecialYLaw, CommunicationComplexity.Functions.Disjointness.RandomizedLowerBound.uniformBool, CommunicationComplexity.Functions.Disjointness.RandomizedLowerBound.uniformBoolPair_eq_prod, CommunicationComplexity.Functions.Disjointness.RandomizedLowerBound.xDistance, CommunicationComplexity.Functions.Disjointness.RandomizedLowerBound.yDistance, CommunicationComplexity.Functions.Disjointness.RandomizedLowerBound.zDistance, CommunicationComplexity.TVDistance.tvDistance_prod_le}
\difficulty{2}
Fix a positive integer $n$, a deterministic communication protocol $p$ for length-$n$
disjointness inputs, and a value $z$ of the coarse conditioning variable. Conditioning
on the event that this variable equals $z$ produces a joint law on the
special-coordinate bit-pair $(a,b)$ — where $a$ is Alice's bit and $b$ is Bob's bit at
the special coordinate — together with its two one-bit marginals, the conditional law of
$a$ and the conditional law of $b$. Then the total variation distance between the joint
conditional law and the uniform distribution on $\{0,1\}^2$ is at most the sum of the
total variation distances of the two marginals from the uniform distribution on
$\{0,1\}$:
\[
\mathrm{TV}\big(\mathcal{L}(a,b),\,\mathrm{Unif}(\{0,1\}^2)\big)\;\le\;\mathrm{TV}\big(\mathcal{L}(a),\,\mathrm{Unif}(\{0,1\})\big)\;+\;\mathrm{TV}\big(\mathcal{L}(b),\,\mathrm{Unif}(\{0,1\})\big).
\]
\end{theorem}

\begin{theorem}[Special bit-pair closeness from marginal closeness]
\lean{CommunicationComplexity.Functions.Disjointness.RandomizedLowerBound.mem_goodZEvent_of_xDistance_yDistance_le}
\label{CommunicationComplexity.Functions.Disjointness.RandomizedLowerBound.mem_goodZEvent_of_xDistance_yDistance_le}
\leanok
\uses{CommunicationComplexity.Functions.Disjointness.RandomizedLowerBound.HardSample, CommunicationComplexity.Functions.Disjointness.RandomizedLowerBound.ProtocolType, CommunicationComplexity.Functions.Disjointness.RandomizedLowerBound.goodZEvent, CommunicationComplexity.Functions.Disjointness.RandomizedLowerBound.xDistance, CommunicationComplexity.Functions.Disjointness.RandomizedLowerBound.yDistance, CommunicationComplexity.Functions.Disjointness.RandomizedLowerBound.zDistance, CommunicationComplexity.Functions.Disjointness.RandomizedLowerBound.zDistance_le_xDistance_add_yDistance, CommunicationComplexity.Functions.Disjointness.RandomizedLowerBound.zVariable}
\difficulty{3}
Let $p$ be a deterministic communication protocol for disjointness inputs of length $n$,
let $\gamma$ be a real number, and let $\omega$ be a sample of the hard distribution.
Associated with $\omega$ is a value $z$ of the observation statistic, which records the
transcript of $p$ on the input pair induced by $\omega$, the special coordinate, Alice's
input bits at the coordinates before it, and Bob's input bits at the coordinates after
it; condition the hard distribution on the fiber of samples producing this same $z$.
Suppose that, under this conditioning, the total variation distance between the law of
Alice's bit at the special coordinate and the uniform law on one bit is at most
$\gamma$, and likewise the total variation distance between the law of Bob's bit at the
special coordinate and the uniform law on one bit is at most $\gamma$. Then the total
variation distance between the conditional law of the special-coordinate bit-pair and
the uniform law on a bit-pair is at most $2\gamma$.
\end{theorem}

\begin{definition}[Aliceinfoterm]
\lean{CommunicationComplexity.Functions.Disjointness.RandomizedLowerBound.aliceInfoTerm}
\label{CommunicationComplexity.Functions.Disjointness.RandomizedLowerBound.aliceInfoTerm}
\leanok
\uses{CommunicationComplexity.Functions.Disjointness.RandomizedLowerBound.ProtocolType, CommunicationComplexity.Functions.Disjointness.RandomizedLowerBound.aliceClaimConditioning, CommunicationComplexity.Functions.Disjointness.RandomizedLowerBound.disjointCondMeasure, CommunicationComplexity.Functions.Disjointness.RandomizedLowerBound.message, CommunicationComplexity.Functions.Disjointness.RandomizedLowerBound.specialX}
The corrected Alice special-coordinate information term: $I(X_{T} : M | T, X_{<T}, Y_{≥T}, D)$.
\end{definition}

\begin{definition}[Bobinfoterm]
\lean{CommunicationComplexity.Functions.Disjointness.RandomizedLowerBound.bobInfoTerm}
\label{CommunicationComplexity.Functions.Disjointness.RandomizedLowerBound.bobInfoTerm}
\leanok
\uses{CommunicationComplexity.Functions.Disjointness.RandomizedLowerBound.ProtocolType, CommunicationComplexity.Functions.Disjointness.RandomizedLowerBound.bobClaimConditioning, CommunicationComplexity.Functions.Disjointness.RandomizedLowerBound.disjointCondMeasure, CommunicationComplexity.Functions.Disjointness.RandomizedLowerBound.message, CommunicationComplexity.Functions.Disjointness.RandomizedLowerBound.specialY}
The corrected Bob special-coordinate information term: $I(Y_{T} : M | T, X_{≤T}, Y_{>T}, D)$.
\end{definition}

\begin{definition}[Claiminfo]
\lean{CommunicationComplexity.Functions.Disjointness.RandomizedLowerBound.claimInfo}
\label{CommunicationComplexity.Functions.Disjointness.RandomizedLowerBound.claimInfo}
\leanok
\uses{CommunicationComplexity.Functions.Disjointness.RandomizedLowerBound.ProtocolType, CommunicationComplexity.Functions.Disjointness.RandomizedLowerBound.aliceInfoTerm, CommunicationComplexity.Functions.Disjointness.RandomizedLowerBound.bobInfoTerm}
The total corrected special-coordinate information.
\end{definition}

\begin{theorem}[Nonnegativity of Bob's special-coordinate information term]
\lean{CommunicationComplexity.Functions.Disjointness.RandomizedLowerBound.bobInfoTerm_nonneg}
\label{CommunicationComplexity.Functions.Disjointness.RandomizedLowerBound.bobInfoTerm_nonneg}
\leanok
\uses{CommunicationComplexity.Functions.Disjointness.RandomizedLowerBound.ProtocolType, CommunicationComplexity.Functions.Disjointness.RandomizedLowerBound.bobClaimConditioning, CommunicationComplexity.Functions.Disjointness.RandomizedLowerBound.bobInfoTerm, CommunicationComplexity.Functions.Disjointness.RandomizedLowerBound.disjointCondMeasure, CommunicationComplexity.Functions.Disjointness.RandomizedLowerBound.message, CommunicationComplexity.Functions.Disjointness.RandomizedLowerBound.specialY}
\difficulty{2}
Fix a positive integer $n$ and a deterministic communication protocol $p$ for the
disjointness problem on inputs indexed by $\{0,1,\dots,n-1\}$. Draw a sample from the
hard distribution over inputs, conditioned on the event that Alice's set and Bob's set
are disjoint, and consider the following random variables: the special coordinate $T$;
Alice's input bits $X_{\le T}$ at coordinates up to and including $T$ (padded by $0$
elsewhere); Bob's input bits $Y_{>T}$ at coordinates strictly above $T$ (padded by $0$
elsewhere); Bob's bit $Y_T$ at the special coordinate; and the transcript $M$ obtained
by running $p$ on the generated input pair. Then the conditional mutual information of
$Y_T$ and $M$ given the tuple $(T, X_{\le T}, Y_{>T})$ is nonnegative:
\[
I\bigl(Y_T : M \mid T, X_{\le T}, Y_{>T}\bigr) \ge 0.
\]
\end{theorem}

\begin{theorem}[Alice's information term is bounded by the total claim information]
\lean{CommunicationComplexity.Functions.Disjointness.RandomizedLowerBound.aliceInfoTerm_le_claimInfo}
\label{CommunicationComplexity.Functions.Disjointness.RandomizedLowerBound.aliceInfoTerm_le_claimInfo}
\leanok
\uses{CommunicationComplexity.Functions.Disjointness.RandomizedLowerBound.ProtocolType, CommunicationComplexity.Functions.Disjointness.RandomizedLowerBound.aliceInfoTerm, CommunicationComplexity.Functions.Disjointness.RandomizedLowerBound.bobInfoTerm_nonneg, CommunicationComplexity.Functions.Disjointness.RandomizedLowerBound.claimInfo}
\difficulty{2}
Fix a positive integer $n$ and a deterministic communication protocol $p$ for
disjointness inputs of length $n$, and give the hard sample space the distribution
conditioned on the event that the generated input pair $(X,Y)$ is disjoint. Write $T$
for the special coordinate, $X_T$ and $Y_T$ for Alice's and Bob's bits at $T$, and $M$
for the transcript produced by running $p$ on the input $(X,Y)$. Let Alice's information
term be the conditional mutual information
\[
I\bigl(X_T \,;\, M \,\big|\, T,\, X_{<T},\, Y_{\ge T}\bigr),
\]
Bob's information term be $I\bigl(Y_T \,;\, M \,\big|\, T,\, X_{\le T},\, Y_{>T}\bigr)$,
and the total claim information be the sum of the two. Then Alice's information term is
at most the total claim information.
\end{theorem}

\begin{theorem}[Alice–Bob duality of the disjointness information terms]
\lean{CommunicationComplexity.Functions.Disjointness.RandomizedLowerBound.aliceInfoTerm_dualProtocol_eq_bobInfoTerm}
\label{CommunicationComplexity.Functions.Disjointness.RandomizedLowerBound.aliceInfoTerm_dualProtocol_eq_bobInfoTerm}
\leanok
\uses{CommunicationComplexity.Functions.Disjointness.RandomizedLowerBound.HardSample, CommunicationComplexity.Functions.Disjointness.RandomizedLowerBound.ProtocolType, CommunicationComplexity.Functions.Disjointness.RandomizedLowerBound.aliceClaimConditioning, CommunicationComplexity.Functions.Disjointness.RandomizedLowerBound.aliceClaimConditioning_dualHardSample, CommunicationComplexity.Functions.Disjointness.RandomizedLowerBound.aliceInfoTerm, CommunicationComplexity.Functions.Disjointness.RandomizedLowerBound.bobClaimConditioning, CommunicationComplexity.Functions.Disjointness.RandomizedLowerBound.bobInfoTerm, CommunicationComplexity.Functions.Disjointness.RandomizedLowerBound.disjointCondMeasure, CommunicationComplexity.Functions.Disjointness.RandomizedLowerBound.disjointCondMeasure_measurePreserving_dualHardSample, CommunicationComplexity.Functions.Disjointness.RandomizedLowerBound.dualConditioningValue, CommunicationComplexity.Functions.Disjointness.RandomizedLowerBound.dualConditioningValue_injective, CommunicationComplexity.Functions.Disjointness.RandomizedLowerBound.dualHardSample, CommunicationComplexity.Functions.Disjointness.RandomizedLowerBound.dualProtocol, CommunicationComplexity.Functions.Disjointness.RandomizedLowerBound.dualProtocolTranscriptMap, CommunicationComplexity.Functions.Disjointness.RandomizedLowerBound.dualProtocolTranscriptMap_injective, CommunicationComplexity.Functions.Disjointness.RandomizedLowerBound.identDistrib_self_comp_measurePreserving, CommunicationComplexity.Functions.Disjointness.RandomizedLowerBound.message, CommunicationComplexity.Functions.Disjointness.RandomizedLowerBound.message_dualProtocol_dualHardSample, CommunicationComplexity.Functions.Disjointness.RandomizedLowerBound.specialX, CommunicationComplexity.Functions.Disjointness.RandomizedLowerBound.specialX_dualHardSample, CommunicationComplexity.Functions.Disjointness.RandomizedLowerBound.specialY, ProbabilityTheory.IdentDistrib.condMutualInfo_eq_finite, ProbabilityTheory.condMutualInfo_comp_right_conditioning_of_injective}
\difficulty{5}
Fix a length $n \ge 1$ and a deterministic communication protocol $p$ for the
set-disjointness problem, whose inputs are subsets of $\{0, \dots, n-1\}$. Write $\mu$
for the hard-distribution measure on samples conditioned on the event that Alice's set
$X$ and Bob's set $Y$ are disjoint, let $T$ be the special coordinate, and let $M_q$
denote the transcript produced by a protocol $q$ on the input $(X, Y)$. Let $p^\ast$ be
the protocol dual, obtained from $p$ by exchanging the roles of Alice and Bob and
reversing the coordinate order of both input sets. Then the conditional mutual
information (under $\mu$) between Alice's bit $x_T$ at the special coordinate and the
transcript of $p^\ast$, given the special coordinate together with Alice's bits strictly
before $T$ and Bob's bits from $T$ onward, equals the conditional mutual information
between Bob's bit $y_T$ at the special coordinate and the transcript of $p$, given the
special coordinate together with Alice's bits up to and including $T$ and Bob's bits
strictly after $T$:
\[
I\!\left[\,x_T : M_{p^\ast} \mid (T,\, x_{<T},\, y_{\ge T})\,;\, \mu\,\right]
\;=\;
I\!\left[\,y_T : M_{p} \mid (T,\, x_{\le T},\, y_{>T})\,;\, \mu\,\right].
\]
\end{theorem}

\begin{theorem}[Duality of Alice's and Bob's information terms]
\lean{CommunicationComplexity.Functions.Disjointness.RandomizedLowerBound.bobInfoTerm_dualProtocol_eq_aliceInfoTerm}
\label{CommunicationComplexity.Functions.Disjointness.RandomizedLowerBound.bobInfoTerm_dualProtocol_eq_aliceInfoTerm}
\leanok
\uses{CommunicationComplexity.Functions.Disjointness.RandomizedLowerBound.HardSample, CommunicationComplexity.Functions.Disjointness.RandomizedLowerBound.ProtocolType, CommunicationComplexity.Functions.Disjointness.RandomizedLowerBound.aliceClaimConditioning, CommunicationComplexity.Functions.Disjointness.RandomizedLowerBound.aliceInfoTerm, CommunicationComplexity.Functions.Disjointness.RandomizedLowerBound.bobClaimConditioning, CommunicationComplexity.Functions.Disjointness.RandomizedLowerBound.bobClaimConditioning_dualHardSample, CommunicationComplexity.Functions.Disjointness.RandomizedLowerBound.bobInfoTerm, CommunicationComplexity.Functions.Disjointness.RandomizedLowerBound.disjointCondMeasure, CommunicationComplexity.Functions.Disjointness.RandomizedLowerBound.disjointCondMeasure_measurePreserving_dualHardSample, CommunicationComplexity.Functions.Disjointness.RandomizedLowerBound.dualConditioningValue, CommunicationComplexity.Functions.Disjointness.RandomizedLowerBound.dualConditioningValue_injective, CommunicationComplexity.Functions.Disjointness.RandomizedLowerBound.dualHardSample, CommunicationComplexity.Functions.Disjointness.RandomizedLowerBound.dualProtocol, CommunicationComplexity.Functions.Disjointness.RandomizedLowerBound.dualProtocolTranscriptMap, CommunicationComplexity.Functions.Disjointness.RandomizedLowerBound.dualProtocolTranscriptMap_injective, CommunicationComplexity.Functions.Disjointness.RandomizedLowerBound.identDistrib_self_comp_measurePreserving, CommunicationComplexity.Functions.Disjointness.RandomizedLowerBound.message, CommunicationComplexity.Functions.Disjointness.RandomizedLowerBound.message_dualProtocol_dualHardSample, CommunicationComplexity.Functions.Disjointness.RandomizedLowerBound.specialX, CommunicationComplexity.Functions.Disjointness.RandomizedLowerBound.specialY, CommunicationComplexity.Functions.Disjointness.RandomizedLowerBound.specialY_dualHardSample, ProbabilityTheory.IdentDistrib.condMutualInfo_eq_finite, ProbabilityTheory.condMutualInfo_comp_right_conditioning_of_injective}
\difficulty{5}
Fix a positive integer $n$ and write $\mu$ for the hard distribution on samples
$\omega$, namely the uniform measure on the finite sample space conditioned on the
disjointness event $X(\omega) \cap Y(\omega) = \emptyset$; here $X(\omega), Y(\omega)
\subseteq \{0, 1, \dots, n-1\}$ are Alice's and Bob's input sets determined by $\omega$,
with distinguished special coordinate $T(\omega)$ and bits $x_T, y_T$ of Alice and Bob
at that coordinate. For a deterministic two-party protocol $p$ on these inputs, let
$M_p(\omega)$ be the transcript produced by running $p$ on the input pair $(X(\omega),
Y(\omega))$, and let $p^\ast$ be its dual, obtained by exchanging the roles of Alice and
Bob and reversing the coordinate order. Then, with all conditional mutual informations
computed under $\mu$,
\[
I\big(y_T \,;\, M_{p^\ast} \;\big|\; T,\, x_{\le T},\, y_{>T}\big) \;=\; I\big(x_T \,;\,
M_p \;\big|\; T,\, x_{<T},\, y_{\ge T}\big),
\]
where $x_{<T}$ and $x_{\le T}$ denote Alice's input bits strictly before, respectively
up to and including, the special coordinate (padded with $\mathrm{false}$ at all other
coordinates), and $y_{>T}$ and $y_{\ge T}$ denote Bob's input bits strictly after,
respectively from, the special coordinate onward (padded with $\mathrm{false}$
elsewhere).
\end{theorem}

\begin{theorem}[Duality identity for conditional transcript information]
\lean{CommunicationComplexity.Functions.Disjointness.RandomizedLowerBound.xVector_message_info_dualProtocol_eq_yVector_message_info}
\label{CommunicationComplexity.Functions.Disjointness.RandomizedLowerBound.xVector_message_info_dualProtocol_eq_yVector_message_info}
\leanok
\uses{CommunicationComplexity.Functions.Disjointness.RandomizedLowerBound.HardSample, CommunicationComplexity.Functions.Disjointness.RandomizedLowerBound.ProtocolType, CommunicationComplexity.Functions.Disjointness.RandomizedLowerBound.disjointCondMeasure, CommunicationComplexity.Functions.Disjointness.RandomizedLowerBound.disjointCondMeasure_measurePreserving_dualHardSample, CommunicationComplexity.Functions.Disjointness.RandomizedLowerBound.dualHardSample, CommunicationComplexity.Functions.Disjointness.RandomizedLowerBound.dualProtocol, CommunicationComplexity.Functions.Disjointness.RandomizedLowerBound.dualProtocolTranscriptMap, CommunicationComplexity.Functions.Disjointness.RandomizedLowerBound.dualProtocolTranscriptMap_injective, CommunicationComplexity.Functions.Disjointness.RandomizedLowerBound.identDistrib_self_comp_measurePreserving, CommunicationComplexity.Functions.Disjointness.RandomizedLowerBound.message, CommunicationComplexity.Functions.Disjointness.RandomizedLowerBound.message_dualProtocol_dualHardSample, CommunicationComplexity.Functions.Disjointness.RandomizedLowerBound.reverseBoolVector, CommunicationComplexity.Functions.Disjointness.RandomizedLowerBound.reverseBoolVector_injective, CommunicationComplexity.Functions.Disjointness.RandomizedLowerBound.xVector, CommunicationComplexity.Functions.Disjointness.RandomizedLowerBound.xVector_dualHardSample, CommunicationComplexity.Functions.Disjointness.RandomizedLowerBound.yVector, CommunicationComplexity.Functions.Disjointness.RandomizedLowerBound.yVector_dualHardSample, ProbabilityTheory.IdentDistrib.condMutualInfo_eq_finite}
\difficulty{5}
Fix a positive integer $n$, and consider the length-$n$ set-disjointness problem in
which Alice and Bob each hold a subset of $\{0,1,\dots,n-1\}$. Let $p$ be a
deterministic two-party communication protocol on such inputs with Boolean output, and
let $p^\ast$ be the dual protocol formed from $p$ by exchanging the roles of Alice and
Bob and reversing the coordinate order of both input sets. Draw a random sample $\omega$
from the hard distribution conditioned on the event that the two generated sets are
disjoint; write $x(\omega), y(\omega) \in \{0,1\}^n$ for Alice's and Bob's full
length-$n$ input vectors at $\omega$, and, for a protocol $q$, write $\Pi_q(\omega)$ for
the transcript produced by running $q$ on the pair $(x(\omega), y(\omega))$. Then, with
every information quantity taken over $\omega$ under the disjointness-conditioned
distribution, the conditional mutual information between Alice's input vector and the
transcript of $p^\ast$ given Bob's input vector equals the conditional mutual
information between Bob's input vector and the transcript of $p$ given Alice's input
vector: \[ I\bigl(x(\omega) \,;\, \Pi_{p^\ast}(\omega) \mid y(\omega)\bigr) =
I\bigl(y(\omega) \,;\, \Pi_{p}(\omega) \mid x(\omega)\bigr). \]
\end{theorem}

\begin{theorem}[Duality invariance of the claim information]
\lean{CommunicationComplexity.Functions.Disjointness.RandomizedLowerBound.claimInfo_dualProtocol}
\label{CommunicationComplexity.Functions.Disjointness.RandomizedLowerBound.claimInfo_dualProtocol}
\leanok
\uses{CommunicationComplexity.Functions.Disjointness.RandomizedLowerBound.ProtocolType, CommunicationComplexity.Functions.Disjointness.RandomizedLowerBound.aliceInfoTerm_dualProtocol_eq_bobInfoTerm, CommunicationComplexity.Functions.Disjointness.RandomizedLowerBound.bobInfoTerm_dualProtocol_eq_aliceInfoTerm, CommunicationComplexity.Functions.Disjointness.RandomizedLowerBound.claimInfo, CommunicationComplexity.Functions.Disjointness.RandomizedLowerBound.dualProtocol}
\difficulty{2}
Fix a positive integer $n$ and identify the coordinate set with $\{0,1,\dots,n-1\}$.
Draw a sample $\omega$ from the hard distribution on inputs, conditioned on the event
that Alice's set $X(\omega)$ and Bob's set $Y(\omega)$ are disjoint, and write $T$ for
the special coordinate of $\omega$ and $X_T, Y_T$ for Alice's and Bob's bits at $T$. For
a deterministic protocol $p$ whose two inputs both range over subsets of
$\{0,\dots,n-1\}$, let $\Pi_p$ be the transcript that $p$ produces on the pair
$\bigl(X(\omega), Y(\omega)\bigr)$. Let $X_{<T}$ and $X_{\le T}$ be the vectors of
Alice's bits at the coordinates below $T$, respectively at the coordinates at most $T$,
each set to $0$ outside that range, and let $Y_{\ge T}$ and $Y_{>T}$ be the vectors of
Bob's bits at the coordinates at least $T$, respectively above $T$, each set to $0$
outside that range. Define the claim information of $p$ by
\[
\mathcal{I}(p) \;=\; I\!\left(X_T \,;\, \Pi_p \,\middle|\, T,\, X_{<T},\, Y_{\ge
T}\right) \;+\; I\!\left(Y_T \,;\, \Pi_p \,\middle|\, T,\, X_{\le T},\, Y_{>T}\right),
\]
both conditional mutual informations being taken under the conditioned distribution. Let
$p^{\dagger}$ be the protocol obtained from $p$ by exchanging the roles of Alice and Bob
throughout and reversing the coordinate order $i \mapsto n-1-i$ on both parties' input
sets. Then $\mathcal{I}(p^{\dagger}) = \mathcal{I}(p)$.
\end{theorem}

\begin{definition}[Fixedaliceinfoterm]
\lean{CommunicationComplexity.Functions.Disjointness.RandomizedLowerBound.fixedAliceInfoTerm}
\label{CommunicationComplexity.Functions.Disjointness.RandomizedLowerBound.fixedAliceInfoTerm}
\leanok
\uses{CommunicationComplexity.Functions.Disjointness.RandomizedLowerBound.ProtocolType, CommunicationComplexity.Functions.Disjointness.RandomizedLowerBound.disjointCondMeasure, CommunicationComplexity.Functions.Disjointness.RandomizedLowerBound.fixedAliceConditioning, CommunicationComplexity.Functions.Disjointness.RandomizedLowerBound.fixedXBit, CommunicationComplexity.Functions.Disjointness.RandomizedLowerBound.message}
The fixed-coordinate summand $I(X_{i} : M | X_{<i}, Y_{≥i})$.
\end{definition}

\begin{definition}[Fixedalicefullyinfoterm]
\lean{CommunicationComplexity.Functions.Disjointness.RandomizedLowerBound.fixedAliceFullYInfoTerm}
\label{CommunicationComplexity.Functions.Disjointness.RandomizedLowerBound.fixedAliceFullYInfoTerm}
\leanok
\uses{CommunicationComplexity.Functions.Disjointness.RandomizedLowerBound.ProtocolType, CommunicationComplexity.Functions.Disjointness.RandomizedLowerBound.disjointCondMeasure, CommunicationComplexity.Functions.Disjointness.RandomizedLowerBound.fixedAliceFullYConditioning, CommunicationComplexity.Functions.Disjointness.RandomizedLowerBound.fixedXBit, CommunicationComplexity.Functions.Disjointness.RandomizedLowerBound.message}
The fixed-coordinate term with the full $Y$ vector in the conditioning.
\end{definition}

\begin{definition}[Fixedalicecrossinfoterm]
\lean{CommunicationComplexity.Functions.Disjointness.RandomizedLowerBound.fixedAliceCrossInfoTerm}
\label{CommunicationComplexity.Functions.Disjointness.RandomizedLowerBound.fixedAliceCrossInfoTerm}
\leanok
\uses{CommunicationComplexity.Functions.Disjointness.RandomizedLowerBound.disjointCondMeasure, CommunicationComplexity.Functions.Disjointness.RandomizedLowerBound.fixedAliceConditioning, CommunicationComplexity.Functions.Disjointness.RandomizedLowerBound.fixedXBit, CommunicationComplexity.Functions.Disjointness.RandomizedLowerBound.fixedYBefore}
The zero cross-information term $I(X_{i} : Y_{<i} | X_{<i}, Y_{≥i})$.
\end{definition}

\begin{theorem}[Vanishing cross-information for Alice's fixed coordinate]
\lean{CommunicationComplexity.Functions.Disjointness.RandomizedLowerBound.fixedAliceCrossInfoTerm_eq_zero}
\label{CommunicationComplexity.Functions.Disjointness.RandomizedLowerBound.fixedAliceCrossInfoTerm_eq_zero}
\leanok
\uses{CommunicationComplexity.Functions.Disjointness.RandomizedLowerBound.coordinateAliceConditioning, CommunicationComplexity.Functions.Disjointness.RandomizedLowerBound.coordinateXBit, CommunicationComplexity.Functions.Disjointness.RandomizedLowerBound.coordinateYBefore, CommunicationComplexity.Functions.Disjointness.RandomizedLowerBound.disjointCondMeasure, CommunicationComplexity.Functions.Disjointness.RandomizedLowerBound.fixedAliceConditioning, CommunicationComplexity.Functions.Disjointness.RandomizedLowerBound.fixedAliceCrossInfoTerm, CommunicationComplexity.Functions.Disjointness.RandomizedLowerBound.fixedXBit, CommunicationComplexity.Functions.Disjointness.RandomizedLowerBound.fixedYBefore, CommunicationComplexity.Functions.Disjointness.RandomizedLowerBound.identDistrib_fixedAliceCrossInfoTriple_uniform, CommunicationComplexity.Functions.Disjointness.RandomizedLowerBound.uniformDisjointCoordinateVector, CommunicationComplexity.Functions.Disjointness.RandomizedLowerBound.uniformDisjointCoordinateVector_crossInfo_eq_zero, ProbabilityTheory.IdentDistrib.condMutualInfo_eq_finite}
\difficulty{3}
Fix a positive integer $n$, write the coordinates as $i \in \{0,\dots,n-1\}$, and
consider the hard distribution on samples conditioned on the event that Alice's input
set and Bob's input set are disjoint. For any coordinate $i$, under this conditioned
distribution the conditional mutual information between Alice's bit $X_i$ at coordinate
$i$ and Bob's prefix $Y_{<i}$ (his bits at all coordinates below $i$), given the pair
$(X_{<i}, Y_{\ge i})$ formed by Alice's prefix and Bob's suffix, is zero:
\[
  I\!\left(X_i : Y_{<i} \,\middle|\, X_{<i},\, Y_{\ge i}\right) = 0.
\]
\end{theorem}

\begin{theorem}[Fixed-coordinate conditioning inequality for Alice's information term]
\lean{CommunicationComplexity.Functions.Disjointness.RandomizedLowerBound.fixedAliceInfoTerm_le_fixedAliceFullYInfoTerm}
\label{CommunicationComplexity.Functions.Disjointness.RandomizedLowerBound.fixedAliceInfoTerm_le_fixedAliceFullYInfoTerm}
\leanok
\uses{CommunicationComplexity.Functions.Disjointness.RandomizedLowerBound.HardSample, CommunicationComplexity.Functions.Disjointness.RandomizedLowerBound.ProtocolType, CommunicationComplexity.Functions.Disjointness.RandomizedLowerBound.TranscriptType, CommunicationComplexity.Functions.Disjointness.RandomizedLowerBound.disjointCondMeasure, CommunicationComplexity.Functions.Disjointness.RandomizedLowerBound.fixedAliceChainConditioningValue_fixedAliceFullYConditioning, CommunicationComplexity.Functions.Disjointness.RandomizedLowerBound.fixedAliceChainConditioningValue_injective, CommunicationComplexity.Functions.Disjointness.RandomizedLowerBound.fixedAliceConditioning, CommunicationComplexity.Functions.Disjointness.RandomizedLowerBound.fixedAliceCrossInfoTerm, CommunicationComplexity.Functions.Disjointness.RandomizedLowerBound.fixedAliceCrossInfoTerm_eq_zero, CommunicationComplexity.Functions.Disjointness.RandomizedLowerBound.fixedAliceFullYConditioning, CommunicationComplexity.Functions.Disjointness.RandomizedLowerBound.fixedAliceFullYInfoTerm, CommunicationComplexity.Functions.Disjointness.RandomizedLowerBound.fixedAliceInfoTerm, CommunicationComplexity.Functions.Disjointness.RandomizedLowerBound.fixedXBit, CommunicationComplexity.Functions.Disjointness.RandomizedLowerBound.fixedYBefore, CommunicationComplexity.Functions.Disjointness.RandomizedLowerBound.message, ProbabilityTheory.condMutualInfo_le_prod_right_snd, ProbabilityTheory.condMutualInfo_prod_right_eq_add}
\difficulty{5}
Fix a positive integer $n$, a deterministic communication protocol $p$ for disjointness
inputs on the coordinate set $\{0,\dots,n-1\}$, and a coordinate $i$. Draw a hard sample
and condition on the event that the induced input sets $X$ and $Y$ are disjoint; under
this conditioned distribution let $X_i$ denote Alice's bit at coordinate $i$, let $M$
denote the transcript produced by executing $p$ on the induced input pair, let $X_{<i}$
denote Alice's bits at the coordinates below $i$, let $Y_{\ge i}$ denote Bob's bits at
the coordinates at least $i$, and let $Y$ denote Bob's full bit vector. Then the
conditional mutual information between $X_i$ and $M$, given the partial pair $(X_{<i},
Y_{\ge i})$, is at most the conditional mutual information given the fuller pair
$(X_{<i}, Y)$:
\[
I\!\left(X_i : M \mid X_{<i},\, Y_{\ge i}\right) \;\le\; I\!\left(X_i : M \mid X_{<i},\,
Y\right).
\]
\end{theorem}

\begin{theorem}[Alice's information sum under fuller conditioning on Bob's input]
\lean{CommunicationComplexity.Functions.Disjointness.RandomizedLowerBound.sum_fixedAliceInfoTerm_le_sum_fixedAliceFullYInfoTerm}
\label{CommunicationComplexity.Functions.Disjointness.RandomizedLowerBound.sum_fixedAliceInfoTerm_le_sum_fixedAliceFullYInfoTerm}
\leanok
\uses{CommunicationComplexity.Functions.Disjointness.RandomizedLowerBound.ProtocolType, CommunicationComplexity.Functions.Disjointness.RandomizedLowerBound.fixedAliceFullYInfoTerm, CommunicationComplexity.Functions.Disjointness.RandomizedLowerBound.fixedAliceInfoTerm, CommunicationComplexity.Functions.Disjointness.RandomizedLowerBound.fixedAliceInfoTerm_le_fixedAliceFullYInfoTerm}
\difficulty{1}
Fix a positive integer $n$ and a deterministic communication protocol $p$ for the
disjointness problem on inputs of length $n$. Let $\mu$ be the hard distribution on
samples conditioned on the event that Alice's input set and Bob's input set are
disjoint, and, for a sample governed by $\mu$, write $M$ for the transcript produced by
running $p$ on the sample's input pair, write $X_i$ for Alice's bit at coordinate $i$,
write $X_{<i}$ for Alice's bits at all coordinates below $i$, write $Y_{\ge i}$ for
Bob's bits at the coordinates from $i$ onward, and write $Y$ for Bob's full bit vector.
With every mutual information computed under $\mu$,
\[
\sum_{i=0}^{n-1} I\!\left(X_i ; M \mid X_{<i},\, Y_{\ge i}\right) \;\le\;
\sum_{i=0}^{n-1} I\!\left(X_i ; M \mid X_{<i},\, Y\right).
\]
\end{theorem}

\begin{theorem}[Chain rule for the mutual information between Alice's input and the transcript]
\lean{CommunicationComplexity.Functions.Disjointness.RandomizedLowerBound.sum_fixedAliceFullYInfoTerm_eq_xVector_info}
\label{CommunicationComplexity.Functions.Disjointness.RandomizedLowerBound.sum_fixedAliceFullYInfoTerm_eq_xVector_info}
\leanok
\uses{CommunicationComplexity.Functions.Disjointness.RandomizedLowerBound.ProtocolType, CommunicationComplexity.Functions.Disjointness.RandomizedLowerBound.disjointCondMeasure, CommunicationComplexity.Functions.Disjointness.RandomizedLowerBound.fixedAliceFullYConditioning, CommunicationComplexity.Functions.Disjointness.RandomizedLowerBound.fixedAliceFullYInfoTerm, CommunicationComplexity.Functions.Disjointness.RandomizedLowerBound.fixedXBit, CommunicationComplexity.Functions.Disjointness.RandomizedLowerBound.fixedXStrictPrefix, CommunicationComplexity.Functions.Disjointness.RandomizedLowerBound.message, CommunicationComplexity.Functions.Disjointness.RandomizedLowerBound.xVector, CommunicationComplexity.Functions.Disjointness.RandomizedLowerBound.yVector, ProbabilityTheory.boolVectorStrictPrefix, ProbabilityTheory.condMutualInfo_boolVector_eq_sum_strictPrefix}
\difficulty{3}
Fix a positive integer $n$ and a deterministic protocol $p$ for length-$n$ disjointness
inputs, and work throughout under the hard distribution on samples conditioned on the
event that the generated input pair is disjoint. Under this distribution write $X =
(X_1, \ldots, X_n)$ for Alice's input vector, $Y$ for Bob's input vector, $M$ for the
transcript produced by running $p$ on the input pair $(X, Y)$, and $X_{<i}$ for the
strict prefix $(X_1, \ldots, X_{i-1})$ of Alice's vector. Then
\[ \sum_{i=1}^{n} I\bigl(X_i : M \mid X_{<i},\, Y\bigr) = I\bigl(X : M \mid Y\bigr), \]
where each $I(\,\cdot : \cdot \mid \cdot\,)$ denotes conditional mutual information
taken with respect to the conditioned distribution.
\end{theorem}

\begin{theorem}[Chain-rule bound on Alice's per-coordinate information terms]
\lean{CommunicationComplexity.Functions.Disjointness.RandomizedLowerBound.sum_fixedAliceInfoTerm_le_xVector_info}
\label{CommunicationComplexity.Functions.Disjointness.RandomizedLowerBound.sum_fixedAliceInfoTerm_le_xVector_info}
\leanok
\uses{CommunicationComplexity.Functions.Disjointness.RandomizedLowerBound.ProtocolType, CommunicationComplexity.Functions.Disjointness.RandomizedLowerBound.disjointCondMeasure, CommunicationComplexity.Functions.Disjointness.RandomizedLowerBound.fixedAliceFullYInfoTerm, CommunicationComplexity.Functions.Disjointness.RandomizedLowerBound.fixedAliceInfoTerm, CommunicationComplexity.Functions.Disjointness.RandomizedLowerBound.message, CommunicationComplexity.Functions.Disjointness.RandomizedLowerBound.sum_fixedAliceFullYInfoTerm_eq_xVector_info, CommunicationComplexity.Functions.Disjointness.RandomizedLowerBound.sum_fixedAliceInfoTerm_le_sum_fixedAliceFullYInfoTerm, CommunicationComplexity.Functions.Disjointness.RandomizedLowerBound.xVector, CommunicationComplexity.Functions.Disjointness.RandomizedLowerBound.yVector}
\difficulty{2}
Fix a positive integer $n$ and a deterministic protocol $p$ for the $n$-coordinate
disjointness problem. Draw an input from the hard distribution conditioned on the event
that Alice's and Bob's sets are disjoint, and write $X = (X_1,\dots,X_n)$ and $Y =
(Y_1,\dots,Y_n)$ for Alice's and Bob's Boolean input vectors and $M$ for the transcript
produced by running $p$ on that input pair; for each coordinate $i$ let $X_i$ be Alice's
bit there, $X_{<i}$ her vector of bits at earlier coordinates, and $Y_{\ge i}$ Bob's
vector of bits at $i$ and later coordinates. Then, with every mutual information taken
under this disjointness-conditioned distribution,
\[
\sum_{i=1}^{n} I\bigl(X_i : M \mid X_{<i},\, Y_{\ge i}\bigr) \;\le\; I\bigl(X : M \mid
Y\bigr).
\]
\end{theorem}

\begin{theorem}[Fiber decomposition of Alice's special-coordinate information term]
\lean{CommunicationComplexity.Functions.Disjointness.RandomizedLowerBound.aliceInfoTerm_eq_sum_specialCoordinate_fiber_info}
\label{CommunicationComplexity.Functions.Disjointness.RandomizedLowerBound.aliceInfoTerm_eq_sum_specialCoordinate_fiber_info}
\leanok
\uses{CommunicationComplexity.Functions.Disjointness.RandomizedLowerBound.ProtocolType, CommunicationComplexity.Functions.Disjointness.RandomizedLowerBound.aliceClaimConditioning_eq_specialCoordinate_prod_dynamic, CommunicationComplexity.Functions.Disjointness.RandomizedLowerBound.aliceDynamicConditioning, CommunicationComplexity.Functions.Disjointness.RandomizedLowerBound.aliceInfoTerm, CommunicationComplexity.Functions.Disjointness.RandomizedLowerBound.disjointCondMeasure, CommunicationComplexity.Functions.Disjointness.RandomizedLowerBound.disjointCondMeasure_measureReal_specialCoordinate_preimage_singleton, CommunicationComplexity.Functions.Disjointness.RandomizedLowerBound.message, CommunicationComplexity.Functions.Disjointness.RandomizedLowerBound.specialCoordinate, CommunicationComplexity.Functions.Disjointness.RandomizedLowerBound.specialX, ProbabilityTheory.condMutualInfo_prod_conditioning_eq_sum}
\difficulty{3}
Let $n$ be a positive integer and let $p$ be a deterministic protocol for disjointness
inputs of length $n$, and write $M$ for the transcript produced by running $p$ on the
input pair generated by a hard-distribution sample. Work throughout under the hard
distribution conditioned on the generated pair being disjoint, and write $T$ for the
special coordinate, $X_T$ for Alice's bit at that coordinate, and $X_{<T}$ and $Y_{\ge
T}$ for Alice's bits before the special coordinate and Bob's bits from the special
coordinate onward. Then Alice's special-coordinate information term is the uniform
average, over the $n$ possible values of $T$, of the corresponding conditional mutual
information computed on the fiber where $T$ is fixed:
\[
I\!\left(X_T : M \mid T,\, X_{<T},\, Y_{\ge T}\right)
= \sum_{i=1}^{n} \frac{1}{n}\, I_{\,T = i}\!\left(X_T : M \mid X_{<T},\, Y_{\ge
T}\right),
\]
where $I_{\,T=i}$ denotes conditional mutual information taken under the same
disjoint-conditioned distribution further conditioned on the event $T = i$.
\end{theorem}

\begin{theorem}[Dynamic Alice information at a fixed special coordinate]
\lean{CommunicationComplexity.Functions.Disjointness.RandomizedLowerBound.aliceDynamicConditioning_fiber_info_eq_fixedAliceInfoTerm}
\label{CommunicationComplexity.Functions.Disjointness.RandomizedLowerBound.aliceDynamicConditioning_fiber_info_eq_fixedAliceInfoTerm}
\leanok
\uses{CommunicationComplexity.Functions.Disjointness.RandomizedLowerBound.HardSample, CommunicationComplexity.Functions.Disjointness.RandomizedLowerBound.ProtocolType, CommunicationComplexity.Functions.Disjointness.RandomizedLowerBound.aliceDynamicConditioning, CommunicationComplexity.Functions.Disjointness.RandomizedLowerBound.coordinateAliceConditioning, CommunicationComplexity.Functions.Disjointness.RandomizedLowerBound.coordinateMessage, CommunicationComplexity.Functions.Disjointness.RandomizedLowerBound.coordinateXBit, CommunicationComplexity.Functions.Disjointness.RandomizedLowerBound.disjointCondMeasure, CommunicationComplexity.Functions.Disjointness.RandomizedLowerBound.disjointCondMeasure_measureReal_specialCoordinate_preimage_singleton, CommunicationComplexity.Functions.Disjointness.RandomizedLowerBound.disjointCoordinateVector, CommunicationComplexity.Functions.Disjointness.RandomizedLowerBound.fixedAliceConditioning, CommunicationComplexity.Functions.Disjointness.RandomizedLowerBound.fixedAliceConditioning_ae_eq_coordinateAliceConditioning, CommunicationComplexity.Functions.Disjointness.RandomizedLowerBound.fixedAliceInfoTerm, CommunicationComplexity.Functions.Disjointness.RandomizedLowerBound.fixedXBefore, CommunicationComplexity.Functions.Disjointness.RandomizedLowerBound.fixedXBit, CommunicationComplexity.Functions.Disjointness.RandomizedLowerBound.fixedXBit_ae_eq_coordinateXBit, CommunicationComplexity.Functions.Disjointness.RandomizedLowerBound.fixedYGe, CommunicationComplexity.Functions.Disjointness.RandomizedLowerBound.identDistrib_disjointCoordinateVector_uniform, CommunicationComplexity.Functions.Disjointness.RandomizedLowerBound.identDistrib_disjointCoordinateVector_uniform_cond_specialCoordinate, CommunicationComplexity.Functions.Disjointness.RandomizedLowerBound.message, CommunicationComplexity.Functions.Disjointness.RandomizedLowerBound.message_ae_eq_coordinateMessage, CommunicationComplexity.Functions.Disjointness.RandomizedLowerBound.specialCoordinate, CommunicationComplexity.Functions.Disjointness.RandomizedLowerBound.specialX, CommunicationComplexity.Functions.Disjointness.RandomizedLowerBound.uniformDisjointCoordinateVector, CommunicationComplexity.Functions.Disjointness.RandomizedLowerBound.xBeforeSpecial, CommunicationComplexity.Functions.Disjointness.RandomizedLowerBound.xBit, CommunicationComplexity.Functions.Disjointness.RandomizedLowerBound.yGeSpecial, ProbabilityTheory.IdentDistrib.condMutualInfo_eq_finite, ProbabilityTheory.condMutualInfo_congr_ae_finite}
\difficulty{6}
Let $n$ be a positive integer and let $p$ be a deterministic communication protocol for
the disjointness problem on the $n$ coordinates, in which each player holds a subset of
$\{0,\dots,n-1\}$ and the protocol outputs a Boolean value. Running $p$ on the input
pair $(X(\omega),Y(\omega))$ determined by a hard sample $\omega$ produces a transcript
$M(\omega)$. Write $\mu$ for the hard distribution conditioned on the event that
$X(\omega)$ and $Y(\omega)$ are disjoint, and for a sample $\omega$ let $T$ be its
special coordinate and $x_T$ Alice's bit there; the dynamic conditioning of $\omega$ is
the pair $(x_{<T},\,y_{\ge T})$ recording Alice's bits at coordinates below $T$ and
Bob's bits at coordinates at or above $T$, with every remaining coordinate set to
$\mathrm{false}$. Then for each coordinate $i$,
\[
I\!\left(x_T ; M \,\middle|\, (x_{<T},\,y_{\ge T})\right)_{\mu(\,\cdot\,\mid T=i)}
\;=\;
I\!\left(x_i ; M \,\middle|\, (x_{<i},\,y_{\ge i})\right)_{\mu},
\]
where the left-hand conditional mutual information is taken under $\mu$ further
conditioned on the event $\{T=i\}$, and the right-hand one under $\mu$ with the
fixed-coordinate conditioning $(x_{<i},\,y_{\ge i})$ (Alice's bits at coordinates below
$i$ and Bob's bits at coordinates at or above $i$, all other coordinates
$\mathrm{false}$) and Alice's bit $x_i$ at coordinate $i$.
\end{theorem}

\begin{theorem}[Alice information term as an average over coordinates]
\lean{CommunicationComplexity.Functions.Disjointness.RandomizedLowerBound.aliceInfoTerm_eq_average_fixedAliceInfoTerm}
\label{CommunicationComplexity.Functions.Disjointness.RandomizedLowerBound.aliceInfoTerm_eq_average_fixedAliceInfoTerm}
\leanok
\uses{CommunicationComplexity.Functions.Disjointness.RandomizedLowerBound.ProtocolType, CommunicationComplexity.Functions.Disjointness.RandomizedLowerBound.aliceDynamicConditioning_fiber_info_eq_fixedAliceInfoTerm, CommunicationComplexity.Functions.Disjointness.RandomizedLowerBound.aliceInfoTerm, CommunicationComplexity.Functions.Disjointness.RandomizedLowerBound.aliceInfoTerm_eq_sum_specialCoordinate_fiber_info, CommunicationComplexity.Functions.Disjointness.RandomizedLowerBound.fixedAliceInfoTerm}
\difficulty{2}
Fix a positive integer $n$ and a deterministic communication protocol $p$ for
disjointness inputs of length $n$, and let $M$ denote the transcript that $p$ produces
on the input pair generated by a hard-distribution sample. All quantities below are
computed under the hard distribution conditioned on the event that the generated pair of
sets is disjoint. Writing $X_i$ for Alice's bit at coordinate $i$, $X_{<i}$ for Alice's
input bits at the coordinates below $i$, $Y_{\ge i}$ for Bob's input bits at the
coordinates $i$ and above, and $T$ for the special coordinate, one has
\[
I\!\left(X_T ; M \mid T,\, X_{<T},\, Y_{\ge T}\right)
  \;=\; \frac{1}{n} \sum_{i=1}^{n} I\!\left(X_i ; M \mid X_{<i},\, Y_{\ge i}\right),
\]
so the conditional mutual information between Alice's bit at the special coordinate and
the transcript, given the triple $(T, X_{<T}, Y_{\ge T})$, is the average over the $n$
coordinates of the corresponding fixed-coordinate information terms.
\end{theorem}

\begin{theorem}[Averaged information bound for Alice's special-coordinate bit]
\lean{CommunicationComplexity.Functions.Disjointness.RandomizedLowerBound.aliceInfoTerm_le_average_xVector_info}
\label{CommunicationComplexity.Functions.Disjointness.RandomizedLowerBound.aliceInfoTerm_le_average_xVector_info}
\leanok
\uses{CommunicationComplexity.Functions.Disjointness.RandomizedLowerBound.ProtocolType, CommunicationComplexity.Functions.Disjointness.RandomizedLowerBound.aliceInfoTerm, CommunicationComplexity.Functions.Disjointness.RandomizedLowerBound.aliceInfoTerm_eq_average_fixedAliceInfoTerm, CommunicationComplexity.Functions.Disjointness.RandomizedLowerBound.disjointCondMeasure, CommunicationComplexity.Functions.Disjointness.RandomizedLowerBound.message, CommunicationComplexity.Functions.Disjointness.RandomizedLowerBound.sum_fixedAliceInfoTerm_le_xVector_info, CommunicationComplexity.Functions.Disjointness.RandomizedLowerBound.xVector, CommunicationComplexity.Functions.Disjointness.RandomizedLowerBound.yVector}
\difficulty{2}
Fix an integer $n \ge 1$ and a deterministic two-party communication protocol $p$ for
set disjointness on inputs of length $n$, in which Alice and Bob each hold a subset of
$\{1,\dots,n\}$. Let $\mu$ denote the hard distribution over such input pairs
conditioned on the event that Alice's and Bob's sets are disjoint, and take all mutual
information $I(\,\cdot\,;\,\cdot\mid\,\cdot\,)$ below with respect to $\mu$. Under
$\mu$, let $X, Y \colon \{1,\dots,n\} \to \{0,1\}$ be the input vectors of Alice and
Bob, let $T$ be the special coordinate, let $X_T$ be Alice's bit at coordinate $T$, and
let $M$ be the transcript produced by running $p$ on the pair $(X,Y)$; write $X_{<T}$
for the vector of Alice's bits at coordinates strictly below $T$, padded with $0$ at
every other coordinate, and $Y_{\ge T}$ for the vector of Bob's bits at coordinates at
least $T$, padded with $0$ at every other coordinate. Then
\[
I\big(X_T \,;\, M \mid T,\, X_{<T},\, Y_{\ge T}\big) \;\le\; \frac{1}{n}\, I\big(X \,;\,
M \mid Y\big).
\]
\end{theorem}

\begin{theorem}[Bob's single-coordinate information is at most the average full-vector information]
\lean{CommunicationComplexity.Functions.Disjointness.RandomizedLowerBound.bobInfoTerm_le_average_yVector_info}
\label{CommunicationComplexity.Functions.Disjointness.RandomizedLowerBound.bobInfoTerm_le_average_yVector_info}
\leanok
\uses{CommunicationComplexity.Functions.Disjointness.RandomizedLowerBound.ProtocolType, CommunicationComplexity.Functions.Disjointness.RandomizedLowerBound.aliceInfoTerm_dualProtocol_eq_bobInfoTerm, CommunicationComplexity.Functions.Disjointness.RandomizedLowerBound.aliceInfoTerm_le_average_xVector_info, CommunicationComplexity.Functions.Disjointness.RandomizedLowerBound.bobInfoTerm, CommunicationComplexity.Functions.Disjointness.RandomizedLowerBound.disjointCondMeasure, CommunicationComplexity.Functions.Disjointness.RandomizedLowerBound.dualProtocol, CommunicationComplexity.Functions.Disjointness.RandomizedLowerBound.message, CommunicationComplexity.Functions.Disjointness.RandomizedLowerBound.xVector, CommunicationComplexity.Functions.Disjointness.RandomizedLowerBound.xVector_message_info_dualProtocol_eq_yVector_message_info, CommunicationComplexity.Functions.Disjointness.RandomizedLowerBound.yVector}
\difficulty{2}
Fix an integer $n \ge 1$ and a deterministic communication protocol $p$ for the
disjointness problem on $n$ coordinates, in which Alice and Bob each hold a subset of
$\{1,\dots,n\}$. Work under the hard distribution over input pairs conditioned on the
event that Alice's set $X$ and Bob's set $Y$ are disjoint, and let $M$ be the transcript
that $p$ produces on a sampled input pair. Let $T$ be the special coordinate, let $Y_T$
be Bob's bit at coordinate $T$, let $X_{\le T}$ be the vector recording Alice's bits at
all coordinates up to and including $T$ and $0$ beyond, let $Y_{>T}$ be the vector
recording Bob's bits at all coordinates strictly greater than $T$ and $0$ elsewhere, and
let $X$ and $Y$ be the full input vectors of Alice and Bob. With every mutual
information taken with respect to the hard distribution conditioned on disjointness of
$X$ and $Y$,
\[
I\!\left(Y_T ; M \mid T,\, X_{\le T},\, Y_{>T}\right) \;\le\; \frac{1}{n}\, I\!\left(Y ;
M \mid X\right).
\]
\end{theorem}

\begin{theorem}[Information cost bound for the disjointness transcript]
\lean{CommunicationComplexity.Functions.Disjointness.RandomizedLowerBound.xVector_message_info_le_complexity_mul_log_two}
\label{CommunicationComplexity.Functions.Disjointness.RandomizedLowerBound.xVector_message_info_le_complexity_mul_log_two}
\leanok
\uses{CommunicationComplexity.Deterministic.Protocol.complexity, CommunicationComplexity.Functions.Disjointness.RandomizedLowerBound.ProtocolType, CommunicationComplexity.Functions.Disjointness.RandomizedLowerBound.disjointCondMeasure, CommunicationComplexity.Functions.Disjointness.RandomizedLowerBound.entropy_message_le_complexity_mul_log_two_of_measure, CommunicationComplexity.Functions.Disjointness.RandomizedLowerBound.message, CommunicationComplexity.Functions.Disjointness.RandomizedLowerBound.xVector, CommunicationComplexity.Functions.Disjointness.RandomizedLowerBound.yVector, ProbabilityTheory.condMutualInfo_le_entropy_right}
\difficulty{2}
Let $n$ be a positive integer and let $p$ be a deterministic two‑party communication
protocol on disjointness inputs of length $n$. Sample from the hard distribution
conditioned on the event that the two generated sets are disjoint, and write $X, Y
\colon \{1,\dots,n\} \to \{0,1\}$ for Alice's and Bob's resulting input vectors and $M$
for the transcript that $p$ produces on the corresponding input pair. Then the
conditional mutual information between $X$ and $M$ given $Y$, computed under this
conditioned distribution, satisfies
\[
  I(X ; M \mid Y) \le \mathrm{complexity}(p)\cdot \log 2,
\]
where $\mathrm{complexity}(p)$ is the worst‑case number of bits exchanged by $p$ and
$\log$ denotes the natural logarithm.
\end{theorem}

\begin{theorem}[Conditional information between Bob's input and the transcript]
\lean{CommunicationComplexity.Functions.Disjointness.RandomizedLowerBound.yVector_message_info_le_complexity_mul_log_two}
\label{CommunicationComplexity.Functions.Disjointness.RandomizedLowerBound.yVector_message_info_le_complexity_mul_log_two}
\leanok
\uses{CommunicationComplexity.Deterministic.Protocol.complexity, CommunicationComplexity.Functions.Disjointness.RandomizedLowerBound.ProtocolType, CommunicationComplexity.Functions.Disjointness.RandomizedLowerBound.disjointCondMeasure, CommunicationComplexity.Functions.Disjointness.RandomizedLowerBound.entropy_message_le_complexity_mul_log_two_of_measure, CommunicationComplexity.Functions.Disjointness.RandomizedLowerBound.message, CommunicationComplexity.Functions.Disjointness.RandomizedLowerBound.xVector, CommunicationComplexity.Functions.Disjointness.RandomizedLowerBound.yVector, ProbabilityTheory.condMutualInfo_le_entropy_right}
\difficulty{2}
Let $n$ be a positive integer, and let $p$ be a deterministic two-party communication
protocol for the length-$n$ set-disjointness problem, in which Alice holds a set $X
\subseteq \{0,\dots,n-1\}$ and Bob holds a set $Y \subseteq \{0,\dots,n-1\}$. Draw an
input pair from the hard distribution conditioned on the disjointness event $X \cap Y =
\varnothing$; write $x, y \in \{0,1\}^{n}$ for the indicator vectors of $X$ and $Y$
respectively, and write $M$ for the transcript produced by executing $p$ on the pair
$(X, Y)$. The conditional mutual information, measured in nats, between $y$ and $M$
given $x$ then satisfies \[ I(y \,;\, M \mid x) \;\le\; \mathrm{complexity}(p)\cdot\log
2, \] where $\mathrm{complexity}(p)$ denotes the worst-case number of bits exchanged by
$p$.
\end{theorem}

\begin{theorem}[Bounding Alice's special-coordinate information by the protocol cost]
\lean{CommunicationComplexity.Functions.Disjointness.RandomizedLowerBound.aliceInfoTerm_le_average_entropy_bound}
\label{CommunicationComplexity.Functions.Disjointness.RandomizedLowerBound.aliceInfoTerm_le_average_entropy_bound}
\leanok
\uses{CommunicationComplexity.Deterministic.Protocol.complexity, CommunicationComplexity.Functions.Disjointness.RandomizedLowerBound.ProtocolType, CommunicationComplexity.Functions.Disjointness.RandomizedLowerBound.aliceInfoTerm, CommunicationComplexity.Functions.Disjointness.RandomizedLowerBound.aliceInfoTerm_le_average_xVector_info, CommunicationComplexity.Functions.Disjointness.RandomizedLowerBound.xVector_message_info_le_complexity_mul_log_two}
\difficulty{3}
Let $n$ be a positive integer and let $p$ be a deterministic two-party communication
protocol for disjointness on inputs of length $n$, with communication complexity $\ell$.
Draw a hard sample and condition on the event $D$ that Alice's generated set $X$ and
Bob's generated set $Y$ are disjoint; let $T$ be the special coordinate, $X_T$ Alice's
bit at $T$, and $M$ the transcript obtained by running $p$ on the generated input pair
$(X,Y)$. Then, under this disjointness-conditioned hard distribution, the conditional
mutual information between Alice's special-coordinate bit and the transcript, given the
special coordinate together with Alice's bits before it and Bob's bits from it onward,
satisfies
\[
I\bigl(X_T : M \mid T,\, X_{<T},\, Y_{\ge T}\bigr) \;\le\; \frac{\ell \,\log 2}{n}.
\]
\end{theorem}

\begin{theorem}[Information bound for Bob's special-coordinate bit]
\lean{CommunicationComplexity.Functions.Disjointness.RandomizedLowerBound.bobInfoTerm_le_average_entropy_bound}
\label{CommunicationComplexity.Functions.Disjointness.RandomizedLowerBound.bobInfoTerm_le_average_entropy_bound}
\leanok
\uses{CommunicationComplexity.Deterministic.Protocol.complexity, CommunicationComplexity.Functions.Disjointness.RandomizedLowerBound.ProtocolType, CommunicationComplexity.Functions.Disjointness.RandomizedLowerBound.bobInfoTerm, CommunicationComplexity.Functions.Disjointness.RandomizedLowerBound.bobInfoTerm_le_average_yVector_info, CommunicationComplexity.Functions.Disjointness.RandomizedLowerBound.yVector_message_info_le_complexity_mul_log_two}
\difficulty{3}
Fix a positive integer $n$ and a deterministic communication protocol $p$ for
disjointness inputs of length $n$, with communication complexity $\ell$. Sample an input
from the hard distribution conditioned on the event that Alice's and Bob's generated
sets are disjoint; let $M$ denote the transcript that $p$ produces on the sampled input
pair, let $T$ denote the special coordinate, let $Y_T$ denote Bob's bit at the special
coordinate, and let $(T, X_{\le T}, Y_{>T})$ denote Bob's special-coordinate
conditioning variable. Then the conditional mutual information between $Y_T$ and $M$
given $(T, X_{\le T}, Y_{>T})$ satisfies
\[
I\!\left(Y_T : M \mid T,\, X_{\le T},\, Y_{>T}\right) \le \frac{\ell \, \log 2}{n}.
\]
\end{theorem}

\begin{theorem}[Claim-information bound for disjointness protocols]
\lean{CommunicationComplexity.Functions.Disjointness.RandomizedLowerBound.claimInfo_le_average_info_upper}
\label{CommunicationComplexity.Functions.Disjointness.RandomizedLowerBound.claimInfo_le_average_info_upper}
\leanok
\uses{CommunicationComplexity.Deterministic.Protocol.complexity, CommunicationComplexity.Functions.Disjointness.RandomizedLowerBound.ProtocolType, CommunicationComplexity.Functions.Disjointness.RandomizedLowerBound.aliceInfoTerm, CommunicationComplexity.Functions.Disjointness.RandomizedLowerBound.aliceInfoTerm_le_average_entropy_bound, CommunicationComplexity.Functions.Disjointness.RandomizedLowerBound.bobInfoTerm, CommunicationComplexity.Functions.Disjointness.RandomizedLowerBound.bobInfoTerm_le_average_entropy_bound, CommunicationComplexity.Functions.Disjointness.RandomizedLowerBound.claimInfo}
\difficulty{3}
Fix a positive integer $n$ and let $p$ be a deterministic two-party communication
protocol whose inputs are the pairs of subsets of $\{1,\dots,n\}$ held by Alice and Bob
in the set-disjointness problem. Draw a sample from the hard distribution on such inputs
conditioned on the event that Alice's set $X$ and Bob's set $Y$ are disjoint, let $M$ be
the transcript that $p$ produces on the resulting input pair, let $T$ be the special
coordinate carried by the sample, and write $X_T$ and $Y_T$ for Alice's and Bob's bits
at $T$; for a threshold coordinate, write $X_{<T}$, $X_{\le T}$, $Y_{>T}$, $Y_{\ge T}$
for the tuples of input bits at the coordinates strictly before, up to and including,
strictly after, and from $T$ onward. Then, writing $C(p)$ for the communication
complexity of $p$ (the worst-case number of bits exchanged) and $\ln$ for the natural
logarithm, the sum of the conditional mutual information of $X_T$ with $M$ given $(T,
X_{<T}, Y_{\ge T})$ and the conditional mutual information of $Y_T$ with $M$ given $(T,
X_{\le T}, Y_{>T})$ satisfies
\[
I\bigl(X_T ; M \mid T, X_{<T}, Y_{\ge T}\bigr) + I\bigl(Y_T ; M \mid T, X_{\le T},
Y_{>T}\bigr) \;\le\; \frac{2\, C(p)\, \ln 2}{n}.
\]
\end{theorem}

\begin{theorem}[Equality of the two suffix vectors when Bob's special bit is false]
\lean{CommunicationComplexity.Functions.Disjointness.RandomizedLowerBound.yGeSpecial_eq_yAfterSpecial_of_specialY_false}
\label{CommunicationComplexity.Functions.Disjointness.RandomizedLowerBound.yGeSpecial_eq_yAfterSpecial_of_specialY_false}
\leanok
\uses{CommunicationComplexity.Functions.Disjointness.RandomizedLowerBound.HardSample, CommunicationComplexity.Functions.Disjointness.RandomizedLowerBound.specialY, CommunicationComplexity.Functions.Disjointness.RandomizedLowerBound.yAfterSpecial, CommunicationComplexity.Functions.Disjointness.RandomizedLowerBound.yBit, CommunicationComplexity.Functions.Disjointness.RandomizedLowerBound.yGeSpecial}
\difficulty{4}
Let $\omega$ be a sample of the hard distribution, with special coordinate $T$, and
suppose Bob's bit at the special coordinate is false. Then Bob's bit-vector from the
special coordinate onward — his bits at all coordinates $i \ge T$, set to false at every
coordinate $i < T$ — coincides, as a function $\mathrm{Fin}\,n \to \{\text{false},
\text{true}\}$, with Bob's bit-vector strictly after the special coordinate — his bits
at all coordinates $i > T$, set to false elsewhere.
\end{theorem}

\begin{theorem}[Coincidence of Alice's and coarse conditioning when Bob's special bit is off]
\lean{CommunicationComplexity.Functions.Disjointness.RandomizedLowerBound.aliceClaimConditioning_eq_coarseConditioning_of_specialY_false}
\label{CommunicationComplexity.Functions.Disjointness.RandomizedLowerBound.aliceClaimConditioning_eq_coarseConditioning_of_specialY_false}
\leanok
\uses{CommunicationComplexity.Functions.Disjointness.RandomizedLowerBound.HardSample, CommunicationComplexity.Functions.Disjointness.RandomizedLowerBound.aliceClaimConditioning, CommunicationComplexity.Functions.Disjointness.RandomizedLowerBound.coarseConditioning, CommunicationComplexity.Functions.Disjointness.RandomizedLowerBound.specialY, CommunicationComplexity.Functions.Disjointness.RandomizedLowerBound.yGeSpecial_eq_yAfterSpecial_of_specialY_false}
\difficulty{1}
Let $\omega$ be a sample of the hard distribution, with special coordinate $T$, and let
$X_{<T}$ denote Alice's input bits at the coordinates strictly before $T$ (padded with
$\mathrm{false}$ elsewhere), $Y_{\ge T}$ Bob's input bits at the coordinates from $T$
onward (padded with $\mathrm{false}$ elsewhere), and $Y_{>T}$ Bob's input bits at the
coordinates strictly after $T$ (padded with $\mathrm{false}$ elsewhere). If Bob's bit at
the special coordinate is $\mathrm{false}$, then Alice's corrected special-coordinate
conditioning variable and the coarse conditioning variable coincide:
\[
(T,\; X_{<T},\; Y_{\ge T}) \;=\; (T,\; X_{<T},\; Y_{>T}).
\]
\end{theorem}

\begin{theorem}[Special bit vanishes under the special-bit-false conditioning]
\lean{CommunicationComplexity.Functions.Disjointness.RandomizedLowerBound.disjointSpecialYFalseMeasure_ae_specialY_false}
\label{CommunicationComplexity.Functions.Disjointness.RandomizedLowerBound.disjointSpecialYFalseMeasure_ae_specialY_false}
\leanok
\uses{CommunicationComplexity.Functions.Disjointness.RandomizedLowerBound.disjointSpecialYFalseMeasure, CommunicationComplexity.Functions.Disjointness.RandomizedLowerBound.specialY}
\difficulty{2}
Fix a positive integer $n$, and take the hard distribution on samples, conditioned first
on the event that Alice's and Bob's input sets are disjoint and then further conditioned
on the event that Bob's special bit $y_T$ is false. Under this doubly-conditioned
distribution, Bob's special bit $y_T$ equals $\mathtt{false}$ almost surely.
\end{theorem}

\begin{theorem}[Almost-sure agreement of Alice's and the coarse conditioning variable]
\lean{CommunicationComplexity.Functions.Disjointness.RandomizedLowerBound.aliceClaimConditioning_ae_eq_coarseConditioning_disjointSpecialYFalse}
\label{CommunicationComplexity.Functions.Disjointness.RandomizedLowerBound.aliceClaimConditioning_ae_eq_coarseConditioning_disjointSpecialYFalse}
\leanok
\uses{CommunicationComplexity.Functions.Disjointness.RandomizedLowerBound.aliceClaimConditioning, CommunicationComplexity.Functions.Disjointness.RandomizedLowerBound.aliceClaimConditioning_eq_coarseConditioning_of_specialY_false, CommunicationComplexity.Functions.Disjointness.RandomizedLowerBound.coarseConditioning, CommunicationComplexity.Functions.Disjointness.RandomizedLowerBound.disjointSpecialYFalseMeasure, CommunicationComplexity.Functions.Disjointness.RandomizedLowerBound.disjointSpecialYFalseMeasure_ae_specialY_false}
\difficulty{2}
Fix a positive integer $n$, and work on the sample space of the hard distribution for
set disjointness, where each sample determines a special coordinate $T \in
\{0,\dots,n-1\}$ together with Alice's and Bob's input bits. Write $X_{<T}$ for Alice's
bits at the coordinates strictly below $T$ (padded with $\mathrm{false}$ elsewhere),
$Y_{\ge T}$ for Bob's bits at the coordinates from $T$ onward (padded with
$\mathrm{false}$ elsewhere), and $Y_{>T}$ for Bob's bits at the coordinates strictly
above $T$ (padded with $\mathrm{false}$ elsewhere). Then, under the hard distribution
conditioned on the event that Alice's and Bob's input sets are disjoint and further on
Bob's bit at the special coordinate being $\mathrm{false}$, the triple $(T, X_{<T},
Y_{\ge T})$ equals the triple $(T, X_{<T}, Y_{>T})$ almost everywhere.
\end{theorem}

\begin{definition}[Aliceinfotermspecialyfalse]
\lean{CommunicationComplexity.Functions.Disjointness.RandomizedLowerBound.aliceInfoTermSpecialYFalse}
\label{CommunicationComplexity.Functions.Disjointness.RandomizedLowerBound.aliceInfoTermSpecialYFalse}
\leanok
\uses{CommunicationComplexity.Functions.Disjointness.RandomizedLowerBound.ProtocolType, CommunicationComplexity.Functions.Disjointness.RandomizedLowerBound.aliceClaimConditioning, CommunicationComplexity.Functions.Disjointness.RandomizedLowerBound.disjointSpecialYFalseMeasure, CommunicationComplexity.Functions.Disjointness.RandomizedLowerBound.message, CommunicationComplexity.Functions.Disjointness.RandomizedLowerBound.specialX}
The Alice information term after additionally conditioning on $Y_{T}=false$: $I(X_{T} : M | T, X_{<T}, Y_{≥T}, Y_{T}=0, D)$.
\end{definition}

\begin{definition}[Alicecoarseinfotermspecialyfalse]
\lean{CommunicationComplexity.Functions.Disjointness.RandomizedLowerBound.aliceCoarseInfoTermSpecialYFalse}
\label{CommunicationComplexity.Functions.Disjointness.RandomizedLowerBound.aliceCoarseInfoTermSpecialYFalse}
\leanok
\uses{CommunicationComplexity.Functions.Disjointness.RandomizedLowerBound.ProtocolType, CommunicationComplexity.Functions.Disjointness.RandomizedLowerBound.coarseConditioning, CommunicationComplexity.Functions.Disjointness.RandomizedLowerBound.disjointSpecialYFalseMeasure, CommunicationComplexity.Functions.Disjointness.RandomizedLowerBound.message, CommunicationComplexity.Functions.Disjointness.RandomizedLowerBound.specialX}
The same Alice $Y_{T}=false$ information term using the coarse conditioning $T, X_{<T}, Y_{>T}$.
\end{definition}

\begin{theorem}[Conditional KL to the uniform bit equals mutual information]
\lean{CommunicationComplexity.Functions.Disjointness.RandomizedLowerBound.condKLDiv_specialX_zVariable_eq_mutualInfo_zVariable}
\label{CommunicationComplexity.Functions.Disjointness.RandomizedLowerBound.condKLDiv_specialX_zVariable_eq_mutualInfo_zVariable}
\leanok
\uses{CommunicationComplexity.Functions.Disjointness.RandomizedLowerBound.ProtocolType, CommunicationComplexity.Functions.Disjointness.RandomizedLowerBound.disjointSpecialYFalseMeasure, CommunicationComplexity.Functions.Disjointness.RandomizedLowerBound.disjointSpecialYFalseMeasure_isProbabilityMeasure, CommunicationComplexity.Functions.Disjointness.RandomizedLowerBound.disjointSpecialYFalseMeasure_specialX_law_eq_uniformBool, CommunicationComplexity.Functions.Disjointness.RandomizedLowerBound.specialX, CommunicationComplexity.Functions.Disjointness.RandomizedLowerBound.uniformBool, CommunicationComplexity.Functions.Disjointness.RandomizedLowerBound.uniformBool_toPMF_ne_zero, CommunicationComplexity.Functions.Disjointness.RandomizedLowerBound.zVariable}
\difficulty{4}
Fix a length $n \ge 1$ and a deterministic two-party protocol $p$ for the disjointness
problem on inputs of that length, and equip the hard-distribution sample space with the
law obtained by conditioning on the generated input pair being disjoint and, in
addition, on Bob's bit $Y_T$ at the special coordinate $T$ being $\mathrm{false}$. Write
$X_T$ for Alice's bit at the special coordinate, and let $Z = (M, T, X_{<T}, Y_{>T})$ be
the random variable that records the protocol transcript $M$ produced by running $p$,
the special coordinate $T$, Alice's input bits strictly before $T$, and Bob's input bits
strictly after $T$. Then the average over $Z$ of the Kullback–Leibler divergence between
the conditional law of $X_T$ given $Z$ and the uniform distribution on a single bit
equals the mutual information $I(X_T \, ; Z)$.
\end{theorem}

\begin{theorem}[Equality of fine and coarse Alice information terms]
\lean{CommunicationComplexity.Functions.Disjointness.RandomizedLowerBound.aliceInfoTermSpecialYFalse_eq_aliceCoarseInfoTermSpecialYFalse}
\label{CommunicationComplexity.Functions.Disjointness.RandomizedLowerBound.aliceInfoTermSpecialYFalse_eq_aliceCoarseInfoTermSpecialYFalse}
\leanok
\uses{CommunicationComplexity.Functions.Disjointness.RandomizedLowerBound.ProtocolType, CommunicationComplexity.Functions.Disjointness.RandomizedLowerBound.aliceClaimConditioning, CommunicationComplexity.Functions.Disjointness.RandomizedLowerBound.aliceClaimConditioning_ae_eq_coarseConditioning_disjointSpecialYFalse, CommunicationComplexity.Functions.Disjointness.RandomizedLowerBound.aliceCoarseInfoTermSpecialYFalse, CommunicationComplexity.Functions.Disjointness.RandomizedLowerBound.aliceInfoTermSpecialYFalse, CommunicationComplexity.Functions.Disjointness.RandomizedLowerBound.coarseConditioning, CommunicationComplexity.Functions.Disjointness.RandomizedLowerBound.disjointSpecialYFalseMeasure, CommunicationComplexity.Functions.Disjointness.RandomizedLowerBound.disjointSpecialYFalseMeasure_isProbabilityMeasure, CommunicationComplexity.Functions.Disjointness.RandomizedLowerBound.message, CommunicationComplexity.Functions.Disjointness.RandomizedLowerBound.specialX, ProbabilityTheory.condMutualInfo_congr_ae_finite}
\difficulty{3}
Fix a positive integer $n$ and a deterministic communication protocol $p$ for
disjointness inputs of length $n$, and let $\mu$ be the hard distribution on samples,
conditioned on the generated input pair $(X,Y)$ being disjoint and on Bob's special bit
$Y_T$ being false. Write $X_T$ for Alice's bit at the special coordinate $T$ and $\Pi$
for the transcript that $p$ produces on the generated input. Then the conditional mutual
information of $X_T$ and $\Pi$ given the triple $(T,\,X_{<T},\,Y_{\ge T})$ equals the
conditional mutual information of $X_T$ and $\Pi$ given the coarser triple
$(T,\,X_{<T},\,Y_{>T})$, both computed under $\mu$; that is,
\[
I\big(X_T ; \Pi \mid T, X_{<T}, Y_{\ge T}\big) \;=\; I\big(X_T ; \Pi \mid T, X_{<T},
Y_{>T}\big).
\]
Here $X_{<T}$ denotes Alice's bits strictly before the special coordinate (padded with
false elsewhere), $Y_{>T}$ Bob's bits strictly after it, and $Y_{\ge T}$ Bob's bits from
the special coordinate onward.
\end{theorem}

\begin{theorem}[Fiber decomposition of the averaged Alice KL cost]
\lean{CommunicationComplexity.Functions.Disjointness.RandomizedLowerBound.integral_xFiberKL_disjointSpecialYFalse_eq_sum_zVariable}
\label{CommunicationComplexity.Functions.Disjointness.RandomizedLowerBound.integral_xFiberKL_disjointSpecialYFalse_eq_sum_zVariable}
\leanok
\uses{CommunicationComplexity.FiniteMeasureSpace.integral_comp_eq_sum_measureReal_fibers, CommunicationComplexity.Functions.Disjointness.RandomizedLowerBound.ProtocolType, CommunicationComplexity.Functions.Disjointness.RandomizedLowerBound.ZType, CommunicationComplexity.Functions.Disjointness.RandomizedLowerBound.disjointSpecialYFalseMeasure, CommunicationComplexity.Functions.Disjointness.RandomizedLowerBound.disjointSpecialYFalseMeasure_isProbabilityMeasure, CommunicationComplexity.Functions.Disjointness.RandomizedLowerBound.xFiberKL, CommunicationComplexity.Functions.Disjointness.RandomizedLowerBound.zFiber, CommunicationComplexity.Functions.Disjointness.RandomizedLowerBound.zVariable}
\difficulty{2}
Fix a positive integer $n$ and a deterministic protocol $p$ for set disjointness on
subsets of $\{1,\dots,n\}$. Let $\mu$ denote the hard distribution over samples,
conditioned both on the two generated input sets being disjoint and on Bob's bit at the
special coordinate being $\mathrm{false}$. Let $Z$ be the statistic sending each sample
$\omega$ to the tuple consisting of the transcript of $p$ on its input pair, the special
coordinate, Alice's input bits strictly before that coordinate, and Bob's input bits
strictly after it; the map $Z$ takes only finitely many values. For each such value $z$,
let $\mathrm{KL}(z)$ be the real-valued Kullback–Leibler divergence, from the uniform
distribution on one bit, of the conditional law of Alice's special bit given $Z = z$.
Then the $\mu$-average of the KL cost, evaluated along $Z$, splits as a finite
fiber-weighted sum:
\[
\int \mathrm{KL}\bigl(Z(\omega)\bigr)\, d\mu(\omega)
  \;=\; \sum_{z}\, \mu\bigl(Z^{-1}(z)\bigr)\,\mathrm{KL}(z),
\]
where $z$ ranges over all values of $Z$ and $\mu\bigl(Z^{-1}(z)\bigr)$ is the
$\mu$-measure of the fiber $\{\omega : Z(\omega) = z\}$.
\end{theorem}

\begin{theorem}[Fiber-wise decomposition of the averaged Alice-bit KL cost]
\lean{CommunicationComplexity.Functions.Disjointness.RandomizedLowerBound.integral_xFiberKL_disjointSpecialYFalse_eq_sum_zVariable_klDiv}
\label{CommunicationComplexity.Functions.Disjointness.RandomizedLowerBound.integral_xFiberKL_disjointSpecialYFalse_eq_sum_zVariable_klDiv}
\leanok
\uses{CommunicationComplexity.Functions.Disjointness.RandomizedLowerBound.ProtocolType, CommunicationComplexity.Functions.Disjointness.RandomizedLowerBound.ZType, CommunicationComplexity.Functions.Disjointness.RandomizedLowerBound.disjointSpecialYFalseMeasure, CommunicationComplexity.Functions.Disjointness.RandomizedLowerBound.integral_xFiberKL_disjointSpecialYFalse_eq_sum_zVariable, CommunicationComplexity.Functions.Disjointness.RandomizedLowerBound.specialX, CommunicationComplexity.Functions.Disjointness.RandomizedLowerBound.uniformBool, CommunicationComplexity.Functions.Disjointness.RandomizedLowerBound.xFiberKL, CommunicationComplexity.Functions.Disjointness.RandomizedLowerBound.xFiberKL_eq_disjointSpecialYFalseMeasure_cond_zVariable_klDiv_of_ne_zero, CommunicationComplexity.Functions.Disjointness.RandomizedLowerBound.zFiber, CommunicationComplexity.Functions.Disjointness.RandomizedLowerBound.zVariable}
\difficulty{3}
Fix a positive integer $n$ and a deterministic communication protocol $p$ whose inputs
are pairs of subsets of an $n$-element coordinate set. Let $\mu$ be the hard
distribution conditioned both on the two generated input sets being disjoint and on
Bob's bit at the special coordinate being $\mathrm{false}$. For a sample $\omega$, let
$Z(\omega)$ record the transcript of $p$ on the generated input, the special coordinate
$T$, Alice's input bits at the coordinates strictly before $T$, and Bob's input bits at
the coordinates strictly after $T$, and let $\mathcal Z$ be the finite set of values it
takes; for each $z \in \mathcal Z$ write $F_z = \{\omega : Z(\omega) = z\}$ for the
corresponding fiber. For a value $z$, let $\kappa(z)$ be the Kullback–Leibler
divergence, taken as a real number, of the law of Alice's bit at the special coordinate
— obtained from the ambient sample measure conditioned on $F_z$ — from the uniform law
on one bit. Then
\[ \int \kappa\big(Z(\omega)\big)\, d\mu(\omega) \;=\; \sum_{z \in \mathcal Z} \mu(F_z)\; D_{\mathrm{KL}}\!\Big( \mathcal L_\mu\big(\text{Alice's special bit} \mid Z = z\big) \,\Big\|\, \mathrm{Unif}\{0,1\} \Big), \]
where $\mathcal L_\mu(\,\cdot \mid Z = z)$ denotes the law under $\mu$ conditioned on
the fiber $F_z$.
\end{theorem}

\begin{theorem}[Averaged fiberwise KL as a conditional KL divergence]
\lean{CommunicationComplexity.Functions.Disjointness.RandomizedLowerBound.integral_xFiberKL_disjointSpecialYFalse_eq_condKLDiv_zVariable}
\label{CommunicationComplexity.Functions.Disjointness.RandomizedLowerBound.integral_xFiberKL_disjointSpecialYFalse_eq_condKLDiv_zVariable}
\leanok
\uses{CommunicationComplexity.Functions.Disjointness.RandomizedLowerBound.ProtocolType, CommunicationComplexity.Functions.Disjointness.RandomizedLowerBound.disjointSpecialYFalseMeasure, CommunicationComplexity.Functions.Disjointness.RandomizedLowerBound.disjointSpecialYFalseMeasure_isProbabilityMeasure, CommunicationComplexity.Functions.Disjointness.RandomizedLowerBound.integral_xFiberKL_disjointSpecialYFalse_eq_sum_zVariable_klDiv, CommunicationComplexity.Functions.Disjointness.RandomizedLowerBound.specialX, CommunicationComplexity.Functions.Disjointness.RandomizedLowerBound.uniformBool, CommunicationComplexity.Functions.Disjointness.RandomizedLowerBound.uniformBool_toPMF_ne_zero, CommunicationComplexity.Functions.Disjointness.RandomizedLowerBound.xFiberKL, CommunicationComplexity.Functions.Disjointness.RandomizedLowerBound.zFiber, CommunicationComplexity.Functions.Disjointness.RandomizedLowerBound.zVariable, ProbabilityTheory.toReal_klDiv_map_bool_eq_KLDiv_of_measureReal_ne_zero}
\difficulty{4}
Fix a positive integer $n$ and a deterministic communication protocol $p$ for the
$n$-coordinate disjointness problem, and let $\mu$ be the hard distribution on samples
$\omega$ conditioned on the two generated input sets being disjoint and on Bob's bit at
the special coordinate $T$ being false. Write $X_T(\omega)$ for Alice's bit at the
special coordinate, and let $Z(\omega)$ be the summary variable recording the transcript
of $p$ on the generated input pair together with the special coordinate $T$, Alice's
coordinates strictly before $T$, and Bob's coordinates strictly after $T$. For a value
$z$ of $Z$, let $\mathcal{L}(X_T \mid Z = z)$ denote the law of $X_T$ under the hard
distribution conditioned on the fiber $\{Z = z\}$, and let $U$ be the uniform law on a
single bit. Then
\[
\int D_{\mathrm{KL}}\!\big(\mathcal{L}(X_T \mid Z = Z(\omega)) \,\big\|\, U\big)\,
d\mu(\omega)
  \;=\; D_{\mathrm{KL}}\big(X_T \,\big\|\, U \,\big|\, Z\big)_{\mu},
\]
so that the $\mu$-average of the fiberwise Kullback–Leibler divergences from uniform
equals the conditional Kullback–Leibler divergence of $X_T$ given $Z$, measured from
uniform, under $\mu$.
\end{theorem}

\begin{definition}[Flipspecialx]
\lean{CommunicationComplexity.Functions.Disjointness.RandomizedLowerBound.flipSpecialX}
\label{CommunicationComplexity.Functions.Disjointness.RandomizedLowerBound.flipSpecialX}
\leanok
\uses{CommunicationComplexity.Functions.Disjointness.RandomizedLowerBound.HardSample}
Flip only Alice's special-coordinate bit.  This preserves the coarse conditioning data $T, X_{<T}, Y_{>T}$ and toggles $X_{T}$.
\end{definition}

\begin{theorem}[Double bit-flip is the identity]
\lean{CommunicationComplexity.Functions.Disjointness.RandomizedLowerBound.flipSpecialX_flipSpecialX}
\label{CommunicationComplexity.Functions.Disjointness.RandomizedLowerBound.flipSpecialX_flipSpecialX}
\leanok
\uses{CommunicationComplexity.Functions.Disjointness.RandomizedLowerBound.HardSample, CommunicationComplexity.Functions.Disjointness.RandomizedLowerBound.flipSpecialX}
\difficulty{2}
Let $\omega$ be a sample of the hard distribution, recording a special coordinate $T$,
Alice's and Bob's bits $x_T, y_T$ at that coordinate, and an assignment of the remaining
coordinates. Toggling Alice's special-coordinate bit twice — that is, applying to
$\omega$ the map that negates $x_T$ while leaving $T$, $y_T$, and the other coordinates
unchanged, and then applying that same map again — recovers $\omega$ exactly.
\end{theorem}

\begin{theorem}[Preimage of a singleton under the special-bit flip]
\lean{CommunicationComplexity.Functions.Disjointness.RandomizedLowerBound.flipSpecialX_preimage_singleton}
\label{CommunicationComplexity.Functions.Disjointness.RandomizedLowerBound.flipSpecialX_preimage_singleton}
\leanok
\uses{CommunicationComplexity.Functions.Disjointness.RandomizedLowerBound.HardSample, CommunicationComplexity.Functions.Disjointness.RandomizedLowerBound.flipSpecialX, CommunicationComplexity.Functions.Disjointness.RandomizedLowerBound.flipSpecialX_flipSpecialX}
\difficulty{3}
Fix $n \ge 1$, and let $\phi$ be the map on hard samples that flips only Alice's
special-coordinate bit, sending a sample $\omega = (T, x_T, y_T, \mathrm{other})$ to
$(T, \lnot x_T, y_T, \mathrm{other})$. Then for every hard sample $\omega$, the preimage
of the singleton $\{\omega\}$ under $\phi$ is exactly the singleton $\{\phi(\omega)\}$;
that is, $\phi^{-1}(\{\omega\}) = \{\phi(\omega)\}$.
\end{theorem}

\begin{theorem}[Flipping Alice's special bit preserves singleton mass]
\lean{CommunicationComplexity.Functions.Disjointness.RandomizedLowerBound.volume_measureReal_singleton_flipSpecialX}
\label{CommunicationComplexity.Functions.Disjointness.RandomizedLowerBound.volume_measureReal_singleton_flipSpecialX}
\leanok
\uses{CommunicationComplexity.Functions.Disjointness.RandomizedLowerBound.HardSample, CommunicationComplexity.Functions.Disjointness.RandomizedLowerBound.flipSpecialX, CommunicationComplexity.uniformOn_univ_measureReal_eq_card_subtype}
\difficulty{2}
Fix a positive integer $n$, and let $\omega$ be a point of the hard sample space, i.e. a
choice of special coordinate $T \in \{0,\dots,n-1\}$, bits $X_T, Y_T \in \{0,1\}$, and
an assignment of one of the three disjointness states to each coordinate. Write
$\omega'$ for the point obtained from $\omega$ by flipping only Alice's
special-coordinate bit $X_T$, leaving $T$, $Y_T$, and all the other-coordinate data
unchanged. Then the ambient volume measure assigns the singleton $\{\omega'\}$ the same
real-valued mass as the singleton $\{\omega\}$.
\end{theorem}

\begin{theorem}[Bit-flip on Alice's special coordinate preserves volume]
\lean{CommunicationComplexity.Functions.Disjointness.RandomizedLowerBound.volume_measurePreserving_flipSpecialX}
\label{CommunicationComplexity.Functions.Disjointness.RandomizedLowerBound.volume_measurePreserving_flipSpecialX}
\leanok
\uses{CommunicationComplexity.Functions.Disjointness.RandomizedLowerBound.flipSpecialX, CommunicationComplexity.Functions.Disjointness.RandomizedLowerBound.flipSpecialX_preimage_singleton, CommunicationComplexity.Functions.Disjointness.RandomizedLowerBound.volume_measureReal_singleton_flipSpecialX}
\difficulty{3}
Fix a positive integer $n$, and consider the finite sample space whose points are the
hard samples $\omega = (T, x_T, y_T, \mathrm{other})$, with $T \in \{0, \dots, n-1\}$,
bits $x_T, y_T \in \{\text{false}, \text{true}\}$, and a coordinate labelling
$\mathrm{other} \colon \{0,\dots,n-1\} \to \{\text{neither}, \text{leftOnly},
\text{rightOnly}\}$. Let $f$ be the self-map of this space that toggles Alice's
special-coordinate bit, sending $\omega$ to $(T, \lnot x_T, y_T, \mathrm{other})$ while
leaving $T$, $y_T$, and $\mathrm{other}$ fixed. Then $f$ is measure-preserving for the
volume (uniform counting) measure on the finite sample space: it is measurable and
pushes the volume measure forward to itself.
\end{theorem}

\begin{theorem}[Flipping Alice's special bit negates it]
\lean{CommunicationComplexity.Functions.Disjointness.RandomizedLowerBound.specialX_flipSpecialX}
\label{CommunicationComplexity.Functions.Disjointness.RandomizedLowerBound.specialX_flipSpecialX}
\leanok
\uses{CommunicationComplexity.Functions.Disjointness.RandomizedLowerBound.HardSample, CommunicationComplexity.Functions.Disjointness.RandomizedLowerBound.flipSpecialX, CommunicationComplexity.Functions.Disjointness.RandomizedLowerBound.specialX}
\difficulty{1}
Let $\omega$ be a sample from the hard distribution, and write $X_T(\omega)$ for Alice's
bit at the special coordinate of $\omega$. Let $\omega'$ be the sample obtained from
$\omega$ by toggling Alice's special-coordinate bit while leaving all other data
unchanged. Then $X_T(\omega') = \lnot\, X_T(\omega)$; that is, reading Alice's special
bit after the flip yields the negation of Alice's special bit before it.
\end{theorem}

\begin{theorem}[Bob's special bit is unaffected by flipping Alice's]
\lean{CommunicationComplexity.Functions.Disjointness.RandomizedLowerBound.specialY_flipSpecialX}
\label{CommunicationComplexity.Functions.Disjointness.RandomizedLowerBound.specialY_flipSpecialX}
\leanok
\uses{CommunicationComplexity.Functions.Disjointness.RandomizedLowerBound.HardSample, CommunicationComplexity.Functions.Disjointness.RandomizedLowerBound.flipSpecialX, CommunicationComplexity.Functions.Disjointness.RandomizedLowerBound.specialY}
\difficulty{1}
Let $\omega$ be a sample of the hard distribution, and let $\omega'$ be obtained from
$\omega$ by flipping only Alice's special-coordinate bit (toggling $X_T$ while leaving
$T$, the coordinate types, and Bob's bits unchanged). Then Bob's bit at the special
coordinate is the same for both samples: the special-$Y$ value of $\omega'$ equals the
special-$Y$ value of $\omega$.
\end{theorem}

\begin{theorem}[Flipping Alice's special bit preserves her prefix]
\lean{CommunicationComplexity.Functions.Disjointness.RandomizedLowerBound.xBeforeSpecial_flipSpecialX}
\label{CommunicationComplexity.Functions.Disjointness.RandomizedLowerBound.xBeforeSpecial_flipSpecialX}
\leanok
\uses{CommunicationComplexity.Functions.Disjointness.RandomizedLowerBound.HardSample, CommunicationComplexity.Functions.Disjointness.RandomizedLowerBound.flipSpecialX, CommunicationComplexity.Functions.Disjointness.RandomizedLowerBound.xBeforeSpecial, CommunicationComplexity.Functions.Disjointness.RandomizedLowerBound.xBit}
\difficulty{3}
Let $n$ be a positive integer and let $\omega$ be a hard-distribution sample, consisting
of a special coordinate $T \in \{0,\dots,n-1\}$, Alice's special bit $x_T$, Bob's
special bit $y_T$, and a coordinate assignment. Then flipping only Alice's special bit
leaves her padded prefix before the special coordinate unchanged; that is, the function
$i \mapsto \bigl(x\text{-bit of }\omega\text{ at }i\text{ if }i < T,\ \text{else
}\mathrm{false}\bigr)$ is the same whether it is computed from $\omega$ or from the
sample obtained by toggling $x_T$ alone.
\end{theorem}

\begin{theorem}[Flipping Alice's special bit fixes Bob's post-special input]
\lean{CommunicationComplexity.Functions.Disjointness.RandomizedLowerBound.yAfterSpecial_flipSpecialX}
\label{CommunicationComplexity.Functions.Disjointness.RandomizedLowerBound.yAfterSpecial_flipSpecialX}
\leanok
\uses{CommunicationComplexity.Functions.Disjointness.RandomizedLowerBound.HardSample, CommunicationComplexity.Functions.Disjointness.RandomizedLowerBound.flipSpecialX, CommunicationComplexity.Functions.Disjointness.RandomizedLowerBound.yAfterSpecial, CommunicationComplexity.Functions.Disjointness.RandomizedLowerBound.yBit}
\difficulty{2}
Let $n$ be a positive integer, and let $\omega$ be a sample of the hard distribution,
consisting of a distinguished coordinate $T \in \{0,\dots,n-1\}$, Alice's special bit
$x_T \in \{0,1\}$, Bob's special bit $y_T \in \{0,1\}$, and an assignment $i \mapsto
c_i$ of a coordinate type to each index $i$. Let Bob's bit at coordinate $i$ be $y_T$
when $i = T$ and otherwise the bit determined by the coordinate type $c_i$, and let
$\beta(\omega) \in \{0,1\}^n$ be Bob's post-special vector, whose $i$-th entry equals
Bob's bit at $i$ when $T < i$ and is $0$ otherwise. Writing $\omega'$ for the sample
obtained from $\omega$ by replacing $x_T$ with its negation and leaving $T$, $y_T$, and
the coordinate-type assignment unchanged, we have \[ \beta(\omega') = \beta(\omega); \]
that is, toggling only Alice's special-coordinate bit leaves Bob's input bits after the
special coordinate unchanged.
\end{theorem}

\begin{theorem}[Invariance of the coarse conditioning under flipping Alice's special bit]
\lean{CommunicationComplexity.Functions.Disjointness.RandomizedLowerBound.coarseConditioning_flipSpecialX}
\label{CommunicationComplexity.Functions.Disjointness.RandomizedLowerBound.coarseConditioning_flipSpecialX}
\leanok
\uses{CommunicationComplexity.Functions.Disjointness.RandomizedLowerBound.HardSample, CommunicationComplexity.Functions.Disjointness.RandomizedLowerBound.coarseConditioning, CommunicationComplexity.Functions.Disjointness.RandomizedLowerBound.flipSpecialX, CommunicationComplexity.Functions.Disjointness.RandomizedLowerBound.specialCoordinate, CommunicationComplexity.Functions.Disjointness.RandomizedLowerBound.xBeforeSpecial_flipSpecialX, CommunicationComplexity.Functions.Disjointness.RandomizedLowerBound.yAfterSpecial_flipSpecialX}
\difficulty{2}
Fix $n \ge 1$, and let $\omega$ be a sample of the hard distribution: a special
coordinate $T \in \{0,1,\dots,n-1\}$, bits $x_T$ and $y_T$ for Alice and Bob at that
coordinate, and a disjointness type assigned to every coordinate. From $\omega$ form the
coarse conditioning $(T,\, X_{<T},\, Y_{>T})$, where $X_{<T}$ is Alice's input bits at
the coordinates strictly below $T$ (padded with $\mathrm{false}$ elsewhere) and $Y_{>T}$
is Bob's input bits at the coordinates strictly above $T$ (padded with $\mathrm{false}$
elsewhere). Let $\omega'$ be obtained from $\omega$ by toggling only Alice's special bit
$x_T$, leaving $T$, $y_T$, and all coordinate types unchanged. Then $\omega$ and
$\omega'$ produce the same coarse conditioning: $(T,\, X_{<T},\, Y_{>T})$ is identical
for $\omega'$ and $\omega$.
\end{theorem}

\begin{theorem}[Uniformity of Alice's special bit within a coarse fiber]
\lean{CommunicationComplexity.Functions.Disjointness.RandomizedLowerBound.volume_measureReal_specialYFalse_inter_coarseConditioning_inter_specialX}
\label{CommunicationComplexity.Functions.Disjointness.RandomizedLowerBound.volume_measureReal_specialYFalse_inter_coarseConditioning_inter_specialX}
\leanok
\uses{CommunicationComplexity.Functions.Disjointness.RandomizedLowerBound.HardSample, CommunicationComplexity.Functions.Disjointness.RandomizedLowerBound.coarseConditioning, CommunicationComplexity.Functions.Disjointness.RandomizedLowerBound.coarseConditioning_flipSpecialX, CommunicationComplexity.Functions.Disjointness.RandomizedLowerBound.flipSpecialX, CommunicationComplexity.Functions.Disjointness.RandomizedLowerBound.specialX, CommunicationComplexity.Functions.Disjointness.RandomizedLowerBound.specialX_flipSpecialX, CommunicationComplexity.Functions.Disjointness.RandomizedLowerBound.specialY, CommunicationComplexity.Functions.Disjointness.RandomizedLowerBound.specialY_flipSpecialX, CommunicationComplexity.Functions.Disjointness.RandomizedLowerBound.volume_measurePreserving_flipSpecialX}
\difficulty{5}
Work with the volume (counting) measure on the finite space of hard-distribution
samples, in which each sample records a special coordinate $T$, Alice's and Bob's
special bits $X_T$ and $Y_T$, and the remaining per-coordinate data. Fix a bit $b$ and a
triple $c=(t,u,v)$, where $t$ is a coordinate and $u,v$ are bit-vectors, playing the
role of a value of the coarse conditioning $(T,\,X_{<T},\,Y_{>T})$. Then the mass of the
event that Bob's special bit satisfies $Y_T=\mathrm{false}$, the coarse conditioning
equals $c$, and Alice's special bit satisfies $X_T=b$ is exactly $\tfrac12$ times the
mass of the event that $Y_T=\mathrm{false}$ and the coarse conditioning equals $c$.
\end{theorem}

\begin{theorem}[Uniformity of Alice's special bit over coarse fibers]
\lean{CommunicationComplexity.Functions.Disjointness.RandomizedLowerBound.disjointSpecialYFalseMeasure_measureReal_specialX_inter_coarseConditioning}
\label{CommunicationComplexity.Functions.Disjointness.RandomizedLowerBound.disjointSpecialYFalseMeasure_measureReal_specialX_inter_coarseConditioning}
\leanok
\uses{CommunicationComplexity.Functions.Disjointness.RandomizedLowerBound.HardSample, CommunicationComplexity.Functions.Disjointness.RandomizedLowerBound.coarseConditioning, CommunicationComplexity.Functions.Disjointness.RandomizedLowerBound.disjointCondMeasure, CommunicationComplexity.Functions.Disjointness.RandomizedLowerBound.disjointEvent, CommunicationComplexity.Functions.Disjointness.RandomizedLowerBound.disjointSpecialYFalseMeasure, CommunicationComplexity.Functions.Disjointness.RandomizedLowerBound.mem_disjointEvent_of_specialY_eq_false, CommunicationComplexity.Functions.Disjointness.RandomizedLowerBound.specialX, CommunicationComplexity.Functions.Disjointness.RandomizedLowerBound.specialY, CommunicationComplexity.Functions.Disjointness.RandomizedLowerBound.volume_measureReal_specialYFalse_inter_coarseConditioning_inter_specialX}
\difficulty{4}
Fix a positive integer $n$, and let $\mu$ denote the hard distribution on samples
conditioned on the event that Alice's and Bob's generated input sets are disjoint and
further conditioned on Bob's special bit being $\mathsf{false}$. Write $X_T$ for Alice's
bit at the special coordinate and write $(T, X_{<T}, Y_{>T})$ for the coarse variable
recording the special coordinate together with Alice's bits strictly before it and Bob's
bits strictly after it. Then for every bit $b \in \{\mathsf{false}, \mathsf{true}\}$ and
every value $c$ of the coarse variable,
\[
\mu\big(\{X_T = b\} \cap \{(T, X_{<T}, Y_{>T}) = c\}\big) = \tfrac{1}{2}\,\mu\big(\{(T,
X_{<T}, Y_{>T}) = c\}\big),
\]
where the measures are taken as real numbers; that is, under $\mu$, Alice's special bit
is uniformly distributed over each fiber of the coarse variable.
\end{theorem}

\begin{theorem}[Independence of Alice's special bit from the coarse variable]
\lean{CommunicationComplexity.Functions.Disjointness.RandomizedLowerBound.mutualInfo_specialX_coarseConditioning_disjointSpecialYFalse_eq_zero}
\label{CommunicationComplexity.Functions.Disjointness.RandomizedLowerBound.mutualInfo_specialX_coarseConditioning_disjointSpecialYFalse_eq_zero}
\leanok
\uses{CommunicationComplexity.Functions.Disjointness.RandomizedLowerBound.HardSample, CommunicationComplexity.Functions.Disjointness.RandomizedLowerBound.coarseConditioning, CommunicationComplexity.Functions.Disjointness.RandomizedLowerBound.disjointSpecialYFalseMeasure, CommunicationComplexity.Functions.Disjointness.RandomizedLowerBound.disjointSpecialYFalseMeasure_isProbabilityMeasure, CommunicationComplexity.Functions.Disjointness.RandomizedLowerBound.disjointSpecialYFalseMeasure_measureReal_specialX_inter_coarseConditioning, CommunicationComplexity.Functions.Disjointness.RandomizedLowerBound.disjointSpecialYFalseMeasure_measureReal_specialX_singleton, CommunicationComplexity.Functions.Disjointness.RandomizedLowerBound.specialX, ProbabilityTheory.indepFun_of_measureReal_inter_preimage_singleton_eq_mul}
\difficulty{3}
Fix a positive integer $n$ and take the hard distribution on samples conditioned on the
event that Alice's set $X$ and Bob's set $Y$ are disjoint, and further conditioned on
Bob's special bit $Y_T$ being false. Under this conditional distribution the mutual
information between Alice's special bit $X_T$ and the coarse variable $(T,\, X_{<T},\,
Y_{>T})$ is zero:
\[
I\bigl(X_T \,;\, (T,\, X_{<T},\, Y_{>T})\bigr) = 0,
\]
so that $X_T$ is independent of $(T,\, X_{<T},\, Y_{>T})$.
\end{theorem}

\begin{theorem}[Full Z-information equals the coarse-conditioned Alice term]
\lean{CommunicationComplexity.Functions.Disjointness.RandomizedLowerBound.mutualInfo_specialX_zVariable_eq_aliceCoarseInfoTermSpecialYFalse}
\label{CommunicationComplexity.Functions.Disjointness.RandomizedLowerBound.mutualInfo_specialX_zVariable_eq_aliceCoarseInfoTermSpecialYFalse}
\leanok
\uses{CommunicationComplexity.Functions.Disjointness.RandomizedLowerBound.HardSample, CommunicationComplexity.Functions.Disjointness.RandomizedLowerBound.ProtocolType, CommunicationComplexity.Functions.Disjointness.RandomizedLowerBound.TranscriptType, CommunicationComplexity.Functions.Disjointness.RandomizedLowerBound.ZType, CommunicationComplexity.Functions.Disjointness.RandomizedLowerBound.ZType.ext, CommunicationComplexity.Functions.Disjointness.RandomizedLowerBound.ZType.specialCoordinate, CommunicationComplexity.Functions.Disjointness.RandomizedLowerBound.ZType.transcript, CommunicationComplexity.Functions.Disjointness.RandomizedLowerBound.ZType.xBefore, CommunicationComplexity.Functions.Disjointness.RandomizedLowerBound.ZType.yAfter, CommunicationComplexity.Functions.Disjointness.RandomizedLowerBound.aliceCoarseInfoTermSpecialYFalse, CommunicationComplexity.Functions.Disjointness.RandomizedLowerBound.coarseConditioning, CommunicationComplexity.Functions.Disjointness.RandomizedLowerBound.disjointSpecialYFalseMeasure, CommunicationComplexity.Functions.Disjointness.RandomizedLowerBound.disjointSpecialYFalseMeasure_isProbabilityMeasure, CommunicationComplexity.Functions.Disjointness.RandomizedLowerBound.message, CommunicationComplexity.Functions.Disjointness.RandomizedLowerBound.mutualInfo_specialX_coarseConditioning_disjointSpecialYFalse_eq_zero, CommunicationComplexity.Functions.Disjointness.RandomizedLowerBound.rawZVariable, CommunicationComplexity.Functions.Disjointness.RandomizedLowerBound.specialX, CommunicationComplexity.Functions.Disjointness.RandomizedLowerBound.zVariable, ProbabilityTheory.mutualInfo_comp_right_of_injective, ProbabilityTheory.mutualInfo_prod_right_eq_add}
\difficulty{5}
Let $n$ be a positive integer and let $p$ be a deterministic two-party communication
protocol for disjointness inputs of length $n$. Let $\mu$ denote the hard distribution
on inputs, conditioned both on the event that Alice's and Bob's generated input sets are
disjoint and on Bob's bit at the special coordinate being $\mathrm{false}$. Write $T$
for the special coordinate, $X_T$ for Alice's bit at coordinate $T$, $X_{<T}$ for the
string of Alice's input bits at coordinates below $T$ (set to $\mathrm{false}$ at all
other coordinates), $Y_{>T}$ for the string of Bob's input bits at coordinates above $T$
(set to $\mathrm{false}$ at all other coordinates), and $M$ for the transcript that $p$
produces on the sampled input pair. Then, with every information quantity computed under
$\mu$,
\[
I\big(X_T ; (M,\,T,\,X_{<T},\,Y_{>T})\big) \;=\; I\big(X_T ; M \mid
(T,\,X_{<T},\,Y_{>T})\big).
\]
\end{theorem}

\begin{theorem}[Averaged special-bit KL equals conditional mutual information]
\lean{CommunicationComplexity.Functions.Disjointness.RandomizedLowerBound.integral_xFiberKL_disjointSpecialYFalse_eq_aliceCoarseInfoTermSpecialYFalse}
\label{CommunicationComplexity.Functions.Disjointness.RandomizedLowerBound.integral_xFiberKL_disjointSpecialYFalse_eq_aliceCoarseInfoTermSpecialYFalse}
\leanok
\uses{CommunicationComplexity.Functions.Disjointness.RandomizedLowerBound.ProtocolType, CommunicationComplexity.Functions.Disjointness.RandomizedLowerBound.aliceCoarseInfoTermSpecialYFalse, CommunicationComplexity.Functions.Disjointness.RandomizedLowerBound.condKLDiv_specialX_zVariable_eq_mutualInfo_zVariable, CommunicationComplexity.Functions.Disjointness.RandomizedLowerBound.disjointSpecialYFalseMeasure, CommunicationComplexity.Functions.Disjointness.RandomizedLowerBound.integral_xFiberKL_disjointSpecialYFalse_eq_condKLDiv_zVariable, CommunicationComplexity.Functions.Disjointness.RandomizedLowerBound.mutualInfo_specialX_zVariable_eq_aliceCoarseInfoTermSpecialYFalse, CommunicationComplexity.Functions.Disjointness.RandomizedLowerBound.xFiberKL, CommunicationComplexity.Functions.Disjointness.RandomizedLowerBound.zVariable}
\difficulty{2}
Fix a positive integer $n$ and a deterministic protocol $p$ for the disjointness problem
on inputs of length $n$. Let $\mu$ denote the hard distribution on samples, conditioned
on the event that Alice's set and Bob's set are disjoint and on the event that Bob's bit
at the special coordinate $T$ equals $\mathtt{false}$; let $\E_\mu$ denote expectation
under $\mu$, and let $I_\mu$ denote conditional mutual information under $\mu$. Let $M$
denote the transcript produced by running $p$ on the generated input pair, let $X_T$
denote Alice's bit at the special coordinate, and let $Z = (M,\,T,\,X_{<T},\,Y_{>T})$
denote the fine variable consisting of the transcript together with the coarse data $(T,
X_{<T}, Y_{>T})$. Then the $\mu$-average of the Kullback–Leibler divergence
$D(\cdot\,\|\,\cdot)$ from the conditional law of $X_T$ given $Z$ to the uniform
distribution on $\{0,1\}$ equals the conditional mutual information between $X_T$ and
$M$ given the coarse data:
\[
\E_\mu\!\left[\, D\!\big(\operatorname{Law}_\mu(X_T \mid Z)\ \big\|\
\mathrm{Unif}\{0,1\}\big)\,\right]
\;=\; I_\mu\big(X_T \,;\, M \,\big|\, T,\, X_{<T},\, Y_{>T}\big).
\]
\end{theorem}

\begin{theorem}[Averaged fiber KL equals Alice's conditional mutual information]
\lean{CommunicationComplexity.Functions.Disjointness.RandomizedLowerBound.integral_xFiberKL_disjointSpecialYFalse_eq_aliceInfoTermSpecialYFalse}
\label{CommunicationComplexity.Functions.Disjointness.RandomizedLowerBound.integral_xFiberKL_disjointSpecialYFalse_eq_aliceInfoTermSpecialYFalse}
\leanok
\uses{CommunicationComplexity.Functions.Disjointness.RandomizedLowerBound.ProtocolType, CommunicationComplexity.Functions.Disjointness.RandomizedLowerBound.aliceInfoTermSpecialYFalse, CommunicationComplexity.Functions.Disjointness.RandomizedLowerBound.aliceInfoTermSpecialYFalse_eq_aliceCoarseInfoTermSpecialYFalse, CommunicationComplexity.Functions.Disjointness.RandomizedLowerBound.disjointSpecialYFalseMeasure, CommunicationComplexity.Functions.Disjointness.RandomizedLowerBound.integral_xFiberKL_disjointSpecialYFalse_eq_aliceCoarseInfoTermSpecialYFalse, CommunicationComplexity.Functions.Disjointness.RandomizedLowerBound.xFiberKL, CommunicationComplexity.Functions.Disjointness.RandomizedLowerBound.zVariable}
\difficulty{2}
Fix a positive integer $n$ and a deterministic communication protocol $p$ for
disjointness inputs of length $n$, and work under the hard distribution conditioned on
Alice's and Bob's generated input sets being disjoint and on Bob's bit $Y_T$ at the
special coordinate $T$ being $\mathrm{false}$. Write $M$ for the transcript that $p$
produces on the generated input, let $Z = (M,\, T,\, X_{<T},\, Y_{>T})$ be the
associated fiber variable, and let $U$ be the uniform law on one bit. Then the
expectation, taken over this conditioned distribution, of the Kullback–Leibler
divergence $D_{\mathrm{KL}}\!\big(P_{X_T \mid Z}\,\big\|\,U\big)$ of the conditional law
of Alice's special bit $X_T$ given $Z$ relative to $U$ equals the conditional mutual
information $I\big(X_T ; M \mid T,\, X_{<T},\, Y_{\ge T}\big)$ between $X_T$ and $M$
given Alice's conditioning variable $(T,\, X_{<T},\, Y_{\ge T})$, computed under the
same conditioned distribution.
\end{theorem}

\begin{theorem}[Alice-side reweighting of conditional mutual information]
\lean{CommunicationComplexity.Functions.Disjointness.RandomizedLowerBound.two_thirds_mul_aliceInfoTermSpecialYFalse_le_aliceInfoTerm}
\label{CommunicationComplexity.Functions.Disjointness.RandomizedLowerBound.two_thirds_mul_aliceInfoTermSpecialYFalse_le_aliceInfoTerm}
\leanok
\uses{CommunicationComplexity.Functions.Disjointness.RandomizedLowerBound.HardSample, CommunicationComplexity.Functions.Disjointness.RandomizedLowerBound.ProtocolType, CommunicationComplexity.Functions.Disjointness.RandomizedLowerBound.aliceClaimConditioning, CommunicationComplexity.Functions.Disjointness.RandomizedLowerBound.aliceClaimConditioningYFalseValues, CommunicationComplexity.Functions.Disjointness.RandomizedLowerBound.aliceInfoTerm, CommunicationComplexity.Functions.Disjointness.RandomizedLowerBound.aliceInfoTermSpecialYFalse, CommunicationComplexity.Functions.Disjointness.RandomizedLowerBound.disjointCondMeasure, CommunicationComplexity.Functions.Disjointness.RandomizedLowerBound.disjointCondMeasure_measureReal_specialY_false, CommunicationComplexity.Functions.Disjointness.RandomizedLowerBound.disjointSpecialYFalseMeasure, CommunicationComplexity.Functions.Disjointness.RandomizedLowerBound.message, CommunicationComplexity.Functions.Disjointness.RandomizedLowerBound.specialX, CommunicationComplexity.Functions.Disjointness.RandomizedLowerBound.specialY, CommunicationComplexity.Functions.Disjointness.RandomizedLowerBound.specialY_false_eq_preimage_aliceClaimConditioningYFalseValues, ProbabilityTheory.measureReal_mul_cond_condMutualInfo_le_condMutualInfo_of_event_eq_preimage}
\difficulty{4}
Fix a positive integer $n$ and a deterministic communication protocol $p$ for
disjointness inputs of length $n$. Consider the hard distribution over samples
conditioned on the generated input pair $(X, Y)$ being disjoint, and let $M$ be the
transcript that $p$ produces on that input, $X_T$ Alice's bit at the special coordinate
$T$, and $(T, X_{<T}, Y_{\ge T})$ Alice's corrected special-coordinate conditioning
variable. Writing $I_D(X_T : M \mid T, X_{<T}, Y_{\ge T})$ for the conditional mutual
information under this disjoint-conditioned distribution, and $I_{D_0}(X_T : M \mid T,
X_{<T}, Y_{\ge T})$ for the same conditional mutual information under the distribution
obtained by further conditioning on Bob's special bit $Y_T$ being false, one has
\[
\tfrac{2}{3}\, I_{D_0}(X_T : M \mid T, X_{<T}, Y_{\ge T}) \;\le\; I_D(X_T : M \mid T,
X_{<T}, Y_{\ge T}).
\]
\end{theorem}

\begin{theorem}[Average one-bit KL cost bounded by Alice's information term]
\lean{CommunicationComplexity.Functions.Disjointness.RandomizedLowerBound.integral_xFiberKL_disjointSpecialYFalse_le_three_halves_mul_aliceInfoTerm}
\label{CommunicationComplexity.Functions.Disjointness.RandomizedLowerBound.integral_xFiberKL_disjointSpecialYFalse_le_three_halves_mul_aliceInfoTerm}
\leanok
\uses{CommunicationComplexity.Functions.Disjointness.RandomizedLowerBound.ProtocolType, CommunicationComplexity.Functions.Disjointness.RandomizedLowerBound.aliceInfoTerm, CommunicationComplexity.Functions.Disjointness.RandomizedLowerBound.disjointSpecialYFalseMeasure, CommunicationComplexity.Functions.Disjointness.RandomizedLowerBound.integral_xFiberKL_disjointSpecialYFalse_eq_aliceInfoTermSpecialYFalse, CommunicationComplexity.Functions.Disjointness.RandomizedLowerBound.two_thirds_mul_aliceInfoTermSpecialYFalse_le_aliceInfoTerm, CommunicationComplexity.Functions.Disjointness.RandomizedLowerBound.xFiberKL, CommunicationComplexity.Functions.Disjointness.RandomizedLowerBound.zVariable}
\difficulty{2}
Fix a positive integer $n$ and a deterministic communication protocol $p$ for
set-disjointness, whose inputs are the pairs of subsets $X, Y \subseteq \{0, 1, \dots,
n-1\}$ produced by the hard distribution. Let $\mu$ be the hard distribution conditioned
both on the two generated sets being disjoint and on Bob's bit $Y_T$ at the special
coordinate $T$ being $\mathrm{false}$, and for a sample $\omega$ write $\Pi(\omega)$ for
the transcript of $p$ run on the input pair $(X, Y)$ generated by $\omega$. Let
$Z(\omega) = \big(\Pi(\omega),\, T,\, X_{<T},\, Y_{>T}\big)$ be the statistic recording
the transcript, the special coordinate, Alice's bits strictly before $T$, and Bob's bits
strictly after $T$. For each value $z$ that $Z$ can take, let $D(z)$ be the
Kullback–Leibler divergence, from the uniform distribution on one bit, of the
conditional law of Alice's special bit $X_T$ given $Z = z$, where the conditioning is
taken with respect to the ambient uniform law on samples. Then
\[
\int D\big(Z(\omega)\big)\, d\mu(\omega) \;\le\; \tfrac{3}{2}\, I\!\big(X_T \,;\, \Pi
\,\big|\, T,\, X_{<T},\, Y_{\ge T}\big),
\]
where the conditional mutual information on the right is computed under the hard
distribution conditioned only on the two generated sets being disjoint.
\end{theorem}

\begin{theorem}[Average KL cost of Alice's special bit under small claimed information]
\lean{CommunicationComplexity.Functions.Disjointness.RandomizedLowerBound.integral_xFiberKL_disjointSpecialYFalse_le_two_mul_gamma_pow_four}
\label{CommunicationComplexity.Functions.Disjointness.RandomizedLowerBound.integral_xFiberKL_disjointSpecialYFalse_le_two_mul_gamma_pow_four}
\leanok
\uses{CommunicationComplexity.Functions.Disjointness.RandomizedLowerBound.ProtocolType, CommunicationComplexity.Functions.Disjointness.RandomizedLowerBound.aliceInfoTerm_le_claimInfo, CommunicationComplexity.Functions.Disjointness.RandomizedLowerBound.claimInfo, CommunicationComplexity.Functions.Disjointness.RandomizedLowerBound.disjointSpecialYFalseMeasure, CommunicationComplexity.Functions.Disjointness.RandomizedLowerBound.integral_xFiberKL_disjointSpecialYFalse_le_three_halves_mul_aliceInfoTerm, CommunicationComplexity.Functions.Disjointness.RandomizedLowerBound.xFiberKL, CommunicationComplexity.Functions.Disjointness.RandomizedLowerBound.zVariable}
\difficulty{2}
Fix a positive integer $n$ and a deterministic two-party communication protocol $p$ for
set-disjointness on $n$ coordinates, and let $\gamma > 0$. For a sample $\omega$ of the
hard distribution, write $T$ for its special coordinate, $X_T$ and $Y_T$ for Alice's and
Bob's bits at $T$, $\Pi$ for the transcript produced by running $p$ on the induced input
pair $(X(\omega), Y(\omega))$, and $Z = (\Pi,\, T,\, X_{<T},\, Y_{>T})$ for the
transcript together with the special coordinate, Alice's input bits strictly before $T$,
and Bob's input bits strictly after $T$. Suppose that, under the hard distribution
conditioned on the two input sets being disjoint, the total claimed information
satisfies
\[
I\big(X_T ; \Pi \,\big|\, T,\, X_{<T},\, Y_{\ge T}\big) \;+\; I\big(Y_T ; \Pi \,\big|\,
T,\, X_{\le T},\, Y_{>T}\big) \;\le\; \tfrac{2}{3}\,\gamma^4 .
\]
Then, taking the expectation over the same disjoint-conditioned hard distribution
further conditioned on $Y_T$ being $\mathrm{false}$, the mean Kullback–Leibler
divergence of the conditional law of $X_T$ given the observed value of $Z$ from the
uniform distribution on one bit satisfies
\[
\E\Big[\, D_{\mathrm{KL}}\big(\operatorname{law}(X_T \mid Z) \,\big\|\,
\mathrm{Unif}\{0,1\}\big) \,\Big] \;\le\; 2\,\gamma^4 .
\]
\end{theorem}

\begin{theorem}[Averaged squared bias of Alice's special bit under disjointness]
\lean{CommunicationComplexity.Functions.Disjointness.RandomizedLowerBound.integral_xDistance_sq_disjointSpecialYFalse_le_gamma_pow_four}
\label{CommunicationComplexity.Functions.Disjointness.RandomizedLowerBound.integral_xDistance_sq_disjointSpecialYFalse_le_gamma_pow_four}
\leanok
\uses{CommunicationComplexity.Functions.Disjointness.RandomizedLowerBound.ProtocolType, CommunicationComplexity.Functions.Disjointness.RandomizedLowerBound.claimInfo, CommunicationComplexity.Functions.Disjointness.RandomizedLowerBound.disjointSpecialYFalseMeasure, CommunicationComplexity.Functions.Disjointness.RandomizedLowerBound.integral_xFiberKL_disjointSpecialYFalse_le_two_mul_gamma_pow_four, CommunicationComplexity.Functions.Disjointness.RandomizedLowerBound.two_mul_integral_xDistance_sq_le_integral_xFiberKL_disjointSpecialYFalse, CommunicationComplexity.Functions.Disjointness.RandomizedLowerBound.xDistance, CommunicationComplexity.Functions.Disjointness.RandomizedLowerBound.zVariable}
\difficulty{2}
Fix an input length $n \ge 1$, a deterministic communication protocol $p$ for
disjointness on inputs of length $n$, and a real number $\gamma > 0$. Draw a sample
$\omega$ from the hard distribution $\mathbb{P}$; write $M$ for the transcript produced
by $p$ on the associated input pair, write $T$ for the special coordinate, and set $Z =
(M,\, T,\, X_{<T},\, Y_{>T})$, collecting the transcript together with $T$, Alice's bits
strictly before $T$, and Bob's bits strictly after $T$. For each value $z$ of $Z$, let
$\Delta(z)$ denote the total variation distance between the conditional law of Alice's
special bit $X_T$ given $Z = z$ under $\mathbb{P}$ and the uniform law on one bit.
Suppose the combined conditional mutual information
\[
I\!\left(X_T : M \,\middle|\, T,\, X_{<T},\, Y_{\ge T}\right) + I\!\left(Y_T : M
\,\middle|\, T,\, X_{\le T},\, Y_{>T}\right) \;\le\; \tfrac{2}{3}\,\gamma^4,
\]
where both terms are computed under $\mathbb{P}$ conditioned on the disjointness event
$D$ (that Alice's and Bob's input sets are disjoint). Then, with the expectation taken
under $\mathbb{P}$ conditioned on $D$ together with $Y_T = 0$,
\[
  \E\!\left[\Delta(Z)^2 \,\middle|\, D,\ Y_T = 0\right] \;\le\; \gamma^4.
\]
\end{theorem}

\begin{theorem}[Mass of the Alice-bit deviation event on the zero–zero slice]
\lean{CommunicationComplexity.Functions.Disjointness.RandomizedLowerBound.measureReal_specialZeroZero_inter_xDistance_bad_le}
\label{CommunicationComplexity.Functions.Disjointness.RandomizedLowerBound.measureReal_specialZeroZero_inter_xDistance_bad_le}
\leanok
\uses{CommunicationComplexity.Functions.Disjointness.RandomizedLowerBound.HardSample, CommunicationComplexity.Functions.Disjointness.RandomizedLowerBound.ProtocolType, CommunicationComplexity.Functions.Disjointness.RandomizedLowerBound.claimInfo, CommunicationComplexity.Functions.Disjointness.RandomizedLowerBound.disjointSpecialYFalseMeasure, CommunicationComplexity.Functions.Disjointness.RandomizedLowerBound.disjointSpecialYFalseMeasure_integral_xDistance_le_of_integral_sq_le, CommunicationComplexity.Functions.Disjointness.RandomizedLowerBound.disjointSpecialYFalseMeasure_xDistance_bad_le_of_integral_le, CommunicationComplexity.Functions.Disjointness.RandomizedLowerBound.integral_xDistance_sq_disjointSpecialYFalse_le_gamma_pow_four, CommunicationComplexity.Functions.Disjointness.RandomizedLowerBound.measureReal_specialZeroZero, CommunicationComplexity.Functions.Disjointness.RandomizedLowerBound.specialZeroZero, CommunicationComplexity.Functions.Disjointness.RandomizedLowerBound.volume_cond_specialZeroZero_measureReal_le_two_mul_disjointSpecialYFalseMeasure, CommunicationComplexity.Functions.Disjointness.RandomizedLowerBound.xDistance, CommunicationComplexity.Functions.Disjointness.RandomizedLowerBound.zVariable}
\difficulty{5}
Fix a positive integer $n$ and a deterministic communication protocol $p$ for the
length-$n$ set-disjointness problem, and let $\gamma > 0$. Work under the hard
distribution on samples $\omega$, each of which determines a special coordinate $T$,
Alice's and Bob's input sets $X$ and $Y$ (with special bits $x_T$ and $y_T$), and the
transcript $M$ produced by $p$ on input $(X,Y)$. For a value $z$ of the variable $Z =
(M, T, X_{<T}, Y_{>T})$, let $\delta_X(z)$ be the total variation distance between the
conditional law of $x_T$ given $\{Z = z\}$ and the uniform distribution on one bit.
Suppose that, under the hard distribution conditioned on the disjointness event $X \cap
Y = \varnothing$,
\[ I\big(x_T ; M \mid T,\, X_{<T},\, Y_{\ge T}\big) \;+\; I\big(y_T ; M \mid T,\, X_{\le T},\, Y_{>T}\big) \;\le\; \frac{2\gamma^4}{3}. \]
Then the hard distribution assigns mass at most $\gamma/2$ to the set of samples for
which both special bits vanish, $x_T = y_T = 0$, and $\delta_X(Z) > \gamma$.
\end{theorem}

\begin{theorem}[Probability bound for a biased special y-bit]
\lean{CommunicationComplexity.Functions.Disjointness.RandomizedLowerBound.measureReal_specialZeroZero_inter_yDistance_bad_le}
\label{CommunicationComplexity.Functions.Disjointness.RandomizedLowerBound.measureReal_specialZeroZero_inter_yDistance_bad_le}
\leanok
\uses{CommunicationComplexity.Functions.Disjointness.RandomizedLowerBound.ProtocolType, CommunicationComplexity.Functions.Disjointness.RandomizedLowerBound.claimInfo, CommunicationComplexity.Functions.Disjointness.RandomizedLowerBound.claimInfo_dualProtocol, CommunicationComplexity.Functions.Disjointness.RandomizedLowerBound.dualProtocol, CommunicationComplexity.Functions.Disjointness.RandomizedLowerBound.measureReal_specialZeroZero_inter_xDistance_bad_le, CommunicationComplexity.Functions.Disjointness.RandomizedLowerBound.specialZeroZero, CommunicationComplexity.Functions.Disjointness.RandomizedLowerBound.volume_specialZeroZero_inter_xDistance_dualProtocol_eq_yDistance, CommunicationComplexity.Functions.Disjointness.RandomizedLowerBound.yDistance, CommunicationComplexity.Functions.Disjointness.RandomizedLowerBound.zVariable}
\difficulty{3}
Fix a positive integer $n$ and a deterministic two-party protocol $p$ for the
set-disjointness problem in which Alice and Bob hold subsets $X, Y \subseteq \{1, \dots,
n\}$. Draw a sample $\omega$ uniformly from the finite hard-sample space; a draw fixes a
special coordinate $T$, Alice's and Bob's bits $x_T$ and $y_T$ at the special
coordinate, and the two input sets $X(\omega)$ and $Y(\omega)$. Let $M(\omega)$ be the
transcript produced by $p$ on the input $(X(\omega), Y(\omega))$, and let
\[ Z(\omega) = \bigl(M(\omega),\, T,\, x_{<T},\, y_{>T}\bigr) \]
record the transcript, the special coordinate, Alice's bits at coordinates strictly
below $T$, and Bob's bits at coordinates strictly above $T$. For a value $z$ taken by
$Z$, let $\delta(z)$ be the total-variation distance between the uniform law on a single
bit and the conditional law of $y_T$ given $Z = z$, the latter formed under the uniform
law on the hard-sample space. Write $I(\,\cdot\,;\,\cdot \mid \cdot\,)$ for conditional
mutual information computed under the measure obtained by conditioning the uniform law
on the event that $X(\omega)$ and $Y(\omega)$ are disjoint. Let $\gamma > 0$ be a real
number for which
\[ I\bigl(x_T ; M \mid T,\, x_{<T},\, y_{\ge T}\bigr) + I\bigl(y_T ; M \mid T,\, x_{\le T},\, y_{>T}\bigr) \le \tfrac{2}{3}\gamma^{4}. \]
Then the probability, under the uniform law on the hard-sample space, that $x_T = y_T =
0$ and $\delta(Z(\omega)) > \gamma$ is at most $\gamma/2$.
\end{theorem}

\begin{theorem}[Good-fiber probability under small claim information]
\lean{CommunicationComplexity.Functions.Disjointness.RandomizedLowerBound.one_div_four_mul_one_sub_four_mul_le_measureReal_goodZEvent}
\label{CommunicationComplexity.Functions.Disjointness.RandomizedLowerBound.one_div_four_mul_one_sub_four_mul_le_measureReal_goodZEvent}
\leanok
\uses{CommunicationComplexity.Functions.Disjointness.RandomizedLowerBound.HardSample, CommunicationComplexity.Functions.Disjointness.RandomizedLowerBound.ProtocolType, CommunicationComplexity.Functions.Disjointness.RandomizedLowerBound.claimInfo, CommunicationComplexity.Functions.Disjointness.RandomizedLowerBound.goodZEvent, CommunicationComplexity.Functions.Disjointness.RandomizedLowerBound.measureReal_specialZeroZero, CommunicationComplexity.Functions.Disjointness.RandomizedLowerBound.measureReal_specialZeroZero_inter_xDistance_bad_le, CommunicationComplexity.Functions.Disjointness.RandomizedLowerBound.measureReal_specialZeroZero_inter_yDistance_bad_le, CommunicationComplexity.Functions.Disjointness.RandomizedLowerBound.mem_goodZEvent_of_xDistance_yDistance_le, CommunicationComplexity.Functions.Disjointness.RandomizedLowerBound.specialZeroZero, CommunicationComplexity.Functions.Disjointness.RandomizedLowerBound.xDistance, CommunicationComplexity.Functions.Disjointness.RandomizedLowerBound.yDistance, CommunicationComplexity.Functions.Disjointness.RandomizedLowerBound.zVariable}
\difficulty{5}
Fix a positive integer $n$ and a deterministic two-party protocol $p$ for the
set-disjointness problem on subsets of $\{1,\dots,n\}$. The hard distribution is the
uniform probability measure on the finite sample space; a sample $\omega$ determines a
special coordinate $T \in \{1,\dots,n\}$, Alice's input set $X$ and Bob's input set $Y$,
and the two special-coordinate bits $X_T$ and $Y_T$. Let $M$ denote the transcript that
$p$ produces on the pair $(X,Y)$, and let $Z = (M,\,T,\,X_{<T},\,Y_{>T})$ be the
statistic recording the transcript, the special coordinate, Alice's input bits strictly
before the special coordinate, and Bob's input bits strictly after the special
coordinate. Say that a value $z$ is \emph{good} when the law of $(X_T,Y_T)$ under the
hard distribution conditioned on the fiber $\{Z = z\}$ lies within total variation
distance $2\gamma$ of the uniform distribution on $\{0,1\}^2$. Suppose $\gamma > 0$ and
that, under the hard distribution conditioned on the event that $X$ and $Y$ are
disjoint,
\[
I\!\left(X_T ; M \mid T,\,X_{<T},\,Y_{\ge T}\right) + I\!\left(Y_T ; M \mid T,\,X_{\le
T},\,Y_{>T}\right) \le \frac{2\gamma^{4}}{3}.
\]
Then the hard-distribution probability that $Z$ is good is at least
$\tfrac14\,(1-4\gamma)$.
\end{theorem}

\begin{theorem}[Conditioning on a fiber is supported on that fiber]
\lean{CommunicationComplexity.Functions.Disjointness.RandomizedLowerBound.zFiberMeasure_inter_fiber}
\label{CommunicationComplexity.Functions.Disjointness.RandomizedLowerBound.zFiberMeasure_inter_fiber}
\leanok
\uses{CommunicationComplexity.Functions.Disjointness.RandomizedLowerBound.HardSample, CommunicationComplexity.Functions.Disjointness.RandomizedLowerBound.ProtocolType, CommunicationComplexity.Functions.Disjointness.RandomizedLowerBound.ZType, CommunicationComplexity.Functions.Disjointness.RandomizedLowerBound.zFiber, CommunicationComplexity.Functions.Disjointness.RandomizedLowerBound.zFiberMeasure}
\difficulty{2}
Fix a length $n \ge 1$ and a deterministic protocol $p$ for the disjointness problem on
inputs indexed by $\{0,\dots,n-1\}$, and let $z$ be an achievable value of the coarse
variable $Z = (M, T, X_{<T}, Y_{>T})$ recorded by the protocol. Write $F_z$ for the
fiber of $Z$ above $z$, that is, the set of hard-distribution samples $\omega$ with
$Z(\omega) = z$, and let $\mu_z(\cdot) = \operatorname{vol}(\cdot \mid F_z)$ denote the
hard distribution conditioned on the event $F_z$. Then for every set $S$ of samples,
\[
  \mu_z(F_z \cap S) = \mu_z(S),
\]
where both sides are taken as real numbers. In other words, restricting to the
conditioning event $F_z$ leaves the conditional measure unchanged, so $\mu_z$ is
supported on $F_z$.
\end{theorem}

\begin{theorem}[Singleton masses of a good conditional special-pair law]
\lean{CommunicationComplexity.Functions.Disjointness.RandomizedLowerBound.abs_conditionalSpecialPairLaw_singleton_sub_quarter_le}
\label{CommunicationComplexity.Functions.Disjointness.RandomizedLowerBound.abs_conditionalSpecialPairLaw_singleton_sub_quarter_le}
\leanok
\uses{CommunicationComplexity.Functions.Disjointness.RandomizedLowerBound.ProtocolType, CommunicationComplexity.Functions.Disjointness.RandomizedLowerBound.ZType, CommunicationComplexity.Functions.Disjointness.RandomizedLowerBound.conditionalSpecialPairLaw, CommunicationComplexity.Functions.Disjointness.RandomizedLowerBound.goodZ, CommunicationComplexity.Functions.Disjointness.RandomizedLowerBound.uniformBoolPair, CommunicationComplexity.Functions.Disjointness.RandomizedLowerBound.uniformBoolPair_singleton, CommunicationComplexity.Functions.Disjointness.RandomizedLowerBound.zDistance, CommunicationComplexity.TVDistance.abs_measureReal_sub_le_tvDistance, CommunicationComplexity.tvDistance}
\difficulty{3}
Fix a positive integer $n$ and a deterministic communication protocol $p$ for the
disjointness problem on inputs of length $n$, drawn under the hard distribution. Write
$(x_T, y_T)$ for the pair of Alice's and Bob's bits at the special coordinate $T$, and
for a value $z$ of the conditioning statistic let $\mu_z$ denote the conditional law of
this bit-pair on $\{0,1\}^2$ given that the statistic equals $z$. Suppose $z$ is
$\gamma$-good, in the sense that the total variation distance between $\mu_z$ and the
uniform distribution on $\{0,1\}^2$ is at most $2\gamma$. Then for every bit-pair $b \in
\{0,1\}^2$,
\[
  \abs{\mu_z(\{b\}) - \tfrac{1}{4}} \le 2\gamma .
\]
\end{theorem}

\begin{theorem}[Uniform lower bound on singleton masses of a good conditional law]
\lean{CommunicationComplexity.Functions.Disjointness.RandomizedLowerBound.quarter_sub_two_mul_le_conditionalSpecialPairLaw_singleton}
\label{CommunicationComplexity.Functions.Disjointness.RandomizedLowerBound.quarter_sub_two_mul_le_conditionalSpecialPairLaw_singleton}
\leanok
\uses{CommunicationComplexity.Functions.Disjointness.RandomizedLowerBound.ProtocolType, CommunicationComplexity.Functions.Disjointness.RandomizedLowerBound.ZType, CommunicationComplexity.Functions.Disjointness.RandomizedLowerBound.abs_conditionalSpecialPairLaw_singleton_sub_quarter_le, CommunicationComplexity.Functions.Disjointness.RandomizedLowerBound.conditionalSpecialPairLaw, CommunicationComplexity.Functions.Disjointness.RandomizedLowerBound.goodZ}
\difficulty{2}
Fix a deterministic disjointness protocol $p$ on inputs of length $n$, a real number
$\gamma$, and a value $z$ of the transcript-statistic. Suppose $z$ is *good* for
$\gamma$, meaning that the conditional special-pair law given $z$ — the distribution of
the special-coordinate bit-pair $(x_T, y_T)$ conditioned on the value $z$ — lies within
total variation distance $2\gamma$ of the uniform law on bit-pairs. Then for every pair
$b \in \{0,1\} \times \{0,1\}$, this conditional special-pair law assigns to the
singleton $\{b\}$ a mass of at least $\tfrac14 - 2\gamma$.
\end{theorem}

\begin{theorem}[Diagonal special-pair samples in a fixed transcript fiber are protocol errors]
\lean{CommunicationComplexity.Functions.Disjointness.RandomizedLowerBound.zFiber_inter_diag_specialPair_subset_protocolErrorEvent}
\label{CommunicationComplexity.Functions.Disjointness.RandomizedLowerBound.zFiber_inter_diag_specialPair_subset_protocolErrorEvent}
\leanok
\uses{CommunicationComplexity.Deterministic.Protocol.run, CommunicationComplexity.Functions.Disjointness.RandomizedLowerBound.ProtocolType, CommunicationComplexity.Functions.Disjointness.RandomizedLowerBound.X, CommunicationComplexity.Functions.Disjointness.RandomizedLowerBound.Y, CommunicationComplexity.Functions.Disjointness.RandomizedLowerBound.ZType, CommunicationComplexity.Functions.Disjointness.RandomizedLowerBound.disjoint_X_Y_iff, CommunicationComplexity.Functions.Disjointness.RandomizedLowerBound.protocolErrorEvent, CommunicationComplexity.Functions.Disjointness.RandomizedLowerBound.run_eq_zOutput_of_zVariable_eq, CommunicationComplexity.Functions.Disjointness.RandomizedLowerBound.specialPair, CommunicationComplexity.Functions.Disjointness.RandomizedLowerBound.specialX, CommunicationComplexity.Functions.Disjointness.RandomizedLowerBound.specialY, CommunicationComplexity.Functions.Disjointness.RandomizedLowerBound.zFiber, CommunicationComplexity.Functions.Disjointness.RandomizedLowerBound.zOutput, CommunicationComplexity.Functions.Disjointness.disjointness}
\difficulty{4}
Fix a positive integer $n$ and a deterministic two-party communication protocol $p$ that
outputs a Boolean on pairs of subsets of $[n]$, intended to compute set disjointness.
For a sample $\omega$ of the hard distribution, let $Z(\omega)$ denote the statistic
recording the transcript produced by running $p$ on the induced input pair $(X_\omega,
Y_\omega)$, together with the special coordinate $T$, Alice's input bits at coordinates
before $T$, and Bob's input bits at coordinates after $T$. Fix one value $z$ of this
statistic, and let $b$ be the Boolean output attached to the transcript stored in $z$.
Then every sample $\omega$ with $Z(\omega) = z$ whose bits at the special coordinate
satisfy $(\text{Alice's special bit}, \text{Bob's special bit}) = (b, b)$ belongs to the
protocol error event; that is, the value $p$ computes on $(X_\omega, Y_\omega)$ differs
from the true disjointness value of $X_\omega$ and $Y_\omega$.
\end{theorem}

\begin{theorem}[Conditional protocol-error probability on a good fiber]
\lean{CommunicationComplexity.Functions.Disjointness.RandomizedLowerBound.quarter_sub_two_mul_le_zFiberMeasure_protocolErrorEvent}
\label{CommunicationComplexity.Functions.Disjointness.RandomizedLowerBound.quarter_sub_two_mul_le_zFiberMeasure_protocolErrorEvent}
\leanok
\uses{CommunicationComplexity.Functions.Disjointness.RandomizedLowerBound.ProtocolType, CommunicationComplexity.Functions.Disjointness.RandomizedLowerBound.ZType, CommunicationComplexity.Functions.Disjointness.RandomizedLowerBound.conditionalSpecialPairLaw_singleton, CommunicationComplexity.Functions.Disjointness.RandomizedLowerBound.goodZ, CommunicationComplexity.Functions.Disjointness.RandomizedLowerBound.protocolErrorEvent, CommunicationComplexity.Functions.Disjointness.RandomizedLowerBound.quarter_sub_two_mul_le_conditionalSpecialPairLaw_singleton, CommunicationComplexity.Functions.Disjointness.RandomizedLowerBound.specialPair, CommunicationComplexity.Functions.Disjointness.RandomizedLowerBound.zFiber, CommunicationComplexity.Functions.Disjointness.RandomizedLowerBound.zFiberMeasure, CommunicationComplexity.Functions.Disjointness.RandomizedLowerBound.zFiberMeasure_inter_fiber, CommunicationComplexity.Functions.Disjointness.RandomizedLowerBound.zFiber_inter_diag_specialPair_subset_protocolErrorEvent, CommunicationComplexity.Functions.Disjointness.RandomizedLowerBound.zOutput}
\difficulty{4}
Fix a positive integer $n$ and consider the set-disjointness problem on subsets of $[n]$
under the hard distribution: a sample $\omega$ specifies a special coordinate $T \in
[n]$, a bit-pair $(x_T, y_T)$ at $T$, and at each remaining coordinate one of the three
options *neither*, *Alice only*, or *Bob only*; from $\omega$ one reads off Alice's set
$X(\omega)$ and Bob's set $Y(\omega)$. Let $p$ be a deterministic communication protocol
with Boolean output on such inputs, and let $Z(\omega)$ be the statistic recording the
transcript of $p$ on $(X(\omega), Y(\omega))$, the special coordinate $T$, Alice's bits
at the coordinates below $T$, and Bob's bits at the coordinates above $T$. Let $\gamma
\in \bbr$, and suppose the value $z$ of $Z$ is *good*, meaning that the total variation
distance between the conditional law of the special bit-pair $(x_T, y_T)$ given $Z = z$
and the uniform distribution on $\{0,1\}^2$ is at most $2\gamma$. Then, conditioning on
the event $Z = z$, the probability that $p$ errs on disjointness — that its output on
$(X(\omega), Y(\omega))$ disagrees with whether $X(\omega)$ and $Y(\omega)$ are disjoint
— is at least $\tfrac{1}{4} - 2\gamma$.
\end{theorem}

\begin{theorem}[Error lower bound from the good-summary event]
\lean{CommunicationComplexity.Functions.Disjointness.RandomizedLowerBound.goodZEvent_mul_quarter_sub_two_mul_le_protocolErrorEvent}
\label{CommunicationComplexity.Functions.Disjointness.RandomizedLowerBound.goodZEvent_mul_quarter_sub_two_mul_le_protocolErrorEvent}
\leanok
\uses{CommunicationComplexity.FiniteMeasureSpace.measureReal_eq_sum_cond_fiber_real, CommunicationComplexity.FiniteMeasureSpace.measureReal_preimage_eq_sum_fibers, CommunicationComplexity.Functions.Disjointness.RandomizedLowerBound.ProtocolType, CommunicationComplexity.Functions.Disjointness.RandomizedLowerBound.goodZ, CommunicationComplexity.Functions.Disjointness.RandomizedLowerBound.goodZEvent, CommunicationComplexity.Functions.Disjointness.RandomizedLowerBound.protocolErrorEvent, CommunicationComplexity.Functions.Disjointness.RandomizedLowerBound.quarter_sub_two_mul_le_zFiberMeasure_protocolErrorEvent, CommunicationComplexity.Functions.Disjointness.RandomizedLowerBound.zVariable}
\difficulty{4}
Fix a positive integer $n$ and a deterministic two-party communication protocol $p$ that
attempts to decide set disjointness on pairs of subsets of $[n]=\{1,\dots,n\}$, and let
$\gamma\in\bbr$. Draw a sample $\omega$ from the hard distribution, which produces
inputs $X(\omega),Y(\omega)\subseteq[n]$ together with a distinguished special
coordinate $T$, and let $Z(\omega)$ be the summary statistic consisting of the
transcript of $p$ on $(X(\omega),Y(\omega))$, the coordinate $T$, Alice's input bits
strictly before $T$, and Bob's input bits strictly after $T$. Call $\omega$ \emph{good}
when the conditional law of the special bit-pair $(x_T,y_T)$, given the value of
$Z(\omega)$, lies within total variation distance $2\gamma$ of the uniform law on
$\{0,1\}^2$, and call $\omega$ \emph{erroneous} when $p$ outputs a value different from
the true disjointness of $X(\omega)$ and $Y(\omega)$. Then, with all probabilities taken
under the hard distribution,
\[
\Pr[\omega \text{ good}]\cdot\left(\tfrac14 - 2\gamma\right)\ \le\ \Pr[\omega \text{
erroneous}].
\]
\end{theorem}

\begin{theorem}[Good-event lower bound on distributional error for set disjointness]
\lean{CommunicationComplexity.Functions.Disjointness.RandomizedLowerBound.goodZEvent_mul_quarter_sub_two_mul_le_distributionalError}
\label{CommunicationComplexity.Functions.Disjointness.RandomizedLowerBound.goodZEvent_mul_quarter_sub_two_mul_le_distributionalError}
\leanok
\uses{CommunicationComplexity.Deterministic.Protocol.distributionalError, CommunicationComplexity.Functions.Disjointness.RandomizedLowerBound.ProtocolType, CommunicationComplexity.Functions.Disjointness.RandomizedLowerBound.distributionalError_inputDist_eq_protocolErrorEvent, CommunicationComplexity.Functions.Disjointness.RandomizedLowerBound.goodZEvent, CommunicationComplexity.Functions.Disjointness.RandomizedLowerBound.goodZEvent_mul_quarter_sub_two_mul_le_protocolErrorEvent, CommunicationComplexity.Functions.Disjointness.RandomizedLowerBound.inputDist, CommunicationComplexity.Functions.Disjointness.disjointness}
\difficulty{2}
Fix a positive integer $n$, a deterministic protocol $p$ that computes set disjointness
on pairs of subsets of $[n]$ (Alice and Bob each holding one subset), and a real number
$\gamma$. Draw a sample $\omega$ from the hard distribution, and let $Z(\omega)$ be the
summary consisting of the transcript that $p$ produces on the associated input pair
$(X(\omega),Y(\omega))$, the special coordinate $T$, Alice's input bits at the
coordinates below $T$, and Bob's input bits at the coordinates above $T$. Call $\omega$
\emph{good} if, conditioned on the value of $Z(\omega)$, the law of the special bit-pair
$(x_T,y_T)$ lies within total variation distance $2\gamma$ of the uniform distribution
on $\{0,1\}^2$. Then
\[
\Pr[\omega \text{ is good}]\cdot\left(\tfrac14 - 2\gamma\right) \le \mathrm{err}(p),
\]
where $\mathrm{err}(p)$ is the distributional error of $p$ against the disjointness
function under the input distribution induced by the hard distribution.
\end{theorem}

\begin{theorem}[Information lower bound for low-error disjointness protocols]
\lean{CommunicationComplexity.Functions.Disjointness.RandomizedLowerBound.one_div_32768_sq_lt_claimInfo_of_distributionalError_le}
\label{CommunicationComplexity.Functions.Disjointness.RandomizedLowerBound.one_div_32768_sq_lt_claimInfo_of_distributionalError_le}
\leanok
\uses{CommunicationComplexity.Deterministic.Protocol.distributionalError, CommunicationComplexity.Functions.Disjointness.RandomizedLowerBound.ProtocolType, CommunicationComplexity.Functions.Disjointness.RandomizedLowerBound.claimInfo, CommunicationComplexity.Functions.Disjointness.RandomizedLowerBound.goodZEvent_mul_quarter_sub_two_mul_le_distributionalError, CommunicationComplexity.Functions.Disjointness.RandomizedLowerBound.inputDist, CommunicationComplexity.Functions.Disjointness.RandomizedLowerBound.one_div_four_mul_one_sub_four_mul_le_measureReal_goodZEvent, CommunicationComplexity.Functions.Disjointness.disjointness}
\difficulty{3}
Fix a positive integer $n$, and let $p$ be a deterministic two-party communication
protocol whose inputs are subsets $X, Y \subseteq [n]$ (held by Alice and Bob,
respectively) and whose output is Boolean. Consider the hard input distribution
generated as follows: choose a special coordinate $T \in [n]$, bits $x_T, y_T \in
\{0,1\}$, and, for every coordinate $i \ne T$, a label in $\{\text{neither}, \text{left
only}, \text{right only}\}$, all uniformly and independently; a "left only" (resp.
"right only") label places $i$ in $X$ (resp. $Y$) and in no other set, "neither" places
it in neither, while at $T$ the bits $x_T, y_T$ decide membership of $T$ in $X$ and in
$Y$. Suppose that against this distribution the protocol $p$ computes set-disjointness —
outputting whether $X \cap Y = \emptyset$ — with distributional error at most $1/32$,
i.e.\ the probability that the output of $p$ on $(X, Y)$ disagrees with the disjointness
indicator is at most $1/32$. Then, writing $\Pi$ for the transcript produced by running
$p$ on $(X, Y)$ and conditioning throughout on the event that $X$ and $Y$ are disjoint,
\[
I\bigl(x_T \,;\, \Pi \,\bigm|\, T,\, X_{<T},\, Y_{\ge T}\bigr) \;+\; I\bigl(y_T \,;\,
\Pi \,\bigm|\, T,\, X_{\le T},\, Y_{>T}\bigr) \;>\; \Bigl(\tfrac{1}{32768}\Bigr)^{2},
\]
where $X_{<T}$ (resp. $X_{\le T}$) records Alice's membership bits at the coordinates
strictly below $T$ (resp. at most $T$), $Y_{\ge T}$ (resp. $Y_{>T}$) records Bob's
membership bits at the coordinates at least $T$ (resp. strictly above $T$), and all
remaining coordinates are set to $0$.
\end{theorem}

\begin{theorem}[Distributional communication lower bound for set disjointness]
\lean{CommunicationComplexity.Functions.Disjointness.RandomizedLowerBound.const_mul_n_le_complexity_of_distributionalError_le}
\statementsource{roughgarden-cc-bootcamp}{
  pdf-d54d402b16e2-p017-b001,
  pdf-d54d402b16e2-p017-b004
}
\label{CommunicationComplexity.Functions.Disjointness.RandomizedLowerBound.const_mul_n_le_complexity_of_distributionalError_le}
\leanok
\uses{CommunicationComplexity.Deterministic.Protocol.complexity, CommunicationComplexity.Deterministic.Protocol.distributionalError, CommunicationComplexity.Functions.Disjointness.RandomizedLowerBound.ProtocolType, CommunicationComplexity.Functions.Disjointness.RandomizedLowerBound.claimInfo_le_average_info_upper, CommunicationComplexity.Functions.Disjointness.RandomizedLowerBound.inputDist, CommunicationComplexity.Functions.Disjointness.RandomizedLowerBound.one_div_32768_sq_lt_claimInfo_of_distributionalError_le, CommunicationComplexity.Functions.Disjointness.disjointness}
\difficulty{4}
Fix a positive integer $n$, and let $\mu_n$ be the hard input distribution on pairs of
subsets $(X, Y)$ of $[n] = \{0, 1, \dots, n-1\}$. Let $p$ be a deterministic two-party
communication protocol for set disjointness on inputs from $[n]$, and suppose its
distributional error against the disjointness function under $\mu_n$ is at most
$\tfrac{1}{32}$, i.e.\ the $\mu_n$-probability that the output of $p$ disagrees with
$\mathrm{disjointness}(n, X, Y)$ is at most $\tfrac{1}{32}$. Then the communication
complexity of $p$ satisfies
\[
\frac{\left(1/32768\right)^2 \, n}{3 \ln 2} \;\le\; \mathrm{complexity}(p).
\]
\end{theorem}

\begin{theorem}[Linear lower bound for randomized disjointness communication]
\lean{CommunicationComplexity.Functions.Disjointness.RandomizedLowerBound.lt_publicCoin_communicationComplexity_disjointness_of_lt_const_mul_n}
\statementsource{roughgarden-cc-oneway}{
  pdf-043896d6545f-p012-b004,
  pdf-043896d6545f-p012-b005,
  pdf-043896d6545f-p013-b001
}
\label{CommunicationComplexity.Functions.Disjointness.RandomizedLowerBound.lt_publicCoin_communicationComplexity_disjointness_of_lt_const_mul_n}
\leanok
\uses{CommunicationComplexity.Deterministic.Protocol.complexity, CommunicationComplexity.Deterministic.Protocol.distributionalError, CommunicationComplexity.Functions.Disjointness.RandomizedLowerBound.const_mul_n_le_complexity_of_distributionalError_le, CommunicationComplexity.Functions.Disjointness.RandomizedLowerBound.inputDist, CommunicationComplexity.Functions.Disjointness.disjointness, CommunicationComplexity.PublicCoin.communicationComplexity, CommunicationComplexity.PublicCoin.lt_communicationComplexity_of_forall_distributionalError_gt}
\difficulty{3}
Fix a positive integer $n$, and let $\mathrm{DISJ}_n$ denote the set-disjointness
function that sends a pair of subsets $X, Y \subseteq [n]$ to $\mathtt{true}$ precisely
when $X \cap Y = \emptyset$. Write $\mathrm{R}^{\mathrm{pub}}_{1/32}(\mathrm{DISJ}_n)
\in \bbn \cup \{\infty\}$ for the public-coin randomized communication complexity of
$\mathrm{DISJ}_n$ at error $\tfrac{1}{32}$, i.e.\ the least worst-case number of bits
exchanged by any public-coin protocol that computes $\mathrm{DISJ}_n$ correctly with
probability at least $\tfrac{31}{32}$ on every input. Then for every natural number $k$
satisfying
\[
  k \;<\; \Bigl(\tfrac{1}{32768}\Bigr)^{2}\,\frac{n}{3\ln 2},
\]
one has $k < \mathrm{R}^{\mathrm{pub}}_{1/32}(\mathrm{DISJ}_n)$.
\end{theorem}

\begin{theorem}[Linear randomized lower bound for set disjointness]
\lean{CommunicationComplexity.Functions.Disjointness.RandomizedLowerBound.floor_div_pow_lt_publicCoin_communicationComplexity_disjointness}
\statementsource{roughgarden-cc-bootcamp}{pdf-d54d402b16e2-p014-b003}
\statementsource{roughgarden-cc-oneway}{
  pdf-043896d6545f-p012-b002,
  pdf-043896d6545f-p016-b001,
  pdf-043896d6545f-p013-b001
}
\label{CommunicationComplexity.Functions.Disjointness.RandomizedLowerBound.floor_div_pow_lt_publicCoin_communicationComplexity_disjointness}
\leanok
\uses{CommunicationComplexity.Functions.Disjointness.RandomizedLowerBound.lt_publicCoin_communicationComplexity_disjointness_of_lt_const_mul_n, CommunicationComplexity.Functions.Disjointness.disjointness, CommunicationComplexity.PublicCoin.communicationComplexity}
\difficulty{3}
Let $n$ be a positive integer, and consider the set-disjointness function on subsets of
$[n] = \{0, 1, \dots, n-1\}$, which sends a pair $X, Y \subseteq [n]$ to $\mathtt{true}$
precisely when $X \cap Y = \emptyset$. Then its $\tfrac{1}{32}$-error public-coin
randomized communication complexity $R^{\mathrm{pub}}_{1/32}$ satisfies
\[
\left\lfloor \frac{n}{2^{32}} \right\rfloor \;<\;
R^{\mathrm{pub}}_{1/32}\bigl(\mathrm{disjointness}_n\bigr),
\]
where the left-hand side is the floor of the real number $n / 2^{32}$ and the right-hand
side is valued in $\bbn_\infty$.
\end{theorem}
